% Topological dynamics of topological groups: a living, hypertextual encyclopaedia.
% Main file for both outputs:
%   * web: plasTeX + plastexdepgraph  (see build.sh)
%   * PDF: pdflatex + bibtex
% Conventions for authors (human or AI) are in ../AGENTS.md.
\documentclass[11pt]{article}
\usepackage[utf8]{inputenc}
\usepackage[T1]{fontenc}
\usepackage{amsmath,amssymb,amsthm}
\usepackage{hyperref}
\usepackage{depgraph}
\usepackage{graphicx}
% Figures: one SVG file per figure in ../figures, drawn by figures/make_figures.py.
% build.sh copies them into _figures/ as PDF (for pdflatex) and PNG (for plasTeX);
% the website then shows the SVG. Figures are numbered within sections.
\graphicspath{{_figures/}}
\numberwithin{figure}{section}
% \illustrates{label,...}, inside a figure: the statements whose proofs it
% illustrates. It prints nothing; on both websites their links name the figure.
\newcommand{\illustrates}[1]{}

\theoremstyle{plain}
\newtheorem{theorem}{Theorem}[section]
\newtheorem{proposition}[theorem]{Proposition}
\newtheorem{corollary}[theorem]{Corollary}
\newtheorem{lemma}[theorem]{Lemma}
\newtheorem{conjecture}[theorem]{Conjecture}
\theoremstyle{definition}
\newtheorem{definition}[theorem]{Definition}
\newtheorem{example}[theorem]{Example}
\theoremstyle{remark}
\newtheorem{remark}[theorem]{Remark}

\title{Topological dynamics of topological groups}
\author{A living, hypertextual encyclopaedia. Editor: LogicArtificialis}
\date{Prototype, version of 26 November 2026}

\begin{document}
\maketitle

\section{About this encyclopaedia}\label{sec:about}

This is a prototype of a living, hypertextual reference work on the topological
dynamics of topological groups, with an emphasis on universal minimal flows. It
is meant to grow: entries are added and revised over time, and every
definition, theorem and remark keeps a permanent \emph{tag}, a four-character
identifier that never changes and is never reused, even when the item moves,
is renumbered or is reworded. Cite results by tag, not by number.

Each statement lists the earlier statements it depends on; together these
dependencies form the dependency graph of the encyclopaedia, which can be
browsed from the navigation menu. The interactive proof tree at
\url{https://topologicaldynamics.org/proof-tree/} draws every statement with
what it rests on; selecting a result shows its statement, its proof and its
proof tree. The proof trees of the main results, the theorems that no other
statement uses, are collected at
\url{https://topologicaldynamics.org/proof-tree/main/}.

\paragraph{Conventions.} Topological groups and compact spaces are Hausdorff.
Actions are left actions. Conventions on left and right uniform structures
differ across the literature; the ones used here are fixed in
Definition~\ref{def:ruc}.

\paragraph{Sources and contributing.} The encyclopaedia is built from the
public repository
\url{https://github.com/eulogica/topdyn-encyclopaedia}. Every section links
to its source files there, with their history. Corrections and new entries are
proposed by pull request and reviewed by the editor; the repository's
\texttt{CONTRIBUTING.md} file explains how.

\paragraph{Status of the seed entries.} The first entries were drafted with
the help of an AI assistant from the cited literature. Statements and
attributions were checked against the original papers' abstracts and
statements and against standard secondary sources, and every bibliographic
record was checked against the publisher or arXiv. They have not yet been
reviewed by the editor or by experts. Sections~\ref{sec:kpt-umf}
to~\ref{sec:ea-catalogue} were added later in the same way, from the full
texts of the cited papers where available and otherwise from their abstracts,
and have not been reviewed either. Sections~\ref{sec:maps-transformations}
to~\ref{sec:inductive-limits} report results from preprints
posted on Zenodo, which have not been refereed. Sections~\ref{sec:proofs-maps-transformations}
to~\ref{sec:proofs-inductive-limits} reproduce the proofs from those preprints,
in the versions cited in the bibliography; the bibliographic data of the works
they cite were checked against Crossref, zbMATH or arXiv.

\paragraph{The editor.} The encyclopaedia is edited by LogicArtificialis
(\url{https://logicartificialis.org}). Its seal and one of its emblems,
reproduced from its website, are shown below.

\begin{center}
\emblem{seal}{A gold circular seal with a sapphire turtle, a corgi, a capybara
  and the Latin motto Organum rationis, sapphirus mentis, above the wordmark
  LogicArtificialis.}\hspace{2em}%
\emblem{meaning-guides}{A corgi and a capybara riding a green turtle above the
  wordmark LogicArtificialis and the motto Meaning guides. Syntax decides.}
\end{center}

\paragraph{Funding.} This work is supported by ERC grant DAT (Definable
Algebraic Topology), 101077154.

\begin{center}
\emblem[2.9cm]{dat-logo}{The DAT logo: the words Definable Algebraic Topology,
  with large magenta initials D, A and T.}\hspace{2em}%
\emblem[2.9cm]{erc-eu-logo}{Funded by the European Union, with the flag of the
  European Union, beside the logo of the European Research Council, established
  by the European Commission.}
\end{center}

Funded by the European Union. Views and opinions expressed are however those of
the author(s) only and do not necessarily reflect those of the European Union
or the European Research Council Executive Agency. Neither the European Union
nor the granting authority can be held responsible for them.

\section{Flows}\label{sec:flows}

Abstract topological dynamics studies continuous actions of topological groups
on compact spaces. The basic objects and morphisms are the following.

\begin{definition}[Flow]\label{def:flow}
Let $G$ be a topological group. A \emph{$G$-flow} is a compact space $X$
together with a continuous action $G\times X\to X$, $(g,x)\mapsto gx$. If $X$
and $Y$ are $G$-flows, a \emph{$G$-map} is a continuous map
$\varphi\colon X\to Y$ such that $\varphi(gx)=g\varphi(x)$ for all $g\in G$
and $x\in X$. A \emph{subflow} of $X$ is a nonempty closed $G$-invariant
subset of $X$. A $G$-flow $Y$ is a \emph{factor} of $X$ if there is a
surjective $G$-map $X\to Y$, and $X$ and $Y$ are \emph{isomorphic} if there is
a bijective $G$-map $X\to Y$.
\end{definition}

\begin{remark}\label{rem:flow-terminology}
\uses{def:flow}
By compactness, the inverse of a bijective $G$-map is continuous, so
isomorphic flows are isomorphic in every reasonable sense. Flows are also
called compact $G$-spaces, $G$-systems or dynamical systems in the literature.
\end{remark}

\begin{example}\label{ex:flows}
\uses{def:flow}
\begin{enumerate}
\item A compact group $G$ acting on itself by left translation.
\item A $\mathbb{Z}$-flow is the same thing as a compact space together with
  a homeomorphism $T$, the action being $(n,x)\mapsto T^n x$.
\item The group $\mathrm{Homeo}_+(\mathbb{S}^1)$ of orientation-preserving
  homeomorphisms of the circle, with the compact-open topology, acting on
  $\mathbb{S}^1$; see Section~\ref{sec:pestov}.
\item The group $S_\infty$ of all permutations of $\mathbb{N}$ acting on the
  space of linear orders on $\mathbb{N}$; see Section~\ref{sec:glasner-weiss}.
\end{enumerate}
\end{example}

\section{Minimal flows}\label{sec:minimal}

Minimal flows are the irreducible objects of topological dynamics: they cannot
be shrunk to a smaller subflow.

\begin{definition}[Minimal flow]\label{def:minimal}
\uses{def:flow}
A $G$-flow $X$ is \emph{minimal} if its only subflow is $X$ itself;
equivalently, the orbit $Gx$ is dense in $X$ for every $x\in X$.
\end{definition}

\begin{proposition}\label{prop:minimal-subflow}
\uses{def:minimal}
Every $G$-flow contains a minimal subflow.
\end{proposition}

\begin{proof}
Order the subflows of $X$ by reverse inclusion. The intersection of a chain of
subflows is closed and invariant, and it is nonempty by compactness, since the
chain has the finite intersection property. By Zorn's lemma there is a
minimal subflow.
\end{proof}

\begin{example}\label{ex:minimal}
\uses{def:minimal}
A fixed point is a minimal subflow. The rotation of the circle by an
irrational multiple of $2\pi$ is a minimal $\mathbb{Z}$-flow, since every orbit
is dense.
\end{example}

\begin{figure}[htbp]
\centering
\includegraphics[width=\linewidth]{irrational-rotation}
\caption{Example~\ref{ex:minimal}: the first $12$ and the first $100$ points $x$, $Rx$, $R^2x$, \dots{} of an orbit of the rotation $R$ of the circle by $2\pi\alpha$, where $\alpha=(\sqrt5-1)/2$ is irrational. Every orbit is dense, so the rotation is a minimal $\mathbb{Z}$-flow.}\label{fig:irrational-rotation}
\illustrates{ex:minimal}
\end{figure}

\section{The greatest ambit}\label{sec:ambit}

Every topological group has a largest flow with a distinguished point of dense
orbit. It is the Samuel compactification of the group, and it is where the
universal minimal flow lives.

\begin{definition}[Ambit]\label{def:ambit}
\uses{def:flow}
A \emph{$G$-ambit} is a pair $(X,x_0)$ consisting of a $G$-flow $X$ and a point
$x_0\in X$ whose orbit $Gx_0$ is dense in $X$. A \emph{morphism of ambits}
$(X,x_0)\to(Y,y_0)$ is a $G$-map sending $x_0$ to $y_0$. There is at most one
such morphism, since $Gx_0$ is dense, and it is surjective.
\end{definition}

\begin{definition}[Right uniformly continuous functions]\label{def:ruc}
A bounded function $f\colon G\to\mathbb{C}$ is \emph{right uniformly
continuous} if for every $\varepsilon>0$ there is a neighbourhood $V$ of the
identity such that $|f(vg)-f(g)|<\varepsilon$ for all $g\in G$ and $v\in V$.
These functions form a unital C*-subalgebra $\mathrm{RUC}_b(G)$ of
$\ell^\infty(G)$.
\end{definition}

\begin{remark}[Conventions]\label{rem:ruc-conventions}
\uses{def:ruc}
These are the functions that are uniformly continuous for the uniform
structure with basic entourages $\{(x,y) : xy^{-1}\in V\}$. This structure is
called the right uniformity by Bourbaki and Pestov and the left uniformity by
Ellis and several other authors, so the same algebra appears in the literature
under both names. Check each source's convention before comparing statements.
\end{remark}

\begin{definition}[Samuel compactification]\label{def:samuel}
\uses{def:ruc, def:ambit}
The \emph{Samuel compactification} $S(G)$ of $G$ is the Gelfand spectrum of
$\mathrm{RUC}_b(G)$: the space of characters of $\mathrm{RUC}_b(G)$ with the
weak* topology. For $f\in\mathrm{RUC}_b(G)$ and $h\in G$ the function
$f^h(g)=f(hg)$ is again right uniformly continuous, and
$(h\varphi)(f)=\varphi(f^h)$ defines a continuous action of $G$ on $S(G)$.
Evaluation at points embeds $G$ equivariantly as a dense subset of $S(G)$, so
$(S(G),1)$ is a $G$-ambit.
\end{definition}

\begin{proposition}[Universal property]\label{prop:greatest-ambit}
\uses{def:samuel}
For every $G$-ambit $(X,x_0)$ there is a unique morphism of ambits
$(S(G),1)\to(X,x_0)$. For this reason $S(G)$ is called the \emph{greatest
ambit} of $G$.
\end{proposition}

\begin{proof}[Sketch]
Continuity of the action and compactness of $X$ give, for every
$f\in C(X)$ and $\varepsilon>0$, a neighbourhood $V$ of the identity with
$|f(vx)-f(x)|<\varepsilon$ for all $x\in X$ and $v\in V$. Hence
$g\mapsto f(gx_0)$ lies in $\mathrm{RUC}_b(G)$. This defines a unital
*-homomorphism $C(X)\to\mathrm{RUC}_b(G)=C(S(G))$, whose dual map
$S(G)\to X$ is the required morphism. Uniqueness holds because $G$ is dense in
$S(G)$. See \cite{Auslander1988,Pestov2006}.
\end{proof}

\section{The universal minimal flow}\label{sec:umf}

Every topological group has a canonical minimal flow that maps onto all
others. Computing it, or deciding whether it is trivial or metrizable, is one
of the central problems of the subject.

\begin{definition}[Universal minimal flow]\label{def:umf}
\uses{def:minimal}
A minimal $G$-flow $M$ is \emph{universal} if every minimal $G$-flow is a
factor of $M$.
\end{definition}

\begin{theorem}[Ellis \cite{Ellis1960}]\label{thm:ellis}
\uses{def:umf, prop:greatest-ambit, prop:minimal-subflow}
Every topological group $G$ has a universal minimal flow, and any two
universal minimal flows of $G$ are isomorphic. It is denoted $M(G)$.
\end{theorem}

\begin{proof}[Proof of existence]
By Proposition~\ref{prop:minimal-subflow}, $S(G)$ contains a minimal subflow
$M$. Let $Y$ be a minimal $G$-flow and $y\in Y$. Then $(Y,y)$ is an ambit, so
Proposition~\ref{prop:greatest-ambit} gives a $G$-map $\pi\colon S(G)\to Y$.
The image $\pi(M)$ is a subflow of $Y$, so $\pi(M)=Y$ by minimality.
Uniqueness is the subtle part; see \cite{Ellis1960}, and \cite{GutmanLi2013}
for a short proof.
\end{proof}

\begin{corollary}\label{cor:umf-in-ambit}
\uses{thm:ellis}
Every minimal subflow of the greatest ambit $S(G)$ is isomorphic to $M(G)$.
\end{corollary}

\begin{example}\label{ex:umf-compact}
\uses{thm:ellis}
If $G$ is compact, then $M(G)$ is $G$ acting on itself by left translation.
Indeed, this flow is minimal, and for a minimal $G$-flow $Y$ and $y\in Y$ the
map $g\mapsto gy$ is a $G$-map onto the orbit $Gy$, which is compact and dense,
hence equal to $Y$.
\end{example}

\begin{remark}[Ellis's problem]\label{rem:ellis-problem}
\uses{thm:ellis, prop:greatest-ambit}
The greatest ambit maps canonically onto the enveloping (Ellis) semigroup
$E(M(G))$ of the universal minimal flow. Ellis asked, in his 1969
\emph{Lectures on Topological Dynamics}, whether this map is always an
isomorphism. Pestov observed that every nontrivial extremely amenable group
(Section~\ref{sec:ea}) answers this in the negative, since its universal
minimal flow, and hence its enveloping semigroup, is trivial
\cite{Pestov1998}.
\end{remark}

\section{Extreme amenability}\label{sec:ea}

The simplest possible universal minimal flow is a single point. Groups with
this property are far from locally compact, and they arise from concentration
of measure and from structural Ramsey theory.

\begin{definition}[Extremely amenable group]\label{def:ea}
\uses{def:flow}
A topological group $G$ is \emph{extremely amenable} if every $G$-flow has a
fixed point, that is, a point $x$ with $gx=x$ for all $g\in G$. This is also
called the \emph{fixed point on compacta property}.
\end{definition}

\begin{proposition}\label{prop:ea-trivial-umf}
\uses{def:ea, thm:ellis, prop:minimal-subflow}
A topological group $G$ is extremely amenable if and only if $M(G)$ is a
single point.
\end{proposition}

\begin{proof}
If $G$ is extremely amenable, a fixed point of $M(G)$ is a subflow, so
$M(G)$ is a point by minimality. Conversely, if $M(G)$ is a point, let $X$ be
a $G$-flow. By Proposition~\ref{prop:minimal-subflow} it contains a minimal
subflow $Y$, which is a factor of $M(G)$ and hence a single fixed point.
\end{proof}

\begin{remark}\label{rem:ea-amenable}
\uses{def:ea}
A topological group is \emph{amenable} if every $G$-flow carries a
$G$-invariant regular Borel probability measure. Extremely amenable groups are
amenable: the Dirac measure at a fixed point is invariant.
\end{remark}

\begin{remark}[The first examples]\label{rem:herer-christensen}
\uses{def:ea, rem:ea-amenable}
Historically, the first extremely amenable groups were the groups
$L^0(\varphi,\mathbb{R})$ of Herer and Christensen \cite{HererChristensen1975},
built from a pathological submeasure $\varphi$. They have no nontrivial
continuous unitary representations; being abelian, they are amenable, and
together the two properties make them extremely amenable. Natural examples
came later: see Sections~\ref{sec:gromov-milman}, \ref{sec:pestov}
and~\ref{sec:kpt}, and the monograph \cite{Pestov2006}.
\end{remark}

\section{Veech's theorem}\label{sec:veech}

For locally compact groups the universal minimal flow is as large as it can
be: the group acts on it freely.

\begin{theorem}[Veech \cite{Veech1977}]\label{thm:veech}
\uses{prop:greatest-ambit, cor:umf-in-ambit}
Every locally compact group $G$ acts freely on its greatest ambit $S(G)$:
$gx\neq x$ for all $x\in S(G)$ and all $g\neq 1$. In particular, $G$ acts
freely on $M(G)$.
\end{theorem}

\begin{remark}\label{rem:veech-history}
\uses{thm:veech}
Veech's theorem shows in particular that every locally compact group admits a
free action on a compact space; for discrete groups this was proved earlier by
Ellis. Pestov's paper \cite{Pestov1998} starts from the question of which
other topological groups admit such actions.
\end{remark}

\begin{corollary}\label{cor:lc-not-ea}
\uses{thm:veech, prop:ea-trivial-umf}
A nontrivial locally compact group is not extremely amenable.
\end{corollary}

\begin{proof}
If $M(G)$ is a single point and $G$ acts freely on it, then $G$ is trivial.
\end{proof}

\begin{remark}\label{rem:lc-metrizable}
\uses{thm:veech}
Kechris, Pestov and Todorcevic observed that a locally compact Polish group
with metrizable universal minimal flow is compact \cite{KPT2005}; see also
\cite{BYMT2017}. Compare Section~\ref{sec:metrizable}, where metrizable
universal minimal flows of Polish groups are described.
\end{remark}

\section{The Gromov--Milman theorem}\label{sec:gromov-milman}

The first natural example of an extremely amenable group came from asymptotic
geometric analysis: concentration of measure on high-dimensional spheres and
groups forces fixed points.

\begin{theorem}[Gromov--Milman \cite{GromovMilman1983}]\label{thm:gromov-milman}
\uses{def:ea}
The unitary group $U(\ell^2)$ of an infinite-dimensional separable Hilbert
space, with the strong operator topology, is extremely amenable.
\end{theorem}

\begin{remark}[Lévy groups]\label{rem:levy}
\uses{thm:gromov-milman}
The proof shows that $U(\ell^2)$ is a \emph{Lévy group}: roughly, it contains
an increasing chain of compact subgroups with dense union whose Haar measures
concentrate, in the sense of Lévy's concentration of measure phenomenon.
In \cite{GromovMilman1983} the fixed point theorem is stated for increasing
unions of compact groups whose normalised Haar measures form a Lévy family,
acting equicontinuously on a compact space (Theorem~7.1 there), and it is
applied to the union of the groups $SO(n)$, to the group of finitary
permutations of $\mathbb{N}$ and to the union of the groups
$(\mathbb{Z}/2\mathbb{Z})^n$. Glasner observed that equicontinuity is not
needed \cite{Glasner1998}, so every Lévy group is extremely amenable.
Theorem~\ref{thm:gromov-milman} is the form in which the result is usually
quoted \cite{Uspenskij2000,GiordanoPestov2007}. This mechanism, which Pestov
calls the Ramsey--Dvoretzky--Milman phenomenon, gives many further examples
\cite{Pestov2006}.
\end{remark}

\section{Pestov's theorems}\label{sec:pestov}

Pestov found the first extremely amenable groups of homeomorphisms and
automorphisms, and with them the first nontrivial universal minimal flow
computed explicitly for a non-locally compact group.

\begin{theorem}[Pestov \cite{Pestov1998}]\label{thm:pestov-ea}
\uses{def:ea}
The group $\mathrm{Aut}(\mathbb{Q},<)$ of order-preserving bijections of the
rationals, with the topology of pointwise convergence, and the group
$\mathrm{Homeo}_+[0,1]$ of orientation-preserving homeomorphisms of the unit
interval, with the compact-open topology, are extremely amenable.
\end{theorem}

\begin{remark}\label{rem:pestov-ramsey}
\uses{thm:pestov-ea}
Extreme amenability of $\mathrm{Aut}(\mathbb{Q},<)$ is a consequence of the
finite Ramsey theorem. This is the prototype of the Kechris--Pestov--Todorcevic
correspondence of Section~\ref{sec:kpt}.
\end{remark}

\begin{theorem}[Pestov \cite{Pestov1998}]\label{thm:pestov-circle}
\uses{thm:pestov-ea, def:umf}
The universal minimal flow of $\mathrm{Homeo}_+(\mathbb{S}^1)$ is the circle
$\mathbb{S}^1$ with the natural action.
\end{theorem}

\begin{proof}[Idea]
The stabiliser of a point $\theta\in\mathbb{S}^1$ is isomorphic to
$\mathrm{Homeo}_+[0,1]$, hence extremely amenable. If $X$ is a
$\mathrm{Homeo}_+(\mathbb{S}^1)$-flow, the stabiliser therefore fixes a point
$x'\in X$, and one checks that $g\theta\mapsto gx'$ is a well-defined
continuous $G$-map $\mathbb{S}^1\to X$. Since $\mathbb{S}^1$ is minimal, it
maps onto every minimal flow.
\end{proof}

\begin{remark}[Uspenskij]\label{rem:uspenskij}
\uses{thm:pestov-circle}
Uspenskij showed that the circle cannot be replaced here by the Hilbert cube
or by a compact manifold of dimension greater than one, answering a question of
Pestov, and that the action of a topological group on its universal minimal
flow is never 3-transitive \cite{Uspenskij2000}.
\end{remark}

\section{The universal minimal flow of $S_\infty$}\label{sec:glasner-weiss}

The group of all permutations of the integers has a metrizable, explicitly
described universal minimal flow: the space of linear orders.

\begin{definition}[The flow of linear orders]\label{def:lo}
\uses{def:flow}
Let $S_\infty$ be the group of all permutations of $\mathbb{N}$, with the
topology of pointwise convergence. Let $\mathrm{LO}(\mathbb{N})$ be the set of
linear orders on $\mathbb{N}$, viewed as a closed, hence compact, subset of
$\{0,1\}^{\mathbb{N}\times\mathbb{N}}$. The group $S_\infty$ acts on
$\mathrm{LO}(\mathbb{N})$ by transporting orders: $a\mathrel{(g\cdot{<})}b$ if
and only if $g^{-1}(a)<g^{-1}(b)$. This makes $\mathrm{LO}(\mathbb{N})$ an
$S_\infty$-flow.
\end{definition}

\begin{theorem}[Glasner--Weiss \cite{GlasnerWeiss2002}]\label{thm:glasner-weiss}
\uses{def:lo, def:umf}
The universal minimal flow of $S_\infty$ is $\mathrm{LO}(\mathbb{N})$.
\end{theorem}

\begin{remark}\label{rem:glasner-weiss}
\uses{thm:glasner-weiss}
Glasner and Weiss also showed that this flow is uniquely ergodic
\cite{GlasnerWeiss2002}. A short modern proof derives the theorem from
extreme amenability of $\mathrm{Aut}(\mathbb{Q},<)$; see
Example~\ref{ex:lo-coprecompact}.
\end{remark}

\section{The Kechris--Pestov--Todorcevic correspondence}\label{sec:kpt}

For automorphism groups of countable homogeneous structures, extreme
amenability is equivalent to a purely combinatorial property of the finite
substructures: the Ramsey property. For simplicity we restrict to relational
signatures.

\begin{definition}[Fraïssé classes]\label{def:fraisse}
Let $L$ be a relational signature. A class $\mathcal{K}$ of finite
$L$-structures, closed under isomorphism, is a \emph{Fraïssé class} if it
contains only countably many structures up to isomorphism, contains structures
of arbitrarily large finite cardinality, and has the hereditary property, the
joint embedding property and the amalgamation property. By Fraïssé's theorem
there is a unique, up to isomorphism, countable ultrahomogeneous structure
$\mathbf{F}=\mathrm{Flim}(\mathcal{K})$ whose finite substructures are, up to
isomorphism, exactly the members of $\mathcal{K}$. Its automorphism group
$\mathrm{Aut}(\mathbf{F})$, with the topology of pointwise convergence, is a
closed subgroup of $S_\infty$. The class $\mathcal{K}$ is a \emph{Fraïssé order
class} if $L$ contains a binary relation symbol $<$ that is interpreted as a
linear order in every member of $\mathcal{K}$.
\end{definition}

\begin{definition}[Ramsey property]\label{def:ramsey}
\uses{def:fraisse}
For $\mathbf{A},\mathbf{B}\in\mathcal{K}$, write
$\binom{\mathbf{B}}{\mathbf{A}}$ for the set of substructures of $\mathbf{B}$
isomorphic to $\mathbf{A}$. The class $\mathcal{K}$ has the \emph{Ramsey
property} if for all $\mathbf{A},\mathbf{B}\in\mathcal{K}$ such that
$\mathbf{A}$ embeds in $\mathbf{B}$, and every $k\geq 1$, there is
$\mathbf{C}\in\mathcal{K}$ such that for every colouring
$c\colon\binom{\mathbf{C}}{\mathbf{A}}\to\{1,\dots,k\}$ there is
$\mathbf{B}'\in\binom{\mathbf{C}}{\mathbf{B}}$ on which $c$ is constant, in
the sense that $c$ is constant on $\binom{\mathbf{B}'}{\mathbf{A}}$.
\end{definition}

\begin{theorem}[Kechris--Pestov--Todorcevic \cite{KPT2005}]\label{thm:kpt}
\uses{def:fraisse, def:ramsey, def:ea}
Let $\mathcal{K}$ be a Fraïssé order class with Fraïssé limit $\mathbf{F}$.
Then $\mathrm{Aut}(\mathbf{F})$ is extremely amenable if and only if
$\mathcal{K}$ has the Ramsey property.
\end{theorem}

\begin{example}\label{ex:kpt-rationals}
\uses{thm:kpt, thm:pestov-ea}
The finite linear orders form a Fraïssé order class with Fraïssé limit
$(\mathbb{Q},<)$, and their Ramsey property is the finite Ramsey theorem.
Theorem~\ref{thm:kpt} thus recovers the extreme amenability of
$\mathrm{Aut}(\mathbb{Q},<)$.
\end{example}

\begin{remark}\label{rem:kpt-general}
\uses{thm:kpt}
More generally, for any Fraïssé structure $\mathbf{F}$, the group
$\mathrm{Aut}(\mathbf{F})$ is extremely amenable if and only if the class of
finite substructures of $\mathbf{F}$ has the Ramsey property and consists of
rigid structures; see \cite{NguyenVanThe2014}. Since every closed subgroup of
$S_\infty$ is the automorphism group of a Fraïssé structure, this
characterises the extremely amenable closed subgroups of $S_\infty$.
Kechris, Pestov and Todorcevic also developed a method for computing universal
minimal flows from order expansions with the Ramsey property and the ordering
property \cite{KPT2005}.
\end{remark}

\section{Metrizable universal minimal flows}\label{sec:metrizable}

For Polish groups, a metrizable universal minimal flow is always of a very
concrete form: the completion of a homogeneous space by an extremely amenable
subgroup.

\begin{definition}[Co-precompact subgroups]\label{def:coprecompact}
\uses{def:flow, def:ruc}
Let $H$ be a closed subgroup of $G$. Equip $G/H$ with the uniformity whose
basic entourages are the sets $\{(gH,vgH) : g\in G,\ v\in V\}$, for $V$ a
neighbourhood of the identity, and let $\widehat{G/H}$ be its completion. The
subgroup $H$ is \emph{co-precompact} if $\widehat{G/H}$ is compact. In that
case the action of $G$ on $G/H$ extends to $\widehat{G/H}$, which becomes a
$G$-flow.
\end{definition}

\begin{proposition}\label{prop:coprecompact-ea}
\uses{def:coprecompact, def:ea, thm:ellis}
If $H\leq G$ is closed, co-precompact and extremely amenable, then every
minimal subflow of $\widehat{G/H}$ is isomorphic to $M(G)$.
\end{proposition}

\begin{proof}
Let $X$ be a minimal $G$-flow. Since $H$ is extremely amenable, it fixes a
point $x_0\in X$. The map $gH\mapsto gx_0$ is well defined and uniformly
continuous, so it extends to a $G$-map $\widehat{G/H}\to X$, which maps every
minimal subflow onto $X$. Thus every minimal subflow of $\widehat{G/H}$ is a
universal minimal flow, and Theorem~\ref{thm:ellis} identifies it with
$M(G)$.
\end{proof}

\begin{figure}[htbp]
\centering
\includegraphics[width=\linewidth]{completion-pattern}
\caption{The argument of Proposition~\ref{prop:coprecompact-ea}. The extremely amenable subgroup $H$ fixes a point $x_0$ of a minimal flow $X$. The map $gH\mapsto gx_0$ is uniformly continuous, so it extends to a $G$-map from the completion $\widehat{G/H}$ to $X$. The image of every minimal subflow of $\widehat{G/H}$ is closed and invariant, hence all of $X$. Several other proofs use the same argument; their links name this figure.}\label{fig:completion-pattern}
\illustrates{prop:coprecompact-ea, thm:pestov-circle, thm:kpt-umf, lem:pw-umf-principles, thm:pw-metrizable-fibre-wreath-umf, thm:pw-full-group-atomic-umf, lem:cu-completion-criterion}
\end{figure}

\begin{example}\label{ex:lo-coprecompact}
\uses{prop:coprecompact-ea, thm:pestov-ea, thm:glasner-weiss}
Identify $\mathbb{N}$ with $\mathbb{Q}$. Then $H=\mathrm{Aut}(\mathbb{Q},<)$ is
a closed, co-precompact, extremely amenable subgroup of $S_\infty$, and
$\widehat{S_\infty/H}$ is $\mathrm{LO}(\mathbb{N})$, which is minimal. This
recovers Theorem~\ref{thm:glasner-weiss}.
\end{example}

\begin{theorem}[Melleray--Nguyen Van Thé--Tsankov \cite{MNT2016}]\label{thm:mnt}
\uses{def:coprecompact, def:ea, thm:ellis}
Let $G$ be a Polish group. If $M(G)$ is metrizable and contains a $G_\delta$
orbit $Gx_0$, then the stabiliser $G_{x_0}$ is closed, co-precompact and
extremely amenable, and $M(G)$ is isomorphic to $\widehat{G/G_{x_0}}$.
\end{theorem}

\begin{theorem}[Ben Yaacov--Melleray--Tsankov \cite{BYMT2017}]\label{thm:bymt}
\uses{thm:mnt}
Let $G$ be a Polish group with metrizable universal minimal flow. Then $M(G)$
has a comeagre orbit. Consequently, there is a closed, co-precompact,
extremely amenable subgroup $G^*\leq G$ such that
$M(G)\cong\widehat{G/G^*}$.
\end{theorem}

\begin{remark}[Zucker]\label{rem:zucker}
For closed subgroups of $S_\infty$, Zucker had earlier characterised
combinatorially when $M(G)$ is metrizable, and showed that in that case every
minimal $G$-flow has a point with comeagre orbit, solving the Generic Point
Problem of Angel, Kechris and Lyons for these groups \cite{Zucker2016}. He
later gave a direct proof of Theorem~\ref{thm:bymt} \cite{Zucker2018}.
\end{remark}

\section{Computing universal minimal flows with order expansions}\label{sec:kpt-umf}

Kechris, Pestov and Todorcevic turned the correspondence of
Section~\ref{sec:kpt} into a method for computing universal minimal flows:
expand a Fraïssé class by linear orders that have the Ramsey property and
the ordering property, and the universal minimal flow is the space of all
admissible orders. Throughout, $L$ is a relational signature that does not
contain the symbol $<$.

\begin{definition}[Order expansions]\label{def:order-expansion}
\uses{def:fraisse}
Let $\mathcal{K}$ be a Fraïssé class in $L$ with Fraïssé limit $\mathbf{F}$.
An \emph{order expansion} of $\mathcal{K}$ is a class $\mathcal{K}^*$ of
structures in $L\cup\{<\}$, closed under isomorphism, in which $<$ is
interpreted as a linear order, and such that the $L$-reducts of the members
of $\mathcal{K}^*$ are exactly the members of $\mathcal{K}$. For
$\mathbf{A}\in\mathcal{K}$, an \emph{admissible order} on $\mathbf{A}$ is a
linear order $\prec$ with $(\mathbf{A},\prec)\in\mathcal{K}^*$.
\begin{enumerate}
\item $\mathcal{K}^*$ is \emph{reasonable} if for all
  $\mathbf{A},\mathbf{B}\in\mathcal{K}$ with $\mathbf{A}$ a substructure of
  $\mathbf{B}$, every admissible order on $\mathbf{A}$ extends to an
  admissible order on $\mathbf{B}$.
\item $\mathcal{K}^*$ has the \emph{ordering property} if for every
  $\mathbf{A}\in\mathcal{K}$ there is $\mathbf{B}\in\mathcal{K}$ such that
  $(\mathbf{A},\prec)$ embeds into $(\mathbf{B},\prec')$ for every admissible
  order $\prec$ on $\mathbf{A}$ and every admissible order $\prec'$ on
  $\mathbf{B}$.
\end{enumerate}
\end{definition}

\begin{definition}[The flow of admissible orders]\label{def:admissible-orders}
\uses{def:order-expansion, def:lo}
Let $\mathcal{K}^*$ be an order expansion of $\mathcal{K}$. Let
$X_{\mathcal{K}^*}$ be the set of linear orders $\prec$ on the universe of
$\mathbf{F}$ such that $(\mathbf{A},{\prec}|_{\mathbf{A}})\in\mathcal{K}^*$
for every finite substructure $\mathbf{A}$ of $\mathbf{F}$. It is a closed
subset of the space of all linear orders on $\mathbf{F}$, defined as in
Definition~\ref{def:lo}, and it is invariant under the action of
$\mathrm{Aut}(\mathbf{F})$ by transport of orders. Hence, when it is
nonempty, it is an $\mathrm{Aut}(\mathbf{F})$-flow.
\end{definition}

\begin{theorem}[Kechris--Pestov--Todorcevic \cite{KPT2005}]\label{thm:kpt-umf}
\uses{def:admissible-orders, def:ramsey, thm:kpt, prop:coprecompact-ea}
Let $\mathcal{K}$ be a Fraïssé class with Fraïssé limit $\mathbf{F}$, and let
$\mathcal{K}^*$ be a reasonable order expansion of $\mathcal{K}$ that is itself
a Fraïssé order class.
\begin{enumerate}
\item The Fraïssé limit of $\mathcal{K}^*$ is of the form
  $(\mathbf{F},\prec_0)$ with $\prec_0\in X_{\mathcal{K}^*}$.
\item The $\mathrm{Aut}(\mathbf{F})$-flow $X_{\mathcal{K}^*}$ is minimal if
  and only if $\mathcal{K}^*$ has the ordering property.
\item If $\mathcal{K}^*$ has the Ramsey property and the ordering property,
  then $M(\mathrm{Aut}(\mathbf{F}))$ is $X_{\mathcal{K}^*}$. In particular it
  is metrizable.
\end{enumerate}
\end{theorem}

\begin{proof}[Proof of (3) from (1) and (2)]
Let $G=\mathrm{Aut}(\mathbf{F})$ and let $H\leq G$ be the stabiliser of
$\prec_0$, that is, the automorphism group of the Fraïssé limit of
$\mathcal{K}^*$. It is a closed subgroup, and it is extremely amenable by
Theorem~\ref{thm:kpt}. The map $gH\mapsto g{\prec_0}$ is well defined and
injective. For a finite set $A$, let $V_A$ be its pointwise stabiliser in $G$.
Two orders $g{\prec_0}$ and $h{\prec_0}$ agree on $A$ if and only if
$hH=vgH$ for some $v\in V_A$: if they agree, the map $g^{-1}a\mapsto h^{-1}a$
($a\in A$) is an isomorphism between finite substructures of
$(\mathbf{F},\prec_0)$, so ultrahomogeneity gives $k\in H$ extending it, and
$v=hkg^{-1}$ works. So the map is a uniform isomorphism of $G/H$,
with the uniformity of Definition~\ref{def:coprecompact}, onto the orbit of
$\prec_0$. The orbit is dense: given $\prec\in X_{\mathcal{K}^*}$ and a
finite $A$, embed $(A,{\prec}|_A)$ into $(\mathbf{F},\prec_0)$ and use
ultrahomogeneity of $\mathbf{F}$ to find $g\in G$ such that $g{\prec_0}$
agrees with $\prec$ on $A$. Since $X_{\mathcal{K}^*}$ is compact, it is the
completion of $G/H$, so $H$ is co-precompact. By (2) the flow
$X_{\mathcal{K}^*}$ is minimal, and Proposition~\ref{prop:coprecompact-ea}
identifies it with $M(G)$.
\end{proof}

\begin{example}[Finite sets]\label{ex:kpt-sets}
\uses{thm:kpt-umf, thm:glasner-weiss, ex:kpt-rationals}
Let $\mathcal{K}$ be the class of finite sets, in the empty signature, so that
$\mathbf{F}=\mathbb{N}$ and $\mathrm{Aut}(\mathbf{F})=S_\infty$, and let
$\mathcal{K}^*$ be the class of finite linear orders. Every linear order is
admissible, so $\mathcal{K}^*$ is reasonable and $X_{\mathcal{K}^*}$ is
$\mathrm{LO}(\mathbb{N})$. The ordering property holds with
$\mathbf{B}=\mathbf{A}$, since any two linear orders on a finite set are
isomorphic, and the Ramsey property is the finite Ramsey theorem
(Example~\ref{ex:kpt-rationals}). Theorem~\ref{thm:kpt-umf} therefore
recovers Theorem~\ref{thm:glasner-weiss}.
\end{example}

\begin{example}[Computations of Kechris, Pestov and Todorcevic]\label{ex:kpt-computations}
\uses{thm:kpt-umf}
The following universal minimal flows are computed in
\cite[Section~8]{KPT2005} with Theorem~\ref{thm:kpt-umf}.
\begin{enumerate}
\item For the random graph, the random $K_n$-free graphs, and the random
  $r$-uniform hypergraphs, $M(\mathrm{Aut}(\mathbf{F}))$ is the space of
  \emph{all} linear orders on $\mathbf{F}$. The same holds for the rational
  Urysohn space and its isometry group.
\item For the countable equivalence relation with infinitely many classes,
  all of them infinite, $M(\mathrm{Aut}(\mathbf{F}))$ is the space of
  \emph{convex} linear orders, those in which every class is an interval.
\item For the countable atomless Boolean algebra and for the vector space of
  countably infinite dimension over a finite field,
  $M(\mathrm{Aut}(\mathbf{F}))$ is the space of \emph{natural} orders, defined
  in \cite{KPT2005} from antilexicographic orders. For the Boolean algebra this
  gives, by Stone duality, the universal minimal flow of the homeomorphism
  group of the Cantor set; see Section~\ref{sec:zero-dim}.
\end{enumerate}
In each case the automorphism group of the Fraïssé limit of $\mathcal{K}^*$
is extremely amenable (Theorem~\ref{thm:kpt}). See also the survey
\cite{NguyenVanThe2014}.
\end{example}

\section{The flow of maximal chains}\label{sec:maximal-chains}

Uspenskij attached to every flow a second flow: the space of maximal chains of
closed sets. It shows that the action of a group on its universal minimal flow
is never highly transitive, and for several homeomorphism groups it is the
universal minimal flow itself.

\begin{definition}[Hyperspace]\label{def:hyperspace}
\uses{def:flow}
For a compact space $X$, let $\mathrm{Exp}(X)$ be the set of nonempty closed
subsets of $X$ with the \emph{Vietoris topology}. Its basic open sets are
the sets $\langle U_1,\dots,U_k\rangle$, for $U_1,\dots,U_k$ open in $X$,
consisting of all $F\in\mathrm{Exp}(X)$ such that
$F\subseteq U_1\cup\dots\cup U_k$ and $F$ meets every $U_j$. It is a compact
space, metrizable if $X$ is.
If $X$ is a $G$-flow, then $gF=\{gx : x\in F\}$ makes $\mathrm{Exp}(X)$ a
$G$-flow.
\end{definition}

\begin{definition}[Maximal chains]\label{def:maximal-chain}
\uses{def:hyperspace}
A \emph{chain} in a compact space $X$ is a family $c\subseteq\mathrm{Exp}(X)$
that is linearly ordered by inclusion. It is \emph{maximal} if it is not
properly contained in another chain. The set of maximal chains in $X$ is
denoted $\Phi(X)$. A maximal chain is \emph{connected} if all its members are
connected.
\end{definition}

\begin{proposition}\label{prop:maximal-chains}
\uses{def:maximal-chain}
Let $X$ be a compact space.
\begin{enumerate}
\item Every maximal chain $c$ is a closed subset of $\mathrm{Exp}(X)$. It
  contains $X$, and it has a least element, which is a singleton
  $\{r(c)\}$. The point $r(c)$ is the \emph{root} of $c$.
\item $\Phi(X)$ is a closed subset of $\mathrm{Exp}(\mathrm{Exp}(X))$, hence
  compact \cite[Proposition~2.1]{Uspenskij2000}.
\item If $X$ is a $G$-flow, then $\Phi(X)$ is a subflow of
  $\mathrm{Exp}(\mathrm{Exp}(X))$, and the root map $r\colon\Phi(X)\to X$ is a
  surjective $G$-map.
\end{enumerate}
\end{proposition}

\begin{proof}
(1) The relation $A\subseteq B$ is closed in
$\mathrm{Exp}(X)\times\mathrm{Exp}(X)$, so the closure of a chain is a chain,
and a maximal chain is closed. The set $X$ is comparable with every member of
$c$, so $X\in c$. The intersection $F_0$ of all members of $c$ is nonempty by
compactness and comparable with every member, so $F_0\in c$. If $x\in F_0$,
then $\{x\}$ is contained in every member of $c$, so $\{x\}\in c$ and
$F_0=\{x\}$. Statement (2) is due to Uspenskij; see also
\cite[Section~1]{GlasnerWeiss2003}.

(3) Homeomorphisms of $X$ preserve inclusion, so $\Phi(X)$ is invariant, and
$r(gc)=gr(c)$. If $c_i\to c$ and $U$ is a neighbourhood of $r(c)$, then
eventually some member $F_i$ of $c_i$ is contained in $U$, because
$\{r(c)\}\in c$; hence $r(c_i)\in F_i\subseteq U$, and $r$ is continuous. For
$x\in X$, Zorn's lemma extends the chain $\{\{x\},X\}$ to a maximal chain,
whose root is $x$.
\end{proof}

\begin{theorem}[Uspenskij \cite{Uspenskij2000}]\label{thm:uspenskij-chains}
\uses{prop:maximal-chains}
Let $X$ be a $G$-flow with at least three points on which $G$ acts
$3$-transitively, that is, transitively on triples of distinct points. Then
there is no $G$-map $X\to\Phi(X)$.
\end{theorem}

\begin{proof}
Suppose $f\colon X\to\Phi(X)$ is a $G$-map and fix $x\in X$. The stabiliser
$G_x$ fixes $f(x)$, hence the root of $f(x)$. By $3$-transitivity, $G_x$ acts
transitively on $X\setminus\{x\}$, which has at least two points, so the root
of $f(x)$ is $x$. The chain $\{\{x\},X\}$ is not maximal, since a set
$\{x,y\}$ with $y\neq x$ can be added to it; so $f(x)$ has a member $F$ with
$\{x\}\subsetneq F\subsetneq X$. Pick $y\in F\setminus\{x\}$ and
$z\in X\setminus F$, and $g\in G$ with $g(x)=x$, $g(y)=z$ and $g(z)=y$. Then
$gF\in gf(x)=f(x)$, but $z\in gF\setminus F$ and $y\in F\setminus gF$, so $F$
and $gF$ are not comparable, a contradiction.
\end{proof}

\begin{corollary}[Uspenskij \cite{Uspenskij2000}]\label{cor:not-3-transitive}
\uses{thm:uspenskij-chains, def:umf, prop:minimal-subflow}
If $M(G)$ has at least three points, then $G$ does not act $3$-transitively on
$M(G)$.
\end{corollary}

\begin{proof}
By Proposition~\ref{prop:maximal-chains}, $\Phi(M(G))$ is a $G$-flow. It
contains a minimal subflow, and $M(G)$ maps onto it, which contradicts
Theorem~\ref{thm:uspenskij-chains} if the action is $3$-transitive.
\end{proof}

\begin{remark}[Uspenskij's formulation]\label{rem:uspenskij-proof}
\uses{thm:uspenskij-chains, cor:not-3-transitive, thm:pestov-circle}
Uspenskij proves Corollary~\ref{cor:not-3-transitive}, which is Theorem~1.1
of \cite{Uspenskij2000}, by the argument above, phrased with the stabiliser
$H$ of a point $p\in M(G)$: every $G$-flow receives a $G$-map from $M(G)$, so
it has an $H$-fixed point, and in particular some maximal chain in $M(G)$ is
fixed by $H$ \cite[Proposition~2.2]{Uspenskij2000}.
Theorem~\ref{thm:uspenskij-chains} is the formulation of the argument in
\cite{GutmanTsankovZucker2021}. The corollary cannot be improved to
$2$-transitivity: $\mathrm{Homeo}_+(\mathbb{S}^1)$ acts $2$-transitively on
its universal minimal flow $\mathbb{S}^1$, and the stabiliser of a point fixes
exactly two maximal chains \cite[Example~2.3]{Uspenskij2000}.
\end{remark}

\begin{remark}\label{rem:not-natural-action}
\uses{cor:not-3-transitive, rem:uspenskij}
Consequently, if $X$ is a compact manifold of dimension greater than one, the
Hilbert cube or the Cantor set, the universal minimal flow of
$\mathrm{Homeo}(X)$ is not the natural action on $X$
\cite{Uspenskij2000,GlasnerWeiss2003}. This is the result recalled in
Remark~\ref{rem:uspenskij}.
\end{remark}

\begin{remark}[Connected chains]\label{rem:connected-chains}
\uses{prop:maximal-chains}
If $X$ is connected, the connected maximal chains form a nonempty, closed,
invariant subset of $\Phi(X)$, so $\Phi(X)$ is not minimal. For the interval,
the circle and the Lelek fan, the universal minimal flow of the homeomorphism
group can be realised inside this subspace \cite{GutmanTsankovZucker2021}; see
Section~\ref{sec:continua}. For zero-dimensional spaces, $\Phi(X)$ itself can
be the universal minimal flow; see Section~\ref{sec:zero-dim}.
\end{remark}

\section{Homeomorphism groups of zero-dimensional spaces}\label{sec:zero-dim}

For the Cantor set, and more generally for every compact zero-dimensional
space whose nonempty clopen subsets are all homeomorphic, the flow of maximal
chains is the universal minimal flow of the homeomorphism group.

\begin{definition}[Proximal flows]\label{def:proximal}
\uses{def:flow}
Two points $x,y$ of a $G$-flow $X$ are \emph{proximal} if there is a net
$(g_i)$ in $G$ such that $(g_ix)$ and $(g_iy)$ converge to the same point. The
flow is \emph{proximal} if every pair of its points is proximal. A factor of
a proximal flow is proximal.
\end{definition}

\begin{theorem}[Glasner--Weiss \cite{GlasnerWeiss2003}]\label{thm:gw-cantor}
\uses{prop:maximal-chains, def:umf, def:proximal, thm:ellis}
Let $E$ be the Cantor set and let $G=\mathrm{Homeo}(E)$, with the topology of
uniform convergence. Then $M(G)$ is isomorphic to the flow $\Phi(E)$ of
maximal chains in $E$. In particular $M(G)$ is metrizable, and in fact
homeomorphic to the Cantor set. Moreover, $\Phi(E)$ is proximal, so every
minimal $G$-flow is proximal.
\end{theorem}

\begin{example}[The standard chain]\label{ex:standard-chain}
\uses{thm:gw-cantor, thm:pestov-ea}
Realise $E$ as the ternary Cantor set in $[0,1]$. The chain
$c_0=\{E\cap[0,t] : t\in[0,1]\}$ is maximal, with root $0$. Its stabiliser is
the group $P$ of order-preserving homeomorphisms of $E$, and $P$ is isomorphic
to $\mathrm{Aut}(\mathbb{Q},<)$, through its action on the complementary open
intervals of $E$, which are ordered like the rationals. Hence $P$ is extremely
amenable by Theorem~\ref{thm:pestov-ea}. Moreover, $P$ has exactly six fixed
points in $\Phi(E)$, and the action of $G$ on $\Phi(E)$ is minimal
\cite[Theorems~1.2 and~1.3]{GlasnerWeiss2003}.
\end{example}

\begin{proof}[Proof of proximality, after Glasner and Weiss]
Let $c_1,c_2\in\Phi(E)$ and let $Z$ be the closure of the orbit of
$(c_1,c_2)$ in $\Phi(E)\times\Phi(E)$. Since $P$ is extremely amenable, $Z$
contains a pair $(a_1,a_2)$ of $P$-fixed points, and Glasner and Weiss check
directly that any two of the six $P$-fixed chains are proximal. So $Z$
contains a diagonal point $(a,a)$, which means that $c_1$ and $c_2$ are
proximal.
\end{proof}

\begin{remark}[Stone duality]\label{rem:stone-duality}
\uses{thm:gw-cantor, ex:kpt-computations, thm:ellis}
The clopen subsets of $E$ form the countable atomless Boolean algebra
$\mathbf{B}_\infty$, and $\mathrm{Homeo}(E)$ is isomorphic, as a topological
group, to $\mathrm{Aut}(\mathbf{B}_\infty)$ with the topology of pointwise
convergence. Kechris, Pestov and Todorcevic computed the universal minimal flow
of $\mathrm{Aut}(\mathbf{B}_\infty)$ as the space of natural orders on
$\mathbf{B}_\infty$ \cite[Section~8]{KPT2005}
(Example~\ref{ex:kpt-computations}). By Theorem~\ref{thm:ellis}, this flow is
isomorphic to $\Phi(E)$.
\end{remark}

\begin{definition}[h-homogeneous spaces]\label{def:h-homogeneous}
A compact zero-dimensional space $X$ is \emph{h-homogeneous} if all its
nonempty clopen subsets are homeomorphic to each other; equivalently, by Stone
duality, if its algebra of clopen sets is a homogeneous Boolean algebra. A
metrizable h-homogeneous space with more than one point is homeomorphic to the
Cantor set.
\end{definition}

\begin{theorem}[Glasner--Gutman \cite{GlasnerGutman2012}]\label{thm:glasner-gutman}
\uses{def:h-homogeneous, prop:maximal-chains, thm:ellis}
Let $X$ be an h-homogeneous compact zero-dimensional space and let
$G=\mathrm{Homeo}(X)$, with the compact-open topology. Then $M(G)$ is isomorphic
to $\Phi(X)$.
\end{theorem}

\begin{remark}\label{rem:glasner-gutman}
\uses{thm:glasner-gutman}
For metrizable $X$ this is Theorem~\ref{thm:gw-cantor}, although the proof of
Glasner and Gutman is new even in that case, since it uses no fixed point
theorem. As for the Cantor set, the main combinatorial tool is the dual Ramsey theorem,
a consequence of the
Graham--Rothschild theorem on $n$-parameter sets \cite{GlasnerWeiss2003}.
Non-metrizable h-homogeneous spaces give many new examples, among them the
generalised Cantor cubes $\{0,1\}^\kappa$ for uncountable $\kappa$ and the
corona $\beta\omega\setminus\omega$ \cite{GlasnerGutman2012}.
\end{remark}

\section{Homeomorphism groups of continua}\label{sec:continua}

For connected spaces the flow of all maximal chains is not minimal. For the
interval, the circle and the Lelek fan, the universal minimal flow of the
homeomorphism group can be realised as a flow of connected chains, while for
manifolds of dimension at least two it is not metrizable. A \emph{continuum} is a compact connected metrizable space;
homeomorphism groups of compact metrizable spaces carry the compact-open
topology, which is the topology of uniform convergence.

\begin{remark}[The interval and the circle]\label{rem:interval-circle}
\uses{thm:pestov-ea, thm:pestov-circle}
The group $\mathrm{Homeo}_+[0,1]$ is extremely amenable, and
$M(\mathrm{Homeo}_+(\mathbb{S}^1))$ is the circle (Section~\ref{sec:pestov}).
\end{remark}

\begin{definition}[The Lelek fan]\label{def:lelek-fan}
Let $C$ be the Cantor set. The \emph{Cantor fan} is the quotient of
$C\times[0,1]$ obtained by collapsing $C\times\{0\}$ to a point $v$, the
\emph{top} of the fan. A point $x$ of a continuum $X$ is an \emph{endpoint} if
$x\in\{h(0),h(1)\}$ for every embedding $h\colon[0,1]\to X$ whose image
contains $x$. The \emph{Lelek fan} $L$ is the unique subcontinuum of the
Cantor fan with more than one point whose endpoints are dense in it
\cite{BartosovaKwiatkowska2019}.
\end{definition}

\begin{theorem}[Bartošová--Kwiatkowska \cite{BartosovaKwiatkowska2019}]\label{thm:lelek}
\uses{def:lelek-fan, prop:maximal-chains, def:umf}
Let $Y^*\subseteq\Phi(L)$ be the set of connected maximal chains in the Lelek
fan $L$ all of whose members contain the top $v$. Then $Y^*$ is a subflow of
$\Phi(L)$, and the universal minimal flow of $\mathrm{Homeo}(L)$ is $Y^*$. In
particular it is metrizable.
\end{theorem}

\begin{remark}[Method]\label{rem:lelek-method}
\uses{thm:lelek, prop:coprecompact-ea, thm:kpt}
Bartošová and Kwiatkowska show that the stabiliser $H$ of a generic chain in
$Y^*$ is extremely amenable and co-precompact, and that the universal minimal
flow is the completion of $\mathrm{Homeo}(L)/H$; compare
Proposition~\ref{prop:coprecompact-ea}. They realise $L$ as a quotient of a
projective Fraïssé limit, compute the universal minimal flow of its
automorphism group with the methods of Kechris, Pestov and Todorcevic and a
new Ramsey theorem proved with the dual Ramsey theorem, and then pass from the
projective limit to $L$ \cite{BartosovaKwiatkowska2019}.
\end{remark}

\begin{definition}[Generalised Ważewski dendrites]\label{def:wazewski}
A \emph{dendrite} is a locally connected continuum containing no simple
closed curve. The \emph{order} of a point $x$ is the number of connected
components of $X\setminus\{x\}$, and $x$ is a \emph{ramification point} if its
order is at least $3$. For $P\subseteq\{3,4,\dots,\omega\}$, a
\emph{generalised Ważewski dendrite} $W_P$ is a dendrite in which every
ramification point has order in $P$ and, for every $p\in P$, every arc contains
a ramification point of order $p$. It is unique up to homeomorphism; see
\cite{Kwiatkowska2018}.
\end{definition}

\begin{theorem}[Kwiatkowska \cite{Kwiatkowska2018}]\label{thm:wazewski}
\uses{def:wazewski, def:umf}
Let $P\subseteq\{3,4,\dots,\omega\}$ be nonempty.
\begin{enumerate}
\item If $P$ is finite, the universal minimal flow of $\mathrm{Homeo}(W_P)$ is
  metrizable, and it is computed explicitly. This answers a question of
  Duchesne.
\item If $P$ is infinite, the universal minimal flow of $\mathrm{Homeo}(W_P)$
  is not metrizable. These groups are Roelcke precompact, and their universal
  minimal flows have a comeagre orbit.
\end{enumerate}
\end{theorem}

\begin{remark}\label{rem:wazewski}
\uses{thm:wazewski, thm:bymt}
By Theorem~\ref{thm:bymt}, a metrizable universal minimal flow of a Polish
group has a comeagre orbit. Part (2) shows that the converse fails
\cite{GutmanTsankovZucker2021}. Kwiatkowska works with an identification,
which follows from work of Duchesne and Monod, of $\mathrm{Homeo}(W_P)$ with
the automorphism group of a countable structure whose universe is the set of
ramification points of $W_P$ \cite{Kwiatkowska2018}.
\end{remark}

\begin{theorem}[Gutman--Tsankov--Zucker \cite{GutmanTsankovZucker2021}]\label{thm:gtz}
\uses{def:umf, prop:maximal-chains, thm:bymt}
Let $X$ be a closed manifold of dimension at least $2$, or the Hilbert cube.
Then the universal minimal flow of $\mathrm{Homeo}(X)$ is not metrizable. If
$X$ is a closed manifold of dimension at least $3$, the connected maximal
chains in $X$ form a minimal $\mathrm{Homeo}(X)$-flow all of whose orbits are
meagre.
\end{theorem}

\begin{remark}\label{rem:gtz}
\uses{thm:gtz, thm:bymt, rem:uspenskij}
This answers a question of Uspenskij. As explained in
\cite{GutmanTsankovZucker2021}, Theorem~\ref{thm:bymt} implies that if the
universal minimal flow of a Polish group $G$ is metrizable, then every minimal
$G$-flow is metrizable and has a comeagre orbit. So a single metrizable minimal
flow with meagre orbits shows that $M(G)$ is not metrizable.
\end{remark}

\begin{remark}[The pseudo-arc]\label{rem:pseudo-arc}
\uses{thm:gtz, rem:connected-chains}
Uspenskij asked whether the universal minimal flow of the homeomorphism group
of the pseudo-arc $\mathbf{P}$ is metrizable, and in particular whether it is
the natural action on $\mathbf{P}$ \cite{Uspenskij2000}. Since $\mathbf{P}$ is
hereditarily indecomposable, its space of connected maximal chains is
isomorphic to $\mathbf{P}$ itself \cite{GutmanTsankovZucker2021}. The question
was still open when \cite{BartosovaKwiatkowska2019} and
\cite{GutmanTsankovZucker2021} were written.
\end{remark}

\section{The semigroup structure of the greatest ambit}\label{sec:semigroup}

The universal property of the greatest ambit turns it into a compact
semigroup. Minimal flows inside it are minimal ideals, and comparing it with
the enveloping semigroup of the universal minimal flow leads to Ellis's
problem and to a conjecture of Pestov. Some sources, such as
\cite{BartosovaZucker2019}, work with right actions; their semigroups are then
left topological rather than right topological.

\begin{definition}[Right topological semigroups]\label{def:rt-semigroup}
A \emph{compact right topological semigroup} is a semigroup $S$ with a compact
topology in which, for every $q\in S$, the map $p\mapsto pq$ is continuous. An
element $u\in S$ is an \emph{idempotent} if $uu=u$. A \emph{left ideal} is a
nonempty set $I\subseteq S$ with $SI\subseteq I$. Terminology varies: such
semigroups are called left topological in \cite{Uspenskij2000}.
\end{definition}

\begin{proposition}\label{prop:ambit-semigroup}
\uses{def:rt-semigroup, prop:greatest-ambit}
For $x\in S(G)$, let $\rho_x\colon(S(G),1)\to(\overline{Gx},x)$ be the morphism
of ambits given by Proposition~\ref{prop:greatest-ambit}, and put
$yx=\rho_x(y)$. This makes $S(G)$ a compact right topological semigroup with
identity $1$, in which left multiplication by $g\in G$ is the action of $g$
\cite[Theorem~3.1]{Uspenskij2000}.
\end{proposition}

\begin{proof}
Right multiplication by $x$ is $\rho_x$, which is continuous. Since $\rho_x$ is
a $G$-map with $\rho_x(1)=x$, we get $gx=\rho_x(g)$ for $g\in G$; also
$\rho_1$ is the identity, so $y1=y$. For associativity, $\rho_z\circ\rho_x$ is
a $G$-map sending $1$ to $xz$, hence a morphism of ambits
$(S(G),1)\to(\overline{Gxz},xz)$, so it equals $\rho_{xz}$ by uniqueness;
that is, $(yx)z=y(xz)$.
\end{proof}

\begin{proposition}\label{prop:minimal-ideals}
\uses{prop:ambit-semigroup, cor:umf-in-ambit}
\begin{enumerate}
\item The subflows of $S(G)$ are exactly its closed left ideals, and its
  minimal subflows are exactly its minimal left ideals. Each of them is
  isomorphic to $M(G)$.
\item Every minimal left ideal $M$ contains an idempotent $u$, and
  $p\mapsto pu$ is a $G$-map $S(G)\to M$ that restricts to the identity on
  $M$. In particular $M(G)$ is a retract of $S(G)$.
\end{enumerate}
\end{proposition}

\begin{proof}
(1) A left ideal is invariant under left multiplication by $G$. Conversely, if
$Y$ is a subflow, $x\in Y$ and $y=\lim g_i$ with $g_i\in G$, then
$yx=\lim g_ix\in Y$, by continuity of $\rho_x$. Every left ideal contains a
left ideal $S(G)x=\rho_x(S(G))$, which is closed, so minimal left ideals are
closed and are the minimal subflows. They are isomorphic to $M(G)$ by
Corollary~\ref{cor:umf-in-ambit}.

(2) A minimal left ideal $M$ is a closed subsemigroup, and every compact right
topological semigroup contains an idempotent, by a theorem of Ellis; see
\cite[Theorem~3.2]{Uspenskij2000} for a short proof. Let $u\in M$ be an idempotent. Then
$S(G)u$ is a left ideal contained in $M$, so $M=S(G)u$, and $pu=p$ for every
$p=qu\in M$. The map $p\mapsto pu$ is continuous, and it is a $G$-map by
associativity.
\end{proof}

\begin{definition}[Enveloping semigroup]\label{def:enveloping}
\uses{def:flow, def:ambit}
The \emph{enveloping semigroup} $E(X)$ of a $G$-flow $X$ is the closure, in
the product space $X^X$, of the set of maps $x\mapsto gx$ for $g\in G$. It is a
compact right topological semigroup under composition, and $G$ acts on it by
$(gp)(x)=g\,p(x)$; with the identity map as base point, it is a $G$-ambit. By
Proposition~\ref{prop:greatest-ambit} there is a canonical morphism of ambits
$S(G)\to E(M(G))$.
\end{definition}

\begin{remark}\label{rem:ellis-question-semigroup}
\uses{def:enveloping, rem:ellis-problem, prop:ea-trivial-umf}
Ellis's problem (Remark~\ref{rem:ellis-problem}) asks whether the canonical
map $S(G)\to E(M(G))$ is always an isomorphism. For a nontrivial extremely
amenable group it is not, since $E(M(G))$ is a point while $S(G)$ is not
\cite{Pestov1998}.
\end{remark}

\begin{theorem}[Glasner \cite{Glasner1998}]\label{thm:glasner-monothetic}
\uses{rem:ellis-question-semigroup, def:ea}
\begin{enumerate}
\item There is an infinite monothetic Polish group (one with a dense cyclic
  subgroup) that is extremely amenable. This answers a question of Mitchell.
\item For $\mathbb{Z}$ with the discrete topology, the greatest ambit
  $\beta\mathbb{Z}$ is not isomorphic to $E(M(\mathbb{Z}))$.
\end{enumerate}
\end{theorem}

\begin{conjecture}[Pestov]\label{conj:pestov}
\uses{rem:ellis-question-semigroup}
For a topological group $G$, the canonical map $S(G)\to E(M(G))$ is an
isomorphism if and only if $G$ is precompact, that is, isomorphic to a
subgroup of a compact group \cite[page~4163]{Pestov1998}.
\end{conjecture}

\begin{remark}\label{rem:pestov-conjecture}
\uses{conj:pestov, ex:umf-compact}
If $G$ is precompact, $S(G)$, $M(G)$ and $E(M(G))$ are all isomorphic to the
completion of $G$, which is a compact group \cite{BartosovaZucker2019}; compare
Example~\ref{ex:umf-compact}. So the conjecture concerns non-precompact
groups. Pestov showed that the canonical map is an isomorphism if and only if
the retractions of $S(G)$ onto a fixed minimal subflow, as in
Proposition~\ref{prop:minimal-ideals}, separate the points of $S(G)$
\cite{BartosovaZucker2019}.
\end{remark}

\begin{theorem}[Glasner--Tsankov--Weiss--Zucker \cite{GTWZ2021}]\label{thm:gtwz}
\uses{conj:pestov}
Pestov's conjecture holds for discrete groups: if $G$ is an infinite discrete
group, then $\beta G$ is not isomorphic to $E(M(G))$.
\end{theorem}

\begin{theorem}[Bartošová--Zucker \cite{BartosovaZucker2019}]\label{thm:bartosova-zucker}
\uses{conj:pestov, ex:kpt-computations}
Let $G$ be $S_\infty$, the automorphism group of a random $K_n$-free graph, or
the automorphism group of a random $r$-uniform hypergraph. Then $S(G)$ is not
isomorphic to $E(M(G))$.
\end{theorem}

\begin{remark}\label{rem:bartosova-zucker}
\uses{thm:bartosova-zucker, thm:gtwz}
The proof uses the explicit universal minimal flows of
Section~\ref{sec:kpt-umf} and only works when $M(G)$ is metrizable. The
methods of \cite{GTWZ2021} are of a different kind: they rely on a
disjointness property that groups with metrizable universal minimal flow
never have \cite{BartosovaZucker2019}.
\end{remark}

\section{Discrete groups}\label{sec:discrete}

For locally compact groups that are not compact, and in particular for
infinite discrete groups, the universal minimal flow is a large,
non-metrizable space. For countable discrete groups its topology is
nevertheless determined up to one cardinal parameter.

\begin{remark}[Size]\label{rem:umf-size}
\uses{thm:veech}
If $G$ is locally compact and not compact, then $M(G)$ is not metrizable
\cite{GlasnerWeiss2003,GiordanoPestov2007}; compare
Remark~\ref{rem:lc-metrizable}. For infinite discrete groups this is due to
Ellis, and Kechris, Pestov and Todorcevic extended it to locally compact
groups that are not compact \cite{GlasnerGutman2012}. By Theorem~\ref{thm:veech}, $G$ acts freely
on it.
\end{remark}

\begin{example}\label{ex:discrete-ambit}
\uses{def:samuel}
If $G$ is discrete, every bounded function on $G$ is right uniformly
continuous, so $\mathrm{RUC}_b(G)=\ell^\infty(G)$, and $S(G)$ is the
Stone--Čech compactification $\beta G$.
\end{example}

\begin{proposition}\label{prop:extremely-disconnected}
\uses{ex:discrete-ambit, prop:minimal-ideals}
If $G$ is discrete, then $M(G)$ is extremely disconnected: the closure of
every open subset is open.
\end{proposition}

\begin{proof}
The space $\beta G$ is extremely disconnected, and by
Proposition~\ref{prop:minimal-ideals} there is a retraction
$r\colon\beta G\to M$ onto a minimal subflow $M\cong M(G)$. Let $U\subseteq M$
be open and $V=r^{-1}(U)$, so that $U\subseteq V$ and $r(V)=U$. Then
$\overline{U}\subseteq\overline{V}\cap M\subseteq r(\overline{V})
\subseteq\overline{r(V)}=\overline{U}$, where the closure of $U$ is taken in
$M$. Hence $\overline{U}=\overline{V}\cap M$, which is open in $M$ because
$\overline{V}$ is open in $\beta G$.
\end{proof}

\begin{theorem}[Balcar--Błaszczyk \cite{BalcarBlaszczyk1990}]\label{thm:balcar-blaszczyk}
Let $B$ be an infinite Boolean algebra and $H$ a countable group of
automorphisms of $B$ such that $(B,H)$ is \emph{minimal}: for every nonzero
$u\in B$ there are $h_1,\dots,h_n\in H$ with $h_1(u)\vee\dots\vee h_n(u)=1$.
Then the completion of $B$ is isomorphic to the completion of the free Boolean
algebra on $\pi(B)$ generators, where $\pi(B)$ is the least cardinality of an
order-dense subset of $B\setminus\{0\}$.
\end{theorem}

\begin{remark}\label{rem:balcar-blaszczyk}
\uses{thm:balcar-blaszczyk}
Minimality of $(B,H)$ means that the dual action of $H$ on the Stone space of
$B$ is minimal. The countability of $H$ cannot be dropped: a homogeneous
Boolean algebra with all its automorphisms is a minimal system
\cite{BalcarBlaszczyk1990}.
\end{remark}

\begin{corollary}\label{cor:discrete-cohen}
\uses{thm:balcar-blaszczyk, prop:extremely-disconnected, thm:veech}
Let $G$ be a countably infinite discrete group. Then $M(G)$ is homeomorphic to
the Stone space of the completion of a free Boolean algebra; equivalently, to
the Gleason cover (absolute) of a Cantor cube $\{0,1\}^\kappa$.
\end{corollary}

\begin{proof}
By Proposition~\ref{prop:extremely-disconnected}, the clopen subsets of $M(G)$
form a complete Boolean algebra $B$, whose Stone space is $M(G)$. It is
infinite, because $G$ is infinite and acts freely on $M(G)$
(Theorem~\ref{thm:veech}). The group $G$ acts on $B$ by automorphisms, and
minimality of $M(G)$ says that for every nonzero $b\in B$ finitely many
translates of $b$ cover $M(G)$, which is minimality of $(B,G)$. Since $B$ is
complete, Theorem~\ref{thm:balcar-blaszczyk} applies to $B$ itself. The
completion of the free Boolean algebra on $\kappa$ generators is the algebra
of regular open subsets of $\{0,1\}^\kappa$, whose Stone space is the Gleason
cover of $\{0,1\}^\kappa$.
\end{proof}

\begin{example}[The integers]\label{ex:umf-integers}
\uses{cor:discrete-cohen, prop:minimal-ideals}
For $G=\mathbb{Z}$, the cardinal is the continuum. Balcar and Błaszczyk show
that every minimal subflow of $\beta\mathbb{Z}$, with the action extending
the shift, is homeomorphic to the Gleason cover of $\{0,1\}^{\mathfrak{c}}$
\cite{BalcarBlaszczyk1990}. By Proposition~\ref{prop:minimal-ideals}, so is
$M(\mathbb{Z})$: it is the Stone space of the Boolean algebra of regular open
subsets of $\{0,1\}^{\mathfrak{c}}$; see also \cite{BartosovaKwiatkowska2019}.
\end{example}

\section{A catalogue of extremely amenable groups}\label{sec:ea-catalogue}

Extremely amenable groups arise from concentration of measure, from structural
Ramsey theory, and as stabilisers inside groups whose universal minimal flow
is metrizable. This section collects the main examples; groups and
topologies are as in the cited sources.

\begin{theorem}[Pestov \cite{Pestov2002}]\label{thm:pestov-urysohn}
\uses{def:ea}
The isometry group $\mathrm{Iso}(\mathbb{U})$ of the Urysohn universal
complete separable metric space $\mathbb{U}$, with the compact-open topology, is
extremely amenable. The same holds for sufficiently homogeneous generalised
Urysohn metric spaces, and together with a theorem of Uspenskij this shows that
every topological group embeds into an extremely amenable group.
\end{theorem}

\begin{theorem}[Giordano--Pestov \cite{GiordanoPestov2007}]\label{thm:giordano-pestov}
\uses{def:ea}
\begin{enumerate}
\item The group of measure-preserving automorphisms of a standard Lebesgue
  space with a non-atomic measure, with the weak topology, is extremely
  amenable.
\item The isometry groups of the spaces $L^p(0,1)$, $1\leq p<\infty$, are
  extremely amenable. This extends the case $p=2$ of
  Theorem~\ref{thm:gromov-milman}.
\item A measure class preserving equivalence relation $R$ on a standard Borel
  space is amenable if and only if its full group $[R]$, with the uniform
  topology, is extremely amenable.
\item A von Neumann algebra is approximately finite-dimensional if and only if
  its unitary group, with the strong topology, is the product of an extremely
  amenable group and a compact group.
\end{enumerate}
\end{theorem}

\begin{remark}[Groups of measurable maps]\label{rem:l0}
\uses{def:ea, thm:glasner-monothetic}
The group $L^0(I,U(1))$ of measurable maps from the unit interval to the circle
group, with the topology of convergence in measure, is extremely amenable. This
was proved by Glasner and, independently and unpublished, by Furstenberg and
Weiss; see \cite{GiordanoPestov2007}. Glasner also constructed an
infinite monothetic extremely amenable Polish group
(Theorem~\ref{thm:glasner-monothetic}).
\end{remark}

\begin{remark}[Summary of the examples in this encyclopaedia]\label{rem:ea-summary}
\uses{rem:herer-christensen, thm:gromov-milman, rem:levy, thm:pestov-ea, thm:kpt, ex:kpt-computations, ex:standard-chain, rem:lelek-method, cor:lc-not-ea}
\begin{enumerate}
\item The groups of Herer and Christensen (Remark~\ref{rem:herer-christensen}),
  historically the first examples.
\item The unitary group of $\ell^2$ with the strong operator topology
  (Theorem~\ref{thm:gromov-milman}), and more generally Lévy groups
  (Remark~\ref{rem:levy}).
\item $\mathrm{Aut}(\mathbb{Q},<)$ and $\mathrm{Homeo}_+[0,1]$
  (Theorem~\ref{thm:pestov-ea}).
\item Automorphism groups of Fraïssé limits of Fraïssé order classes with the
  Ramsey property (Theorem~\ref{thm:kpt}), such as the order expansions in
  Example~\ref{ex:kpt-computations}.
\item The stabiliser of the standard chain in $\mathrm{Homeo}(E)$, for $E$ the
  Cantor set (Example~\ref{ex:standard-chain}), and the stabiliser of a
  generic connected chain in the homeomorphism group of the Lelek fan
  (Remark~\ref{rem:lelek-method}).
\end{enumerate}
By Corollary~\ref{cor:lc-not-ea}, none of these groups is locally compact;
according to \cite{GiordanoPestov2007}, this was first shown by Granirer and
Lau.
\end{remark}

\section{Groups of maps and of measure-space transformations}\label{sec:maps-transformations}

Two elementary principles compute many universal minimal flows: a group
that is dense in a compact group has that compact group as its universal
minimal flow, and a closed normal extremely amenable subgroup can be divided
out. The computations in this section, apart from
Proposition~\ref{prop:ea-normal}, are taken from a preprint of the editor
posted on Zenodo \cite{LupiniPointwise2026}; it has not been refereed.

\begin{proposition}\label{prop:ea-normal}
\uses{def:ea, thm:ellis, def:samuel, cor:umf-in-ambit}
Let $N$ be a closed normal subgroup of a topological group $G$.
\begin{enumerate}
\item If $N$ is extremely amenable, then $N$ acts trivially on every minimal
  $G$-flow, and $M(G)$ is $M(G/N)$ with $G$ acting through the quotient map.
\item If $G$ is a dense subgroup of a compact group $Q$, with the induced
  topology, then $M(G)$ is $Q$ with $G$ acting by left translation.
\end{enumerate}
\end{proposition}

\begin{proof}
(1) In a minimal $G$-flow $X$, the set of $N$-fixed points is nonempty by
extreme amenability, closed, and $G$-invariant because $N$ is normal, so it
is all of $X$. The action of $G$ therefore factors through $G/N$, and the
factored action is continuous because the quotient map $G\to G/N$ is open.
Thus minimal $G$-flows and minimal $G/N$-flows are the same. (2) The right
uniform structure of $G$ is the one induced by the unique uniform structure
of $Q$, so every bounded right uniformly continuous function on $G$ extends
uniquely to a continuous function on $Q$. Hence $S(G)=Q$, which is a minimal
$G$-flow because $G$ is dense, and Corollary~\ref{cor:umf-in-ambit} gives
$M(G)=Q$. See also \cite[Lemma~2.1]{LupiniPointwise2026}.
\end{proof}

\begin{theorem}[Lupini \cite{LupiniPointwise2026}]\label{thm:lupini-cp}
\uses{prop:ea-normal}
Let $X$ be a Tychonoff space and $K$ a compact path-connected topological
group, and let $C_p(X,K)$ be the group of continuous maps $X\to K$, with
pointwise multiplication and the topology of pointwise convergence. Then
$M(C_p(X,K))\cong K^X$, with $C_p(X,K)$ acting by left translation.
\end{theorem}

\begin{remark}\label{rem:lupini-cp}
\uses{thm:lupini-cp}
Continuous maps can take any finitely many prescribed values, so
$C_p(X,K)$ is dense in $K^X$ and Proposition~\ref{prop:ea-normal} applies.
The same argument computes, for the pointwise topology, the universal
minimal flows of path, loop and gauge groups of compact Lie groups, and of
groups of Lipschitz, Hölder, real-analytic and holomorphic maps. In the finer
natural topologies of path and loop groups the same compact group is still a
factor of the universal minimal flow, which is therefore not metrizable when
the target group has nontrivial identity component
\cite[Section~2]{LupiniPointwise2026}.
\end{remark}

\begin{theorem}[Lupini \cite{LupiniPointwise2026}]\label{thm:lupini-aut-measure}
\uses{prop:ea-normal, thm:giordano-pestov, thm:glasner-weiss}
Let $(Y,\mu)$ be a standard probability space, let $A_\mu$ be its at most
countable set of atoms, let $R_\mu=\{\mu(\{a\}) : a\in A_\mu\}$, and for
$r\in R_\mu$ let $A_r$ be the finite set of atoms of mass $r$. With the weak
topologies,
\[
\begin{gathered}
M(\mathrm{Aut}(Y,\mu))\cong\prod_{r\in R_\mu}\mathrm{Sym}(A_r),\\
M(\mathrm{Aut}^*(Y,\mu))\cong\mathrm{LO}(A_\mu),
\end{gathered}
\]
where $\mathrm{Aut}(Y,\mu)$ is the group of measure-preserving automorphisms,
$\mathrm{Aut}^*(Y,\mu)$ is the group of nonsingular automorphisms, and
$\mathrm{LO}(A_\mu)$ is the space of linear orders on $A_\mu$.
\end{theorem}

\begin{remark}\label{rem:lupini-aut-measure}
\uses{thm:lupini-aut-measure}
Both groups split as the product of the corresponding group of the atomless
part, which is extremely amenable by Giordano and Pestov, and a permutation
group of the atoms; Proposition~\ref{prop:ea-normal} and
Theorem~\ref{thm:glasner-weiss} do the rest. The same decomposition gives the
universal minimal flows of measurable wreath products, of full groups of
measured equivalence relations whose atomless part has an extremely amenable
full group, which are products of spaces of linear orders over the atomic
classes, and of automorphism groups of von Neumann algebras such as the
hyperfinite factors \cite[Sections~4 and~5]{LupiniPointwise2026}.
\end{remark}

\begin{theorem}[Lupini \cite{LupiniPointwise2026}]\label{thm:lupini-lp}
\uses{thm:lupini-aut-measure, thm:giordano-pestov}
Let $(Y,\mu)$ be a standard probability space with set of atoms $A_\mu$, let
$\mathbb{F}$ be $\mathbb{R}$ or $\mathbb{C}$, let $S_{\mathbb{F}}$ be the
group of scalars of modulus one, and let $1\leq p<\infty$ with $p\neq2$. The
universal minimal flow of the group of surjective linear isometries of
$L^p(Y,\mu;\mathbb{F})$, with the strong operator topology, has phase space
$S_{\mathbb{F}}^{A_\mu}\times\mathrm{LO}(A_\mu)$, with the skew action of the
group $S_{\mathbb{F}}^{A_\mu}\rtimes\mathrm{Sym}(A_\mu)$ of monomial
isometries of the atomic part; the isometries of the atomless part act
trivially.
In particular, the isometry group is extremely amenable when $\mu$ has no
atoms.
\end{theorem}

\begin{theorem}[Lupini \cite{LupiniPointwise2026}]\label{thm:lupini-wreath}
\uses{thm:bymt, thm:pestov-ea, prop:coprecompact-ea}
Let $F$ be a Polish group whose universal minimal flow $X=M(F)$ is
metrizable, and let $\Omega$ be a countable set. Give $F^\Omega$ the product
topology and $\mathrm{Sym}(\Omega)$ the topology of pointwise convergence.
Then
\[
M\bigl(F^\Omega\rtimes\mathrm{Sym}(\Omega)\bigr)\cong X^\Omega\times\mathrm{LO}(\Omega),
\]
with the skew action
$((f_a)_a,\sigma)\cdot((x_a)_a,<)=((f_a x_{\sigma^{-1}(a)})_a,\sigma\cdot{<})$.
\end{theorem}

\begin{remark}\label{rem:lupini-wreath}
\uses{thm:lupini-wreath}
The proof combines Theorem~\ref{thm:bymt}, which provides a co-precompact
extremely amenable subgroup $F^*\leq F$, with the extremely amenable subgroup
$(F^*)^\Omega\rtimes\mathrm{Aut}(\Omega,<_0)$ for an order $<_0$ of type
$\mathbb{Q}$. Applied to the automorphism groups of the Gurarij space and of
the noncommutative Gurarij space, which are extremely amenable, and of the
Poulsen simplex and the noncommutative Poulsen simplex, whose universal
minimal flows are known, it gives explicit wreath-product formulas; the
preprint makes no separate priority claim for them
\cite[Theorem~4.5 and Corollary~4.6]{LupiniPointwise2026}.
\end{remark}

\begin{theorem}[Lupini \cite{LupiniPointwise2026}]\label{thm:lupini-nuclear}
\uses{prop:ea-normal}
Let $A$ be a unital separable nuclear C*-algebra, and give its unitary group
$U(A)$ the relative weak topology $\sigma(A,A^*)$. Then
\[
M(U(A))\cong\prod_{[\pi]}U(H_\pi),
\]
where $[\pi]$ ranges over the unitary equivalence classes of
finite-dimensional irreducible representations $\pi\colon A\to B(H_\pi)$, and
$u\in U(A)$ acts through the coordinates $\pi(u)$. In particular, $U(A)$ is
extremely amenable in this topology when $A$ has no finite-dimensional
irreducible representations, and $M(U(C(X)))\cong\mathbb{T}^X$ for every
compact metrizable space $X$.
\end{theorem}

\section{Automorphism groups of C*-algebras}\label{sec:operator-algebras}

Automorphism groups of separable C*-algebras, with the point-norm topology
($\alpha_i\to\alpha$ if and only if $\alpha_i(a)\to\alpha(a)$ in norm for
every $a$), are Polish groups. For Kirchberg algebras and AF algebras their
universal minimal flows reduce to classification invariants, because a large
normal subgroup of approximately inner automorphisms is extremely amenable
and can be divided out (Proposition~\ref{prop:ea-normal}). A \emph{Kirchberg
algebra} is a separable, simple, nuclear, purely infinite C*-algebra. The
results in this section are taken from preprints of the editor posted on
Zenodo \cite{LupiniCuntz2026,LupiniStable2026,LupiniUnital2026}; they have
not been refereed.

\begin{theorem}[Lupini \cite{LupiniCuntz2026}]\label{thm:lupini-cuntz}
\uses{def:ea, prop:coprecompact-ea}
Let $\mathcal{O}_2$ be the Cuntz algebra, $\mathcal{K}$ the algebra of
compact operators and $D=\mathcal{O}_2\otimes\mathcal{K}$, with minimal
unitisation $D^+$.
\begin{enumerate}
\item The group $U_1(D)=\{u\in U(D^+) : u-1\in D\}$, with the norm topology,
  is extremely amenable.
\item For every pure state $\varphi$ of $\mathcal{O}_2$, the stabiliser
  $\{\alpha\in\mathrm{Aut}(\mathcal{O}_2) : \varphi\circ\alpha=\varphi\}$, with
  the point-norm topology, is extremely amenable.
\item With the point-norm topology, $M(\mathrm{Aut}(\mathcal{O}_2))$ is the
  state space $S(\mathcal{O}_2)$, with the weak* topology and the action
  $\alpha\cdot\rho=\rho\circ\alpha^{-1}$.
\item With the norm topology, $M(U(D^+))$ is the circle $\mathbb{T}$, on which
  $U(D^+)$ acts through the scalar character of $D^+$.
\end{enumerate}
\end{theorem}

The proof is in Section~\ref{sec:proofs-cuntz}.

\begin{remark}\label{rem:lupini-cuntz}
\uses{thm:lupini-cuntz}
Since $\mathcal{O}_2$ is strongly self-absorbing and its state space is not a
point, part~(3) answers negatively a question of Eagle, Farah, Hart, Kadets,
Kalashnyk and Lupini, whether the automorphism group of every strongly
self-absorbing C*-algebra is extremely amenable. Alekseev, Schmidt and Thom
asked, as a case of a conjecture of theirs, whether
$U((\mathcal{O}_n\otimes\mathcal{K})^+)$ is amenable in the norm topology;
parts~(1) and~(4) answer this affirmatively for $n=2$ \cite{LupiniCuntz2026},
and Theorem~\ref{thm:lupini-stable-norm} below answers it for every $n$.
Part~(3) is obtained from part~(2) as in
Proposition~\ref{prop:coprecompact-ea}: the stabiliser of a pure state is a
co-precompact subgroup whose completion is the whole state space.
\end{remark}

\begin{theorem}[Lupini \cite{LupiniStable2026,LupiniCuntz2026}]\label{thm:lupini-stable}
\uses{prop:ea-normal}
Let $A$ be a stable Kirchberg algebra, and let $N_A$ be the closure in
$\mathrm{Aut}(A)$ of the automorphisms $\mathrm{Ad}(u)$, for $u$ a unitary of
the multiplier algebra of $A$.
\begin{enumerate}
\item The group $N_A$ is extremely amenable, and the map
  $\mathrm{Aut}(A)\to KL(A,A)^\times$ to the group of invertible elements of
  $KL(A,A)$ is continuous, open and surjective, with kernel $N_A$. Hence
  $M(\mathrm{Aut}(A))$ is $M(KL(A,A)^\times)$ with $\mathrm{Aut}(A)$ acting
  through this map, and the two groups have the same Bohr compactification.
  No UCT assumption is needed.
\item If $A$ satisfies the UCT, the same holds with $KL(A,A)^\times$ replaced
  by the automorphism group of total K-theory $\underline{K}(A)$ with its
  coefficient operations, with the topology of pointwise convergence.
\item If $A$ satisfies the UCT and the groups $G_i=K_i(A)$ are finitely
  generated, then the quotient group in~(2) is countable and discrete, and it
  is finite exactly when the free ranks of $G_0$ and $G_1$ are at most one. In
  that case $M(\mathrm{Aut}(A))$ is the translation flow of this finite group,
  of order
  $|\mathrm{Aut}(G_0)|\,|\mathrm{Aut}(G_1)|\,|\mathrm{Ext}(G_0,G_1)|\,|\mathrm{Ext}(G_1,G_0)|$;
  otherwise $M(\mathrm{Aut}(A))$ is not metrizable. Moreover,
  $\mathrm{Aut}(A)$ is amenable if and only if both free ranks are at most
  one.
\end{enumerate}
Part~(3) is proved in \cite{LupiniCuntz2026}, except for the amenability
criterion, which is proved in \cite{LupiniStable2026}.
\end{theorem}

The proof is in Section~\ref{sec:proofs-stable-kirchberg}.

\begin{example}\label{ex:lupini-stable-cuntz}
\uses{thm:lupini-stable}
For $2\leq n<\infty$, the stabilised Cuntz algebra $\mathcal{O}_n\otimes\mathcal{K}$
satisfies the UCT, with $K_0=\mathbb{Z}/(n-1)$ and $K_1=0$. When $K_1(A)=0$
the total K-theory quotient reduces to $\mathrm{Aut}(K_0(A))$, so
$M(\mathrm{Aut}(\mathcal{O}_n\otimes\mathcal{K}))$ is the translation flow of
the group $(\mathbb{Z}/(n-1))^\times$ of units of $\mathbb{Z}/(n-1)$; for
$\mathcal{O}_\infty\otimes\mathcal{K}$, where $K_0=\mathbb{Z}$, it is the
two-point group $\{1,-1\}$ \cite{LupiniCuntz2026}. In particular, the
automorphism groups of $\mathcal{O}_2\otimes\mathcal{K}$ and
$\mathcal{O}_3\otimes\mathcal{K}$ are extremely amenable.
\end{example}

The proof is in Section~\ref{sec:proofs-cuntz}.

\begin{theorem}[Lupini \cite{LupiniStable2026}]\label{thm:lupini-stable-geometric}
\uses{thm:glasner-weiss, thm:gw-cantor}
Let $B$ be a stable Kirchberg algebra whose automorphism group is extremely
amenable, such as $\mathcal{O}_2\otimes\mathcal{K}$. For a nonempty finite set
$I$ and for the Cantor set $X$,
\[
\begin{gathered}
M(\mathrm{Aut}(c_0(I,B)))\cong\mathrm{Sym}(I),\qquad
M(\mathrm{Aut}(c_0(\mathbb{N},B)))\cong\mathrm{LO}(\mathbb{N}),\\
M(\mathrm{Aut}(C(X)\otimes B))\cong\Phi(X),
\end{gathered}
\]
with the actions induced by permutations of the summands and by
homeomorphisms of $X$.
\end{theorem}

The proof is in Section~\ref{sec:proofs-stable-kirchberg}.

\begin{theorem}[Lupini \cite{LupiniStable2026}]\label{thm:lupini-stable-norm}
\uses{prop:ea-normal}
Let $C$ be a unital C*-algebra, let $B=C\otimes\mathcal{K}$, with minimal
unitisation $B^+$, and give $U(B^+)$ the norm topology. Then $U(B^+)$ is
amenable if and only if $B$ is nuclear. If $C$ is nuclear, then
\[
M(U(B^+))\cong\mathbb{T}\times M(K_1(C)),
\]
with $K_1(C)$ discrete, and $U(B^+)$ acting through the scalar character of
$B^+$ and the $K_1$-class. For $C=\mathcal{O}_n$, where $K_1=0$, this gives
$M(U((\mathcal{O}_n\otimes\mathcal{K})^+))\cong\mathbb{T}$ for every $n$.
\end{theorem}

The proof is in Section~\ref{sec:proofs-stable-kirchberg}.

\begin{theorem}[Lupini \cite{LupiniUnital2026}]\label{thm:lupini-unital}
\uses{prop:ea-normal}
Let $A$ be a unital Kirchberg algebra, let $N^0_A$ be the closure of the
group of automorphisms $\mathrm{Ad}(u)$ with $u$ in the identity component
of the unitary group of $A$, and let $q\colon\mathrm{Aut}(A)\to
Q^0_A=\mathrm{Aut}(A)/N^0_A$ be the quotient map. Then
\[
M(\mathrm{Aut}(A))\cong S(A)\times M(Q^0_A),
\qquad \alpha\cdot(\rho,x)=(\rho\circ\alpha^{-1},q(\alpha)x),
\]
and $M(N^0_A)=S(A)$. No UCT assumption is needed. If $A$ satisfies the UCT
and $K_1(A)=0$, then $Q^0_A$ is the group of automorphisms of $K_0(A)$
fixing the class of the unit, with the topology of pointwise convergence.
\end{theorem}

The proof is in Section~\ref{sec:proofs-unital-kirchberg-af}.

\begin{theorem}[Lupini \cite{LupiniUnital2026}]\label{thm:lupini-af}
\uses{prop:ea-normal}
Let $A$ be a separable unital AF algebra with no nonzero finite-dimensional
representations, and let $\Gamma_A$ be the group of automorphisms of the
scaled dimension group $(K_0(A),K_0(A)^+,[1_A])$, with the topology of
pointwise convergence. Then the closure of the group of inner automorphisms of
$A$ is extremely amenable, and $M(\mathrm{Aut}(A))$ is $M(\Gamma_A)$ with
$\mathrm{Aut}(A)$ acting through $\alpha\mapsto K_0(\alpha)$. This applies to
every infinite-dimensional simple separable unital AF algebra. The hypothesis
on representations cannot be dropped: for the matrix algebra $M_d(\mathbb{C})$
with $d\geq2$ the group $\Gamma_A$ is trivial, while $\mathrm{Aut}(A)$ is the
nontrivial compact group $PU(d)$.
\end{theorem}

The proof is in Section~\ref{sec:proofs-unital-kirchberg-af}.

\begin{remark}\label{rem:lupini-unital}
\uses{thm:lupini-unital, thm:lupini-cuntz}
For $A=\mathcal{O}_2$ the group $Q^0_A$ is trivial, and
Theorem~\ref{thm:lupini-unital} recovers part~(3) of
Theorem~\ref{thm:lupini-cuntz}. The same holds for $\mathcal{O}_n$ with
$3\leq n\leq\infty$, and for every purely infinite strongly self-absorbing
C*-algebra $A$: then $\mathrm{Aut}(A)=N^0_A$ and
$M(\mathrm{Aut}(A))\cong S(A)$. For separable stable AF algebras $B$, the map
$\mathrm{Aut}(B)\to\mathrm{Aut}(K_0(B),K_0(B)^+)$ is continuous, open and
surjective with extremely amenable kernel, so $M(\mathrm{Aut}(B))$ is again
computed from the ordered $K_0$-group \cite{LupiniStable2026}. The proofs give more
\cite{LupiniUnital2026}. In the AF case the unitary group, in the norm
topology, and the inner automorphism group are \emph{whirly amenable}, a
property stronger than extreme amenability. In the Kirchberg case $N^0_A$ is
not amenable; the whirly amenable kernel is the stabiliser in $N^0_A$ of a
pure state, and the state space, the universal minimal flow of $N^0_A$,
accounts for the rest.
\end{remark}

\section{Inductive limits of compact Lie groups}\label{sec:inductive-limits}

Increasing unions of compact Lie groups, such as the infinite unitary group,
can carry several natural group topologies, and extreme amenability depends on
the choice. Concentration of measure gives extreme amenability for the
coarser operator topologies (Section~\ref{sec:gromov-milman}); for the finest
topology the answer is negative. The results in this section are taken from
a preprint of the editor posted on Zenodo \cite{LupiniInductive2026}; it has
not been refereed.

\begin{definition}[Final topology]\label{def:final-topology}
Let $K_1\subseteq K_2\subseteq\cdots$ be compact Lie groups, each a closed
subgroup of the next. The \emph{final topology} on $G=\bigcup_nK_n$ declares
$A\subseteq G$ open if and only if $A\cap K_n$ is open in $K_n$ for every $n$.
It is a Hausdorff group topology. For the standard block embeddings
$u\mapsto\mathrm{diag}(u,1)$ this defines the groups $SU(\infty)$,
$SO(\infty)$, $Sp(\infty)$, $U(\infty)$ and $O(\infty)$.
\end{definition}

\begin{theorem}[Lupini \cite{LupiniInductive2026}]\label{thm:lupini-inductive}
\uses{def:final-topology, def:ea}
Let $G=\bigcup_nK_n$ have the final topology, and suppose that some $K_j$ is
not solvable. Then $G$ acts continuously on a nonempty compact metrizable
space $X$ in which no point is fixed by $K_j$. In particular, $G$ is not
extremely amenable.
\end{theorem}

\begin{corollary}[Lupini \cite{LupiniInductive2026}]\label{cor:lupini-inductive}
\uses{thm:lupini-inductive}
With the final topology, none of $SU(\infty)$, $SO(\infty)$, $Sp(\infty)$,
$U(\infty)$ and $O(\infty)$ is extremely amenable.
\end{corollary}

\begin{remark}\label{rem:lupini-inductive}
\uses{cor:lupini-inductive, rem:levy}
For $SU(\infty)$ this answers negatively Question~19 in Pestov's list of open
problems on large topological groups \cite{LupiniInductive2026}. The proof of
Theorem~\ref{thm:lupini-inductive} uses Oliver's smooth actions of compact Lie
groups on disks with prescribed fixed-point sets, joined into an inverse
system of equivariant maps. For $U(\infty)$ and $O(\infty)$ the determinant
already gives a nontrivial compact quotient.
\end{remark}

\begin{theorem}[Lupini \cite{LupiniInductive2026}]\label{thm:lupini-inductive-umf}
\uses{cor:lupini-inductive, rem:umf-size}
For each of $SU(\infty)$, $SO(\infty)$, $Sp(\infty)$, $U(\infty)$ and
$O(\infty)$, with the final topology, the universal minimal flow contains a
homeomorphic copy of $\beta\mathbb{N}$, and is therefore not metrizable. The
Bohr compactifications of $SU(\infty)$, $SO(\infty)$ and $Sp(\infty)$ are
trivial, and those of $U(\infty)$ and $O(\infty)$ are $\mathbb{T}$ and
$\{1,-1\}$, through the determinant.
\end{theorem}

\begin{theorem}[Lupini \cite{LupiniInductive2026}]\label{thm:lupini-trace-class}
\uses{def:ea, rem:levy}
The group
\[
SU_1=\{u\in U(\ell^2) : u-I \text{ is trace class and } \det(u)=1\},
\]
where $\det$ is the Fredholm determinant, is extremely amenable in the
topology of the trace norm $\|u-v\|_1=\mathrm{Tr}|u-v|$, and so is its dense
subgroup $\bigcup_nSU(n)$ with the same topology.
\end{theorem}

\begin{remark}\label{rem:lupini-trace-class}
\uses{thm:lupini-trace-class, thm:lupini-inductive}
The final topology on $\bigcup_nSU(n)$ is strictly finer than the trace-norm
topology. The determinant-one condition is needed in
Theorem~\ref{thm:lupini-trace-class}: on the full group of trace-class
perturbations of the identity, the Fredholm determinant is a nontrivial
compact quotient \cite{LupiniInductive2026}.
\end{remark}

\section{Proofs: groups of maps and of measure-space transformations}\label{sec:proofs-maps-transformations}

This section reproduces, from the preprint of the editor \cite{LupiniPointwise2026} in its version of 17 October 2026, the
proofs of the results stated in Section~\ref{sec:maps-transformations}. The text is that of the
preprint, with British spelling: its sections appear as subsections, its macros
are expanded and its labels renamed, and of its introduction only the
statements of the main results are kept. Citations of the companion preprints
point to their own proof pages where the cited statement is reproduced there.
The preprint has not been refereed. The last subsection derives the statements
of Section~\ref{sec:maps-transformations} from those of the preprint.

\subsection{General principles and pointwise completions}
\label{sec:pw-umf}

We recall the standard completion and extremely-amenable-kernel principles
used throughout; see \cite[Example~6.1.8 and Theorem~6.2.9]{Pestov2006} and
\cite[Lemma~2.3 and Corollary~2.5]{PestovRemarks2002}. Recall that \(G\) is \emph{extremely amenable} when \(M(G)\) is a
singleton.

\begin{lemma}\label{lem:pw-umf-principles}
Let \(G\) be a Hausdorff topological group.
\begin{enumerate}
\item If \(\iota:G\to Q\) is a continuous homomorphism with dense image in a
compact group \(Q\), then left translation through \(\iota\) makes \(Q\) a
minimal \(G\)-flow. Consequently, there is a factor map
\[
M(G)\longrightarrow Q.
\]
\item If \(\iota\) is a topological embedding, so that \(Q\) is the compact
group completion of the precompact group \(G\), then
\[
M(G)\cong Q.
\]
\item If \(N\) is a closed extremely amenable normal subgroup of \(G\), then
\[
M(G)\cong M(G/N).
\]
In particular, if \(G/N\) is compact, then \(M(G)\cong G/N\).
\end{enumerate}
\end{lemma}

\begin{proof}
For the first assertion, every \(G\)-orbit in \(Q\) is dense because
\(\iota(G)\) is dense. The universal property of \(M(G)\) gives the factor
map.

For the second, every compact \(G\)-flow extends uniquely to the completion
\(Q\); see \cite[Fact~2.3]{BassoZucker2021}. If \(X\) is minimal and \(x\in X\), the
extended orbit map \(Q\to X\), \(q\mapsto qx\), is onto. Thus the minimal
\(G\)-flow \(Q\) maps onto every minimal \(G\)-flow.

Finally, let \(X\) be a minimal \(G\)-flow. Extreme amenability supplies an
\(N\)-fixed point. Since \(N\) is normal, the fixed-point set \(X^N\) is
nonempty, closed, and \(G\)-invariant, and hence \(X^N=X\). Thus every
minimal \(G\)-flow factors through \(G/N\), which proves the last assertion.
\end{proof}


\subsubsection{Finite interpolation beyond manifolds}

The pointwise-completion argument only needs the ability to realise arbitrary
finite prescriptions.  This gives a general function-space statement.  We use
the standard notation \(C_p(X,K)\) for the group of continuous maps endowed
with pointwise convergence. The separation property in
\cite[Definition~2.1 and Remark~2.2]{HernandezRodenas2007} supplies finite interpolation;
\cite[Definition~2.2(iii), Proposition~2.3(i), and proof of Theorem~3.3,
equation~(10)]{ShakhmatovSpevak2010} records the resulting completion. The next statement is
its universal-minimal-flow consequence, including the relative version.

\begin{theorem}\label{thm:pw-cp-umf}
Let \(X\) be a Tychonoff space and let \(K\) be a compact path-connected
topological group. Then
\[
M(C_p(X,K))\cong K^X.
\]
More generally, if \(A\subseteq X\) is closed and
\[
C_p(X,A;K)=\{f\in C(X,K):f|_A=e\},
\]
with its pointwise topology, then
\[
M(C_p(X,A;K))\cong K^{X\setminus A}.
\]
In the second formula the omitted coordinates on \(A\) are fixed at the
identity.
\end{theorem}

\begin{proof}
\uses{lem:pw-umf-principles}
Fix distinct points \(x_1,\ldots,x_m\in X\) and prescribed values
\(k_1,\ldots,k_m\in K\). Complete regularity supplies continuous functions
\(\psi_i:X\to[0,1]\) such that
\(\psi_i(x_i)=1\) and \(\psi_i(x_j)=0\) for \(j\ne i\). Choose paths
\(c_i:[0,1]\to K\) from \(e\) to \(k_i\). Then
\[
f(x)=\prod_{i=1}^m c_i(\psi_i(x))
\]
is continuous and satisfies \(f(x_i)=k_i\). Thus the evaluation image of
\(C_p(X,K)\) is dense in \(K^X\).

For the relative group, take the separating functions \(\psi_i\) to vanish on
the closed set \(A\) as well as at the other prescribed points. The same
construction gives density in the closed subgroup of \(K^X\) whose coordinates
on \(A\) are the identity. In both cases this is precisely the compact group
completion, so Lemma~\ref{lem:pw-umf-principles} applies.
\end{proof}

\medskip
The connectedness assumption can be removed if one records the Boolean
information carried by the components. For a Tychonoff space \(X\), write
\(x\sim_{\mathrm{clop}}y\) when every clopen subset of \(X\) contains either
both \(x,y\) or neither of them. Thus the equivalence classes are the
quasi-components of \(X\). If \(A\subseteq X\), put
\[
 \operatorname{bcl}_X(A)
 =\bigcap\{U\subseteq X:U\text{ is clopen and }A\subseteq U\},
\]
with \(\operatorname{bcl}_X(\varnothing)=\varnothing\). We call this the
Boolean hull of \(A\).

Let \(K\) be a compact Lie group, let \(q_0:K\to K/K_0\) be the component
quotient, and define
\begin{align*}
 \mathcal B_X(K)
 &=\{u\in K^X:q_0(u(x))=q_0(u(y))
          \text{ whenever }x\sim_{\mathrm{clop}}y\},\\
 \mathcal B_{X,A}(K)
 &=\{u\in\mathcal B_X(K):u|_A=e,\quad
       q_0(u(x))=K_0\text{ for }x\in\operatorname{bcl}_X(A)\}.
\end{align*}

\begin{theorem}\label{thm:pw-cp-disconnected-umf}
Let \(X\) be Tychonoff and let \(K\) be a compact Lie group, not necessarily
connected. Then
\[
 M(C_p(X,K))\cong\mathcal B_X(K).
\]
If \(A\subseteq X\) is closed, then
\[
 M(C_p(X,A;K))\cong\mathcal B_{X,A}(K).
\]
In particular, the completion is \(K^X\) whenever the quasi-components of
\(X\) are singletons, and it is the subgroup of \(K^X\) with globally
constant component coordinate whenever \(X\) is connected.
\end{theorem}

\begin{proof}
\uses{lem:pw-umf-principles}
Every continuous map from \(X\) to the finite discrete group \(K/K_0\) is
locally constant. Its pointwise closure consists exactly of the functions
\(c:X\to K/K_0\) which are constant on the
\(\sim_{\mathrm{clop}}\)-classes.
Indeed, finitely many distinct classes can be separated by a finite clopen
partition. In the relative case the locally constant maps which equal \(K_0\)
on \(A\) have as their pointwise closure the functions which, in addition,
equal \(K_0\) on \(\operatorname{bcl}_X(A)\). To see the converse, separate
each of the finitely many points outside the Boolean hull from \(A\) by a
clopen set, and refine these sets together with the pairwise separating
clopen sets to a finite clopen partition.

It remains to interpolate inside the components. Fix finitely many points
\(x_1,\ldots,x_m\) and target values prescribed by an element of the displayed
compact group. Choose a locally constant component map with the required
values and, in the relative case, equal to \(K_0\) on \(A\). Choose one
representative of every component of \(K\), taking the representative of
\(K_0\) to be \(e\); composition gives a continuous preliminary map
\(f_0:X\to K\). For each \(x_i\notin A\), the correction
\(f_0(x_i)^{-1}u(x_i)\) lies in the connected compact Lie group \(K_0\).
Choose a path in \(K_0\) from \(e\) to this correction and a continuous scalar
function which is one at \(x_i\) and zero on \(A\) and at the other prescribed
points. Multiplying the resulting corrections into \(f_0\) gives a continuous
map, equal to \(e\) on \(A\) in the relative case, with all the prescribed
values. Thus the evaluation images are dense in the stated compact groups,
which are their compact group completions. Apply
Lemma~\ref{lem:pw-umf-principles}.
\end{proof}

There is also a projective version. If \(Z\leq K\) is closed and central,
write \(\Delta Z\) for the compact subgroup of constant \(Z\)-valued maps.

\begin{corollary}\label{cor:pw-projective-cp-umf}
With the quotient of the pointwise topology,
\[
 M\bigl(C_p(X,K)/\Delta Z\bigr)
 \cong \mathcal B_X(K)/\Delta Z.
\]
For connected \(K\), this reduces to \(K^X/\Delta Z\).
\end{corollary}

\begin{proof}
\uses{lem:pw-umf-principles, thm:pw-cp-disconnected-umf}
The completion of the quotient by the compact subgroup \(\Delta Z\) is the
quotient of the compact completion by \(\Delta Z\). The result follows from
Theorem~\ref{thm:pw-cp-disconnected-umf} and
Lemma~\ref{lem:pw-umf-principles}.
\end{proof}

There is a bundle version. Let \(\mathcal H\to X\) be a locally trivial bundle
of topological groups with typical fibre \(K\), and let
\(\Gamma_p(X,\mathcal H)\) be its continuous-section group with pointwise
convergence.

\begin{corollary}\label{cor:pw-topological-section-umf}
Let \(X\) be paracompact Hausdorff and let \(\mathcal H\to X\) be a locally
trivial bundle of compact path-connected groups. Then
\[
M(\Gamma_p(X,\mathcal H))
 \cong\prod_{x\in X}\mathcal H_x.
\]
If \(A\subseteq X\) is closed, the section group fixed at the identity on
\(A\) has universal minimal flow
\[
\prod_{x\in X\setminus A}\mathcal H_x.
\]
\end{corollary}

\begin{proof}
\uses{lem:pw-umf-principles}
Given finitely many base points, choose pairwise disjoint trivializing
neighbourhoods. In each neighbourhood, a scalar function which is one at the
chosen point and zero off the neighbourhood, composed with a path in the fibre
from the identity to the prescribed value, gives a local section which patches
with the identity section. Multiplying the finitely many local sections proves
density in the displayed fibre product. If the section must be the identity on
\(A\), choose the scalar functions to vanish on \(A\). Apply
Lemma~\ref{lem:pw-umf-principles}.
\end{proof}

The finite-interpolation principle also applies to a regularity class which
does not contain the smooth maps.

\begin{corollary}\label{cor:pw-analytic-mapping-umf}
Let \(X\) be a compact real-analytic manifold and let \(K\) be a connected
compact Lie group. With the topology of pointwise convergence,
\[
M(C^\omega(X,K))\cong K^X.
\]
\end{corollary}

\begin{proof}
\uses{lem:pw-umf-principles}
By the real-analytic embedding theorem \cite{Grauert1958}, choose a real-analytic
embedding \(j:X\to\mathbb R^N\). For distinct
\(x_1,\ldots,x_m\), ordinary polynomial interpolation on the finite subset
\(\{j(x_1),\ldots,j(x_m)\}\) gives polynomials \(p_i\) with
\(p_i(j(x_r))=\delta_{ir}\). Given \(k_i\in K\), choose
\(A_i\in\mathfrak k\) with \(\exp(A_i)=k_i\); the exponential map is
surjective for connected compact Lie groups. The real-analytic map
\[
x\longmapsto\prod_{i=1}^m
 \exp\bigl(p_i(j(x))A_i\bigr)
\]
takes the prescribed values. Hence \(C^\omega(X,K)\) is dense in \(K^X\),
and the conclusion follows from Lemma~\ref{lem:pw-umf-principles}.
\end{proof}

The same argument applies to metric regularity classes which permit finite
scalar interpolation. Fix a compatible bi-invariant Riemannian metric on a
compact connected Lie group \(K\). For a metric space \(X\), write
\(\operatorname{Lip}_p(X,K)\) for all Lipschitz maps with pointwise
convergence, and, for \(0<\gamma\leq1\), define
\(C_p^{0,\gamma}(X,K)\) analogously.

\begin{corollary}\label{cor:pw-lipschitz-mapping-umf}
Let \(X\) be any metric space and let \(K\) be a compact connected Lie group.
Then
\[
 M(\operatorname{Lip}_p(X,K))\cong K^X,
 \qquad
 M(C_p^{0,\gamma}(X,K))\cong K^X
 \quad(0<\gamma\leq1).
\]
\end{corollary}

\begin{proof}
\uses{lem:pw-umf-principles}
On a finite subset \(\{x_1,\ldots,x_m\}\subseteq X\), choose bounded
Lipschitz functions \(\psi_i:X\to[0,1]\) with
\(\psi_i(x_j)=\delta_{ij}\). Given \(k_i\in K\), choose
\(A_i\in\mathfrak k\) with \(\exp(A_i)=k_i\). Then
\[
 x\longmapsto\prod_{i=1}^m\exp(\psi_i(x)A_i)
\]
is Lipschitz and takes the prescribed values. A bounded Lipschitz map is
\(\gamma\)-H\"older for every \(0<\gamma\leq1\), after changing the constant.
Thus both evaluation images are dense in \(K^X\), and
Lemma~\ref{lem:pw-umf-principles} applies.
\end{proof}

There is also a topological compact-support version, extending the smooth
manifold statement above.  Write \(C_{c,p}(X,K)\) for the group of continuous
maps \(f:X\to K\) that equal the identity outside a compact subset of \(X\),
equipped with the topology of pointwise convergence.

\begin{corollary}\label{cor:pw-lch-compact-support-umf}
Let \(X\) be locally compact Hausdorff and let \(K\) be a compact
path-connected topological group. With pointwise convergence,
\[
 M(C_{c,p}(X,K))\cong K^X.
\]
\end{corollary}

\begin{proof}
\uses{thm:pw-cp-umf}
Given finitely many points, local compactness and complete regularity provide
scalar separating functions with compact support. Composing them with paths
in \(K\) and multiplying as in Theorem~\ref{thm:pw-cp-umf} realises every finite
prescription. Hence the compact completion is \(K^X\).
\end{proof}

Finite interpolation has a holomorphic form as well. If \(S\) is a Stein
manifold and \(T\) is a complex Lie group, write
\(\mathcal O_p(S,T)\) for the holomorphic mapping group with pointwise
convergence.

\begin{theorem}\label{thm:pw-holomorphic-mapping-umf}
Let \(S\) be a Stein manifold and let \(T\) be a connected compact complex Lie
group. Then
\[
 M(\mathcal O_p(S,T))\cong T^S.
\]
If \(Z\leq T\) is closed and central, then
\[
 M(\mathcal O_p(S,T)/\Delta Z)\cong T^S/\Delta Z.
\]
\end{theorem}

\begin{proof}
\uses{lem:pw-umf-principles}
A connected compact complex Lie group \(T\) is abelian. Indeed, the
holomorphic matrix entries of its adjoint representation are constant by
compactness, so its Lie algebra is abelian; connectedness then gives the
assertion for \(T\). Its exponential map is onto, since its image is an open
subgroup. Fix distinct \(x_1,\ldots,x_m\in S\). Holomorphic separability of
the Stein manifold supplies functions \(h_j\) with
\(h_j(x_i)=\delta_{ij}\): for each \(i\ne j\), normalise a separating
function to be zero at \(x_i\) and one at \(x_j\), and take the product over
\(i\). Given \(t_j\in T\), choose \(v_j\) in its Lie algebra with
\(\exp(v_j)=t_j\). Then
\[
 x\longmapsto\exp\!\left(\sum_{j=1}^m h_j(x)v_j\right)
\]
is a holomorphic interpolant. Thus \(\mathcal O_p(S,T)\) is dense in
\(T^S\), and Lemma~\ref{lem:pw-umf-principles} and the compact central
quotient give the two formulas. The Oka interpolation theorem
\cite[Theorem~1.1]{ForstnericPrezelj2001} provides a more general route, but is unnecessary
under the present compact-target hypotheses.
\end{proof}

These regularity examples are applications of finite interpolation and the
standard precompact-flow calculation. In particular, the continuous case
uses the explicit earlier completion in \cite{ShakhmatovSpevak2010}; the other displayed
completions are derived directly, without a claim of a new interpolation
principle or historical priority for every formulation.


\subsubsection{One-dimensional NCCW unitary groups}

Let \(E\) and \(F\) be finite-dimensional unital \(C^*\)-algebras, let
\[
\phi=(\phi_0,\phi_1):E\longrightarrow F\oplus F
\]
be a unital \(*\)-homomorphism, and consider the one-dimensional NCCW algebra
\[
A=\{(p,u)\in C(I,F)\oplus E:(p(0),p(1))=\phi(u)\}.
\]
Write
\[
\phi_{\mathrm U}=\phi|_{\mathrm U(E)}:
\mathrm U(E)\longrightarrow\mathrm U(F)\times\mathrm U(F),
\]
and put
\begin{equation}\label{eq:pw-nccw-pointwise-completion}
\mathcal B_A=
\left\{(x,u)\in\mathrm U(F)^I\times\mathrm U(E):
 (x(0),x(1))=\phi_{\mathrm U}(u)\right\}.
\end{equation}
Equip \(\mathrm U(A)\) with the pointwise topology induced by the
\(\mathrm U(E)\)-coordinate and by all evaluations of the path coordinate.

\begin{proposition}\label{prop:pw-nccw-umf}
The compact group \(\mathcal B_A\) is the group completion and universal
minimal flow of pointwise-topological \(\mathrm U(A)\):
\[
M(\mathrm U(A),\tau_{\mathrm{pw}})\cong\mathcal B_A.
\]
For the norm topology, \(\mathcal B_A\) is a compact minimal
\(\mathrm U(A)\)-flow, and hence there is a factor map
\[
M(\mathrm U(A)_{\|\cdot\|})\longrightarrow\mathcal B_A.
\]
If the fibre algebra \(F\) is nonzero, as in the standard unital
one-dimensional NCCW convention, then
\(M(\mathrm U(A)_{\|\cdot\|})\) is nonmetrizable.
\end{proposition}

\begin{proof}
\uses{lem:pw-umf-principles}
The space \(\mathcal B_A\) is the fibre product of the endpoint homomorphism
\(\mathrm U(F)^I\to\mathrm U(F)\times\mathrm U(F)\) and
\(\phi_{\mathrm U}\), so it is a closed subgroup of the compact group
\(\mathrm U(F)^I\times\mathrm U(E)\).

The finite-dimensional group \(\mathrm U(F)\) is a finite product of connected
unitary groups. Given \((x,u)\in\mathcal B_A\) and finitely many coordinates,
add \(0\) and \(1\) to the coordinate set, and join the prescribed values in
their order by continuous paths in \(\mathrm U(F)\), retaining the endpoint
values prescribed by \(\phi_{\mathrm U}(u)\). This gives
\((p,u)\in\mathrm U(A)\) agreeing with \((x,u)\) at the selected coordinates.
Hence \(\mathrm U(A)\) is dense in \(\mathcal B_A\), and
Lemma~\ref{lem:pw-umf-principles} proves both the pointwise calculation and the
norm-topological factor assertion.

Taking \(u=1_E\) and fixing the endpoint coordinates at \(1_F\) gives a closed
copy of \(\mathrm U(F)^{(0,1)}\) in \(\mathcal B_A\). Since \(F\) is a
nonzero unital finite-dimensional \(C^*\)-algebra under the stated
convention, \(\mathrm U(F)\) is
nontrivial, and this product is nonmetrizable. Since a Hausdorff continuous
image of a compact metrizable space is metrizable, the displayed factor map
forces \(M(\mathrm U(A)_{\|\cdot\|})\) to be nonmetrizable.
\end{proof}

\begin{remark}\label{rem:pw-nccw-pointwise-coordinate}
\uses{prop:pw-nccw-umf}
Including the \(\mathrm U(E)\)-coordinate in \(\tau_{\mathrm{pw}}\) is
essential. If the topology is induced only by evaluations of the path
coordinate, then the compact kernel \(\ker(\phi_{\mathrm U})\) is not
separated; the Hausdorff quotient instead completes to
\[
\{x\in\mathrm U(F)^I:
  (x(0),x(1))\in\phi_{\mathrm U}(\mathrm U(E))\}.
\]
Proposition~\ref{prop:pw-nccw-umf} identifies \(\mathcal B_A\) only as a
factor in the norm topology, not as the full norm-topological universal
minimal flow.
\end{remark}


\subsection{Convergence in measure and sampled topologies}

The opposite extreme from pointwise convergence is provided by measurable
maps. Let \((Y,\mu)\) be a standard probability space and let
\[
A_\mu=\{a\in Y:\mu(\{a\})>0\}
\]
be its at most countable set of atomic points.

\begin{proposition}\label{prop:pw-l0-umf}
For every compact metrizable group \(K\), with \(L^0(Y,\mu;K)\) carrying
convergence in measure,
\[
M(L^0(Y,\mu;K))\cong K^{A_\mu}.
\]
In particular, \(L^0(Y,\mu;K)\) is extremely amenable when \(\mu\) is
atomless.
\end{proposition}

\begin{proof}
\uses{lem:pw-umf-principles}
Let \(Y_{\mathrm c}=Y\setminus A_\mu\). Restriction gives a topological group
isomorphism
\[
L^0(Y,\mu;K)
\cong L^0(Y_{\mathrm c},\mu;K)\times K^{A_\mu}.
\]
Indeed, for a bounded compatible bi-invariant metric \(d_K\), the atomic
factor carries the weighted metric
\[
\sum_{a\in A_\mu}\mu(\{a\})d_K(k_a,\ell_a),
\]
which induces the product topology. For the compact target used here, extreme amenability of the atomless
\(L^0\)-group is already \cite[Theorem~2.20]{GiordanoPestov2007}; the circle case appears
in \cite[Theorems~1.2--1.3]{Glasner1998}. Pestov and Schneider prove the more
general amenable-target theorem \cite[Theorem~1.1]{PestovSchneider2017}. If \(\mu(Y_{\mathrm c})>0\), normalise the
restriction of \(\mu\) to \(Y_{\mathrm c}\), which does not change the
convergence-in-measure topology; the resulting standard atomless probability
space is isomorphic modulo null sets to the standard model. If
\(\mu(Y_{\mathrm c})=0\), the first factor is trivial. The compact group \(K\)
is amenable, so Lemma~\ref{lem:pw-umf-principles} applies to the projection onto
\(K^{A_\mu}\).
\end{proof}

The same argument permits the compact group to vary measurably.  Say that a
Borel bundle of compact metrizable groups \(\mathcal H\to Y\) is
\emph{countably trivializable} if there is a countable Borel partition
\(Y=\bigsqcup_jY_j\) and compact metrizable groups \(H_j\) such that
\(\mathcal H|_{Y_j}\) is Borel-isomorphic, as a group bundle, to
\(Y_j\times H_j\).

\begin{theorem}\label{thm:pw-measurable-bundle-umf}
Let \((Y,\mu)\) be a standard probability space and let
\(\mathcal H\to Y\) be a countably trivializable Borel bundle of compact
metrizable groups.  Equip the measurable-section group
\(L^0(Y,\mu;\mathcal H)\) with convergence in measure. Then
\[
M(L^0(Y,\mu;\mathcal H))
 \cong\prod_{a\in A_\mu}\mathcal H_a.
\]
In particular, this measurable-section group is extremely amenable when
\(\mu\) is atomless.
\end{theorem}

\begin{proof}
\uses{lem:pw-umf-principles}
Restriction to the atoms is a continuous surjective homomorphism onto the
displayed compact product. Its kernel is the measurable-section group over
\(Y\setminus A_\mu\). After applying the countable trivializing partition,
this kernel is a product of groups of the form
\(L^0(Y_j\setminus A_\mu,\mu;H_j)\). Each nontrivial factor is extremely
amenable by \cite[Theorem~1.1]{PestovSchneider2017}, and products of extremely amenable
groups are extremely amenable. Here convergence in measure gives the product
topology on the partitioned group: after finitely many pieces the remaining
total measure is arbitrarily small. Lemma~\ref{lem:pw-umf-principles} now
identifies the universal minimal flow with the compact quotient.
\end{proof}

\begin{corollary}\label{cor:pw-measurable-gauge-umf}
Let \(Y\) be a compact manifold with a regular Borel probability measure
\(\mu\), let \(K\) be a compact Lie group, and let \(P\to Y\) be a smooth
principal \(K\)-bundle.  Put
\[
\operatorname{Gau}^0_\mu(P)=L^0(Y,\mu;\operatorname{Ad}P)
\]
be its measurable gauge group, modulo equality almost everywhere, with
convergence in measure. Then
\[
M(\operatorname{Gau}^0_\mu(P))
\cong\prod_{a\in A_\mu}\operatorname{Ad}(P)_a.
\]
In particular, the measurable gauge group is extremely amenable when
\(\mu\) is atomless.
\end{corollary}

\begin{proof}
\uses{thm:pw-measurable-bundle-umf}
A finite or countable Borel partition subordinate to bundle trivializations
shows that \(\operatorname{Ad}(P)\) is countably trivializable. Apply
Theorem~\ref{thm:pw-measurable-bundle-umf}.
\end{proof}

\begin{corollary}\label{cor:pw-measure-topology-umf}
Suppose that \(K\) is connected. Equip any path, loop, twisted loop, finite
endpoint-pullback, or one-dimensional gauge group considered above with the
coarser topology induced solely by convergence of its map or section
coordinate in normalised volume measure. If this topology does not separate
auxiliary coordinates of a pullback, take its Hausdorff quotient. Then the
resulting group is extremely amenable, and its universal minimal flow is a
point. The same conclusion holds after imposing finitely many basepoint or
framing conditions.
\end{corollary}

\begin{proof}
\uses{prop:pw-l0-umf}
The map into the appropriate measurable-map or measurable-section group
induces precisely the stated topology on its Hausdorff image. That image is
dense. For maps, approximate a measurable map by a finite-valued simple map,
approximate its level sets by finite unions of intervals, and join the
finitely many values across transition intervals of arbitrarily small total
measure. For gauge groups, cut the one-manifold into finitely many trivializing
arcs and place the transitions in neighbourhoods of arbitrarily small total
measure; equivalently, use a finite Borel partition subordinate to bundle
trivializations to identify the measurable-section group with an
\(L^0\)-group. Connectedness of \(K\) supplies the joining paths. Endpoint,
derivative-matching, twisting, and framing conditions can be repaired on
neighbourhoods whose measures tend to zero. Thus the completion is an
atomless \(L^0\)-group, which is extremely amenable by
Proposition~\ref{prop:pw-l0-umf}. Universal minimal flows are unchanged on
passing to a dense completion by \cite[Fact~2.3]{BassoZucker2021}.
\end{proof}

Only the map or section coordinate enters this corollary. For example, in an
NCCW pullback a unitary vertex coordinate need not be determined by the path
almost everywhere. If that coordinate is retained with its compact topology,
the completion instead retains it: for connected unitary fibres it is
\(L^0(I;\mathrm U(F))\times\mathrm U(E)\), and its universal minimal flow
is \(\mathrm U(E)\). Endpoint conditions can be repaired on intervals of
arbitrarily small measure, so they impose no relation between these two
completion coordinates. This is a different topology from the one in the
corollary.

\begin{remark}\label{rem:pw-disconnected-measure-umf}
\uses{cor:pw-measure-topology-umf}
For disconnected \(K\), the component restriction does not disappear under
completion. We use the same Hausdorff-quotient convention for auxiliary
coordinates as in Corollary~\ref{cor:pw-measure-topology-umf}. Let \(C=K/K_0\), let \(\overline L\) be the image of
\(L\leq K\times K\) in \(C\times C\), and set
\[
C_L=\{c\in C:(c,c)\in\overline L\}.
\]
For a finite endpoint pullback with convergence in Lebesgue measure, the
completion has an extremely amenable normal \(L^0(I,K_0)\)-kernel and compact
quotient \(C_L\), and hence its universal minimal flow is \(C_L\). By
\cite[Fact~2.3]{BassoZucker2021}, the original dense endpoint pullback has the same
universal minimal flow. Thus unbased paths and ordinary loops give \(C\),
based groups give a point, and an \(\alpha\)-twisted loop group gives
\(\operatorname{Fix}(\pi_0(\alpha))\).
\end{remark}


\subsubsection{Mixed sampling topologies and measure-space symmetries}

There is a continuum of useful group topologies between convergence in measure
and evaluation at every point. Let \(X\) be a compact connected smooth
manifold, let \(\mu\) be an atomless regular Borel probability measure of full
support, let \(K\) be a connected compact Lie group, and fix an arbitrary set
\(D\subseteq X\). Define \(\tau_{\mu,D}\) to be the topology on
\(C^\infty(X,K)\) induced by
\[
\jmath_D:C^\infty(X,K)\longrightarrow
 L^0(X,\mu;K)\times K^D,
 \qquad f\longmapsto([f],(f(d))_{d\in D}).
\]
Full support makes this homomorphism injective.

\begin{theorem}\label{thm:pw-mixed-sampling-umf}
For every \(D\subseteq X\),
\[
M(C^\infty(X,K),\tau_{\mu,D})\cong K^D.
\]
Thus \(D=\varnothing\) gives extreme amenability, while for \(D=X\) the
universal minimal flow agrees with the one for the pointwise topology.
\end{theorem}

\begin{proof}
\uses{lem:pw-umf-principles}
Smooth maps are dense in \(L^0(X,\mu;K)\). One way to see this is to
approximate a measurable map by a finite-valued simple map, approximate its
level sets inside finitely many coordinate charts, and insert smooth transition
regions of arbitrarily small measure; connectedness of \(K\) supplies paths
between the finitely many values.

This approximation can be made simultaneously with any finite prescription on
\(D\). After first choosing a smooth approximation, multiply it by corrections
supported in pairwise disjoint coordinate balls around the finitely many sampled
points. The balls can be chosen with arbitrarily small total measure, so the
corrections impose the prescribed point values without changing the limit in
measure. Hence \(\jmath_D\) has dense image and the two-sided completion is
\[
L^0(X,\mu;K)\times K^D.
\]
The first factor is extremely amenable by \cite[Theorem~1.1]{PestovSchneider2017}; it is a
closed normal subgroup and the quotient is the compact group \(K^D\).
Lemma~\ref{lem:pw-umf-principles} gives the formula.
\end{proof}

The Hausdorff completion underlying the preceding result can be computed for
arbitrary measurable maps, without connectedness or smoothness assumptions.
Let \(\mathcal M(Y,K)\) be the group of Borel maps \(Y\to K\), with no
almost-everywhere identification. For \(D\subseteq Y\), let
\[
 \jmath_{\mu,D}:\mathcal M(Y,K)\longrightarrow
 L^0(Y,\mu;K)\times K^D,
 \qquad f\longmapsto([f],f|_D),
\]
and write \(G_{\mu,D}(K)\) for its image with the subspace topology.
Equivalently, this is the Hausdorff quotient of \(\mathcal M(Y,K)\) by maps
which equal the identity almost everywhere and at every point of \(D\).

\begin{theorem}\label{thm:pw-sampled-measurable-umf}
Let \((Y,\mu)\) be a standard probability space, let \(K\) be a compact
metrizable group, and let \(D\subseteq Y\) be arbitrary. Then
\[
 M(G_{\mu,D}(K))\cong K^{A_\mu\cup D}.
\]
More precisely, the two-sided completion of \(G_{\mu,D}(K)\) is naturally
\[
 L^0(Y_{\mathrm c},\mu;K)\times K^{A_\mu\cup D},
\]
where an atomic coordinate occurring also in \(D\) is included only once.
\end{theorem}

\begin{proof}
\uses{lem:pw-umf-principles}
An element of \(L^0(Y,\mu;K)\) has a well-defined value on every atom. Hence
the closure of the image of \(\jmath_{\mu,D}\) is the closed subgroup
\[
 \{(g,w)\in L^0(Y,\mu;K)\times K^D:
        w_a=g(a)\text{ for }a\in A_\mu\cap D\}.
\]
It is isomorphic to the displayed product. To prove density, choose a Borel
representative of \(g\). A basic neighbourhood specifies only finitely many
coordinates in \(D\); changing the representative at the non-atomic points
among them does not alter its \(L^0\)-class, while compatibility at the atomic
points is already part of the definition of the closed subgroup.

By \cite[Theorem~1.1]{PestovSchneider2017}, the first factor is extremely amenable. Projection
onto \(K^{A_\mu\cup D}\) therefore has an extremely amenable kernel and compact
quotient. Lemma~\ref{lem:pw-umf-principles}, together with invariance under
two-sided completion, proves the formula.
\end{proof}

For an atomless smooth probability space and connected compact Lie target, the
dense smooth subgroup of \(G_{\mu,D}(K)\) is precisely the group in
Theorem~\ref{thm:pw-mixed-sampling-umf}. Thus
Theorem~\ref{thm:pw-sampled-measurable-umf} both explains that result and adds
arbitrary atoms and arbitrary compact metrizable targets.

We next apply the same extremely-amenable-kernel argument to transformations of
a measure space. Let \((Y,\mu)\) be a standard probability space, retain the
notation \(A_\mu\) for its set of atomic points, and put
\[
Y_{\mathrm c}=Y\setminus A_\mu,
\qquad
R_\mu=\{\mu(\{a\}):a\in A_\mu\},
\qquad
A_r=\{a\in A_\mu:\mu(\{a\})=r\}.
\]
Every \(A_r\) is finite. Write \(\operatorname{Aut}(Y,\mu)_{\mathrm w}\) for the
measure-preserving automorphism group with its weak topology and
\(\operatorname{Aut}^*(Y,\mu)_{\mathrm w}\) for the nonsingular automorphism group with
the corresponding weak topology. Equivalently, the latter is the topology
induced by the quasi-regular representation
\[
 U_\tau^{(p)}f=
 \left(\frac{d(\mu\circ\tau^{-1})}{d\mu}\right)^{1/p}
 f\circ\tau^{-1}
\]
in the strong operator topology of \(L^p(Y,\mu)\), for one (and hence every)
\(1\leq p<\infty\); see \cite[Section~6.2]{GiordanoPestov2007}. This convention will also
be used below for vector-valued quasi-regular representations.

\begin{theorem}\label{thm:pw-measure-aut-umf}
There are natural isomorphisms of flows
\begin{align}
M(\operatorname{Aut}(Y,\mu)_{\mathrm w})
 &\cong \prod_{r\in R_\mu}\operatorname{Sym}(A_r),
 \label{eq:pw-measure-preserving-aut-umf}\\
M(\operatorname{Aut}^*(Y,\mu)_{\mathrm w})
 &\cong \operatorname{LO}(A_\mu),
 \label{eq:pw-nonsingular-aut-umf}
\end{align}
where \(\operatorname{LO}(A_\mu)\) is the compact space of all linear orders
on \(A_\mu\), with the natural permutation action. For a finite atomic set this
flow is the regular translation flow of \(\operatorname{Sym}(A_\mu)\), and for
an empty atomic set it is a point.
\end{theorem}

\begin{proof}
\uses{lem:pw-umf-principles}
A measure-preserving automorphism preserves atom masses, while a nonsingular
automorphism may permute all atoms. Restriction to the atomic and nonatomic
parts gives topological direct-product decompositions
\begin{align*}
\operatorname{Aut}(Y,\mu)_{\mathrm w}
 &\cong \operatorname{Aut}(Y_{\mathrm c},\mu)_{\mathrm w}
   \times\prod_{r\in R_\mu}\operatorname{Sym}(A_r),\\
\operatorname{Aut}^*(Y,\mu)_{\mathrm w}
 &\cong \operatorname{Aut}^*(Y_{\mathrm c},\mu)_{\mathrm w}
   \times\operatorname{Sym}(A_\mu).
\end{align*}
The product topology on the atomic coordinates is the induced weak topology:
control of finitely many positive-mass atoms gives a basic coordinate
condition, and the remaining atomic mass can be made arbitrarily small.

Giordano and Pestov prove that both nonatomic factors above are extremely
amenable, already in \cite[Theorems~2.1 and~2.4]{GiordanoPestov2002}; see also the fuller
proofs in \cite[Theorem~4.2 and Theorem~6.1]{GiordanoPestov2007}. The remaining factor in the
first decomposition is compact, which proves
\eqref{eq:pw-measure-preserving-aut-umf}. If \(A_\mu\) is countably infinite,
Glasner and Weiss compute
\(M(\operatorname{Sym}(A_\mu))=\operatorname{LO}(A_\mu)\)
\cite{GlasnerWeiss2002}. For finite \(A_\mu\), the permutation action on its linear orders
is simply transitive, so it is the compact group's own universal minimal flow.
Lemma~\ref{lem:pw-umf-principles} proves
\eqref{eq:pw-nonsingular-aut-umf} in every case.
\end{proof}


\subsection{Wreath products and atomic decompositions}

Let a nonsingular transformation \(T\) act on \(L^0(Y,\mu;K)\) by
\[
 (T\mathbin{\cdot}f)(y)=f(T^{-1}y).
\]
For a compact metrizable group \(K\), form the measurable wreath products
\begin{align*}
 W_\mu(K)
  &=L^0(Y,\mu;K)\rtimes\operatorname{Aut}(Y,\mu)_{\mathrm w},\\
 W_\mu^*(K)
  &=L^0(Y,\mu;K)\rtimes\operatorname{Aut}^*(Y,\mu)_{\mathrm w}.
\end{align*}

\begin{theorem}\label{thm:pw-measurable-wreath-umf}
There are natural isomorphisms of flows
\begin{align}
 M(W_\mu(K))
 &\cong K^{A_\mu}\rtimes
       \prod_{r\in R_\mu}\operatorname{Sym}(A_r),
 \label{eq:pw-preserving-wreath-umf}\\
 M(W_\mu^*(K))
 &\cong K^{A_\mu}\times\operatorname{LO}(A_\mu).
 \label{eq:pw-nonsingular-wreath-umf}
\end{align}
The first compact group acts on itself by left translation. In the second
formula, the atomic subgroup
\(K^{A_\mu}\rtimes\operatorname{Sym}(A_\mu)\) acts by
\[
 (z,\sigma)\mathbin{\cdot}(w,<)
 =\bigl(z(\sigma\mathbin{\cdot}w),\sigma\mathbin{\cdot}<\bigr),
 \qquad
 (\sigma\mathbin{\cdot}w)_a=w_{\sigma^{-1}(a)},
\]
and the nonatomic factor acts trivially. In particular, both measurable wreath
products are extremely amenable when \(\mu\) is atomless.
\end{theorem}

\begin{proof}
\uses{lem:pw-umf-principles}
Restriction to the atomic and nonatomic parts gives
\begin{align*}
 W_\mu(K)
 &\cong
 \bigl(L^0(Y_{\mathrm c},\mu;K)\rtimes
       \operatorname{Aut}(Y_{\mathrm c},\mu)_{\mathrm w}\bigr)
 \times
 \left(K^{A_\mu}\rtimes
       \prod_{r\in R_\mu}\operatorname{Sym}(A_r)\right),\\
 W_\mu^*(K)
 &\cong
 \bigl(L^0(Y_{\mathrm c},\mu;K)\rtimes
       \operatorname{Aut}^*(Y_{\mathrm c},\mu)_{\mathrm w}\bigr)
 \times
 \bigl(K^{A_\mu}\rtimes\operatorname{Sym}(A_\mu)\bigr).
\end{align*}
Each first factor is extremely amenable: its \(L^0\)-kernel is extremely
amenable by \cite[Theorem~1.1]{PestovSchneider2017}, its quotient is extremely amenable by
\cite[Theorem~4.2 and Theorem~6.1]{GiordanoPestov2007}, and
Lemma~\ref{lem:pw-umf-principles} applies twice.

The second factor in the measure-preserving decomposition is compact, which
gives \eqref{eq:pw-preserving-wreath-umf}. In the nonsingular case,
\(K^{A_\mu}\) is compact and normal. The semidirect-product formula
\cite[Proposition~6.1]{KechrisSokic2012}, together with
\(M(\operatorname{Sym}(A_\mu))=\operatorname{LO}(A_\mu)\) from \cite{GlasnerWeiss2002},
gives the phase space and the stated skew action in
\eqref{eq:pw-nonsingular-wreath-umf}. The extension splits, so no cocycle
correction is needed here: the skew action is the one produced by the
splitting, as in the proof of \cite[Proposition~6.1]{KechrisSokic2012}. The extremely amenable direct factors do
not change either universal minimal flow.
\end{proof}

The compact fibre in the preceding theorem can be replaced by a much larger
class without changing the form of the answer.

\begin{definition}\label{def:pw-ea-by-compact}
A Polish group \(F\) is \emph{extremely-amenable-by-compact} if it has a closed
extremely amenable normal subgroup \(F_0\) such that

\[
 Q_F=F/F_0
\]

is compact.
\end{definition}

\begin{theorem}\label{thm:pw-ea-compact-measurable-wreath}
\uses{thm:pw-measurable-wreath-umf}
Let \(F\) be extremely-amenable-by-compact and give \(L^0(Y,\mu;F)\)
the topology of convergence in measure. Then
\begin{align}
 M(L^0(Y,\mu;F))
  &\cong Q_F^{A_\mu},
  \label{eq:pw-ea-compact-l0-umf}\\
 M\!\left(L^0(Y,\mu;F)\rtimes\operatorname{Aut}(Y,\mu)_{\mathrm w}\right)
  &\cong Q_F^{A_\mu}\rtimes
        \prod_{r\in R_\mu}\operatorname{Sym}(A_r),
  \label{eq:pw-ea-compact-preserving-wreath}\\
 M\!\left(L^0(Y,\mu;F)\rtimes\operatorname{Aut}^*(Y,\mu)_{\mathrm w}\right)
  &\cong Q_F^{A_\mu}\times\operatorname{LO}(A_\mu).
  \label{eq:pw-ea-compact-nonsingular-wreath}
\end{align}
The compact group in \eqref{eq:pw-ea-compact-preserving-wreath} acts on itself by
left translation.  In \eqref{eq:pw-ea-compact-nonsingular-wreath} the quotient
\(Q_F^{A_\mu}\rtimes\operatorname{Sym}(A_\mu)\) has the standard skew action
described in Theorem~\ref{thm:pw-measurable-wreath-umf}, and the remaining
nonatomic factor acts trivially.  In particular, all three groups are
extremely amenable when \(\mu\) is atomless.
\end{theorem}

\begin{proof}
\uses{lem:pw-umf-principles}
The group \(F\) is amenable, since both \(F_0\) and its compact quotient are
amenable.  On the nonatomic part, Pestov and Schneider's theorem therefore
makes \(L^0(Y_{\mathrm c},\mu;F)\) extremely amenable
\cite[Theorem~1.1]{PestovSchneider2017}.  On the atoms, the closed normal subgroup
\(F_0^{A_\mu}\) is extremely amenable and the quotient is the compact group
\(Q_F^{A_\mu}\).  Restriction to the two parts, followed by
Lemma~\ref{lem:pw-umf-principles}, proves \eqref{eq:pw-ea-compact-l0-umf}.

After adjoining measure-space transformations, the nonatomic semidirect
factor is still extremely amenable: its kernel is the preceding measurable-map
group and its quotient is extremely amenable by \cite{GiordanoPestov2007}.  On the atomic
part, divide first by \(F_0^{A_\mu}\).  In the measure-preserving case the
remaining group is the compact group displayed in
\eqref{eq:pw-ea-compact-preserving-wreath}.  In the nonsingular case it is

\[
 Q_F^{A_\mu}\rtimes\operatorname{Sym}(A_\mu).
\]

The compact-normal-subgroup formula \cite[Proposition~6.1]{KechrisSokic2012} and the
Glasner--Weiss computation \cite{GlasnerWeiss2002} give
\eqref{eq:pw-ea-compact-nonsingular-wreath} and its skew action.
\end{proof}

The same compact-extension argument gives a general supply of examples over
countable structures.

\begin{corollary}\label{cor:pw-nonarch-wreath-umf}
Let \(\Omega\) be countable, let
\(H\leq\operatorname{Sym}(\Omega)\) be closed in the topology of pointwise
convergence, and let \(K\) be compact metrizable. Then
\[
 M(K^\Omega\rtimes H)\cong K^\Omega\times M(H),
\]
with the standard skew action. Consequently,
\begin{align*}
 M(K^{\mathbb Q}\rtimes\operatorname{Aut}(\mathbb Q,<))
 &\cong K^{\mathbb Q},\\
 M(K^{\mathbb N}\rtimes\operatorname{Aut}(\mathcal R))
 &\cong K^{\mathbb N}\times\operatorname{LO}(\mathbb N),
\end{align*}
where \(\mathcal R\) is the countable random graph.
\end{corollary}

\begin{proof}
The group \(K^\Omega\) is compact and normal, so the first formula is
\cite[Proposition~6.1]{KechrisSokic2012}. The group
\(\operatorname{Aut}(\mathbb Q,<)\) is extremely amenable, while the universal
minimal flow of \(\operatorname{Aut}(\mathcal R)\) is the space of all linear
orders on its vertex set; both statements follow from the Ramsey-theoretic
computations in \cite{KPT2005}.
\end{proof}

The compactness of the fibre in Corollary~\ref{cor:pw-nonarch-wreath-umf} is
unnecessary when its universal minimal flow is metrizable.  This gives a
wreath-product construction which is not restricted to compact homogeneous
spaces. Before giving its stabiliser proof, we record the earlier comparisons.
Soki\'c's relational composition theorem \cite[Theorem~6.5(v)]{Sokic2013}
identifies the universal flow with admissible convex orders: the fibre class
has strong amalgamation, at least one class has strong joint embedding, and
the reasonable ordered expansions have the Ramsey and ordering properties. With finite
sets as the quotient class, a convex order consists of an order on the
quotient and an admissible order in each fibre; the automorphism group and
flow are respectively \(F^{\mathbb N}\rtimes\operatorname{Sym}(\mathbb N)\) and
\(M(F)^{\mathbb N}\times\operatorname{LO}(\mathbb N)\). Compositions of more general
expansions are treated in \cite[Proposition~2.2 of the 2017 preprint]{PawliukSokic2020}.
Jahel and Zucker's \cite[Theorem~5.1]{JahelZucker2026} instead establishes metrizability
and unique ergodicity for Polish wreath products. The next theorem gives the
explicit flow for arbitrary Polish fibres with metrizable flow, using the
standard co-precompact stabiliser construction; it does not require an ordered
relational presentation of the fibre group.

\begin{theorem}\label{thm:pw-metrizable-fibre-wreath-umf}
Let \(F\) be a Polish group whose universal minimal flow \(X=M(F)\) is
metrizable, and let \(\Omega\) be countable.  Give \(F^\Omega\) the product
topology and \(\operatorname{Sym}(\Omega)\) the pointwise topology.  Then
\[
 M\!\left(F^\Omega\rtimes\operatorname{Sym}(\Omega)\right)
 \cong X^\Omega\times\operatorname{LO}(\Omega).
\]
The action is the standard skew action
\[
 ((f_a)_{a\in\Omega},\sigma)\mathbin{\cdot}((x_a)_{a\in\Omega},<)
 =\bigl((f_a\mathbin{\cdot}x_{\sigma^{-1}(a)})_{a\in\Omega},
        \sigma\mathbin{\cdot}<\bigr).
\]
For finite \(\Omega\), the same formula applies and
\(\operatorname{LO}(\Omega)\) is the simply transitive
\(\operatorname{Sym}(\Omega)\)-space of linear orders.
\end{theorem}

\begin{proof}
By \cite[Theorem~1.1]{BYMT2017}, there are a point \(x_0\in X\) and a closed
extremely amenable co-precompact subgroup
\[
 F^*=\operatorname{Stab}_F(x_0)
\]
such that the completion of \(F/F^*\), for the quotient of the right
uniformity, is \(X\).  If \(\Omega\) is infinite, choose a linear order
\(<_0\) of type \(\mathbb Q\) on \(\Omega\); if it is finite, choose any
linear order.  The subgroup
\[
 H=(F^*)^\Omega\rtimes\operatorname{Aut}(\Omega,<_0)
 \leq F^\Omega\rtimes\operatorname{Sym}(\Omega)
\]
is extremely amenable. Indeed, \((F^*)^\Omega\) is extremely amenable, and
so is the quotient \(\operatorname{Aut}(\Omega,<_0)\), by the Ramsey theorem
in the infinite case and trivially in the finite case. Extreme amenability is
preserved under extensions, so the assertion follows.

Put \(G=F^\Omega\rtimes\operatorname{Sym}(\Omega)\) and
\(y_0=((x_0)_{a\in\Omega},<_0)\). The stabiliser of \(y_0\) is exactly \(H\).
The orbit map
\[
 \Theta:G/H\longrightarrow X^\Omega\times\operatorname{LO}(\Omega),
 \qquad \Theta(gH)=g\mathbin{\cdot}y_0,
\]
is one-to-one, and its image is
\[
 (F\mathbin{\cdot}x_0)^\Omega\times
 \bigl(\operatorname{Sym}(\Omega)\mathbin{\cdot}<_0\bigr).
\]
We spell out the uniform-completion point. A basic entourage of the quotient
right uniformity on \(G/H\) is obtained by choosing a finite set
\(D\subseteq\Omega\), requiring the two induced linear orders to agree on
\(D\),
and requiring, for \(a\in D\), the corresponding two \(F/F^*\)-coordinates
to lie in prescribed entourages of the quotient right uniformity. This follows
from the product topology on \(F^\Omega\), the pointwise topology on
\(\operatorname{Sym}(\Omega)\), and the freedom to multiply on the right by
the independent coordinates of \((F^*)^\Omega\) and by
\(\operatorname{Aut}(\Omega,<_0)\). More explicitly, agreement of the two
induced orders on \(D\) gives an order-preserving bijection between their
inverse images under the two permutations. Homogeneity of the rational order
extends this finite bijection to an order automorphism. After making this
correction, the independent \(F^*\)-coordinates realise each of the finitely
many quotient entourages. The converse follows by restricting a small group
displacement to those coordinates. For finite \(\Omega\), agreement on all
of \(\Omega\) replaces the order-extension step. These are precisely the entourages induced
through \(\Theta\) by the product of the quotient right uniformities on finite
coordinate sets.

By the choice of \(F^*\), the completion of \(F/F^*\) is \(X\). The completion
of \(\operatorname{Sym}(\Omega)/\operatorname{Aut}(\Omega,<_0)\) is
\(\operatorname{LO}(\Omega)\), by the Glasner--Weiss computation (and by
finiteness when \(\Omega\) is finite). The finite-coordinate description
therefore shows that \(H\) is co-precompact and that \(\Theta\) extends to a
\(G\)-isomorphism
\[
 \widehat{G/H}\cong X^\Omega\times\operatorname{LO}(\Omega).
\]
The displayed skew action is minimal: first use a permutation to meet any
prescribed finite order condition and then use the independent \(F\)-coordinates
and minimality of \(X\) to meet any prescribed finite family of open subsets of
\(X\). An extremely amenable co-precompact subgroup makes every minimal
subflow of \(\widehat{G/H}\) a copy of \(M(G)\), as recalled in the introduction
of \cite{BYMT2017}; the minimality just proved therefore gives the formula.
\end{proof}

This produces concrete examples from functional analysis that are not
extremely-amenable-by-compact.  Let \(\mathbb G\) and \(\mathbb{NG}\) denote
the Gurarij Banach space and the noncommutative Gurarij operator space.  Let
\(\mathbb P\) be the Poulsen simplex, and let \(\mathbb{NP}_1\) be the compact
space of scalar states of the noncommutative Poulsen simplex.

\begin{corollary}\label{cor:pw-gurarij-poulsen-wreath-umf}
For every countable \(\Omega\),
\begin{align*}
 M\!\left(\operatorname{Aut}(\mathbb G)^\Omega
          \rtimes\operatorname{Sym}(\Omega)\right)
 &\cong\operatorname{LO}(\Omega),\\
 M\!\left(\operatorname{Aut}(\mathbb{NG})^\Omega
          \rtimes\operatorname{Sym}(\Omega)\right)
 &\cong\operatorname{LO}(\Omega),\\
 M\!\left(\operatorname{Aut}(\mathbb P)^\Omega
          \rtimes\operatorname{Sym}(\Omega)\right)
 &\cong\mathbb P^\Omega\times\operatorname{LO}(\Omega),\\
 M\!\left(\operatorname{Aut}(\mathbb{NP})^\Omega
          \rtimes\operatorname{Sym}(\Omega)\right)
 &\cong\mathbb{NP}_1^\Omega\times\operatorname{LO}(\Omega).
\end{align*}
Here automorphisms mean surjective linear isometries for \(\mathbb G\),
surjective complete isometries for \(\mathbb{NG}\), affine homeomorphisms for
\(\mathbb P\), and matrix-affine homeomorphisms for \(\mathbb{NP}\).
\end{corollary}

\begin{proof}
\uses{thm:pw-metrizable-fibre-wreath-umf}
The groups \(\operatorname{Aut}(\mathbb G)\) and
\(\operatorname{Aut}(\mathbb{NG})\) are extremely amenable
\cite{BartosovaLopezAbadLupiniMbombo2022,BartosovaLopezAbadLupiniMbombo2021}.  The canonical actions on \(\mathbb P\) and
\(\mathbb{NP}_1\) are the respective universal minimal flows of
\(\operatorname{Aut}(\mathbb P)\) and \(\operatorname{Aut}(\mathbb{NP})\),
again by \cite{BartosovaLopezAbadLupiniMbombo2022,BartosovaLopezAbadLupiniMbombo2021}.  Apply
Theorem~\ref{thm:pw-metrizable-fibre-wreath-umf}.
\end{proof}

Theorem~\ref{thm:pw-metrizable-fibre-wreath-umf} and
Corollary~\ref{cor:pw-gurarij-poulsen-wreath-umf} are recorded as consequences of
the cited metrizable-flow and Ramsey results; no separate claim of priority is
made for these formulas.


\subsubsection{Full groups of measured equivalence relations}

Let \(\mathcal R\) be a countable nonsingular measured equivalence relation on
the standard probability space \((Y,\mu)\). Its full group \([\mathcal R]\)
is endowed with the uniform metric
\[
 d_{\mathrm u}(S,T)=\mu\{y\in Y:S(y)\ne T(y)\}.
\]
Put \(\mathcal R_{\mathrm c}=\mathcal R|_{Y_{\mathrm c}}\), and let
\[
 \mathcal C_{\mathcal R}=A_\mu/\mathcal R
\]
be the countable family of \(\mathcal R\)-classes contained in the atomic
part. We regard \(\operatorname{LO}(C)\) as a singleton when \(|C|\leq1\).

\begin{theorem}\label{thm:pw-full-group-atomic-umf}
Assume that \([\mathcal R_{\mathrm c}]\), with the uniform topology, is
extremely amenable. Then
\[
 M([\mathcal R])
 \cong \prod_{C\in\mathcal C_{\mathcal R}}\operatorname{LO}(C),
\]
with the coordinatewise order action. In particular, the conclusion applies
when \(\mathcal R_{\mathrm c}\) is ergodic and amenable. If \(\mathcal R\) is
measure preserving, every atomic class \(C\) is finite, and the displayed flow
is equivariantly isomorphic to the left-translation flow of the compact group
\[
 \prod_{C\in\mathcal C_{\mathcal R}}\operatorname{Sym}(C).
\]
\end{theorem}

\begin{proof}
\uses{lem:pw-umf-principles}
Quasi-invariance prevents an atom from being \(\mathcal R\)-equivalent to a
point of \(Y_{\mathrm c}\): the saturation of a null singleton is null and
therefore cannot contain an atom. Restriction consequently gives a topological
direct-product decomposition
\[
 [\mathcal R]
 \cong [\mathcal R_{\mathrm c}]
 \times\prod_{C\in\mathcal C_{\mathcal R}}\operatorname{Sym}(C).
\]
The topology on the second factor is pointwise convergence. Indeed, fixing a
positive-mass atom is a uniform open condition, while outside a sufficiently
large finite set of atoms the total mass is arbitrarily small.

Write \(H=\prod_C\operatorname{Sym}(C)\). For each countably infinite \(C\),
choose an order \(<_C\) of type \(\mathbb Q\), and choose the unique order type
on each finite \(C\). The subgroup
\[
 H^*=\prod_C\operatorname{Aut}(C,<_C)
\]
is extremely amenable by the Ramsey theorem and closure of extreme amenability
under products. It is co-precompact in \(H\), the completion of \(H/H^*\) is
\(\prod_C\operatorname{LO}(C)\), and the resulting action is minimal. The
co-precompact-subgroup description of universal minimal flows
\cite[Theorem~1.1]{BYMT2017} therefore gives
\[
 M(H)=\prod_C\operatorname{LO}(C);
\]
equivalently, this is the coordinatewise product of the Glasner--Weiss
calculation \cite{GlasnerWeiss2002}. The extremely amenable first direct factor does not
change the universal minimal flow by Lemma~\ref{lem:pw-umf-principles}.

If \(\mathcal R\) is measure preserving, atoms in the same class have equal
mass. Such a class is finite, and the action of \(\operatorname{Sym}(C)\) on
\(\operatorname{LO}(C)\) is simply transitive. Finally, Giordano and Pestov
prove the asserted extreme amenability when \(\mathcal R_{\mathrm c}\) is
ergodic and amenable \cite[Theorem~5.7]{GiordanoPestov2007}.
\end{proof}

The same calculation allows compact labels along the equivalence relation.

\begin{corollary}\label{cor:pw-decorated-full-group-umf}
\uses{thm:pw-full-group-atomic-umf}
Let \(K\) be compact metrizable and retain the assumption of
Theorem~\ref{thm:pw-full-group-atomic-umf}. For the action by precomposition, put
\[
 W_{\mathcal R}(K)=L^0(Y,\mu;K)\rtimes[\mathcal R].
\]
Then
\[
 M(W_{\mathcal R}(K))
 \cong K^{A_\mu}\times
       \prod_{C\in\mathcal C_{\mathcal R}}\operatorname{LO}(C).
\]
The atomic group
\[
 K^{A_\mu}\rtimes
 \prod_{C\in\mathcal C_{\mathcal R}}\operatorname{Sym}(C)
\]
acts by the standard skew action, and the nonatomic factor acts trivially.
\end{corollary}

\begin{proof}
\uses{thm:pw-full-group-atomic-umf}
Restriction to \(Y_{\mathrm c}\) and \(A_\mu\) decomposes
\(W_{\mathcal R}(K)\) as the direct product of
\[
 L^0(Y_{\mathrm c},\mu;K)\rtimes[\mathcal R_{\mathrm c}]
\]
and the displayed atomic group. The first factor is extremely amenable: its
measurable-map kernel is extremely amenable by \cite[Theorem~1.1]{PestovSchneider2017}, its
quotient is extremely amenable by assumption, and fixed points lift through
extensions. On the atomic factor, \(K^{A_\mu}\) is compact and normal.
Theorem~\ref{thm:pw-full-group-atomic-umf} and the compact semidirect-product
formula \cite[Proposition~6.1]{KechrisSokic2012} give the result.
\end{proof}

\begin{corollary}\label{cor:pw-ea-compact-full-group-umf}
\uses{thm:pw-full-group-atomic-umf}
Let \(F\) be extremely-amenable-by-compact and retain the hypothesis on
\(\mathcal R\) from Theorem~\ref{thm:pw-full-group-atomic-umf}. Then
\[
 M\!\left(L^0(Y,\mu;F)\rtimes[\mathcal R]\right)
 \cong Q_F^{A_\mu}\times
       \prod_{C\in\mathcal C_{\mathcal R}}\operatorname{LO}(C).
\]
The quotient
\[
 Q_F^{A_\mu}\rtimes
 \prod_{C\in\mathcal C_{\mathcal R}}\operatorname{Sym}(C)
\]
acts by the coordinatewise standard skew action, and the extremely amenable
kernel acts trivially.
\end{corollary}

\begin{proof}
\uses{thm:pw-ea-compact-measurable-wreath, thm:pw-full-group-atomic-umf}
The nonatomic semidirect factor is extremely amenable, as in the proof of
Theorem~\ref{thm:pw-ea-compact-measurable-wreath}.  On the atomic part, divide by
the extremely amenable normal subgroup \(F_0^{A_\mu}\), then apply the compact
semidirect-product formula \cite[Proposition~6.1]{KechrisSokic2012} and
Theorem~\ref{thm:pw-full-group-atomic-umf}.
\end{proof}

The nonatomic extreme-amenability input in
Theorem~\ref{thm:pw-full-group-atomic-umf} and
Corollaries~\ref{cor:pw-decorated-full-group-umf} and
\ref{cor:pw-ea-compact-full-group-umf} is due to Giordano and Pestov. Retaining
the atoms exposes the order flows of the individual atomic equivalence
classes.

\medskip
For scalar \(\ell^p\), the following atomic skew flow is already stated
in \cite[Theorem~4.13 and Remark~4.14]{PestovRemarks2002}. Splitting the atomic and
nonatomic bands gives the corresponding mixed-space formula. We record the
calculation for the isometry groups of \(L^p\)-spaces. Put
\[
\mathbb S_{\mathbb C}=\mathbb T,
\qquad
\mathbb S_{\mathbb R}=\{-1,1\}.
\]

\begin{theorem}\label{thm:pw-lp-isometry-umf}
Let \(\mathbb F\in\{\mathbb R,\mathbb C\}\) and
\(1\leq p<\infty\), \(p\ne2\). Equip the group of surjective linear
isometries of \(L^p(Y,\mu;\mathbb F)\) with the strong operator topology. Then
the phase space of its universal minimal flow is
\[
\mathbb S_{\mathbb F}^{A_\mu}\times\operatorname{LO}(A_\mu).
\]
The action is the standard skew action of the atomic monomial group
\[
\mathbb S_{\mathbb F}^{A_\mu}\rtimes
\operatorname{Sym}(A_\mu),
\]
acting as follows: if \(z,w\in\mathbb S_{\mathbb F}^{A_\mu}\),
\(\sigma\in\operatorname{Sym}(A_\mu)\), and
\(<\in\operatorname{LO}(A_\mu)\), then
\[
 (z,\sigma)\mathbin{\cdot}(w,<)
 =\bigl(z(\sigma\mathbin{\cdot}w),\sigma\mathbin{\cdot}<\bigr),
 \qquad
 (\sigma\mathbin{\cdot}w)_a=w_{\sigma^{-1}(a)},
\]
where \(a\,(\sigma\mathbin{\cdot}<)\,b\) means
\(\sigma^{-1}(a)<\sigma^{-1}(b)\). The isometries of the nonatomic summand
act trivially on this flow. In
particular, the isometry group is extremely amenable when \(\mu\) is atomless.
\end{theorem}

\begin{proof}
\uses{lem:pw-umf-principles, thm:pw-measure-aut-umf}
The Banach--Lamperti description, in the topological form used in
\cite[Section~6.2]{GiordanoPestov2007}, preserves the atomic and nonatomic bands and gives
\[
\operatorname{Iso}(L^p(Y,\mu;\mathbb F))
\cong E_{\mathrm c}\times
 \bigl(\mathbb S_{\mathbb F}^{A_\mu}\rtimes
       \operatorname{Sym}(A_\mu)\bigr).
\]
Here \(E_{\mathrm c}\) is the isometry group of the nonatomic summand and is
extremely amenable by \cite[Theorem~6.6]{GiordanoPestov2007}. On the atomic summand, after
normalising the characteristic function of every atom, the strong operator
topology is the product topology on the scalar coordinates together with
pointwise convergence on the permutation coordinate.

The group \(\mathbb S_{\mathbb F}^{A_\mu}\) is compact and normal. The
semidirect-product formula for universal minimal flows with a compact normal
factor \cite[Proposition~6.1]{KechrisSokic2012}, together with Theorem~\ref{thm:pw-measure-aut-umf},
identifies the phase space with
\(\mathbb S_{\mathbb F}^{A_\mu}\times\operatorname{LO}(A_\mu)\) and supplies
the stated skew action. The extremely amenable direct factor does not change
the universal minimal flow by Lemma~\ref{lem:pw-umf-principles}.
\end{proof}


\subsection{Operator-algebraic applications}


\subsubsection{Weak unitary groups of injective and nuclear algebras}

The operator-algebraic compact factor can be identified explicitly.
The underlying compact-times-extremely-amenable decomposition is already
\cite[Corollary~3.6]{GiordanoPestov2002}; we record its factor coordinates using the fuller
proof in \cite{GiordanoPestov2007}. Let \(N\)
be an injective von Neumann algebra and let
\(\{z_i:i\in I_{\mathrm{fd}}(N)\}\) be the family of minimal central
projections for which \(z_iN\) is finite-dimensional. Write
\(z_iN\cong M_{n_i}(\mathbb C)\) and put
\[
\mathcal C_N=\prod_{i\in I_{\mathrm{fd}}(N)}\mathrm U(n_i).
\]

\begin{theorem}\label{thm:pw-injective-vn-umf}
For the \(s(N,N_{\ast})\)-topology,
\[
M(\mathrm U(N)_{\mathrm s})\cong\mathcal C_N.
\]
Thus \(\mathrm U(N)_{\mathrm s}\) is extremely amenable precisely when \(N\) has no
nonzero finite-dimensional factor summand.
\end{theorem}

\begin{proof}
\uses{lem:pw-umf-principles}
Let \(z=\sum_{i\in I_{\mathrm{fd}}(N)}z_i\). The central decomposition
\[
N=zN\oplus(1-z)N
\]
gives a topological direct product decomposition of the strong-topological
unitary group. The first factor is the compact group \(\mathcal C_N\).
Giordano and Pestov's proof of \cite[Theorem~3.3]{GiordanoPestov2007} shows that the unitary
group of the complementary summand is extremely amenable: the nonatomic
finite type-I, finite continuous, and properly infinite injective parts are
handled in their Proposition~3.4, Proposition~3.7, and Corollary~3.6,
respectively. Lemma~\ref{lem:pw-umf-principles} gives the formula.
\end{proof}

For example, if
\[
N=L^\infty(Y,\mu)\,\overline\otimes\,M_n(\mathbb C),
\]
then \(\mathrm U(N)_{\mathrm s}=L^0(Y,\mu;\mathrm U(n))\)
\cite[Lemma~3.2]{GiordanoPestov2007}, and Theorem~\ref{thm:pw-injective-vn-umf} recovers
\[
M(\mathrm U(N)_{\mathrm s})\cong\mathrm U(n)^{A_\mu}.
\]

The preceding theorem also gives a formula directly at the \(C^*\)-algebra
level. For a unital \(C^*\)-algebra \(A\), let
\(\widehat A_{\mathrm{fd}}\) be the set of unitary-equivalence classes of
finite-dimensional irreducible representations, and choose a representative
\(\pi:A\to\mathcal B(H_\pi)\) in each class.

\begin{lemma}\label{lem:pw-finite-dimensional-envelope-summands}
Let \(A\) be a unital \(C^*\)-algebra. The finite-dimensional factor summands
of its universal von Neumann envelope \(A^{**}\) are in canonical bijection
with \(\widehat A_{\mathrm{fd}}\). If \(z_{\mathrm{fd}}\) is their central
sum, then
\[
 z_{\mathrm{fd}}A^{**}
 \cong
 \bigoplus_{[\pi]\in\widehat A_{\mathrm{fd}}}^{\ell^\infty}
 \mathcal B(H_\pi),
\qquad
 \mathrm U(z_{\mathrm{fd}}A^{**})_{\mathrm s}
 \cong
 \prod_{[\pi]\in\widehat A_{\mathrm{fd}}}\mathrm U(H_\pi),
\]
where the last group has the product topology.
\end{lemma}

\begin{proof}
A finite-dimensional irreducible representation \(\pi\) has a normal
surjective extension
\(\widetilde\pi:A^{**}\to\mathcal B(H_\pi)\). Its kernel is a weak-star
closed ideal, hence has the form \((1-z_\pi)A^{**}\) for a central projection
\(z_\pi\), and
\(z_\pi A^{**}\cong\mathcal B(H_\pi)\). In particular, \(z_\pi\) is a
minimal nonzero central projection. Hence two such central projections are
either orthogonal or equal. If two such central summands agree,
the corresponding finite-dimensional quotient representations differ by an
automorphism of a full matrix algebra and are therefore unitarily equivalent.
Thus inequivalent irreducible representations give orthogonal summands.

Conversely, restricting a finite-dimensional factor-summand representation of
\(A^{**}\) to the weak-star dense copy of \(A\) gives an irreducible
finite-dimensional representation: in finite dimension weak-star density is
norm density, so the image is the full matrix algebra. These two constructions
are inverse. The first displayed decomposition follows by taking the central
sum. The strong-topological unitary group of a von Neumann direct sum is the
product of the strong-topological unitary groups of its summands
\cite[Lemma~3.1]{GiordanoPestov2007}; on each finite-dimensional summand this is the usual
compact topology.
\end{proof}

\begin{theorem}\label{thm:pw-nuclear-weak-umf}
Let \(A\) be a unital separable nuclear \(C^*\)-algebra. Equip
\(\mathrm U(A)\) with the relative weak topology \(\sigma(A,A^*)\). Then
\[
M(\mathrm U(A)_{\mathrm w})
 \cong
 \prod_{[\pi]\in\widehat A_{\mathrm{fd}}}\mathrm U(H_\pi),
\]
where \(u\in\mathrm U(A)\) acts through the coordinates \(\pi(u)\).
In particular, \(\mathrm U(A)_{\mathrm w}\) is extremely amenable when
\(A\) has no finite-dimensional irreducible representations.
\end{theorem}

\begin{proof}
\uses{lem:pw-finite-dimensional-envelope-summands, thm:pw-injective-vn-umf}
The two-sided completion of \(\mathrm U(A)\) in the relative weak topology is
\(\mathrm U(A^{**})\) with the \(s(A^{**},A^*)\)-topology; see the proof of
\cite[Theorem~3.8]{GiordanoPestov2007}. Nuclearity makes \(A^{**}\) injective. Its
finite-dimensional factor summands are indexed precisely by
\(\widehat A_{\mathrm{fd}}\), by
Lemma~\ref{lem:pw-finite-dimensional-envelope-summands}, and their unitary groups
form the displayed compact product. Theorem~\ref{thm:pw-injective-vn-umf} computes the universal
minimal flow of the completion. Universal minimal flows are unchanged on
passing to the two-sided completion by \cite[Fact~2.3]{BassoZucker2021}.
\end{proof}

\begin{corollary}\label{cor:pw-nuclear-weak-examples}
\uses{prop:pw-nccw-umf}
Let \(X\) be compact metrizable.
\begin{enumerate}
\item
\(M(\mathrm U(C(X))_{\mathrm w})\cong\mathbb T^X\).
\item If \(A\) is unital, separable, nuclear, and has no finite-dimensional
representations, then \(M(\mathrm U(A)_{\mathrm w})\) is a point.
\item For the one-dimensional NCCW algebra in
Proposition~\ref{prop:pw-nccw-umf},
\[
M(\mathrm U(A)_{\mathrm w})\cong\mathcal B_A.
\]
Thus \(\mathcal B_A\) is the full universal minimal flow in both the
pointwise and relative weak topologies, but remains only a known factor in the
norm topology.
\end{enumerate}
\end{corollary}

\begin{proof}
\uses{thm:pw-nuclear-weak-umf}
The irreducible representations of \(C(X)\) are its characters, indexed by
\(X\). The second assertion is the empty-product case of
Theorem~\ref{thm:pw-nuclear-weak-umf}. For a one-dimensional NCCW algebra, the
finite-dimensional irreducible representations consist of the interior fibre
representations together with the representations of the finite-dimensional
quotient \(E\). Consequently, the product in
Theorem~\ref{thm:pw-nuclear-weak-umf} is naturally
\(\mathrm U(E)\times\mathrm U(F)^{(0,1)}\), which is isomorphic to
\(\mathcal B_A\) by filling in the endpoint coordinates using
\(\phi_{\mathrm U}\).
\end{proof}

There is a corresponding formula after dividing by scalar unitaries. Put
\[
 \operatorname{PU}(A)=\mathrm U(A)/(\mathbb T1_A),
 \qquad
 \mathcal Q_A=
 \prod_{[\pi]\in\widehat A_{\mathrm{fd}}}\mathrm U(H_\pi),
\]
and, when the product is nonempty, let
\[
 \Delta\mathbb T
 =\{(\lambda1_{H_\pi})_{[\pi]}:\lambda\in\mathbb T\}
 \leq\mathcal Q_A.
\]
For the empty product, \(\mathcal Q_A/\Delta\mathbb T\) is understood to be a
point.

We use the following elementary completion fact. If \(G\) is a Hausdorff
topological group with two-sided completion \(\widehat G\), and \(C\leq G\)
is compact and central, then the two-sided completion of \(G/C\) is
\(\widehat G/C\). Indeed, compactness makes \(C\) closed in \(\widehat G\);
the quotient uniformity on \(G/C\) is the restriction of that on
\(\widehat G/C\), and the image of \(G/C\) is dense. This also shows that no
additional kernel appears when the scalar circle is removed below.

\begin{theorem}\label{thm:pw-projective-nuclear-weak-umf}
Let \(A\) be a unital separable nuclear \(C^*\)-algebra, and equip
\(\operatorname{PU}(A)\) with the quotient of the relative weak topology on
\(\mathrm U(A)\). Then
\[
 M(\operatorname{PU}(A)_{\mathrm w})
 \cong \mathcal Q_A/\Delta\mathbb T.
\]
The action is induced coordinatewise by the finite-dimensional irreducible
representations: the class of \(u\in\mathrm U(A)\) sends the class of
\((v_\pi)_{[\pi]}\) to the class of
\((\pi(u)v_\pi)_{[\pi]}\).
\end{theorem}

\begin{proof}
\uses{lem:pw-umf-principles, thm:pw-injective-vn-umf, thm:pw-nuclear-weak-umf}
Let \(z\in A^{**}\) be the sum of the finite-dimensional factor summands and
set
\[
 E_A=\mathrm U((1-z)A^{**}).
\]
The proof of Theorem~\ref{thm:pw-nuclear-weak-umf} identifies the two-sided
completion of \(\mathrm U(A)_{\mathrm w}\) with
\(E_A\times\mathcal Q_A\), and Theorem~\ref{thm:pw-injective-vn-umf} shows that
\(E_A\) is extremely amenable. Since the scalar circle is compact and central,
the completion of \(\operatorname{PU}(A)_{\mathrm w}\) is
\[
 (E_A\times\mathcal Q_A)/\Delta_A\mathbb T,
 \qquad
 \Delta_A\mathbb T
 =\{(\lambda1_{(1-z)A^{**}},
       (\lambda1_{H_\pi})_{[\pi]}):\lambda\in\mathbb T\}.
\]
If \(\mathcal Q_A\) is nontrivial, projection onto the second coordinate
induces a quotient onto \(\mathcal Q_A/\Delta\mathbb T\) whose kernel is
isomorphic to \(E_A\). If \(\mathcal Q_A\) is trivial, the completion is
\(E_A/(\mathbb T1)\), which is extremely amenable because quotients of
extremely amenable groups are extremely amenable. Thus in either case the
displayed compact quotient has an extremely amenable closed normal kernel.
Lemma~\ref{lem:pw-umf-principles}, followed by invariance of the universal minimal
flow under two-sided completion \cite[Fact~2.3]{BassoZucker2021}, proves the formula.
\end{proof}

\begin{corollary}\label{cor:pw-projective-nuclear-weak-examples}
\uses{prop:pw-nccw-umf}
Let \(X\) be compact metrizable.
\begin{enumerate}
\item
\[
 M(\operatorname{PU}(C(X))_{\mathrm w})
 \cong\mathbb T^X/\Delta\mathbb T.
\]
\item If \(A\) is unital, separable, nuclear, and has no finite-dimensional
irreducible representations, then \(\operatorname{PU}(A)_{\mathrm w}\) is
extremely amenable.
\item For the one-dimensional NCCW algebra in
Proposition~\ref{prop:pw-nccw-umf},
\[
 M(\operatorname{PU}(A)_{\mathrm w})
 \cong\mathcal B_A/\Delta\mathbb T.
\]
\end{enumerate}
\end{corollary}

\noindent\textit{Proof.}
For \(C(X)\), the finite-dimensional irreducible representations are the
characters indexed by \(X\). The second assertion is the empty-product case of
Theorem~\ref{thm:pw-projective-nuclear-weak-umf}. The identification
\(\mathcal Q_A\cong\mathcal B_A\) in the proof of
Corollary~\ref{cor:pw-nuclear-weak-examples} intertwines the two diagonal scalar
circles, which gives the last formula.
\hfill\(\square\)
\par

The weak-topological projective unitary formula is a consequence of the
preceding completion argument. For \(A=M_n(\mathbb C)\),
the formula reduces, as it should, to the compact group
\(\operatorname{PU}(n)=\mathrm U(n)/\mathbb T\).

\subsection{Proofs of the statements in Section~\ref{sec:maps-transformations}}\label{sec:pw-encyclopaedia-proofs}

\begin{proof}[Proof of Theorem~\ref{thm:lupini-cp}]
\proves{thm:lupini-cp}
\uses{thm:pw-cp-umf}
This is the first assertion of Theorem~\ref{thm:pw-cp-umf}.
\end{proof}

\begin{proof}[Proof of Theorem~\ref{thm:lupini-aut-measure}]
\proves{thm:lupini-aut-measure}
\uses{thm:pw-measure-aut-umf}
This is Theorem~\ref{thm:pw-measure-aut-umf}.
\end{proof}

\begin{proof}[Proof of Theorem~\ref{thm:lupini-lp}]
\proves{thm:lupini-lp}
\uses{thm:pw-lp-isometry-umf}
This is Theorem~\ref{thm:pw-lp-isometry-umf}.
\end{proof}

\begin{proof}[Proof of Theorem~\ref{thm:lupini-wreath}]
\proves{thm:lupini-wreath}
\uses{thm:pw-metrizable-fibre-wreath-umf}
This is Theorem~\ref{thm:pw-metrizable-fibre-wreath-umf}.
\end{proof}

\begin{proof}[Proof of Theorem~\ref{thm:lupini-nuclear}]
\proves{thm:lupini-nuclear}
\uses{thm:pw-nuclear-weak-umf, cor:pw-nuclear-weak-examples}
The formula and the extreme amenability are Theorem~\ref{thm:pw-nuclear-weak-umf};
the formula for $U(C(X))$ is part~(1) of Corollary~\ref{cor:pw-nuclear-weak-examples}.
\end{proof}

\section{Proofs: automorphism and unitary groups of Cuntz algebras}\label{sec:proofs-cuntz}

This section reproduces, from the preprint of the editor \cite{LupiniCuntz2026} in its version of 12 September 2026, the
proofs of the results stated in Section~\ref{sec:operator-algebras}. The text is that of the
preprint, with British spelling: its sections appear as subsections, its macros
are expanded and its labels renamed, and of its introduction only the
statements of the main results are kept. Citations of the companion preprints
point to their own proof pages where the cited statement is reproduced there.
The preprint has not been refereed. The last subsection derives the statements
of Section~\ref{sec:operator-algebras} from those of the preprint.

\subsection{Main results of the preprint}\label{sec:cu-main-results}

\begin{theorem}\label{thm:cu-main}
\uses{def:ea, thm:ellis}
 Put $D=\mathcal{O}_2\otimes\mathcal{K}$.

\begin{enumerate}
\item The norm group $U_1(D)$ is extremely amenable.

\item For every pure state $\varphi$ of $\mathcal{O}_2$, the point-norm
group 
\begin{equation*}
H_\varphi=\{\alpha\in\mathrm{Aut}(\mathcal{O}_2):\varphi\circ\alpha=\varphi\}
\end{equation*}
is extremely amenable.

\item With point-norm topology on the automorphism group, 
\begin{equation*}
M(\mathrm{Aut}(\mathcal{O}_2))\cong S(\mathcal{O}_2).
\end{equation*}
The state space has its weak-* topology and the action $\alpha\cdot\rho=\rho \circ\alpha^{-1}$.

\item With norm topology on the unitary group, 
\begin{equation*}
M(U(D^+))\cong\mathbb{T}.
\end{equation*}
The action is multiplication through the scalar character of $D^+$.
\end{enumerate}
\end{theorem}

\begin{theorem}\label{thm:cu-finite-main}
 For every finite tuple $u_1,\ldots,u_m\in U_1(D)$
and every $\varepsilon>0$, there are $w_1,\ldots,w_l\in U_1(D)$ and
permutations $\sigma_j$ of $\{1,\ldots,l\}$ such that 
\begin{equation*}
\frac1l\sum_{a=1}^l \left\|w_{\sigma_j(a)}u_j-w_a\right\|<\varepsilon
\qquad(1\leq j\leq m).
\end{equation*}
The witnessing unitaries $w_a$ can be supported on one finite matrix corner
of $D$, with the identity on its complement.
\end{theorem}

\subsection{Dynamical and operator-algebraic preliminaries}

\label{sec:cu-prelim}


\subsubsection{Topologies and compact actions}

All compact spaces are nonempty Hausdorff spaces and all group actions are
jointly continuous. The group $\mathrm{Aut}(B)$ of a separable C*-algebra
is Polish in the point-norm topology. Norm unitary groups and their compact
central quotients are Polish SIN groups. Multiplier unitary groups, when
used, carry the strict topology. These conventions are not interchangeable.

A group is amenable if it has a left-invariant mean on $\mathrm{RUCB}(G)$,
the bounded right-uniformly continuous functions. Equivalently, every
compact action admits an invariant regular probability. A group is extremely
amenable if every compact action has a fixed point.

\begin{definition}\label{def:cu-definition-1}
We call a group \emph{whirly} if every invariant regular probability in
every compact action has support contained in the fixed-point set.
\end{definition}

We use this fixed-support formulation throughout. Whirliness does not assert
existence of invariant probabilities. A whirly amenable group is extremely
amenable. The term is used here for a property of a topological group and
its compact actions; it is an analogue, not a special case, of the
measure-theoretic whirliness of Glasner--Tsirelson--Weiss \cite{GlasnerTsirelsonWeiss2005} for
measure-preserving actions.

\begin{lemma}\label{lem:cu-permanence}
\uses{def:cu-definition-1, def:ea}
 Extreme amenability passes through continuous
homomorphisms with dense image, through finite products, and through groups
with a dense directed union of extremely amenable subgroups. Whirliness
passes through continuous dense homomorphisms and through groups
topologically generated by continuous images of whirly groups.
\end{lemma}

\begin{proof}
Pull back an action along the homomorphism. A point fixed by its image is
fixed by its closure. For products, first take the fixed points of one
factor and then of the other. For a directed union, the closed fixed-point
sets have the finite intersection property. For the whirly statements, apply
the same argument to the closed support of an invariant probability: this
one set is fixed by every generating image and hence by its generated
closure.
\end{proof}

\begin{lemma}\label{lem:cu-kernel}
\uses{def:ea, thm:ellis}
 Let $q:G\to Q$ be a continuous open surjective
homomorphism with extremely amenable kernel. Then $M(G)\cong q^*M(Q)$. For a
compact group $C$ and any topological group $L$, 
\begin{equation*}
M(C\times L)\cong C\times M(L).
\end{equation*}
\end{lemma}

\begin{proof}
In a minimal $G$-flow, the kernel has a nonempty closed fixed-point set.
Normality makes it $G$-invariant, so it is the entire flow. The action
descends continuously through the open quotient. This identifies minimal
flows and their factor maps. For the second assertion, in a minimal $C\times
L$-flow $X$ choose an $L$-minimal subflow $Y$. The set $CY$ is compact and
invariant, hence equals $X$. An $L$-factor map $f:M(L)\to Y$ gives the
factor map $(c,x)\mapsto cf(x)$. The displayed product is minimal under its
independent actions.
\end{proof}

\begin{lemma}[Homogeneous completion]\label{lem:cu-completion-criterion}
\uses{def:coprecompact, def:ea, thm:ellis}
 Let $G$ be a topological group and $H\leq G$
a closed, extremely amenable, co-precompact subgroup. Let $\widehat{G/H}$ be
the completion of $G/H$ in the quotient right uniformity. If $\widehat{G/H}$
is a minimal $G$-flow, then $M(G)\cong\widehat{G/H}$.
\end{lemma}

\begin{proof}
Let $X$ be a compact minimal $G$-flow and choose an $H$-fixed point $x\in X$. Joint continuity of the action, uniformly over the compact space $X$,
makes $g\mapsto gx$ right-uniformly continuous; it is constant on left
cosets of $H$, so it descends to a uniformly continuous map $G/H\to X$ and
extends to $\widehat{G/H}$. The extension is equivariant, and its image is a
nonempty closed invariant subset of $X$, hence all of $X$. Thus $\widehat{G/H}$ factors onto every minimal flow; being minimal itself, it is the
universal one.
\end{proof}

The converse direction is the substantial one and is due to Melleray--Nguyen
Van Th\'e--Tsankov: a Polish group whose universal minimal flow is
metrizable with a $G_\delta$ orbit admits such a subgroup \cite[Theorem 1.2]{MNT2016}.


\subsubsection{Norm transport and finite means}

For a bounded metric $d$, write 
\begin{equation*}
\mathsf{W}_1^d(\mu,\nu)= \inf_{\gamma\in\mathrm{Cpl}(\mu,\nu)} \int
d(x,y)\,d\gamma(x,y).
\end{equation*}
We omit $d$ when it is the operator-norm metric.

\begin{lemma}\label{lem:cu-transport}
 Let $G$ have a bounded compatible bi-invariant metric $d$. Suppose that for every finite $E\subset G$ and $\varepsilon>0$ there is
a finitely supported probability $\mu$ with 
\begin{equation*}
\max_{u\in E}\mathsf{W}_1^d(\mu,(R_u)_*\mu)<\varepsilon, \qquad R_u(x)=xu.
\end{equation*}
Then $G$ is amenable. A compactly supported law with this property may be
replaced by an equal-weight finite law, at arbitrarily small additional
error. For equal-weight laws, permutation couplings attain the minimum
transport cost.
\end{lemma}

\begin{proof}
For a bounded $L$-Lipschitz function $f$, a coupling gives 
\begin{equation*}
\left|\int f(xu)\,d\mu(x)-\int f(x)\,d\mu(x)\right| \leq L\mathsf{W}_1^d(\mu,(R_u)_*\mu).
\end{equation*}
Take a weak-* cluster point of the averaging functionals as $E$ increases
and the error decreases. Bounded Lipschitz functions are uniformly dense in
the bounded uniformly continuous functions on a metric space (apply infimal
convolution to real and imaginary parts). Thus the cluster point is a
right-invariant mean on that algebra. Inversion converts it to a
left-invariant mean; the two uniformities agree by bi-invariance.

For compact support, push the law to a finite $\delta$-net and then
approximate its weights by rationals. The transport error of this
replacement tends to zero. Right translations are isometries, so replacing a
law within transport distance $\tau$ increases each invariance error by at
most $2\tau$. Rational weights give equal weights after repetition. The
coupling polytope for two equal-weight labelled laws is the polytope of
doubly stochastic matrices. Its extreme points are permutation matrices,
which proves the last assertion.
\end{proof}

The full equivalence with finite matching criteria is due to Schneider--Thom 
\cite{SchneiderThom2018}. The elementary implication above is the one needed in the
construction.


\subsubsection{Two CP tools}

We use the following consequences of \cite{KirchbergPhillips2000}.

\begin{lemma}[Finite-target CP repair]\label{lem:cu-repair}
 Let $\mathcal{S}$ be an operator system in a unital C*-algebra $B$, let $F$ be a finite-dimensional C*-algebra, and let $L:\mathcal{S}\to F$ be unital, selfadjoint, and completely bounded. There is a
UCP map $T:B\to F$ such that 
\begin{equation*}
\left\|T|_{\mathcal{S}}-L\right\|_{\mathrm{cb}}\leq\left\|L\right\|_{\mathrm{cb}}-1.
\end{equation*}
\end{lemma}

\begin{proof}
Represent $F$ faithfully and unitally as a block-diagonal subalgebra of $M_d$. Apply \cite[Lemma 1.6]{KirchbergPhillips2000} with target $B(\mathbb{C}^d)=M_d$, and compose
with the UCP block-diagonal conditional expectation onto $F$.
\end{proof}

\begin{lemma}[Two-way CP comparison]\label{lem:cu-kp}
 Let $B_0,B_1$ be separable exact unital C*-algebras, with
unitaries $x_{ij}\in B_i$, $1\leq j\leq m$, and faithful unital embeddings $\iota_i:B_i\to\mathcal{O}_2$. If there are UCP maps $T_i:B_i\to\mathcal{O}_2$
satisfying 
\begin{equation*}
\left\|T_0(x_{0j})-\iota_1(x_{1j})\right\|\leq\eta,\qquad
\left\|T_1(x_{1j})-\iota_0(x_{0j})\right\|\leq\eta,
\end{equation*}
then for every $\tau>0$ there is $z\in U(\mathcal{O}_2)$ such that 
\begin{equation*}
\max_j\left\|z\iota_0(x_{0j})z^*-\iota_1(x_{1j})\right\| <11\sqrt\eta+\tau.
\end{equation*}
The listed unitaries need not generate $B_0$ or $B_1$.
\end{lemma}

\begin{proof}
Apply \cite[Proposition 1.13]{KirchbergPhillipsFields2000} to the discrete two-point continuous
field with fibers $B_0,B_1$, section algebra $B_0\oplus B_1$, and unitary
sections $(x_{0j},x_{1j})$. A two-point discrete space is compact metric,
and exactness of the section algebra is immediate from that of its fibers,
so the hypotheses of \cite[Lemma 1.10]{KirchbergPhillipsFields2000} hold. The directed
comparison quantity $\rho_0$ of \cite[Proposition 1.11]{KirchbergPhillipsFields2000} is an
infimum over unital completely positive maps and, by part~(1) of that
proposition, does not depend on the choice of the injective unital
homomorphisms; the maps $T_0,T_1$ therefore bound both of its values by $\eta $. The conclusion of \cite[Proposition 1.13]{KirchbergPhillipsFields2000} bounds an
infimum over $z\in U(\mathcal{O}_2)$ by $11\sqrt\eta$, whence the
arbitrarily small $\tau$. The proposition is stated for arbitrary continuous
unitary sections and imposes no generating assumption on them. The exponent $1/2$ cannot be improved \cite[Remark 3.11]{KirchbergPhillipsFields2000}.
\end{proof}

We also use the embedding and uniqueness theorems of \cite{KirchbergPhillips2000}: every
separable exact unital algebra embeds unitally in $\mathcal{O}_2$, and any
two such embeddings are approximately unitarily equivalent. All the
assertions with ambient target $\mathcal{O}_2$ also hold in a matrix algebra
over $\mathcal{O}_2$, since $M_k(\mathcal{O}_2)\cong\mathcal{O}_2$. Lemma~\ref{lem:cu-repair}, however, is used with finite-dimensional target; no
CP-extension property of $\mathcal{O}_2$ is being assumed.


\subsection{Boundary factorizations and a complete inverse estimate}

\label{sec:cu-boundary}

Let $\Gamma=\mathbb{F}_2=\langle g_1,g_2\rangle$, let $X=\partial\Gamma$ be
its boundary of infinite reduced words, and put 
\begin{equation*}
\mathcal{E}=C(X)\rtimes_r\Gamma.
\end{equation*}
Write $\lambda_g$ for the canonical unitary and $\alpha_g(f)(\xi)=f(g^{-1}\xi)$. The boundary crossed-product description is classical; see, for
example, Robertson's preprint \emph{Boundary crossed products from the free group} (1990), as cited in~\cite{LupiniCuntz2026}. We give the approximation maps explicitly,
including a proof of the nuclearity needed below.


\subsubsection{Prefix averages}

For $n\geq1$, set $S_n=\{x\in\Gamma:|x|<n\}$ and $r=|S_n|$. Define $m_x\in
C(X)$ by setting $m_x(\xi)=1/n$ if $x$ is one of the first $n$ vertices of
the ray from $1$ to $\xi$, and zero otherwise. The identity vertex is
included. These are locally constant functions and 
\begin{equation}  \label{eq:cu-prefix}
\sum_{x\in S_n}m_x=1,\qquad 0\leq m_x\leq 1/n.
\end{equation}
Define a column over $\mathcal{E}$ by 
\begin{equation*}
a=(a_x)_{x\in S_n},\qquad a_x=\lambda_x^*\sqrt{m_x}.
\end{equation*}
Thus $a^*a=1$. Compression of the regular representation on the Hilbert $C(X) $-module $\ell^2(\Gamma,C(X))$ gives a UCP map 
\begin{equation*}
\theta_n:\mathcal{E}\longrightarrow M_r(C(X)),\qquad
\theta_n(f\lambda_h)_{x,y} =\mathbf{1}_{\{x=hy\}}\alpha_{x^{-1}}(f).
\end{equation*}
The return map $Y\mapsto a^*Ya$ is also UCP.

\begin{lemma}\label{lem:cu-boundary-recovery}
 For $f\in C(X)$ and $h\in\Gamma$, 
\begin{equation*}
a^*\theta_n(f\lambda_h)a=fc_{n,h}\lambda_h,\qquad c_{n,h}=\sum_x\sqrt{m_x\alpha_h(m_{h^{-1}x})},
\end{equation*}
where missing indices have coefficient zero. Moreover, 
\begin{equation*}
0\leq c_{n,h}\leq1,\qquad \left\|1-c_{n,h}\right\|\leq\min\{|h|/n,1\}.
\end{equation*}
Consequently $a^*\theta_n(y)a\to y$ in norm for every $y\in\mathcal{E}$. The
algebra $\mathcal{E}$ is separable, unital and nuclear.
\end{lemma}

\begin{proof}
Matrix multiplication and the covariance relation give the displayed
formula. At $\xi$, the function $c_{n,h}$ is $1/n$ times the size of the
intersection of the first $n$ vertices of the rays from $1$ and from $h$
toward $\xi$. In a tree their common terminal ray starts after at most $|h|$
vertices have been lost from either initial segment. The intersection
therefore has at least $\max\{n-|h|,0\}$ vertices. The upper bound also
follows from Cauchy--Schwarz and \eqref{eq:cu-prefix}. Norm convergence follows
on the algebraic crossed product, and then on $\mathcal{E}$ by contractivity.

For a finite set of inputs and a fixed $n$, a sufficiently fine finite
clopen partition of $X$ gives a finite-dimensional subalgebra $B_n\subset
C(X)$ and a UCP map $J_n:M_r(C(X))\to M_r(B_n)$ approximating their images
under $\theta_n$. Namely, evaluate at one point of each partition atom and
extend constantly on that atom. The maps 
\begin{equation*}
\mathcal{E}\xrightarrow{J_n\theta_n}M_r(B_n) \xrightarrow{Y\mapsto a^*Ya}\mathcal{E }
\end{equation*}
give finite-dimensional completely positive approximations to the identity.
This proves nuclearity. Separability and unitality follow from the
construction.
\end{proof}

We will use embeddings in both directions between $\mathcal{E}$ and $\mathcal{O}_2$. The embedding $\mathcal{E}\hookrightarrow\mathcal{O}_2$
follows from Kirchberg--Phillips \cite{KirchbergPhillips2000}. The other one has a simple
formula. Let $P_x$, for $x=g_1,g_1^{-1},g_2,g_2^{-1}$, be the first-letter
projections, and set $S_x=\lambda_x(1-P_{x^{-1}})$. Then 
\begin{equation*}
S_x^*S_x=1-P_{x^{-1}},\qquad S_xS_x^*=P_x.
\end{equation*}
The elements 
\begin{equation}  \label{eq:cu-cuntz-boundary}
t_1=S_{g_1}+P_{g_1^{-1}},\qquad t_2=S_{g_2}+P_{g_2^{-1}}
\end{equation}
are isometries. Their final projections are $P_{g_1}+P_{g_1^{-1}}$ and $P_{g_2}+P_{g_2^{-1}}$, respectively, and add to $1$. The cross terms vanish
because the initial and final pieces are orthogonal. Thus \eqref{eq:cu-cuntz-boundary} determines a unital embedding $j:\mathcal{O}_2\hookrightarrow\mathcal{E}$, injective by simplicity.


\subsubsection{Dilations and complementary unitaries}

Fix arbitrary $u_1,\ldots,u_m\in U(\mathcal{E})$. Choose the preceding
clopen approximations so that 
\begin{equation*}
t_{n,j}=J_n\theta_n(u_j^*),\qquad
\epsilon_n=\max_j\left\|a^*t_{n,j}a-u_j^*\right\|\longrightarrow0.
\end{equation*}
Each $t_{n,j}$ is a contraction. Its Halmos dilation is 
\begin{equation*}
Q_{n,j}^0= 
\begin{pmatrix}
t_{n,j} & (1-t_{n,j}t_{n,j}^*)^{1/2} \\ 
(1-t_{n,j}^*t_{n,j})^{1/2} & -t_{n,j}^*\end{pmatrix}\in U(M_{2r}(B_n)).
\end{equation*}
Put $b=(a,0)^t$. Then $b^*b=1$, and expansion of the square gives 
\begin{equation}  \label{eq:cu-column-error}
\left\|Q_{n,j}^0bu_j-b\right\|^2\leq2\epsilon_n.
\end{equation}
For $N=2r+1$, the matrix 
\begin{equation*}
V=V_n=\begin{pmatrix}
b & 1-bb^* \\ 
0 & b^*\end{pmatrix}
\in U(M_N(\mathcal{E}))
\end{equation*}
is unitary. Here its row blocks have sizes $2r,1$ and its column blocks have
sizes $1,2r$. In particular its first column is $\widetilde a=(b,0)^t$. Set 
\begin{equation*}
Q_j=Q_{n,j}^0\oplus1,\quad U_j=\mathrm{diag}(u_j,1_{N-1}),\quad
R_j=V^*Q_jVU_j,\quad p=e_{11},\quad q=1-p.
\end{equation*}
With $\delta_n=\sqrt{2\epsilon_n}$, equation \eqref{eq:cu-column-error} says $\left\|(R_j-1)p\right\|\leq\delta_n$.

\begin{lemma}\label{lem:cu-polar}
 If $\delta_n<1$, then $d_j=qR_jq$ is invertible in the $q$-corner. Its polar unitary $c_j$ satisfies 
\begin{equation*}
\left\|c_j-d_j\right\|\leq s_n:=1-\sqrt{1-\delta_n^2},\qquad
\left\|R_j-(p+c_j)\right\|\leq\beta_n:=\sqrt2\delta_n+s_n.
\end{equation*}
\end{lemma}

\begin{proof}
For any unitary $R$, the two operators $(R-1)p$ and $p(R-1)$ have the same
norm: the squares of these norms both equal $\left\|p(2-R-R^*)p\right\|$.
Thus both off-diagonal blocks of $R_j$ have norm at most $\delta_n$. It
follows that 
\begin{equation*}
\left\|q-d_j^*d_j\right\|,\ \left\|q-d_jd_j^*\right\|\leq\delta_n^2.
\end{equation*}
Functional calculus proves invertibility and the asserted polar bound.
Finally, 
\begin{equation*}
R_j-(p+d_j)=(R_j-1)p+pR_jq.
\end{equation*}
The two summands have orthogonal initial supports, and the norm of their sum
is at most the square root of the sum of their squared norms. Adding the
polar error proves the result.
\end{proof}


\subsubsection{The complete inverse}

Consider the finite-dimensional unital algebra and its UCP compression 
\begin{equation}  \label{eq:cu-compression}
F_n=M_{2r}(B_n)\oplus\mathbb{C}\subset M_N(C(X)),\qquad \Phi_n(Y)=qV^*YVq.
\end{equation}
Its range is in $qM_N(\mathcal{E})q$, whose unit is $q$. Since $U_jq=q$, 
\begin{equation}  \label{eq:cu-compression-tuple}
\Phi_n(Q_j)=d_j.
\end{equation}

\begin{theorem}[Complete inverse estimate]\label{thm:cu-inverse}
 For $n>4$, the map $\Phi_n$ is injective and 
\begin{equation*}
\left\|\Phi_n^{-1}:\Phi_n(F_n)\to F_n\right\|_{\mathrm{cb}} \leq
\gamma_n:=(1-2/\sqrt n)^{-1}.
\end{equation*}
The estimate is independent of the clopen partition, $r$, and the matrix
amplification level.
\end{theorem}

\begin{proof}
For $\xi\in X$, the regular covariant representation on $\ell^2(\Gamma)$ is 
\begin{equation*}
\pi_\xi(f)\delta_h=f(h\xi)\delta_h,\qquad
\pi_\xi(\lambda_x)\delta_h=\delta_{xh}.
\end{equation*}
Write $\Pi=\widetilde a\widetilde a^*$. Let $v\in\mathbb{C}^k\otimes\mathbb{C}^N$. The terms of $(1_k\otimes\widetilde a^*)(v\otimes\delta_h)$
corresponding to different $x\in S_n$ have different group coordinates $xh$.
Indeed $a_x^*=\sqrt{m_x}\lambda_x$. Each coefficient has modulus at most $n^{-1/2}$. Orthogonality and the fact that $\widetilde a$ is an isometry
imply 
\begin{equation}  \label{eq:cu-small-projection}
\left\|(1_k\otimes\Pi)(v\otimes\delta_h)\right\| \leq
n^{-1/2}\left\|v\right\|.
\end{equation}

Take $Y\in M_k(M_N(C(X)))$. Choose $\xi\in X$ and a unit vector $v$ with $\left\|Y(\xi)v\right\|=\left\|Y\right\|$; compactness of $X$ and finite
dimensionality permit this. Elements of the coefficient algebra preserve the
group coordinate, so 
\begin{equation*}
\pi_\xi^{(kN)}(Y)(v\otimes\delta_1)=Y(\xi)v\otimes\delta_1.
\end{equation*}
Suppressing matrix amplifications, put $e=1-\Pi$. The identity 
\begin{equation*}
eYe-Y=-\Pi Y-eY\Pi
\end{equation*}
and \eqref{eq:cu-small-projection}, applied to both $v$ and $Y(\xi)v$, give 
\begin{equation*}
\left\|(eYe-Y)(v\otimes\delta_1)\right\| \leq 2n^{-1/2}\left\|Y\right\|.
\end{equation*}
The norm in the crossed product dominates this represented norm. Therefore 
\begin{equation*}
\left\|eYe\right\|\geq(1-2/\sqrt n)\left\|Y\right\|.
\end{equation*}
Since $VqV^*=1-\Pi$, unitary conjugation identifies this compression with
the amplification of $\Phi_n$. Restriction to $M_k(F_n)$, for every $k$,
proves the theorem.
\end{proof}

The only input from the group is the pair of facts that the prefix weights
satisfy $m_x\leq1/n$ and that the translates $xh$, $x\in S_n$, are pairwise
distinct. Freeness enters exactly here: in the tree, the first $n$ vertices
of a ray are distinct group elements, so the column $\widetilde a$ spreads a
unit vector over $\left|S_n\right|$ orthogonal group coordinates with
weights $n^{-1/2}$. No other property of $\mathbb{F}_2$ is used.

The operator system $\Phi_n(F_n)$ has unit $q$, and its inverse map is
unital and selfadjoint. Lemma~\ref{lem:cu-repair} therefore gives a UCP map 
\begin{equation*}
\Psi_n:qM_N(\mathcal{E})q\longrightarrow F_n,\qquad \left\|\Psi_n\Phi_n-\mathrm{id}_{F_n}\right\|_{\mathrm{cb}}\leq\gamma_n-1.
\end{equation*}
Combining this with Lemma~\ref{lem:cu-polar} gives 
\begin{equation}  \label{eq:cu-reverse-comparison}
\left\|\Psi_n(c_j)-Q_j\right\|\leq s_n+\gamma_n-1.
\end{equation}
Thus the inverse estimate produces a reverse CP comparison for the actual
complementary tuple. It is stronger than CP approximation of the original
inputs alone.


\subsection{Simultaneous corrections and norm amenability}

\label{sec:cu-norm}

Let $D=\mathcal{O}_2\otimes\mathcal{K}$, and first suppose the inputs belong
to a copy of $\mathcal{E}$ in a corner $fDf\cong\mathcal{O}_2$. All matrices
from Section~\ref{sec:cu-boundary} are now taken over this corner. The
complementary corner 
\begin{equation*}
C=qM_N(fDf)q\cong\mathcal{O}_2
\end{equation*}
contains $c_1,\ldots,c_m$. Choose a unital embedding $\iota:F_n\to C$. To
apply Lemma~\ref{lem:cu-kp}, take 
\begin{equation*}
B_0=F_n,\qquad B_1=C^*(q,c_1,\ldots,c_m),
\end{equation*}
embed them in $C\cong\mathcal{O}_2$ by $\iota$ and the inclusion, and use
the UCP maps $\Phi_n:B_0\to C$ and $\iota\Psi_n|_{B_1}:B_1\to C$. Equations \eqref{eq:cu-compression-tuple} and \eqref{eq:cu-reverse-comparison} compare the
tuples $(Q_j)$ and $(c_j)$ in both directions, with error at most 
\begin{equation*}
\eta_n=s_n+\gamma_n-1\longrightarrow0.
\end{equation*}
Both source algebras are separable and exact. Lemma~\ref{lem:cu-kp} therefore
gives a unitary $z_n\in U(C)$ such that the representation 
\begin{equation*}
\pi_n=\mathrm{Ad}(z_n)\circ\iota:F_n\to C
\end{equation*}
satisfies, for any chosen sequence $\tau_n>0$ tending to zero, 
\begin{equation}  \label{eq:cu-one-representation}
\max_j\left\|c_j-\pi_n(Q_j)\right\|<11\sqrt{\eta_n}+\tau_n.
\end{equation}
In particular a single representation works for all inputs. It is defined on 
$F_n$, irrespective of whether the chosen tuple generates that algebra.

\begin{proposition}\label{prop:cu-haar-law}
 For each finite tuple of inputs in $U(\mathcal{E})$,
embedded in $fDf$, there are compactly supported probability laws on $U_1(D)$
whose norm transport costs under the corresponding normalised right
translations tend to zero.
\end{proposition}

\begin{proof}
\uses{lem:cu-polar, lem:cu-kp, lem:cu-repair, thm:cu-inverse}
Let $K_n=U(F_n)$, with Haar probability. For $P\in K_n$, put 
\begin{equation*}
\Theta_P=p+\pi_n(P)^*,\qquad w_P=PV_n\Theta_P.
\end{equation*}
These are unitaries in $M_N(fDf)$, with the first expression written in the $p,q$ decomposition. Each $\Theta_P$ commutes with each $U_j=\mathrm{diag}(u_j,1_{N-1})$. Direct multiplication yields 
\begin{equation}  \label{eq:cu-relative}
w_P^*w_{PQ_j}U_j=\Theta_P^*R_j\Theta_{PQ_j}.
\end{equation}
If $R_j$ is replaced by $p+\pi_n(Q_j)$, the right side is the identity,
since its complementary block becomes 
\begin{equation*}
\pi_n(P)\pi_n(Q_j)\pi_n(PQ_j)^*=q.
\end{equation*}
Hence Lemma~\ref{lem:cu-polar} and \eqref{eq:cu-one-representation} imply 
\begin{equation}  \label{eq:cu-haar-bound}
\sup_{P\in K_n}\max_j\left\|w_{PQ_j}U_j-w_P\right\| \leq\beta_n+11\sqrt{\eta_n}+\tau_n\longrightarrow0.
\end{equation}

Realise $M_N(fDf)$ as a finite matrix corner in $D$, with first diagonal
corner $f$. Extend its unitaries by the identity on the complement. Let $\mu_n$ be the push-forward of Haar probability under $P\mapsto w_P$. For
every $j$, the map 
\begin{equation*}
P\longmapsto(w_P,w_{PQ_j}U_j)
\end{equation*}
is a coupling of $\mu_n$ with its right translate by $U_j$. This uses right
invariance of Haar probability on the same compact group $K_n$. Equation \eqref{eq:cu-haar-bound} bounds the transport cost. The image is norm compact
because the parametrization is norm continuous.
\end{proof}

\begin{proof}[Proof of Theorem~\protect\ref{thm:cu-finite-main}]
\proves{thm:cu-finite-main}
\uses{lem:cu-transport, prop:cu-haar-law}
We transfer Proposition~\ref{prop:cu-haar-law} twice. First fix a finite
unitary tuple in a corner $fDf\cong\mathcal{O}_2$. Choose a unital embedding 
$\mathcal{E}\hookrightarrow fDf$, and compose it with $j:\mathcal{O}_2\hookrightarrow\mathcal{E}$ from \eqref{eq:cu-cuntz-boundary}. This composite
embedding and an isomorphism $\mathcal{O}_2\to fDf$ are approximately
unitarily equivalent by \cite[Theorem 1.13]{KirchbergPhillips2000}. Consequently the prescribed
tuple, after one common corner-unitary conjugation, is arbitrarily close to
a tuple in the chosen copy of $\mathcal{E}$. Apply Proposition~\ref{prop:cu-haar-law} and conjugate the resulting law back. The conjugating corner
unitary, extended by $1-f$, belongs to $U_1(D)$. Conjugation is an isometry,
and changing a translating unitary by less than $t$ changes its transport
cost by at most $t$.

Second, let $u_1,\ldots,u_m\in U_1(D)$ be arbitrary. Let $f_l=1_{\mathcal{O}_2}\otimes p_l$, where $p_l$ are increasing finite-rank projections forming
an approximate identity of $\mathcal{K}$. Since $u_j-1\in D$, 
\begin{equation*}
(1-f_l)+f_lu_jf_l\longrightarrow u_j
\end{equation*}
in norm. For sufficiently large $l$, every compression $f_lu_jf_l$ is
invertible in $f_lDf_l$; its polar unitary, extended by $1-f_l$, still
converges to $u_j$. The corner $f_lDf_l\cong M_l(\mathcal{O}_2)\cong\mathcal{O}_2$ is of the kind already treated. Disjoint rank-$l$ blocks in $\mathcal{K}$ supply as many equivalent diagonal corners as the construction requires.
The first transfer, followed by this truncation, gives laws with arbitrarily
small costs for all prescribed inputs.

Finally apply Lemma~\ref{lem:cu-transport}: compactly supported laws can be
approximated in norm transport by equal-weight finite laws. Right
translations are isometries, so an approximation within $t$ adds at most $2t$
to every cost. For equal-weight atomic laws, a minimum-cost coupling is a
permutation coupling by the doubly stochastic matrix theorem. All sampled
unitaries lie in the same finite corner used in the construction. This is
exactly the assertion of Theorem~\ref{thm:cu-finite-main}.
\end{proof}

\begin{corollary}\label{cor:cu-strong-diagonal}
 Let $D=\mathcal{O}_2\otimes\mathcal{K}$. For
every finite set $E\subset U(D^+)$ and every $\eta>0$, there are $z_1,\ldots,z_M\in U_1(D)$, supported on one finite matrix corner, such that 
\begin{equation*}
\Delta=\frac1M\sum_{t=1}^M z_t^*\otimes z_t \in D^+\widehat\otimes_\pi D^+
\end{equation*}
lies in $\mathrm{co}\{v^*\otimes v:v\in U_1(D)\}$ and satisfies 
\begin{equation*}
m(\Delta)=1,\qquad \left\|\Delta\right\|_\pi\leq1,\qquad \mathrm{flip}(\Delta)=\Delta,\qquad \max_{g\in E}\left\|g\Delta-\Delta g\right\|_\pi<\eta,
\end{equation*}
where $m(x\otimes y)=xy$ and the tensor product has its usual outer bimodule
actions. Consequently $D^+$ has a strong approximate diagonal with
finite-corner witnesses.
\end{corollary}

\begin{proof}
\uses{thm:cu-finite-main}
Write every $g\in E$ as $g=\zeta_g u_g$, where $\zeta_g\in\mathbb{T}$ and $u_g\in U_1(D)$. Apply Theorem~\ref{thm:cu-finite-main} to the finite tuple $(u_g^*)_{g\in E}$, with tolerance $\eta/2$. Thus there are $w_1,\ldots,w_l$,
supported on one finite corner, and permutations $\sigma_g$ such that 
\begin{equation*}
\frac1l\sum_{b=1}^l \left\|w_{\sigma_g(b)}u_g^*-w_b\right\|<\frac\eta2.
\end{equation*}
Index a new list, counting repetitions, by ordered pairs and put $z_{a,b}=w_a^*w_b$, so that $M=l^2$. It stays in the same corner and is
invariant under taking adjoints, since $z_{a,b}^*=z_{b,a}$. Hence the
element 
\begin{equation*}
\Delta=\frac1{l^2}\sum_{a,b=1}^l z_{a,b}^*\otimes z_{a,b}
\end{equation*}
is flip invariant; also $m(\Delta)=1$ and $\left\|\Delta\right\|_\pi\leq1$.

Fix $g\in E$, and reindex the sum by $\rho_g(a,b)=(a,\sigma_g(b))$. If 
\begin{equation*}
e_{a,b}=\left\|z_{\rho_g(a,b)}u_g^*-z_{a,b}\right\|,
\end{equation*}
then $e_{a,b}=\left\|w_{\sigma_g(b)}u_g^*-w_b\right\|$. The two factors in 
\begin{equation*}
u_g z_{\rho_g(a,b)}^*\otimes z_{\rho_g(a,b)} \quad\hbox{and}\quad
z_{a,b}^*\otimes z_{a,b}u_g
\end{equation*}
differ by at most $e_{a,b}$ each. The projective tensor estimate and the
displayed average therefore give 
\begin{equation*}
\left\|u_g\Delta-\Delta u_g\right\|_\pi \leq\frac2{l^2}\sum_{a,b=1}^l
e_{a,b}<\eta.
\end{equation*}
The scalar $\zeta_g$ does not change this bound. Every element of the unital 
C*-algebra $D^+$ is a linear combination of four unitaries. Directing
these elements by finite subsets and decreasing $\eta$ therefore gives
centrality on all of $D^+$, and hence the asserted strong approximate
diagonal.
\end{proof}

\begin{remark}\label{rem:cu-remark-1}
\uses{thm:cu-finite-main}
The qualitative conclusion also follows from norm amenability through \cite[Theorem 2.4]{AlekseevSchmidtThom2025}. The construction above retains the quantitative
centrality bound and the common finite-corner support furnished by Theorem~\ref{thm:cu-finite-main}.
\end{remark}

\begin{corollary}\label{cor:cu-w-amenable}
 The norm group $U_1(\mathcal{O}_2\otimes\mathcal{K})$
is amenable.
\end{corollary}

\begin{proof}
\uses{lem:cu-transport, thm:cu-finite-main}
Apply the transport criterion in Lemma~\ref{lem:cu-transport} to Theorem~\ref{thm:cu-finite-main}.
\end{proof}

\begin{remark}\label{rem:cu-remark-2}
\uses{prop:cu-haar-law}
The corrections above take place in the full $\mathcal{O}_2$ tail corner.
They are not required to remain in the boundary coefficient algebra. This
freedom is essential: the argument uses Kirchberg--Phillips uniqueness after
the reverse CP map has been constructed. Neither a column averaging
statement alone nor independent corrections for individual generators would
establish the common-law assertion \eqref{eq:cu-haar-bound}.
\end{remark}


\subsection{Extreme amenability and pure-state stabilisers}

\label{sec:cu-stabilizers}


\subsubsection{Whirliness of stable identity components}

We first record the concentration input in a form that also applies to the
general Kirchberg discussion.

\begin{lemma}\label{lem:cu-stable-whirly}
 Let $B$ be a unital C*-algebra such that $B\otimes\mathcal{K}$ has real rank zero. The identity component $U_1(B\otimes \mathcal{K})_0$, with its norm topology, is whirly.
\end{lemma}

\begin{proof}
\uses{lem:cu-permanence}
The Hilbert--Schmidt unitary group 
\begin{equation*}
U_2=\{u\in U(\ell^2):u-1\text{ is Hilbert--Schmidt}\}
\end{equation*}
with its unnormalized Hilbert--Schmidt metric is a L\'evy group \cite[Example 2.15]{GiordanoPestov2007}. Its invariant probabilities on compact flows are
supported on fixed points, by \cite[Theorem 1.1 and Remark 1.4]{GlasnerTsirelsonWeiss2005}. The
inclusion into norm $U_1(\mathcal{K})$ is continuous and has dense image,
since finite-rank perturbations of the identity are dense. Thus norm $U_1(\mathcal{K})$ is whirly by Lemma~\ref{lem:cu-permanence}.

Every nonzero projection $p\in B\otimes\mathcal{K}$ can be the first
diagonal matrix unit of a copy of $\mathcal{K}$ in $B\otimes\mathcal{K}$. To
see this, compress $p$ to a sufficiently large canonical finite corner and
use functional calculus to obtain a projection $p^{\prime }$ with $\left\|p-p^{\prime }\right\|<1$. A unitary in the unitisation conjugates $p^{\prime }$ to $p$. Infinitely many disjoint copies of the finite corner
give equivalent copies of $p^{\prime }$, and partial isometries between them
give matrix units. Conjugate these back; the resulting matrix units remain
in the ideal $B\otimes\mathcal{K}$, since a unitary of its unitisation
normalises that ideal. In particular the projection circle 
\begin{equation*}
\{1+(z-1)p:z\in\mathbb{T}\}
\end{equation*}
is contained in a copy of $U_1(\mathcal{K})$.

By real rank zero, every selfadjoint element of $B\otimes\mathcal{K}$ is a
norm limit of finite-spectrum selfadjoint elements. After zero spectral
projections are omitted, their exponentials are products of
projection-circle elements. Such exponentials generate $U_1(B\otimes\mathcal{K})_0$. If a compact action of this group has an invariant probability, its
support is fixed by each of the copies of $U_1(\mathcal{K})$, then by the
group they generate, and finally by its closure. This proves whirliness.
\end{proof}

\begin{theorem}\label{thm:cu-w-ea}
 The norm group $W=U_1(\mathcal{O}_2\otimes\mathcal{K})$ is
extremely amenable.
\end{theorem}

\begin{proof}
\uses{cor:cu-w-amenable, lem:cu-stable-whirly}
The algebra $\mathcal{O}_2\otimes\mathcal{K}$ has real rank zero, and its $K_1$-group vanishes \cite[Theorem 1.2(ii)]{Zhang1992}, \cite{Cuntz1981}. For this
stabilised algebra, stability identifies $W/W_0$ with $K_1(\mathcal{O}_2\otimes\mathcal{K})$ \cite[8.1.1]{Blackadar1998}, so $W=W_0$. Lemma~\ref{lem:cu-stable-whirly} makes $W$ whirly. Corollary~\ref{cor:cu-w-amenable}
supplies an invariant probability in every compact $W$-flow, and its
nonempty support consists of fixed points. Thus $W$ is extremely amenable.
\end{proof}


\subsubsection{The hereditary algebra of a pure state}

Fix a pure state $\varphi$ on $A=\mathcal{O}_2$, with GNS triple $(\pi_\varphi,\mathcal{H}_\varphi,\xi_\varphi)$, and define 
\begin{equation*}
D_\varphi=\{a\in A: \varphi(a^*a)=\varphi(aa^*)=0\}.
\end{equation*}
This is the hereditary subalgebra with open support $1-p_\varphi$, where $p_\varphi$ is the minimal support projection of $\varphi$ in $A^{**}$.

\begin{lemma}\label{lem:cu-hereditary}
 The algebra $D_\varphi$ is isomorphic to $\mathcal{O}_2\otimes\mathcal{K}$. Under its unitisation embedding in $A$, 
\begin{equation*}
U_1(D_\varphi) =\{u\in U(A):\pi_\varphi(u)\xi_\varphi=\xi_\varphi\}.
\end{equation*}
\end{lemma}

\begin{proof}
The hereditary algebra is nonzero: otherwise its support would vanish and $1=p_\varphi$ would be minimal, forcing $A=\mathbb{C}$. It is full by
simplicity. It is nonunital, since a unit for it would equal $1-p_\varphi$
and would put the nonzero minimal projection $p_\varphi$ in $A$, impossible
in a purely infinite simple algebra. It is separable. Zhang's stability
theorem \cite[Theorem 1.2(i)]{Zhang1992} makes this nonunital hereditary algebra
stable, and Brown's full hereditary subalgebra theorem \cite[Theorem 2.8]{Brown1977} identifies its stabilisation with $A\otimes\mathcal{K}$. Thus $D_\varphi\cong A\otimes\mathcal{K}$.

For a unitary $u$, the condition $\pi_\varphi(u)\xi_\varphi =\xi_\varphi$
implies the same condition for $u^*$. It is therefore equivalent to $u-1\in
D_\varphi$, directly from the definition.
\end{proof}

\begin{proposition}\label{prop:cu-stabilizer-density}
 The group $H_\varphi=\mathrm{Aut}(A,\varphi)$
is a closed Polish subgroup of $\mathrm{Aut}(A)$, and the continuous
homomorphism 
\begin{equation*}
\mathrm{Ad}:U_1(D_\varphi)\longrightarrow H_\varphi
\end{equation*}
has dense image.
\end{proposition}

\begin{proof}
Closedness follows from continuity of $\alpha\mapsto\varphi\circ\alpha$ for
the point-norm and weak-* topologies. Thus $H_\varphi$ is Polish. Let $\alpha\in H_\varphi$, and prescribe finite point-norm tests and an error. By 
$\mathcal{O}_2$-uniqueness \cite[Theorem 1.13]{KirchbergPhillips2000}, approximate $\alpha$ on
these tests and on an additional finite state-correction test by $\gamma=\mathrm{Ad} u$, $u\in U(A)$. The state $\psi=\varphi\gamma$ is unitarily
equivalent to $\varphi$ and sufficiently close to it on the additional test,
since $\varphi\alpha=\varphi$.

Property 1.2 of \cite{KishimotoOzawaSakai2003} gives a continuous path of unitaries from the
identity implementing $\psi=\varphi\mathrm{Ad} v$ with $\mathrm{Ad} v$ as
close to the identity on the original tests as prescribed. Applicability
uses that an irreducible representation of $A$ contains no nonzero compact
operators. Nuclearity of $A$ is not required for Property~1.2: \cite[Theorem~4.3]{KishimotoOzawaSakai2003} proves Property~1.3 for every C*-algebra, and \cite[Theorem~3.1]{KishimotoOzawaSakai2003} upgrades Property~1.3 to Property~1.2. Put $z=uv^*$. Then $\varphi\mathrm{Ad} z=\varphi$, and 
\begin{equation*}
\left\|\mathrm{Ad} z(a)-\alpha(a)\right\| \leq\left\|\mathrm{Ad}
v^*(a)-a\right\|+\left\|\gamma(a)-\alpha(a)\right\|.
\end{equation*}
Both terms are small. Equality of the two pure vector states implies $\pi_\varphi(z)\xi_\varphi=\lambda\xi_\varphi$ for some $\lambda\in\mathbb{T}$. Replacing $z$ by $\lambda^{-1}z$ does not change its adjoint action and
puts it in $U_1(D_\varphi)$. Finally, norm continuity of the adjoint map
follows from $\left\|\mathrm{Ad} u(a)-\mathrm{Ad} v(a)\right\|\leq2\left\|a\right\|\left\|u-v\right\|$.
\end{proof}

\begin{corollary}\label{cor:cu-h-ea}
 Every pure-state stabiliser $H_\varphi$ in $\mathrm{Aut}(\mathcal{O}_2)$ is extremely amenable.
\end{corollary}

\begin{proof}
\uses{lem:cu-hereditary, prop:cu-stabilizer-density, thm:cu-w-ea}
Apply Theorem~\ref{thm:cu-w-ea}, Lemma~\ref{lem:cu-hereditary}, Proposition~\ref{prop:cu-stabilizer-density}, and permanence under continuous homomorphisms
with dense image.
\end{proof}

\begin{corollary}\label{cor:cu-scalar-umf}
 The universal minimal flow of norm $U((\mathcal{O}_2\otimes\mathcal{K})^+)$ is $\mathbb{T}$, acted on through the scalar
quotient.
\end{corollary}

\begin{proof}
\uses{lem:cu-kernel}
The scalar character $\chi:D^+\to\mathbb{C}$ gives the topological direct
product 
\begin{equation*}
U(D^+)\longrightarrow\mathbb{T}\times W,\qquad u\longmapsto\bigl(\chi(u),\overline{\chi(u)}u\bigr).
\end{equation*}
The normal subgroup $W$ is extremely amenable. Apply Lemma~\ref{lem:cu-kernel};
the UMF of the compact quotient $\mathbb{T}$ is its translation action.
\end{proof}


\subsection{Uniform pure-state homogeneity and the state-space flow}

\label{sec:cu-state}

The local argument in this section is stated for a larger class of algebras.
Its purpose is to make the uniformity in the base state explicit. All state
neighbourhoods below are weak-* neighbourhoods.


\subsubsection{Approximately central UCP rows}

\begin{lemma}\label{lem:cu-central-row}
 Let $A$ be unital and nuclear. For every finite $F\subset A$ and $\tau>0$, there are $x_1,\ldots,x_n\in A$ such that 
\begin{equation*}
\sum_i x_ix_i^*=1,\qquad
\sup_{\left\|b\right\|\leq1}\left\|[a,\Phi(b)]\right\|<\tau\quad(a\in F),
\qquad \Phi(b)=\sum_i x_i b x_i^*.
\end{equation*}
The row depends only on $F,\tau$.
\end{lemma}

\begin{proof}
Apply \cite[Lemma 4.2]{KishimotoOzawaSakai2003} to the universal representation. Nuclearity
makes its weak closure injective. The lemma gives finite positive-form
tensors 
\begin{equation*}
T=\sum_i y_i\otimes y_i^*,\qquad \sum_i\left\|y_i\right\|^2\leq1,
\end{equation*}
with $\left\|aT-Ta\right\|_\pi$ arbitrarily small on $F$, and with $c=\sum_i
y_iy_i^*$ converging weakly to $1$. Here $\left\|\cdot\right\|_\pi$ is the
Banach projective tensor norm, with the usual outer bimodule actions.
Convexification of tails makes $c\to1$ in norm, preserving the positive form
and the centrality estimates: concatenate rows multiplied by square roots of
the convex weights. Also $0\leq c\leq1$. Choose $c$ invertible and put $d=c^{-1/2}$, $x_i=dy_i$.

For $\Psi(b)=\sum_i y_i b y_i^*$, the projective estimate gives $\left\|\mathrm{ad}(a)\circ\Psi\right\|\leq\left\|aT-Ta\right\|_\pi$. Moreover 
\begin{equation*}
\left\|\Phi-\Psi\right\|\leq(\left\|d\right\|+1)\left\|d-1\right\|.
\end{equation*}
Consequently 
\begin{equation*}
\left\|\mathrm{ad}(a)\circ\Phi\right\| \leq\left\|aT-Ta\right\|_\pi+
2\left\|a\right\|(\left\|d\right\|+1)\left\|d-1\right\|.
\end{equation*}
Both terms can be made arbitrarily small, and $\sum_i x_ix_i^*=dcd=1$.
\end{proof}

\begin{lemma}[Gram comparison]\label{lem:cu-gram}
 Let $X,Y:\mathbb{C}^n\to\mathcal{H}$ have orthogonal
ranges. If their Gram matrices $G_0=X^*X$, $G_1=Y^*Y$ differ entrywise by at
most $\delta$, there is a positive contraction $f$ on the sum of their
ranges such that, on each standard basis vector $e_i$, 
\begin{equation*}
\left\|fXe_i-(Xe_i-Ye_i)/2\right\|,\ \left\|fYe_i-(Ye_i-Xe_i)/2\right\|
\leq\tfrac12\sqrt{n\delta}.
\end{equation*}
\end{lemma}

\begin{proof}
Write $X=U\sqrt{G_0}$, $Y=V\sqrt{G_1}$ in polar form and put $Z=VU^*$. Since
the ranges are orthogonal, $Z+Z^*$ is a selfadjoint contraction on their
sum. Set $f=(1-Z-Z^*)/2$. The support identities give 
\begin{equation*}
ZX-Y=V(\sqrt{G_0}-\sqrt{G_1}),\qquad Z^*Y-X=U(\sqrt{G_1}-\sqrt{G_0}).
\end{equation*}
Also $Z^*X=ZY=0$. The square-root estimate $\left\|\sqrt{G_0}-\sqrt{G_1}\right\|\leq\left\|G_0-G_1\right\|^{1/2} \leq\sqrt{n\delta}$ proves the
claim, including singular Gram matrices.
\end{proof}

\begin{lemma}[Glimm approximation with avoidance]\label{lem:cu-glimm-avoid}
 Let $A$ be separable and unital, and let $\pi:A\to B(\mathcal{H})$ be a faithful irreducible representation with $\pi(A)\cap \mathcal{K}(\mathcal{H})=0$. Given a state $\theta$, a finite set $E\subset
A $, a finite-dimensional subspace $M\subset\mathcal{H}$, and $r>0$, there
is a unit vector $\zeta\in M^\perp$ such that 
\begin{equation*}
\bigl|\langle\pi(a)\zeta,\zeta\rangle-\theta(a)\bigr|<r \qquad(a\in E).
\end{equation*}
\end{lemma}

\begin{proof}
Glimm's vector-state approximation theorem \cite{Glimm1961} gives unit vectors $\zeta_j$ converging weakly to zero whose vector states converge weak-* to $\theta$. If $P_M$ is the projection onto $M$, then $\left\|P_M\zeta_j\right \|\to0$. Consequently the normalised vectors 
\begin{equation*}
\frac{(1-P_M)\zeta_j}{\left\|(1-P_M)\zeta_j\right\|}
\end{equation*}
belong to $M^\perp$, converge in norm to $\zeta_j$, and retain the required
finite weak-* approximation. Here one may use 
\begin{equation*}
\bigl|\langle\pi(a)\xi,\xi\rangle- \langle\pi(a)\eta,\eta\rangle\bigr| \leq
2\left\|a\right\|\left\|\xi-\eta\right\|
\end{equation*}
for unit vectors $\xi,\eta$.
\end{proof}


\subsubsection{Uniform local correction}

\begin{proposition}\label{prop:cu-approx-alignment}
 Let $A$ be simple, separable, unital, infinite
dimensional and nuclear. Given finite $F\subset A$ and $\varepsilon>0$,
there are finite $E\subset A$ and $\delta>0$ with the following property. If
pure states $\omega,\nu$ satisfy $\left|\omega(a)-\nu(a)\right|<\delta$ for $a\in E$, then for every finite $K\subset A$ and $t>0$ there is a unitary
path $(u_s)_{0\leq s\leq1}$, starting at $1$, such that 
\begin{equation*}
\sup_s\max_{a\in F}\left\|\mathrm{Ad} u_s(a)-a\right\|<\varepsilon,\qquad
\max_{a\in K}\left|\omega(\mathrm{Ad} u_1(a))-\nu(a)\right|<t.
\end{equation*}
If the GNS representations of the two pure states are unitarily equivalent, $u_1$ may be chosen so that $\omega\circ\mathrm{Ad} u_1=\nu$.
\end{proposition}

\begin{proof}
\uses{lem:cu-central-row, lem:cu-glimm-avoid, lem:cu-gram}
Let $c=\max\bigl(\{1\}\cup\{\left\|a\right\|:a\in F\}\bigr)$. It suffices to
prove the assertion for $c^{-1}F$ with tolerance $c^{-1}\varepsilon$;
relabeling, assume that $F$ is in the unit ball. Choose $\tau,\kappa>0$ such
that 
\begin{equation*}
\kappa<1/(16\pi),\qquad 2\pi\tau+40\pi\kappa<\varepsilon.
\end{equation*}
Take the row of Lemma~\ref{lem:cu-central-row} with commutator bound $\tau$,
let its length be $n$, and put 
\begin{equation*}
E=\{x_ix_j^*:1\leq i,j\leq n\},\qquad \delta_0=\kappa^2/(2n^2).
\end{equation*}
These choices precede the choice of either state.

First suppose that $\omega,\nu$ are represented by unit vectors $\xi,\eta$
in one irreducible representation $\pi$, and are $\delta_0$-close on $E$.
The representation is faithful and essential: a nonzero intersection of $\pi(A)$ with the compact operators would, by simplicity, make $A$
elementary. Lemma~\ref{lem:cu-glimm-avoid}, applied to the state of $\xi$,
gives a unit vector $\zeta$ whose state is $\delta_0$-close to that of $\xi$
on $E$, and which is orthogonal to 
\begin{equation*}
\mathrm{span}\{\pi(x_ix_j^*)\xi,\pi(x_ix_j^*)\eta:1\leq i,j\leq n\}.
\end{equation*}
The displayed span contains $\xi,\eta$, because $\sum_i x_ix_i^*=1$. The
row-adjoint spans belonging to $\zeta$ and to either $\xi$ or $\eta$ are
orthogonal. The corresponding Gram matrices differ entrywise by less than $2\delta_0$.

Apply Lemma~\ref{lem:cu-gram} to the tuples $(\pi(x_i^*)\xi)_i$ and $(\pi(x_i^*)\zeta)_i$. Let $M$ be their finite-dimensional summed span.
Extend the positive contraction furnished by Lemma~\ref{lem:cu-gram} by zero on 
$M^\perp$. Positive Kadison transitivity \cite[Theorem II.6.1.13]{Blackadar2006} gives $a\in A_+$ whose restriction to $M$ is that contraction.
The subspace $M$ reduces $\pi(a)$, so replacing $a$ by $h=\min(a,1)$ through
continuous functional calculus preserves its restriction and gives $0\leq
h\leq1$. Put $k=\Phi(h)$. Multiplication by the row $(\pi(x_i))_i$, whose
norm is one, shows 
\begin{equation*}
\left\|\pi(k)\xi-(\xi-\zeta)/2\right\|<\kappa,\qquad
\left\|\pi(k)\zeta-(\zeta-\xi)/2\right\|<\kappa.
\end{equation*}
Indeed the bound before enlarging it to $\kappa$ is $\sqrt n\,\sqrt{2n\delta_0}/2=\kappa/2$. Thus 
\begin{equation*}
\left\|\pi(k)(\xi+\zeta)\right\|<2\kappa,\qquad
\left\|(\pi(k)-1)(\xi-\zeta)\right\|<2\kappa.
\end{equation*}
The integral formula for the exponential gives 
\begin{equation*}
\begin{split}
\left\|(\pi(e^{i\pi k})-1)(\xi+\zeta)\right\|&<2\pi\kappa, \\
\left\|(\pi(e^{i\pi k})+1)(\xi-\zeta)\right\|&<2\pi\kappa.
\end{split}\end{equation*}
Taking their half-sum yields 
\begin{equation*}
\left\|\pi(e^{i\pi k})\xi-\zeta\right\|<2\pi\kappa.
\end{equation*}
The path $e^{is\pi k}$ moves each $a\in F$ by less than $\pi\tau$.

For completeness, two unit vectors at distance $d<1/4$ can be joined by a
unitary on their two-dimensional span whose selfadjoint logarithm has norm
at most $4d$. Extend the two-vector isometry to that span and use the
principal logarithm; the estimate also follows from the elementary plane
rotation formula. Norm-controlled selfadjoint Kadison transitivity \cite[Theorem II.6.1.13]{Blackadar2006} lifts the logarithm to an element of $A$ of
norm less than $5d$. Its exponential path therefore has norm distance from $1 $ less than $5d$, and moves each element of the unit ball by less than $10d<20\pi\kappa$. Apply this to correct $\pi(e^{i\pi k})\xi$ exactly to $\zeta$. The concatenated path ends at $U$ with $\pi(U)\xi=\zeta$ and moves $F $ by less than 
\begin{equation*}
B=\pi\tau+20\pi\kappa.
\end{equation*}
The same construction for $\eta,\zeta$ gives $V$ with $\pi(V)\eta=\zeta$ and
the same bound. Concatenate the reversed adjoint path $U_s^*$ with $U^*V_s$.
Its endpoint $w=U^*V$ satisfies $\pi(w^*)\xi=\eta$, so $\omega\mathrm{Ad}
w=\nu $, and its movement on $F$ is less than $2B<\varepsilon$.

For arbitrary pure states take $\delta=\delta_0/3$. In the essential
irreducible representation of $\omega$, Lemma~\ref{lem:cu-glimm-avoid} (with no
avoidance subspace) approximates $\nu$ on $E\cup K$ by a vector state $\nu^{\prime }$, with error smaller than $\min\{\delta_0/3,t\}$. The vector
state is pure and has GNS representation equivalent to that of $\omega$,
since $\pi$ is irreducible. Moreover $\omega,\nu^{\prime }$ are less than $2\delta_0/3$, hence $\delta_0$, apart on $E$. Apply the exact
equivalent-state result just proved.
\end{proof}

\begin{theorem}[Uniform exact alignment]\label{thm:cu-alignment}
\uses{prop:cu-approx-alignment}
 Under the hypotheses of Proposition~\ref{prop:cu-approx-alignment}, for every finite $F\subset A$ and $\varepsilon>0$
there are finite $E\subset A$ and $\delta>0$ such that any two pure states $\rho_0,\rho_1$ that are $\delta$-close on $E$ satisfy 
\begin{equation*}
\rho_1=\rho_0\circ v^{-1},\qquad \max_{a\in
F}\left\|v(a)-a\right\|<\varepsilon
\end{equation*}
for some $v\in\mathrm{Aut}(A)$.
\end{theorem}

\begin{proof}
\uses{prop:cu-approx-alignment}
We give the alternating construction, including inverse control. We write $\rho\gamma=\rho\circ\gamma$; thus $(\rho\mathrm{Ad} R)\mathrm{Ad} u=\rho \mathrm{Ad}(Ru)$. Let $(a_j)$ be dense in the unit ball and choose positive $(\eta_n)_{n\geq1}$ with $\sum_n\eta_n<\varepsilon/4$. Set $P=Q=1$ initially
and $L_0=F\cup\{a_1\}$. Let $(G_0,\delta_0)$ be the test from Proposition~\ref{prop:cu-approx-alignment} for $(L_0,\eta_1)$. This is the required initial
test, independent of the states. Assume the states pass it.

At stage $n\geq1$, update $Q$ if $n$ is odd and $P$ if $n$ is even. Denote
the product to be updated by $R$ and the other one by $S$. Before choosing
its correction, set 
\begin{equation*}
L_n=F\cup\mathrm{Ad}(S^*)(F)\cup \{a_j,\mathrm{Ad}(S^*)(a_j):j\leq n+1\}.
\end{equation*}
Choose the next test $(G_n,\delta_n)$ for control $(L_n,\eta_{n+1})$,
enlarge $G_n$ by $G_{n-1}$ and $\{a_1,\ldots,a_n\}$, and decrease $\delta_n$
below $2^{-n}$. The current states $\rho_0\mathrm{Ad} Q,\rho_1\mathrm{Ad} P$
are $\delta_{n-1}$-close on $G_{n-1}$. Apply the proposition with control $(L_{n-1},\eta_n)$, ordering the two states so that the active one is
changed, and request final closeness on $G_n$ within $\delta_n$. If $u_n$ is
the resulting unitary, replace $R$ by $Ru_n$.

There is no circular choice: $G_n$ is chosen before $u_n$. At stage $n\geq2$, the product $R$ to be updated was the fixed opposite product when $L_{n-1}$
was formed. Hence 
\begin{equation*}
F\cup\{a_j:j\leq n\}\cup \mathrm{Ad}(R^*)\bigl(F\cup\{a_j:j\leq n\}\bigr)\subset L_{n-1}.
\end{equation*}
For $n=1$, the same inclusion follows from $R=1$ and $L_0=F\cup\{a_1\}$.
Therefore, for every $a\in F\cup\{a_j:j\leq n\}$, 
\begin{equation*}
\begin{split}
\left\|\mathrm{Ad}(Ru_n)(a)-\mathrm{Ad} R(a)\right\| &=\left\|\mathrm{Ad}
u_n(a)-a\right\|<\eta_n, \\
\left\|\mathrm{Ad}((Ru_n)^*)(a)-\mathrm{Ad}(R^*)(a)\right\| &=\left\|\mathrm{Ad} u_n^*(\mathrm{Ad}(R^*)(a))-\mathrm{Ad}(R^*)(a)\right\|<\eta_n.
\end{split}\end{equation*}
Here $\left\|\mathrm{Ad} u_n^*(b)-b\right\|=\left\|\mathrm{Ad}
u_n(b)-b\right\|$.

Write $P_n,Q_n$ for the products after stage $n$. Summability and density
now imply point-norm convergence of their adjoint actions and of the adjoint
actions of their inverses. Write 
\begin{equation*}
\mathrm{Ad} P_n\to\alpha,\quad \mathrm{Ad}(P_n^*)\to\widetilde\alpha, \qquad 
\mathrm{Ad} Q_n\to\beta,\quad \mathrm{Ad}(Q_n^*)\to\widetilde\beta.
\end{equation*}
For example, for every $a\in A$, 
\begin{equation*}
\begin{split}
\left\|\alpha(\widetilde\alpha(a))-a\right\| &\leq
\left\|\alpha(\widetilde\alpha(a))- \mathrm{Ad} P_n(\widetilde\alpha(a))\right\| \\
&\quad+\left\|\widetilde\alpha(a)-\mathrm{Ad}(P_n^*)(a)\right\|\longrightarrow0.
\end{split}\end{equation*}
The other compositions are identical, so $\widetilde\alpha=\alpha^{-1}$ and $\widetilde\beta=\beta^{-1}$; in particular $\alpha,\beta$ are automorphisms.

For $n\geq j$, the stage-$n$ state test gives 
\begin{equation*}
\left|\rho_0\mathrm{Ad} Q_n(a_j)-\rho_1\mathrm{Ad} P_n(a_j)\right|
<\delta_n<2^{-n}.
\end{equation*}
Passing to the limit and then using density gives $\rho_0\beta=\rho_1\alpha$. Thus $v=\alpha\beta^{-1}$ has $\rho_0v^{-1}=\rho_1$. Every control set
contains $F$. Applying the isometry $\alpha^{-1}$ gives $\left\|v(a)-a\right \|=\left\|\beta^{-1}(a)-\alpha^{-1}(a)\right\|$, and $\left\|\alpha^{-1}(a)-a\right\|=\left\|\alpha(a)-a\right\|$, so 
\begin{equation*}
\left\|v(a)-a\right\|\leq\left\|\beta^{-1}(a)-a\right\|
+\left\|\alpha(a)-a\right\| \leq2\sum_n\eta_n<\varepsilon\quad(a\in F).
\end{equation*}
Scaling handles finite sets outside the unit ball.
\end{proof}


\subsubsection{Completion and minimality}

\begin{proposition}\label{prop:cu-completion}
\uses{def:coprecompact}
 Let $A$ be simple, separable, unital, infinite
dimensional and nuclear, let $\varphi$ be a pure state, and put $G=\mathrm{Aut}(A)$ and $H=H_\varphi=\mathrm{Aut}(A,\varphi)$. The orbit map 
\begin{equation*}
G/H\longrightarrow P(A),\qquad gH\longmapsto\varphi\circ g^{-1}
\end{equation*}
is a uniform isomorphism from the quotient right uniformity onto the weak-*
subspace uniformity. Its compact completion is $S(A)$. In particular $H$ is
co-precompact.
\end{proposition}

\begin{proof}
\uses{lem:cu-glimm-avoid, thm:cu-alignment}
Pure-state homogeneity \cite{KishimotoOzawaSakai2003} makes the map bijective. The quotient
right uniformity is generated by pairs $(ugH,gH)$ for $u$ in a neighbourhood
of $1$. Their image states have the form $(\rho u^{-1},\rho)$, and 
\begin{equation*}
\left|\rho(u^{-1}(a))-\rho(a)\right|
\leq\left\|u^{-1}(a)-a\right\|=\left\|u(a)-a\right\|.
\end{equation*}
This proves uniform continuity. Theorem~\ref{thm:cu-alignment} proves uniform
continuity of the inverse, with tests independent of $\rho$. Pure states are
weak-* dense by Lemma~\ref{lem:cu-glimm-avoid}, applied in any essential
irreducible representation. Since the state space is compact, it is the
completion.
\end{proof}

\begin{corollary}[State-space criterion]\label{cor:cu-state-criterion}
\uses{prop:cu-completion}
 Under the hypotheses of Proposition~\ref{prop:cu-completion}, if $H_\varphi$ is extremely amenable and the action of $\mathrm{Aut}(A)$ on $S(A)$ is minimal, then 
\begin{equation*}
M(\mathrm{Aut}(A))\cong S(A).
\end{equation*}
\end{corollary}

\begin{proof}
\uses{lem:cu-completion-criterion, prop:cu-completion}
Apply Lemma~\ref{lem:cu-completion-criterion} to Proposition~\ref{prop:cu-completion}.
\end{proof}

\begin{proposition}\label{prop:cu-minimal}
 The action of $\mathrm{Aut}(\mathcal{O}_2)$ on $S(\mathcal{O}_2)$ is minimal and strongly proximal. Its pure-state orbit is
comeager.
\end{proposition}

\begin{proof}
Fix a source state $\rho$, a target pure state $\sigma$, and a finite set $F$
in the unit ball. Pure-state excision \cite[Proposition 2.2]{AkemannAndersonPedersen1986} first
gives a positive contraction $h$ of norm one, with tolerances chosen
sufficiently small, for which $\left\|hah-\sigma(a)h^2\right\|$ is small on $F$. For small $\gamma>0$, real rank zero \cite[Theorem 2.6(iii)]{BrownPedersen1991} gives a nonzero projection $p$ in the hereditary algebra
generated by $(h-(1-\gamma))_+$. Since $\left\|hp-p\right\|\leq\gamma$,
compression and the triangle inequality give 
\begin{equation*}
\left\|pap-\sigma(a)p\right\|<\epsilon\qquad(a\in F).
\end{equation*}
One may ensure properness by passing to a nonzero proper subprojection. For $L$ large, the $2^L$ Cuntz word-range projections sum to $1$; one of them, $z$, satisfies $\rho(z)\leq2^{-L}$. Put $q=1-z$. Nonzero projections in $\mathcal{O}_2$ are Murray--von Neumann equivalent \cite{Cuntz1981}. Applying
this to $p,q$ and to their complements gives a unitary $v$ with $vpv^*=q$.
The state $\rho^{\prime }=\rho\mathrm{Ad} v$ has $\rho^{\prime -L}$.
Cauchy--Schwarz, applied to the two off-diagonal corners of $a$, gives 
\begin{equation*}
\left|\rho^{\prime }(a)-\sigma(a)\right| \leq\epsilon+2\sqrt{2^{-L}}+2\,2^{-L}\qquad(a\in F).
\end{equation*}
Thus every orbit approaches every pure state. Pure-state density proves
minimality.

For strong proximality, let $\mu$ be a probability measure on $S(\mathcal{O}_2)$ and let $\bar\mu$ be its barycenter. By minimality, a net of translates
of $\bar\mu$ converges to a pure state $\sigma$. A cluster point of the
corresponding translated measures has barycenter $\sigma$; extremality
forces that measure to be $\delta_\sigma$. Finally the extreme points of a
metrizable compact convex set form a $G_\delta$. Here they are dense and
form one orbit.
\end{proof}

\begin{proof}[Completion of Theorem~\protect\ref{thm:cu-main}]
\proves{thm:cu-main}
\uses{cor:cu-h-ea, cor:cu-scalar-umf, lem:cu-completion-criterion, prop:cu-completion, prop:cu-minimal, thm:cu-w-ea}
Theorem~\ref{thm:cu-w-ea}, Corollary~\ref{cor:cu-h-ea}, and Corollary~\ref{cor:cu-scalar-umf} prove the first, second, and fourth assertions. For the
third, apply Lemma~\ref{lem:cu-completion-criterion}. Its hypotheses are
precisely Corollary~\ref{cor:cu-h-ea} and Propositions~\ref{prop:cu-completion}
and~\ref{prop:cu-minimal}.
\end{proof}

\begin{corollary}\label{cor:cu-furstenberg}
\uses{def:proximal}
 The state-space action is the universal minimal
strongly proximal flow, or Furstenberg boundary, of $\mathrm{Aut}(\mathcal{O}_2)$. Every minimal compact $\mathrm{Aut}(\mathcal{O}_2)$-flow is metrizable
and strongly proximal; in particular, $S(\mathcal{O}_2)$ is also the
universal minimal proximal flow.
\end{corollary}

\begin{proof}
\uses{prop:cu-minimal, thm:cu-main}
Theorem~\ref{thm:cu-main} identifies the universal minimal flow with $S(\mathcal{O}_2)$, and Proposition~\ref{prop:cu-minimal} makes this flow strongly
proximal. Every minimal flow is a factor of the universal minimal flow;
metrizability and strong proximality pass to factors.
\end{proof}

\begin{corollary}\label{cor:cu-radicals}
 The group $\mathrm{Aut}(\mathcal{O}_2)$ has trivial
Bohr compactification and no nontrivial closed normal amenable subgroup.
\end{corollary}

\begin{proof}
\uses{cor:cu-furstenberg}
A nontrivial compact group compactification of $\mathrm{Aut}(\mathcal{O}_2)$
would give a nontrivial minimal translation flow carrying Haar probability,
contradicting Corollary~\ref{cor:cu-furstenberg}. For the second assertion, let 
$N$ be closed, normal, and amenable. It fixes a probability $\mu$ on $S(\mathcal{O}_2)$. Normality makes every $G$-translate of $\mu$ again $N$-invariant. Strong proximality therefore gives an $N$-fixed state, and
normality plus minimality makes $N$ fix the whole state space. States
separate automorphisms, so $N$ is trivial.
\end{proof}

\begin{remark}\label{rem:cu-aut-nonamenable}
 The state-space flow has no invariant
probability. Such a probability would have an inner-invariant barycenter,
hence a tracial state on $\mathcal{O}_2$. The Cuntz relations would give $1=\tau(s_1s_1^*)+\tau(s_2s_2^*)=2$. Thus $\mathrm{Aut}(\mathcal{O}_2)$ is
nonamenable, although its UMF is metrizable.
\end{remark}


\subsection{Related Cuntz groups}

\label{sec:cu-related}

Throughout this section $A=\mathcal{O}_2$ and $D=A\otimes\mathcal{K}$. The
formulas specify the acting quotient as well as the compact space.


\subsubsection{Stable automorphisms and rigidity}

\begin{corollary}\label{cor:cu-stable-o2}
 Both $\mathrm{Aut}(D)$ and $\mathrm{Aut}(D^+)$, with
point-norm topology, are extremely amenable.
\end{corollary}

\begin{proof}
\uses{thm:cu-w-ea}
We show that $\mathrm{Ad} U_1(D)$ is dense in $\mathrm{Aut}(D)$. Let $\alpha\in\mathrm{Aut}(D)$. Approximate a given finite test set by a
canonical corner $pDp$. The projections $\alpha(p)$ and $p$ can be
conjugated by a unitary in $1+D$. Indeed approximate $\alpha(p)$ by a
projection in a finite canonical corner and first use the close-projection
unitary. Enlarge that corner to contain both projections with nonzero
complements. In this matrix algebra over $\mathcal{O}_2$, the projections
and their respective complements are equivalent; the implementing partial
isometries sum to a corner unitary. Extend it by the identity.

After this conjugation $\alpha$ preserves $pDp\cong\mathcal{O}_2$, where its
restriction is approximately inner by \cite{KirchbergPhillips2000}. Extending the resulting
corner unitary by $1-p$ approximates $\alpha$ on the original tests. Density
and Theorem~\ref{thm:cu-w-ea} imply extreme amenability. Every automorphism of $D^+$ preserves its unique nonzero proper ideal $D$, and every automorphism
of $D$ extends uniquely. Restriction is a point-norm topological isomorphism.
\end{proof}

\begin{proposition}\label{prop:cu-rigidity}
 The only $H_\varphi$-invariant state on $A$ is $\varphi $. In particular $N_{\mathrm{Aut}(A)}(H_\varphi)=H_\varphi$, and the
only continuous $\mathrm{Aut}(A)$-equivariant self-map of $S(A)$ is the
identity.
\end{proposition}

\begin{proof}
An invariant state $\rho$, restricted to $D_\varphi$, is invariant under all
its unitisation inner automorphisms, so its restriction is a bounded trace.
A stable algebra has no nonzero bounded trace: arbitrarily many orthogonal
equivalent copies of a positive element force its trace to vanish. Thus $\rho|_{D_\varphi}=0$. The support of $D_\varphi$ is $1-p_\varphi$, so the
normal extension of $\rho$ is supported on the minimal projection $p_\varphi$, forcing $\rho=\varphi$. If $g$ normalises $H_\varphi$, then $g\cdot\varphi$
is $H_\varphi$-fixed and hence equals $\varphi$. An equivariant self-map
likewise fixes $\varphi$, then the dense pure-state orbit, and therefore the
whole space.
\end{proof}

\begin{remark}\label{rem:cu-remark-3}
The norm groups $U(\mathcal{O}_2)$ and $\mathrm{PU}(\mathcal{O}_2)$ are
nonamenable, by the tracial-barycenter obstruction in Section~\ref{sec:cu-state}. Their UMFs are nonmetrizable: a Polish SIN group with metrizable UMF is
amenable \cite[Corollary 3.4]{BYMT2017}. These groups differ from both the
point-norm automorphism group and the stabilised unitisation group computed
above.
\end{remark}


\subsection{Stable UCT Kirchberg automorphism groups}

\label{sec:cu-kirchberg}

A Kirchberg algebra is nonzero, separable, nuclear, simple and purely
infinite. Here $B$ is stable and satisfies the universal coefficient theorem
(UCT). Write $\underline K(B)$ for total $K$-theory: the ordinary graded
groups, the groups with finite coefficients, and the coefficient and
Bockstein operations. Give 
\begin{equation*}
Q_B=\mathrm{Aut}_\Lambda(\underline K(B))
\end{equation*}
pointwise convergence on this countable discrete invariant. Preservation of
operations is closed in the corresponding permutation group, so $Q_B$ is
Polish. Coefficient information is necessary even for some finite $K$-groups.

\begin{theorem}\label{thm:cu-stable-general}
 The total-$K$-theory homomorphism is a continuous
open surjection 
\begin{equation*}
q_B:\mathrm{Aut}(B)\longrightarrow Q_B
\end{equation*}
whose kernel is extremely amenable. Consequently 
\begin{equation*}
M(\mathrm{Aut}(B))\cong q_B^*M(Q_B),
\end{equation*}
where $q_B^*$ denotes pullback of the action.
\end{theorem}

\begin{proof}
\uses{cor:cu-cuntz-stable, lem:cu-kernel, lem:cu-strict-ea}
Let 
\begin{equation*}
N_B=\overline{\{\mathrm{Ad} u:u\in U(M(B))\}}^{\,\text{point-norm}}.
\end{equation*}
Choose a nonzero projection $p\in B$. It is full, and Brown's theorem \cite[Theorem 2.8]{Brown1977} gives $B\cong B_0\otimes\mathcal{K}$ for the unital
Kirchberg algebra $B_0=pBp$. The adjoint map $\mathrm{Ad}\colon U(M(B))\to \mathrm{Aut}(B)$ is continuous for the strict topology on the source and the
point-norm topology on the target: strict convergence controls $(u_i-u)a$
and $(u_i-u)a^*$, which bound the two terms in $\mathrm{Ad} u_i(a)-\mathrm{Ad} u(a)$. Lemma~\ref{lem:cu-strict-ea} below makes $U(M(B))$ extremely amenable,
and dense-image permanence makes $N_B$ extremely amenable.

With the standard convention for stable algebras, approximate unitary
equivalence is implemented by multiplier unitaries. Thus the group denoted $\overline{\mathrm{Inn}}(B)$ in the classification literature is precisely $N_B$. Kirchberg--Phillips classification together with the universal
multicoefficient theorem \cite{KirchbergPhillips2000,Phillips2000,RosenbergSchochet1987,DadarlatLoring1996,Dadarlat2005} give the exact
sequence 
\begin{equation}  \label{eq:cu-classification-sequence}
1\longrightarrow N_B\longrightarrow\mathrm{Aut}(B) \xrightarrow{q_B}Q_B\longrightarrow1.
\end{equation}
Surjectivity can be read directly from those ingredients. The universal
multicoefficient theorem lifts any $\theta\in Q_B$ to $x\in KK(B,B)$, and
invertibility of $\theta$ on total $K$-theory makes $x$ invertible. Via $B\cong B_0\otimes\mathcal{K}$ and Morita invariance, regard $x$ as a class
in $KK(B_0,B_0)$. Stable Kirchberg classification \cite[Theorem 4.2.1]{Phillips2000} then realises the class by an isomorphism of stabilisations,
hence by an automorphism of $B_0\otimes\mathcal{K}\cong B$. The kernel
assertion in \eqref{eq:cu-classification-sequence} says exactly that trivial
total $K$-theory class is equivalent to approximate unitary equivalence to
the identity. Continuity can be checked on representatives: each coefficient 
$K$-class is represented by a projection or unitary in a fixed coefficient
algebra and amplification, and a sufficiently small perturbation leaves its
class unchanged. Dadarlat identifies the corresponding quotient topology on $KK/\overline 0$ with pointwise convergence on total $K$-theory \cite[Theorem
4.1(b)]{Dadarlat2005}. A continuous surjective homomorphism between Polish
groups is open. Apply Lemma~\ref{lem:cu-kernel} to this topological exact
sequence. The group $Q_B$ is discrete precisely when $\underline K(B)$ is
finitely generated. If in addition $Q_B$ is finite, then $M(\mathrm{Aut}(B))\cong Q_B$ as a translation flow; this is the case for the stabilised
Cuntz algebras in Corollary~\ref{cor:cu-cuntz-stable}. If $Q_B$ is infinite
discrete, $M(\mathrm{Aut}(B))$ is the pullback of the (typically
nonmetrizable) universal minimal flow of that discrete group, not $Q_B$
itself. For infinitely generated invariants, $Q_B$ need not be discrete and
the formula remains a pullback of $M(Q_B)$.
\end{proof}

\begin{lemma}[Strict multiplier unitaries]\label{lem:cu-strict-ea}
\uses{def:ea}
 Let $A$ be a unital separable simple nuclear C*-algebra and put $B=A\otimes\mathcal{K}$. The group $U(M(B))$, equipped with
the strict topology, is extremely amenable. In particular this applies when $B$ is a stable Kirchberg algebra.
\end{lemma}

\begin{proof}
The group is Polish. Nuclearity of $A$ is equivalent to the AF-property of $M(B)$ in the strict topology \cite[Theorem~3.3]{NgRuiz2012}: there is an
increasing sequence of finite-dimensional unital C*-subalgebras $F_n\subset M(B)$ whose union is strictly dense. A Kaplansky-type density
statement in the same paper \cite[Proposition~4.1(vi)]{NgRuiz2012} then makes $\bigcup_n U(F_n)$ strictly dense in $U(M(B))$.

Stability supplies a coherent amplification that keeps the sequence
increasing. Identify $B\cong M_n(B)$ and therefore $M(B)\cong M_n(M(B))$.
Transport $F_n$ to $F_n^{\prime }:=F_n\otimes 1_{M_n}$ under this
isomorphism. Every matrix summand of $F_n^{\prime }$ then has size at least $n$, and one may choose the identifications so that $F_n^{\prime }\subset
F_{n+1}^{\prime }$. Write $U(F_n^{\prime })$ again as $U(F_n)$. Write $F_n\cong\bigoplus_iM_{k_i}$, so that $U(F_n)=\prod_iU(k_i)$, and let $\left\|\cdot\right\|_{\mathrm{HS}}$ denote the unnormalized Hilbert--Schmidt
norm of this finite-dimensional algebra. For $b\in B$ the strict seminorm
obeys 
\begin{equation*}
\left\|(u-v)b\right\|\leq\left\|b\right\|\,\left\|u-v\right\|
\leq\left\|b\right\|\,\left\|u-v\right\|_{\mathrm{HS}},
\end{equation*}
so every strict entourage of $U(F_n)$ contains a Hilbert--Schmidt ball of
fixed radius, and concentration in a coarser metric is implied by
concentration in a finer one. On $U(k)$, Gromov--Milman concentration in the
unnormalized Hilbert--Schmidt metric \cite{GromovMilman1983} bounds the deviation of a $1$-Lipschitz function from its mean by $2e^{-ckt^2}$, and tensorization gives
the same bound on the product with $k=\min_ik_i$. The amplification above
forces $\min_ik_i\geq n$. Hence $\bigl(U(F_n)\bigr)_n$ is a L\'evy family
for the ambient strict entourages. The union $\bigcup_n U(F_n)$ is still
strictly dense, so $U(M(B))$ is a L\'evy group in the sense of
Giordano--Pestov \cite[Section~2]{GiordanoPestov2007} and is therefore extremely amenable.
\end{proof}

\begin{remark}\label{rem:cu-remark-4}
\uses{cor:cu-stable-o2, lem:cu-strict-ea, thm:cu-main}
The strict topology is essential. This argument concerns point-norm
automorphisms and does not prove amenability of norm $U(B^+)$. Ku\v{c}erovsk\'y--Ng \cite[Theorem~3.11]{KucerovskyNg2009} obtain fixed points on compact
metric spaces for the strict unitary group of a stable AH-algebra; that
statement is close to extreme amenability but is restricted to the AH case.
The original citation \cite{KucerovskyNg2010} established nuclearity criteria via
absorbing extensions and does not by itself supply Lemma~\ref{lem:cu-strict-ea}. Ng--Ruiz prove the corresponding structural statement for the projective
group, namely that $U(M(A\otimes\mathcal{K}))/\mathbb{T}$ is the closure of
an increasing union of compact subgroups \cite[Theorem 4.2(ii)]{NgRuiz2012}; the
concentration estimate is what is added here. The lemma is not used for
Theorem~\ref{thm:cu-main}: extreme amenability of $\mathrm{Aut}(D)$ is already
Corollary~\ref{cor:cu-stable-o2}. It is used only for the Kirchberg table in
this section.
\end{remark}


\subsubsection{Ordinary and total \(K\)-theory}

Let $G_i=K_i(B)$, and set 
\begin{equation*}
E_B=\mathrm{Ext}^1_\mathbb{Z}(G_0,G_1)\oplus\mathrm{Ext}^1_\mathbb{Z}(G_1,G_0),\qquad P_B=\mathrm{Pext}^1_\mathbb{Z}(G_0,G_1)\oplus\mathrm{Pext}^1_\mathbb{Z}(G_1,G_0).
\end{equation*}
Put $R_B=\mathrm{End}_\Lambda(\underline K(B))$. The UCT maps the Kasparov
endomorphism ring onto $\mathrm{End}(G_0)\times\mathrm{End}(G_1)$ with
kernel $E_B$, whereas the multicoefficient theorem maps the same ring onto $R_B$ with kernel $P_B\subset E_B$. Passing between these two quotients gives
the exact sequence of rings 
\begin{equation*}
0\longrightarrow E_B/P_B\longrightarrow R_B \longrightarrow\mathrm{End}(G_0)\times\mathrm{End}(G_1)\longrightarrow0.
\end{equation*}
The kernel is square-zero: products of filtration-one UCT classes land in $\mathrm{Ext}^2_\mathbb{Z}=0$. Units lift modulo a square-zero ideal, and the
units in the kernel are exactly $1+E_B/P_B$, with group law given by
addition. Passing to units, with $e\mapsto1+e$, therefore gives the group
exact sequence 
\begin{equation}  \label{eq:cu-total-extension}
1\longrightarrow(E_B/P_B,+)\longrightarrow Q_B \longrightarrow\mathrm{Aut}(G_0)\times\mathrm{Aut}(G_1)\longrightarrow1.
\end{equation}
Quotienting by pure extensions gives the total-$K$-theory ring \cite{RosenbergSchochet1987,DadarlatLoring1996,Dadarlat2005}. No splitting of the displayed sequence is asserted.
Equation \eqref{eq:cu-total-extension} is used algebraically; the topology on $Q_B$ remains pointwise convergence.

\begin{corollary}\label{cor:cu-k1zero}
\uses{thm:ellis}
 If $K_1(B)=0$, then, topologically, 
\begin{equation*}
Q_B\cong\mathrm{Aut}(K_0(B)),\qquad M(\mathrm{Aut}(B))\cong M(\mathrm{Aut}(K_0(B))).
\end{equation*}
The right-hand group carries pointwise convergence on $K_0(B)$. The degrees
may be interchanged.
\end{corollary}

\begin{proof}
Both cross-degree extension terms vanish. Restriction is a continuous
bijection of Polish groups and hence a topological isomorphism.
\end{proof}

\begin{corollary}\label{cor:cu-cuntz-stable}
\uses{thm:ellis}
 For $2\leq n<\infty$, 
\begin{equation*}
M(\mathrm{Aut}(\mathcal{O}_n\otimes\mathcal{K})) \cong(\mathbb{Z}/(n-1)\mathbb{Z})^\times.
\end{equation*}
For $\mathcal{O}_\infty$, the stable automorphism UMF is the two-point group 
$\{1,-1\}$. These are translation actions through $K_0$. For $n=2$ the
quotient is trivial.
\end{corollary}

\begin{proof}
Use $K_0(\mathcal{O}_n)=\mathbb{Z}/(n-1)\mathbb{Z}$, $K_0(\mathcal{O}_\infty)=\mathbb{Z}$, and $K_1=0$ \cite{Cuntz1981}. Their automorphism groups
are finite and equal to their own universal minimal flows. Table~\ref{eq:cu-cuntz} lists small cases.
\end{proof}

\begin{table}[ht]
\centering
\begin{tabular}{lll}
\hline Stable algebra $B$ & $K_0(B)$ & $M(\mathrm{Aut}(B))$ \\ 
\hline $\mathcal{O}_2\otimes\mathcal{K}$ & $0$ & singleton \\ 
$\mathcal{O}_3\otimes\mathcal{K}$ & $\mathbb{Z}/2$ & singleton \\ 
$\mathcal{O}_4\otimes\mathcal{K}$ & $\mathbb{Z}/3$ & $C_2$ \\ 
$\mathcal{O}_6\otimes\mathcal{K}$ & $\mathbb{Z}/5$ & $C_4$ \\ 
$\mathcal{O}_9\otimes\mathcal{K}$ & $\mathbb{Z}/8$ & $C_2\times C_2$ \\ 
$\mathcal{O}_\infty\otimes\mathcal{K}$ & $\mathbb{Z}$ & $C_2$ \\ 
\hline &  & 
\end{tabular}\caption{Stable Cuntz algebras and the universal minimal flows of their
automorphism groups, from Corollary~\protect\ref{cor:cu-k1zero}.}
\label{eq:cu-cuntz}
\end{table}

\begin{example}\label{ex:cu-example-1}
The range theorem \cite[Theorem 4.2.5]{Phillips2000}, followed by stabilisation,
realises arbitrary countable graded abelian $K$-groups by stable UCT
Kirchberg algebras. Let $p$ be prime and $d\geq1$. If $G_0=(\mathbb{Z}/p)^d$
and $G_1=0$, the UMF is the finite translation group $\mathrm{GL}_d(\mathbb{F}_p)$. If $G_0=\mathbb{Z}[1/p]$ and $G_1=0$, then 
\begin{equation*}
Q_B=\{\pm p^k:k\in\mathbb{Z}\}\cong C_2\times\mathbb{Z}.
\end{equation*}
It is discrete because fixing $1$ forces the identity. Its UMF is $C_2\times
M(\mathbb{Z})$, which is nonmetrizable.
\end{example}

\begin{example}\label{ex:cu-coefficient}
 If $G_0=G_1=\mathbb{Z}/2$, the ordinary graded
automorphism group is trivial. However $P_B=0$, since pure extensions of
finitely generated groups split, and $E_B=(\mathbb{Z}/2)^2$. Thus $Q_B=(\mathbb{Z}/2)^2$ and the automorphism UMF is its four-point translation
action. Replacing total $K$-theory by the ordinary graded invariant would
give a wrong answer.
\end{example}

\begin{proposition}\label{prop:cu-finite-k}
\uses{thm:ellis}
 Suppose $G_0,G_1$ are finitely generated. Then $Q_B$
is countable discrete, and finite exactly when both free ranks are at most
one. In the finite case, 
\begin{equation*}
|M(\mathrm{Aut}(B))| =|\mathrm{Aut}(G_0)|\,|\mathrm{Aut}(G_1)|\, |\mathrm{Ext}^1_\mathbb{Z}(G_0,G_1)|\,|\mathrm{Ext}^1_\mathbb{Z}(G_1,G_0)|.
\end{equation*}
If either free rank is at least two, $M(\mathrm{Aut}(B))$ is nonmetrizable.
\end{proposition}

\begin{proof}
Pure extensions vanish, and both extension groups are finite. Equation \eqref{eq:cu-total-extension} makes $Q_B$ countable. A countable Polish group
is discrete by Baire. For $\mathbb{Z}^r\oplus T$, with $T$ finite, the
automorphism group maps onto $\mathrm{GL}_r(\mathbb{Z})$ with finite kernel.
It is finite exactly when $r\leq1$, proving the criterion and the
cardinality formula. Otherwise $Q_B$ is infinite discrete, hence locally
compact and noncompact, and therefore has nonmetrizable universal minimal
flow by Veech's theorem together with \cite[Corollary 3.9 and the remark
following it]{BYMT2017}; for locally compact groups this is the theorem of
Kechris--Pestov--Todor\v{c}evi\'c \cite{KPT2005}.
\end{proof}

For general countable $K$-groups, Theorem~\ref{thm:cu-stable-general} is a
reduction to $Q_B$; it does not assert that every such group's UMF has a
known geometric model.

\begin{example}[A natural-order flow]\label{ex:cu-vector-space}
 For a prime $p$, let $V_p=\bigoplus_{j=1}^{\infty}\mathbb{Z}/p\mathbb{Z}$, viewed as the countable-dimensional vector space
over $\mathbb{F}_p$, and use the range theorem to choose a stable UCT
Kirchberg algebra $B_p$ with $K_0(B_p)=V_p$ and $K_1(B_p)=0$. Then 
\begin{equation*}
M(\mathrm{Aut}(B_p))\cong M(\mathrm{GL}(V_p)) \cong\mathrm{NO}(V_p),
\end{equation*}
where $\mathrm{NO}(V_p)$ is the space of those linear orders on $V_p$ whose
restriction to every finite-dimensional subspace is natural, that is,
antilexicographic with respect to some ordered basis. It is a closed, hence
compact, subspace of $2^{V_p\times V_p}$. The last identification is the
vector-space computation in the KPT correspondence \cite[Theorem 8.2(ii)]{KPT2005}. This gives an infinite, metrizable, nontrivial UMF for a stable
Kirchberg automorphism group.
\end{example}

\begin{example}[Compact $p$-adic flows]\label{ex:cu-prufer}
 For a prime $p$ and $d\geq1$, let $C_{p^\infty}$ be the
Pr\"ufer $p$-group and let $B_{p,d}$ be a stable UCT Kirchberg algebra with $K_0(B_{p,d})=C_{p^\infty}^{\,d}$ and $K_1(B_{p,d})=0$. Then 
\begin{equation*}
\mathrm{Aut}(K_0(B_{p,d}))\cong\mathrm{GL}_d(\mathbb{Z}_p)
\end{equation*}
with its compact $p$-adic topology: fixing the subgroup of elements of order
at most $p^n$ is congruence modulo $p^n$. Indeed, restriction to the $p^n$-torsion subgroups identifies $\mathrm{End}(C_{p^\infty})$ with $\varprojlim_n\mathbb{Z}/p^n\mathbb{Z}=\mathbb{Z}_p$; taking matrices and
units gives the displayed topological group. Hence 
\begin{equation*}
M(\mathrm{Aut}(B_{p,d}))\cong\mathrm{GL}_d(\mathbb{Z}_p)
\end{equation*}
as a translation flow. For $d=1$ this is $\mathbb{Z}_p^\times$.
\end{example}

\subsection{Proofs of statements in Section~\ref{sec:operator-algebras}}\label{sec:cu-encyclopaedia-proofs}

\begin{proof}[Proof of Theorem~\ref{thm:lupini-cuntz}]
\proves{thm:lupini-cuntz}
\uses{thm:cu-main}
This is Theorem~\ref{thm:cu-main}.
\end{proof}

\begin{proof}[Proof of Example~\ref{ex:lupini-stable-cuntz}]
\proves{ex:lupini-stable-cuntz}
\uses{cor:cu-k1zero, cor:cu-cuntz-stable}
By Corollary~\ref{cor:cu-k1zero} the quotient is $\mathrm{Aut}(K_0(A))$ when
$K_1(A)=0$, and Corollary~\ref{cor:cu-cuntz-stable} computes the universal
minimal flows: $(\mathbb{Z}/(n-1))^\times$ for $2\leq n<\infty$, and
$\{1,-1\}$ for $\mathcal{O}_\infty\otimes\mathcal{K}$. For $n=2$ and $n=3$ the
group $(\mathbb{Z}/(n-1))^\times$ is trivial.
\end{proof}

\section{Proofs: stable Kirchberg algebras and norm-unitary groups}\label{sec:proofs-stable-kirchberg}

This section reproduces, from the preprint of the editor \cite{LupiniStable2026} in its version of 15 September 2026, the
proofs of the results stated in Section~\ref{sec:operator-algebras}. The text is that of the
preprint, with British spelling: its sections appear as subsections, its macros
are expanded and its labels renamed, and of its introduction only the
statements of the main results are kept. Citations of the companion preprints
point to their own proof pages where the cited statement is reproduced there.
The preprint has not been refereed. The last subsection derives the statements
of Section~\ref{sec:operator-algebras} from those of the preprint.

\subsection{Main results of the preprint}\label{sec:st-main-results}

\begin{theorem}\label{thm:st-kk}
\uses{def:cu-definition-1}
 Let $A$ be a stable Kirchberg algebra, with no UCT
assumption, and put 
\begin{equation*}
\begin{aligned}
 \mathsf G_A&=\operatorname{Aut}(A), & J_A&=\{\alpha:KK(\alpha)=1\},\\
 \operatorname{Inn}_M(A)&=\{\operatorname{Ad}(u):u\in U(M(A))\}, &
 N_A&=\overline{\operatorname{Inn}_M(A)}.
 \end{aligned}
\end{equation*}
Then $J_A$ is the identity path component of $\mathsf G_A$, it is a Borel Polish
subgroup of $\mathsf G_A$, and $\overline{J_A}=N_A$. The natural map identifies $\mathsf G_A/J_A$ definably with $KK(A,A)^\times$ in its witnessed Cuntz
presentation. The homomorphism 
\begin{equation*}
\kappa_A:\mathsf G_A\longrightarrow R_A^\times,\qquad  \alpha\longmapsto[\alpha]_{KL}, 
\end{equation*}
is continuous, open, and surjective, with whirly amenable kernel $N_A$.
Thus $R_A^\times$ is the maximal Hausdorff quotient of $\mathsf G_A/J_A$ with its
quotient topology, and 
\begin{equation}  \label{eq:st-main-kl}
M(\mathsf G_A)\cong\kappa_A^*M(R_A^\times),\qquad b\mathsf G_A\cong bR_A^\times.
\end{equation}
Moreover, $\mathsf G_A$ is whirly amenable if and only if $R_A^\times$ is whirly
amenable.
\end{theorem}

\begin{theorem}\label{thm:st-nuclear-norm}
\uses{def:cu-definition-1}
If \(C\) is unital and nuclear, \(V_{C\otimes\mathcal K}\) is whirly amenable in
norm, and
\[
 M(U((C\otimes\mathcal K)^+))\cong\mathbb T\times M(K_1(C)).
\]
The action is through scalar evaluation and the stable component class.
For arbitrary unital \(C\), put \(B=C\otimes\mathcal K\). Nuclearity of \(B\) is equivalent to norm
amenability of \(U(B^+)\), and to amenability, extreme amenability, or
whirly amenability of \(U_1(B)_0\).
\end{theorem}

\begin{theorem}\label{thm:st-geometric}
\uses{cor:st-direct-sum-bohr-general, def:lo, def:maximal-chain}
 Let $B$ be a stable Kirchberg algebra such that $\operatorname{Aut} (B)$ is extremely amenable. For a nonempty finite set $I$, for $\mathbb{N}$,
and for the Cantor set $X$, respectively, 
\begin{align*}
M(\operatorname{Aut}(c_0(I,B)))&\cong\operatorname{Sym}(I), \\
M(\operatorname{Aut}(c_0(\mathbb{N},B)))&\cong\operatorname{LO}(\mathbb{N}), \\
M(\operatorname{Aut}(C(X)\otimes B))&\cong\Phi(X).
\end{align*}
Their Bohr compactifications are 
\begin{align*}
b\operatorname{Aut}(c_0(I,B))&\cong\operatorname{Sym}(I), \\
b\operatorname{Aut}(c_0(\mathbb{N},B))&=1, \\
b\operatorname{Aut}(C(X)\otimes B)&=1.
\end{align*}
The actions are induced by the permutations of summand ideals and the
homeomorphisms of $X$. In particular the formulas apply to both $B=\mathcal{O}_2\otimes\mathcal{K}$ and $B=\mathcal{O}_3\otimes\mathcal{K}$. Here $\operatorname{LO}(\mathbb{N})$ is the space of linear orders, and $\Phi(X)$ the space of
maximal chains of nonempty closed subsets. The computations also hold after
unitisation. The groups $\operatorname{Aut}(c_0(I,B))$, for finite nonempty $I$ and for $I=\mathbb{N}$, are amenable and their universal minimal flows are uniquely
ergodic; on $\operatorname{LO}(\mathbb{N})$ the unique measure is the uniform random-order
measure. In contrast, $\operatorname{Aut}(C(X)\otimes B)$ is nonamenable. Thus the two
infinite geometric groups are both minimally almost periodic, although only
the direct-sum group is amenable. The direct-sum Bohr conclusion also extends beyond this class: Corollary~\ref{cor:st-direct-sum-bohr-general}
computes it for every nonzero simple stable $C^*$-algebra $C$ and every index set $I
$. It is $(b\operatorname{Aut}(C))^I\rtimes\operatorname{Sym}(I)$ for finite $I$ and trivial for
infinite $I$, without separability, nuclearity, or amenability assumptions. Moreover, amenability of $\operatorname{Aut}(C)$ implies amenability of the
direct-sum automorphism group.
\end{theorem}

\begin{theorem}\label{thm:st-stable-stabilizer}
\uses{def:cu-definition-1}
 Let $A$ be a stable Kirchberg algebra,
without UCT, and let $\varphi$ be a pure state. Then 
\begin{equation*}
\operatorname{Aut}(A,\varphi)\cap N_A 
\end{equation*}
is whirly amenable. Restriction of the quotient map to $H_\varphi=\operatorname{Aut} (A,\varphi)$ is an open surjection onto $Q_A=\operatorname{Aut}(A)/N_A\cong KL(A,A)^\times$, and 
\begin{equation*}
M(H_\varphi)\cong(q_A|_{H_\varphi})^*M(Q_A),\qquad  bH_\varphi\cong bQ_A. 
\end{equation*}
In particular, restricting $M(\operatorname{Aut}(A))$ to $H_\varphi$ gives its universal
minimal flow.
The groups $H_\varphi$, $\operatorname{Aut}(A)$, and $Q_A$ are simultaneously whirly
amenable. Under UCT, this occurs exactly when $Q_A$ is trivial.
\end{theorem}

\subsection{Minimal flows and whirly amenable kernels}

\label{sec:st-dynamics} All compact spaces are nonempty and Hausdorff. A $G$-flow is a compact space with a jointly continuous $G$-action. It is minimal
if every orbit is dense. Every flow contains a minimal subflow, by
compactness and Zorn's lemma. The universal minimal flow exists and is
unique up to equivariant homeomorphism. A group is extremely amenable if
every compact flow has a fixed point, equivalently if its universal minimal
flow is a singleton.

Call a topological group \emph{whirly} if every invariant regular Borel
probability measure in every continuous action on a compact Hausdorff space
has support contained in the fixed-point set. It is \emph{whirly amenable}
if it is also amenable; the terminology is Schneider's
\cite[Corollary~5.7]{Schneider2019}. Whirly amenability implies extreme amenability,
since amenability provides an invariant probability measure and its
nonempty support consists of fixed points. Whirliness alone does not imply
amenability \cite{HollSchneider2026}.

We use the following \emph{dense-image principle}. If a continuous
homomorphism \(f:H\to G\) has dense image, then amenability, extreme
amenability and whirliness pass from \(H\) to \(G\). Indeed, restrict
any compact \(G\)-flow to \(H\). By continuity and density, every
\(H\)-fixed point is \(G\)-fixed, and every \(H\)-invariant regular
probability is \(G\)-invariant. These observations prove the first two
assertions. If \(H\) is whirly, a \(G\)-invariant regular probability
has support fixed by \(H\), hence by \(G\), proving the third.
In particular, whirly amenability passes to continuous dense images.

\begin{lemma}\label{lem:st-products}
\uses{def:cu-definition-1, def:ea}
 Finite products of extremely amenable groups are
extremely amenable. A group with a dense directed union of extremely
amenable subgroups is extremely amenable. Consequently arbitrary products,
with the product topology, preserve extreme amenability.
The same statements hold for whirly amenability.
\end{lemma}

\begin{proof}
For an $E\times F$-action on $X$, the nonempty closed set $\operatorname{Fix}_E(X)$ is $F$-invariant and contains an $F$-fixed point. Induction gives finite products.
For a directed family $(E_j)$, the nonempty closed sets $\operatorname{Fix}_{E_j}(X)$ have
the finite intersection property. A point in their intersection is fixed by
the dense union, hence by the whole group. Apply this to the subgroups
supported on finitely many coordinates of an arbitrary product.
For whirly amenability, the preceding argument supplies extreme amenability
and hence amenability. If $\mu$ is an invariant regular probability measure
on a compact action, its support is fixed by each whirly amenable factor or
subgroup. It is therefore fixed by the subgroup they generate and its
closure, proving whirliness in every case.
\end{proof}

A Polish group is a \emph{L\'evy group} if it has increasing compact
subgroups with dense union whose Haar probabilities form a concentrating
family for its group uniformity. Every Polish L\'evy group is whirly
amenable by \cite[Theorem 1.1 and Remark 1.4]{GlasnerTsirelsonWeiss2005}, and hence extremely
amenable. We will use the following finite-dimensional input.

The matrix concentration input is the classical Haar concentration
on \(U(n)\). Its Hilbert--Schmidt form is recorded in
Lemma~\ref{lem:st-unitary-levy}, and implies operator-norm concentration.

\begin{lemma}\label{lem:st-nonarch-whirly}
\uses{def:cu-definition-1}
 A Polish non-Archimedean group is whirly
amenable if and only if it is trivial.
\end{lemma}

\begin{proof}
A Polish non-Archimedean group $Q$ has a faithful continuous permutation
action on a countable discrete set $D$: take the disjoint union of the
countable coset spaces for a countable basis of open subgroups. Its induced
action on $\{0,1\}^D$ is jointly continuous and preserves the fair product
probability measure, which has full support. Whirliness would force every
configuration to be fixed, so the faithful action and $Q$ must be trivial.
The converse is immediate.
\end{proof}

\begin{proposition}\label{prop:st-transfer}
\uses{lem:cu-kernel, def:cu-definition-1}
 The minimal-flow assertion is Lemma~\ref{lem:cu-kernel}; the additional consequences are recorded together here. Let $q:G\to Q$ be a continuous open surjective
homomorphism, and suppose $N=\ker q$ is extremely amenable. Every minimal $G$-flow factors through $q$, and 
\begin{equation}  \label{eq:st-transfer}
M(G)\cong q^*M(Q).
\end{equation}
Consequently, $G$ is extremely amenable, respectively amenable, if and only
if $Q$ is, and $M(G)$ is metrizable if and only if $M(Q)$ is metrizable. The
universal minimal flow of $G$ carries a unique invariant probability measure
if and only if that of $Q$ does. Moreover, writing $bG$ for the Bohr
compactification, $q$ induces an isomorphism 
\begin{equation}  \label{eq:st-transfer-bohr}
bG\cong bQ.
\end{equation}
If, in addition, $N$ is whirly amenable, then $G$ is whirly amenable if and
only if $Q$ is whirly amenable.
\end{proposition}

\begin{proof}
For a minimal $G$-flow $X$, $\operatorname{Fix}_N(X)$ is nonempty and closed. Normality
makes it invariant: for $n\in N$, $g\in G$, and $x\in\operatorname{Fix}_N(X)$, $n(gx)=g(g^{-1}ng)x=gx$. Minimality implies $\operatorname{Fix}_N(X)=X$. The descended
action on $Q\times X$ is continuous because $q\times\mathrm{id}_X$ is an
open quotient map. Conversely a minimal $Q$-flow remains minimal after
pullback. Factor maps are unchanged, proving universality in \eqref{eq:st-transfer}. The assertions about extreme amenability and
metrizability are immediate. A group is amenable exactly when its universal
minimal flow carries an invariant regular probability measure: such a
measure maps onto a minimal subflow of any compact action. Under \eqref{eq:st-transfer}, invariance under $G$ is the same as invariance under $Q$, which also proves the unique-measure assertion.

Finally, let $f:G\to C$ be a continuous homomorphism into a compact group.
The closure of $f(N)$, with the left-translation action pulled back to $N$,
is a compact minimal $N$-flow. Extreme amenability forces it to be trivial.
Thus $f$ factors continuously through the open quotient $q$, and \eqref{eq:st-transfer-bohr} follows from the universal property of the Bohr
compactification.
For the additional assertion, the forward implication follows by pulling
compact actions and invariant measures back from $Q$; amenability also
passes to quotients. Conversely, suppose $N$ and $Q$ are whirly amenable.
The preceding amenability equivalence makes $G$ amenable. For an invariant
regular probability measure $\mu$ on a compact $G$-flow $X$, whirliness of
$N$ gives $\operatorname{supp}\mu\subseteq X^N$. This closed $G$-invariant subspace
carries a continuous $Q$-action, and $\mu$, viewed as a regular probability
on $X^N$, is $Q$-invariant. Whirliness of $Q$ therefore gives
$\operatorname{supp}\mu\subseteq (X^N)^Q=X^G$.
\end{proof}

\begin{lemma}\label{lem:st-amenable-extensions}
 Let $N$ be a closed normal subgroup of a
topological group $G$. If $N$ and $G/N$ are amenable, then $G$ is amenable.
\end{lemma}

\begin{proof}
Let $G$ act continuously on a compact space $X$. Amenability of $N$ makes
the compact convex set 
\begin{equation*}
\mathcal{P}(X)^N=\{\mu\in\mathcal{P}(X):n\mu=\mu\text{ for all }n\in N\} 
\end{equation*}
nonempty. Normality makes it $G$-invariant, and $N$ acts trivially on it, so
the continuous affine action factors through $G/N$. Amenability of $G/N$, in
its equivalent affine fixed-point form, gives a fixed point in $\mathcal{P}(X)^N$. This is a $G$-invariant probability measure on $X$.
\end{proof}

\begin{lemma}\label{lem:st-amenable-products}
 A topological group with a dense directed
union of amenable subgroups is amenable. Consequently arbitrary products of
amenable topological groups are amenable.
\end{lemma}

\begin{proof}
\uses{lem:st-amenable-extensions}
Let $(H_j)$ be directed with dense union in $G$, and let $G$ act
continuously on a compact space $X$. The nonempty closed subsets $\mathcal{P}(X)^{H_j}$ of $H_j$-invariant probability measures have the finite
intersection property. A measure in their intersection is fixed by the dense
union and hence by $G$.

Finite products of amenable groups are amenable by repeated application of
Lemma~\ref{lem:st-amenable-extensions}. In an arbitrary product, the
finite-coordinate subgroups form a dense directed union of such finite
products, so the first assertion applies.
\end{proof}

Call a topological group \emph{minimally almost periodic} if its Bohr
compactification is trivial. Equivalently, it has no nontrivial continuous
finite-dimensional unitary representation.

\begin{corollary}\label{cor:st-finite-dimensional-representations}
\uses{prop:st-transfer}
 Under the hypotheses of
Proposition~\ref{prop:st-transfer}, pullback along $q$ gives an equivalence
between the categories of continuous finite-dimensional unitary
representations of $Q$ and of $G$. In particular, 
\begin{equation*}
G\text{ is minimally almost periodic}  \quad\Longleftrightarrow\quad  Q\text{
is minimally almost periodic}. 
\end{equation*}
\end{corollary}

\begin{proof}
If $\pi:G\to U(n)$ is continuous, the closure of $\pi(N)$, with its
left-translation action pulled back to $N$, is a compact minimal $N$-flow.
Thus $\pi(N)=1$, and $\pi$ factors uniquely and continuously through the
open quotient $q$. The converse and the assertion about intertwiners follow
from surjectivity of $q$.
\end{proof}

\begin{lemma}\label{lem:st-minap-tools}
 The following statements hold.

\begin{enumerate}
\item An arbitrary product of minimally almost periodic groups is minimally
almost periodic.

\item Suppose that $H$ is a dense subgroup of $G$ and that, for every $h\in H
$, there are $g_n\in G$ such that $g_nhg_n^{-1}\to1$. Then $G$ is minimally
almost periodic.
\end{enumerate}
\end{lemma}

\begin{proof}
For (i), let $\pi:\prod_iG_i\to U(n)$ be continuous. Choose an identity
neighbourhood in $U(n)$ containing no nontrivial subgroup. The preimage
contains a basic product neighbourhood, so $\pi$ is trivial on every
coordinate group outside a finite set. It therefore factors through a finite
product, where it is trivial on each factor.

For (ii), unitary invariance of the operator norm gives 
\begin{equation*}
\|\pi(h)-1\|=\|\pi(g_nhg_n^{-1})-1\|\longrightarrow0 
\end{equation*}
for every continuous finite-dimensional unitary representation $\pi$. Thus $\pi$ is trivial on $H$, and hence on $G$. The Peter--Weyl theorem identifies
this with triviality of the Bohr compactification.
\end{proof}

\begin{lemma}\label{lem:st-bohr-finite-products}
 For finitely many topological groups $G_1,\ldots,G_n$, the canonical map induces 
\begin{equation*}
b(G_1\times\cdots\times G_n)\cong bG_1\times\cdots\times bG_n. 
\end{equation*}
\end{lemma}

\begin{proof}
It is enough to consider two factors. Given a continuous homomorphism $f:G_1\times G_2\to K$ into a compact group, the restrictions to the two
factors have commuting images. Their extensions to $bG_1$ and $bG_2$ still
have commuting images, by density, and their product is the unique
continuous extension of $f$ to $bG_1\times bG_2$. The universal property and
induction prove the assertion.
\end{proof}

\begin{corollary}\label{cor:st-wreath-minap}
 For every topological group $E$ and every infinite
set $I$, with the product topology on $E^I$ and the pointwise topology on $\operatorname{Sym}(I)$, 
\begin{equation*}
E^I\rtimes\operatorname{Sym}(I) 
\end{equation*}
is minimally almost periodic. Here $\operatorname{Sym}(I)$ permutes the coordinates of $E^I
$.
\end{corollary}

\begin{proof}
\uses{lem:st-minap-tools}
The subgroup $E^{(I)}\rtimes\operatorname{Sym}_{\mathrm{fin}}(I)$ is dense. Each of its
elements is supported on a finite set of coordinates. Conjugating by
permutations that move this support through pairwise disjoint finite subsets
of $I$ makes the conjugates converge to the identity in the product and
pointwise topologies. Apply Lemma~\ref{lem:st-minap-tools}(ii).
\end{proof}

The canonical quotient map by a closed normal subgroup is open. For Polish
groups, the quotient is again Polish.

For compact \(C\), minimal flows are homogeneous translation quotients,
and the compact-factor rule from Lemma~\ref{lem:cu-kernel} gives
\begin{equation}\label{eq:st-compact-product}
 M(C\times L)=C\times M(L).
\end{equation}
The Bohr counterpart follows from Lemma~\ref{lem:st-bohr-finite-products}:
\begin{equation}\label{eq:st-compact-product-bohr}
 b(C\times L)=C\times bL.
\end{equation}

\begin{lemma}\label{lem:st-locally-compact-umf}
\uses{thm:ellis}
 For a locally compact Polish group $L$, the
universal minimal flow $M(L)$ is metrizable if and only if $L$ is compact.
\end{lemma}

\begin{proof}
The action of a locally compact group on its universal minimal flow is free;
see the discussion following Corollary~3.9 in \cite{BYMT2017}. Thus, if $M(L)$
is metrizable, \cite[Corollary~3.9]{BYMT2017} implies that $L$ is compact. The
converse follows from the compact-group facts recalled above.
\end{proof}

\begin{corollary}\label{cor:st-compact-quotient-minimal}
\uses{prop:st-transfer}
 Under the hypotheses of Proposition~\ref{prop:st-transfer}, suppose in addition that $Q$ is compact. Then every minimal 
$G$-flow is, through $q$, a homogeneous space $Q/L$ for a closed subgroup $L\leq Q$. In particular, every minimal $G$-flow is equicontinuous and
uniquely ergodic, and 
\begin{equation*}
M(G)\cong Q,\qquad bG\cong Q. 
\end{equation*}
\end{corollary}

\begin{proof}
\uses{prop:st-transfer}
Proposition~\ref{prop:st-transfer} makes every minimal $G$-flow a minimal $Q$-flow. The compact-group facts recalled above give the dynamical properties and the
formula for $M(G)$; the Bohr formula follows from \eqref{eq:st-transfer-bohr}
and $bQ\cong Q$.
\end{proof}

We also use the Polish open mapping theorem: a continuous surjective
homomorphism between Polish groups is open \cite{BeckerKechris1996}. The
point-norm automorphism group of a separable $C^*$-algebra is Polish; a
complete compatible metric is obtained by testing automorphisms and their
inverses on a countable dense subset of its unit ball. Checking openness
will ensure that the algebraic quotients below carry precisely the required
topology.


\subsection{The approximately inner subgroup}

\label{sec:st-kernel} For nonunital $A$, write $M(A)$ for the multiplier
algebra and $A^+$ for the minimal unitisation. The strict topology on $M(A)$
is generated by 
\begin{equation*}
u\longmapsto\left\|ua\right\|+\left\|au\right\|\qquad(a\in A). 
\end{equation*}
Its restriction makes $U(M(A))$ a topological group. Set 
\begin{equation*}
\operatorname{Inn}_M(A)=\{\operatorname{Ad}(u):u\in U(M(A))\},\qquad  N_A=\overline{\operatorname{Inn}_M(A)}. 
\end{equation*}
This notation fixes the implementer convention for nonunital algebras.

The stable nuclear norm-unitary theorem proved in
Section~\ref{sec:st-unitaries} gives the multiplier conclusion through
strict approximation. Its proof uses stable uniqueness and Haar
concentration and does not depend on the present section.

\begin{theorem}[Strict multiplier unitaries]\label{thm:st-multiplier-whirly}
\uses{def:cu-definition-1, def:ea}
 If $B$ is unital, separable, simple, and
nuclear, put $A=B\otimes\mathcal K$. The inclusion
\[
 \bigl(U_1(A)_0\bigr)_{\|\cdot\|}\longrightarrow U(M(A))_{\mathrm{strict}}
\]
is continuous with dense image. The target is a Polish group and is
whirly amenable, hence extremely amenable.
\end{theorem}

\begin{proof}
\uses{thm:st-v}
The strict multiplier-unitary group is Polish because $A$ is separable,
and norm convergence implies strict convergence. Cuntz--Higson's Kuiper
theorem makes $U(M(A))$ norm contractible \cite{CuntzHigson1987}. Thus every
$u\in U(M(A))$ is a finite product
$u=\prod_{j=1}^k\exp(ih_j)$, with $h_j=h_j^*\in M(A)$.
Let $e_n=1_B\otimes\sum_{\ell=1}^n e_{\ell\ell}$ be the canonical
projection approximate identity of $A$. For each $j$,
$e_nh_je_n\to h_j$ strictly and $\|e_nh_je_n\|\leq\|h_j\|$.
Multiplication is jointly strictly continuous on bounded sets, so
polynomial approximation to the exponential gives
\[
 u_n:=\prod_{j=1}^k\exp(ie_nh_je_n)\longrightarrow u
 \qquad\text{strictly}.
\]
The path $t\mapsto\prod_{j=1}^k\exp(it e_nh_je_n)$ lies in $U_1(A)$
and joins $1$ to $u_n$. Hence $u_n\in U_1(A)_0$, proving density.
Theorem~\ref{thm:st-v} makes $U_1(A)_0$ whirly amenable in norm.
The dense-image principle now gives the assertion for the strict group.
\end{proof}

\begin{remark}\label{rem:st-remark-1}
\uses{thm:st-v}
The strict-density argument above uses stability and the canonical
projection approximate identity; nuclearity enters only through
Theorem~\ref{thm:st-v}. Ng--Ruiz establish strict AF approximation and
unitary density \cite[Theorem~3.3 and Proposition~4.1(vi)]{NgRuiz2012}.
The proof here obtains whirly amenability directly from the norm group.
Whether a concentrating family of compact subgroups realises the strict
multiplier group as a L\'evy group requires a separate construction.
\end{remark}

\begin{lemma}\label{lem:st-adjoint}
 The homomorphism $\operatorname{Ad}:U(M(A))_{\mathrm{strict}}\to\operatorname{Aut}(A)
$ is continuous, and $N_A$ is closed and normal.
\end{lemma}

\begin{proof}
For unitaries $u,v$ and $a\in A$, 
\begin{equation*}
uau^*-vav^*=(u-v)au^*+va(u^*-v^*), 
\end{equation*}
whence 
\begin{equation*}
\left\|uau^*-vav^*\right\| 
\leq\left\|(u-v)a\right\|+\left\|(u-v)a^*\right\|. 
\end{equation*}
This proves continuity. Every $\alpha\in\operatorname{Aut}(A)$ extends uniquely to $\widetilde\alpha\in\operatorname{Aut}(M(A))$, and $\alpha\operatorname{Ad}(u)\alpha^{-1}=\operatorname{Ad} (\widetilde\alpha(u))$. Thus the image and its closure are normal.
\end{proof}

\begin{theorem}\label{thm:st-kernel-ea}
 If $A\cong B\otimes\mathcal{K}$ for a unital separable
simple nuclear algebra $B$, then $N_A$ is whirly amenable. In particular
for every stable Kirchberg algebra, without a UCT assumption, $N_A$ is
whirly amenable, and
\begin{equation}  \label{eq:st-general-quotient}
M(\operatorname{Aut}(A))\cong M(\operatorname{Aut}(A)/N_A)
\end{equation}
with the quotient action.
\end{theorem}

\begin{proof}
\uses{lem:st-adjoint, prop:st-transfer, thm:st-multiplier-whirly}
Theorem~\ref{thm:st-multiplier-whirly} and Lemma~\ref{lem:st-adjoint} and the dense-image principle prove
the first assertion. A stable Kirchberg algebra contains a nonzero
projection $p$. Simplicity makes it full, and $B=pAp$ is unital, separable,
simple, and nuclear. Brown's full hereditary subalgebra theorem \cite[Theorem~2.8]{Brown1977}
gives $B\otimes\mathcal{K}\cong A\otimes\mathcal{K}\cong A$. Apply the first
assertion to obtain whirly amenability of $N_A$.

Finally,
Proposition~\ref{prop:st-transfer} gives the quotient-flow formula. A second
argument for whirliness of $N_A$ in the Kirchberg case is recorded
in Remark~\ref{rem:st-alternative-whirly} below.
\end{proof}

\begin{remark}\label{rem:st-no-norm}
\uses{thm:st-v}
 The adjoint map here has the strict topology on its
multiplier-unitary source and point-norm convergence on its automorphism
target. The norm-topology assertion used in its proof concerns the
scalar-one identity component $U_1(A)_0$, as stated in
Theorem~\ref{thm:st-v}. Passing to a finer group topology enlarges the class
of continuous compact actions.
\end{remark}


\subsection{Definable \(KK\)-units and the Hausdorff
quotient}

\label{sec:st-definable-kk} Throughout this section $A$ is a stable Kirchberg
algebra; no UCT assumption is made. We use the degree-zero Kasparov group
and the composition convention for its product, so that $\alpha\mapsto
KK(\alpha)$ is a homomorphism on automorphism groups. Put 
\begin{equation*}
\mathsf G_A=\operatorname{Aut}(A),\qquad  J_A=\{\alpha\in \mathsf G_A:KK(\alpha)=1\},\qquad 
R_A=KL(A,A)=KK(A,A)/\overline{\{0\}}. 
\end{equation*}
The closure of zero is taken in the natural topology on $KK(A,A)$. It is an
ideal by continuity of multiplication. Dadarlat proves that $R_A$ is a
Polish additive group and that Kasparov products are continuous
\cite[Proposition~2.8 and Theorem~3.5]{Dadarlat2005}. Point-norm convergence
of homomorphisms induces convergence of their classes: if
$\varphi_n\to\varphi$ point-norm, then
$a\mapsto(\varphi_n(a))_{n\in\overline{\mathbb N}}$, with $\varphi(a)$ in
the coordinate at infinity, is a $^*$-homomorphism into
$C(\overline{\mathbb N})\otimes B$, so \cite[Theorem~3.3]{Dadarlat2005} gives
$[\varphi_n]\to[\varphi]$.


\subsubsection{The topology and realisation of \(KL\)-units}

For a unital Polish topological ring $R$, endow $R^\times$ with the topology
of the inverse-pair space 
\begin{equation}  \label{eq:st-inverse-pair}
\{(x,y)\in R\times R:xy=yx=1\}, \qquad x\longleftrightarrow(x,x^{-1}).
\end{equation}
The displayed space is closed in $R^2$. Multiplication is $(x,y)(x^{\prime
},y^{\prime })=(xx^{\prime },y^{\prime }y)$, and inversion interchanges the
two coordinates. Thus $R^\times$ is a Polish group. We use this topology on $R_A^\times$, including when inversion has not separately been shown
continuous for the subspace topology from $R_A$.

We recall the classification inputs needed without UCT. Every invertible $KK$-class of a stable Kirchberg algebra is represented by an automorphism \cite[Theorem 4.2.1]{Phillips2000}, and two such automorphisms with the same $KK$-class are isotopic \cite[Theorem 4.1.4(2)]{Phillips2000}. Here an isotopy is a
point-norm continuous path consisting of automorphisms. For unital Kirchberg
algebras, equality of $KL$-classes of unital homomorphisms is equivalent to
approximate unitary equivalence \cite[Theorem~5.1(a)]{Dadarlat2005}, the
converse implication being \cite[Theorem~3.3]{Dadarlat2005}.

\begin{lemma}\label{lem:st-stable-kl-uniqueness}
 For $\alpha,\beta\in\operatorname{Aut}(A)$, equality of
their $KL$-classes is equivalent to approximate unitary equivalence by
multiplier unitaries. Consequently the continuous homomorphism 
\begin{equation*}
\kappa_A:\mathsf G_A\longrightarrow R_A^\times,\qquad  \alpha\longmapsto[\alpha]_{KL}, 
\end{equation*}
has kernel $N_A$.
\end{lemma}

\begin{proof}
It suffices to compare an automorphism $\gamma$ with the identity. Write $A=B\otimes\mathcal{K}$ for a unital Kirchberg algebra $B$, and let $e_{ij}=1_B\otimes e_{ij}^{\mathcal{K}}$. Equality of $KL$-classes implies $[\gamma(e_{11})]=[e_{11}]$ in $K_0(A)$. Nonzero projections with the same $K_0$-class in a purely infinite simple algebra are Murray--von Neumann
equivalent. We record explicitly why the equivalence extends to a multiplier
unitary. The strictly convergent sum 
\begin{equation*}
s=\sum_{i\geq1}e_{i+1,i}\in M(A) 
\end{equation*}
satisfies $s^*s=1$ and $ss^*=1-e_{11}$. The automorphism $\gamma$ has a
strictly continuous extension $\bar\gamma$ to $M(A)$, so $\bar\gamma(s)\bar\gamma(s)^*=1-\gamma(e_{11})$. Thus the two complementary
multiplier projections are both Murray--von Neumann equivalent to $1$.
Adding a partial isometry between the complements to one between $\gamma(e_{11})$ and $e_{11}$ gives a multiplier unitary carrying $\gamma(e_{11})$ to $e_{11}$. After the corresponding multiplier-inner
perturbation, assume $\gamma(e_{11})=e_{11}$.

The strict sum 
\begin{equation*}
w=\sum_{i\geq1}e_{i1}\gamma(e_{1i}) 
\end{equation*}
is a multiplier unitary, and $\operatorname{Ad}(w)\gamma$ fixes every $e_{ij}$. Indeed the
summands have pairwise orthogonal initial projections $\gamma(e_{ii})$ and
final projections $e_{ii}$, both summing strictly to one. Hence $\operatorname{Ad} (w)\gamma=\theta\otimes\mathrm{id}_{\mathcal{K}}$ for a unital automorphism $\theta$ of $B$. Passing to the full corner preserves equality of $KL$-classes, since the Morita equivalence and its inverse act continuously by
Kasparov product. Dadarlat's unital uniqueness theorem makes $\theta$
approximately inner. The implementing unitaries $v_n\in U(B)$ give
multiplier unitaries $v_n\otimes1$, proving the stable assertion.

Conversely, every multiplier-inner automorphism has $KK$-class one: left
multiplication by the implementing unitary's adjoint intertwines its
Kasparov correspondence with the identity correspondence. Continuity into
the Hausdorff group $R_A$ then implies equality of $KL$-classes under
approximate unitary equivalence. Continuity of $\kappa_A$ for \eqref{eq:st-inverse-pair} follows by applying continuity of the class map to
both $\alpha$ and $\alpha^{-1}$.
\end{proof}

\begin{lemma}\label{lem:st-kl-realization}
 Every unit of $R_A$ is represented by an
automorphism of $A$.
\end{lemma}

\begin{proof}
Choose a full corner $B=pAp$ with $[p]=0$ in $K_0(A)$. To find $p$, take an
infinite nonzero projection $e$ and a proper subprojection $f<e$ equivalent
to $e$; then $p=e-f\ne0$ has class zero and is full. Brown's theorem gives $B\otimes\mathcal{K}\cong A$. Thus $B$ is a unital Kirchberg algebra in
standard form, meaning $[1_B]=0$.

Transport a given unit $x\in R_A$, and its inverse, to $KL(B,B)$. Choose $KK$-lifts of these two classes. The Kirchberg--Phillips existence theorem
\cite[Theorem~8.2.1(i)]{RordamStormer2002} realises these lifts by full
homomorphisms into $B\otimes\mathcal{K}$, and
\cite[Proposition~4.1.4]{RordamStormer2002} cuts them down to unital
homomorphisms $\varphi,\psi:B\to B$, the unit-class condition being
automatic in standard form. This is precisely the existence step used in
\cite[proof of Theorem~5.1(a)]{Dadarlat2005}. Their composites have the identity $KL$-class,
so the uniqueness theorem gives $\psi\varphi\approx_u\mathrm{id}_B$ and $\varphi\psi\approx_u\mathrm{id}_B$.

Here are the intertwining details. Choose summable positive numbers $\varepsilon_n$ and increasing finite selfadjoint subsets $F_n$ of the unit
ball whose union is dense. We enlarge the sets during the construction by
finitely many images, without changing these properties. Starting with $\varphi_1=\varphi$, construct maps $\varphi_n,\psi_n$ such that every $\varphi_n$ is an inner perturbation of $\varphi$, every $\psi_n$ is an inner
perturbation of $\psi$, and 
\begin{align}
\left\|\psi_n\varphi_n(a)-a\right\|&<\varepsilon_n & &(a\in F_n),
\label{eq:st-kl-intertwining-one} \\
\left\|\varphi_{n+1}\psi_n(b)-b\right\|&<\varepsilon_n & &(b\in
F_n\cup\varphi_n(F_n)).  \label{eq:st-kl-intertwining-two}
\end{align}
Indeed, if $\varphi_n=\operatorname{Ad}(u)\varphi$, then $\psi\varphi_n=\operatorname{Ad} (\psi(u))\psi\varphi$ is approximately unitarily equivalent to the identity,
so a post-inner perturbation of $\psi$ gives \eqref{eq:st-kl-intertwining-one}.
Similarly, after $\psi_n=\operatorname{Ad}(v)\psi$ has been chosen, $\varphi\psi_n=\operatorname{Ad} (\varphi(v))\varphi\psi$ is approximately unitarily equivalent to the
identity, and a post-inner perturbation of $\varphi$ gives \eqref{eq:st-kl-intertwining-two}. Before the next step, enlarge $F_{n+1}$ so
that it contains $F_n\cup\psi_n(F_n)$.

For $a\in F_n$, the two displayed estimates give 
\begin{equation*}
\left\|\varphi_{n+1}(a)-\varphi_n(a)\right\|<2\varepsilon_n. 
\end{equation*}
At the next half-step, for $b\in F_n$, applying $\psi_{n+1}$ to \eqref{eq:st-kl-intertwining-two} and using $\psi_n(b)\in F_{n+1}$ gives 
\begin{equation*}
\left\|\psi_{n+1}(b)-\psi_n(b)\right\|  <\varepsilon_n+\varepsilon_{n+1}. 
\end{equation*}
Thus $\varphi_n$ and $\psi_n$ converge point-norm, on a dense union and
hence on all of $B$, to unital homomorphisms $\theta$ and $\eta$. Passing to
the limit in \eqref{eq:st-kl-intertwining-one}--\eqref{eq:st-kl-intertwining-two}
yields $\eta\theta=\mathrm{id}_B$ and $\theta\eta=\mathrm{id}_B$.
Consequently $\theta$ is an automorphism.

Every $\varphi_n$ has the prescribed $KL$-class, and the class map is
continuous with Hausdorff target. Thus $[\theta]_{KL}=x$, after the chosen
transport. Stabilisation and the fixed Morita identification return an
automorphism of $A$ with the original class. The proof does not require the
two initially chosen $KK$-lifts to be inverse to one another.
\end{proof}


\subsubsection{Path components as Polish subgroups}

In this paper, a subgroup $H$ of a Polish group $G$ is called a \emph{Polish
subgroup} if it admits a Polish group topology whose Borel sets are exactly
those inherited from $G$; this topology need not be the subspace topology. A
quotient $G/H$, where $H$ is a normal Borel Polish subgroup, is called a 
\emph{group with a Polish cover}; its quotient topology need not be
Hausdorff \cite[Remark 1.17]{LupiniDefinableKHomology2020}. This category was introduced by
Bergfalk--Lupini--Panagiotopoulos in the first part of their work on
definable homological invariants \cite{BergfalkLupiniPanagiotopoulos2024a} and developed further in the
second \cite{BergfalkLupiniPanagiotopoulos2024b}. Its full subcategory of abelian groups with a Polish
cover is the left heart of the derived category of the quasi-abelian
category of abelian Polish groups \cite{Lupini2024Heart}, while the category of
groups with a Polish cover and Borel-definable homomorphisms is the exact
completion, as a regular category, of the homological category of Polish
groups and continuous homomorphisms \cite{LupiniExactCompletion2026}.

\begin{lemma}\label{lem:st-path-polish}
 The identity path component $G_{\mathrm{path}}$ of
every Polish group $G$ is a normal Borel Polish subgroup.
\end{lemma}

\begin{proof}
With pointwise multiplication and the compact-open topology, 
\begin{equation*}
P_eG=\{\gamma\in C([0,1],G):\gamma(0)=e\} 
\end{equation*}
is a Polish group. Endpoint evaluation is a continuous homomorphism, with
image $G_{\mathrm{path}}$ and closed kernel $L_eG$. The induced continuous
injection $P_eG/L_eG\to G$ has Borel image by the Lusin--Souslin theorem and
is a Borel isomorphism onto that image. Transporting the Polish group
topology of $P_eG/L_eG$ proves that $G_{\mathrm{path}}$ is a Polish
subgroup. Normality follows by conjugating paths.
\end{proof}

\begin{proposition}\label{prop:st-kk-path-cover}
 One has 
\begin{equation*}
J_A=(\mathsf G_A)_{\mathrm{path}},\qquad \overline{J_A}=N_A. 
\end{equation*}
The map $\alpha\mapsto KK(\alpha)$ induces an algebraic isomorphism $\mathsf G_A/J_A\cong KK(A,A)^\times$, giving the latter a group-with-a-Polish-cover
presentation. The maximal Hausdorff quotient of $\mathsf G_A/J_A$, with its quotient
topology, is $\mathsf G_A/N_A$.
\end{proposition}

\begin{proof}
\uses{lem:st-path-polish, lem:st-stable-kl-uniqueness}
The isotopy theorem recalled above and homotopy invariance of $KK$-theory
give $J_A=(\mathsf G_A)_{\mathrm{path}}$. Classification realises every $KK$-unit,
giving the algebraic isomorphism, and Lemma~\ref{lem:st-path-polish} supplies
the Polish cover. Lemma~\ref{lem:st-stable-kl-uniqueness} gives $J_A\subseteq
N_A$. Multiplier inners have $KK$-class one, so $\operatorname{Inn}_M(A)\subseteq J_A$,
whence $\overline{J_A}=N_A$. Finally, the closure of the identity in $\mathsf G_A/J_A
$ is $N_A/J_A$. Quotienting by it gives $\mathsf G_A/N_A$; every continuous
homomorphism to a Hausdorff group annihilates this closure and therefore
factors through that quotient.
\end{proof}


\subsubsection{Compatibility with the Cuntz presentation}

The preceding Polish cover agrees definably with the natural cycle
presentation of the $KK$-unit group. We give the details because this
conclusion does not require a definable theory of all additive groups $KK(A,B)$.

We use the terminology of \cite[Sections 1.3--1.5]{LupiniDefinableKHomology2020}. A \emph{semidefinable set} is a quotient of a standard Borel space by an analytic
equivalence relation. It is \emph{definable} when the relation is Borel and
idealistic. A relation is idealistic if its classes admit an assignment of
proper sigma-filters, with the Borel measurability condition of \cite[Definition 1.6]{LupiniDefinableKHomology2020}; coset relations of Polish subgroups have this
property. A map between quotient sets is \emph{definable} if it admits a
Borel lift to the representative spaces. We use the following transfer
principle:

\begin{theorem}[{\protect\cite[Corollary 1.14]{LupiniDefinableKHomology2020}}]\label{thm:st-definable-transfer}
 A definable bijection from a definable set to
a semidefinable set is a definable isomorphism, and its target is definable.
\end{theorem}

Let $qA=\ker(A*A\to A)$, where the map is the identity on both free-product
copies, and put $C_A=\mathcal{K}\otimes qA$. The stabilised Cuntz picture
identifies 
\begin{equation}  \label{eq:st-cuntz-endomorphism}
KK(A,A)\cong[C_A,C_A],
\end{equation}
where brackets denote homotopy classes of homomorphisms and composition
represents the Kasparov product \cite[Section 5.6]{LupiniDefinableKHomology2020}. The same
source constructs $KK(A,B)$ as a semidefinable group, with Borel-lifted
product, but leaves definability for arbitrary coefficient algebras open.

Write $\mathcal{H}_A=\operatorname{Hom}_*(C_A,C_A)$ for the Polish space of all, possibly
nonunital, $^*$-homomorphisms, with point-norm topology. Composition is
jointly continuous: on a fixed $c$, 
\begin{equation*}
\|f_n(g_n(c))-f(g(c))\|  \leq\|g_n(c)-g(c)\|+\|f_n(g(c))-f(g(c))\|. 
\end{equation*}
Consider the closed Polish subspace 
\begin{equation}  \label{eq:st-unit-witnesses}
\begin{split}
\mathcal{W}_A=\{(f,g,h,k)\in \mathcal{H}_A^2\times C([0,1],\mathcal{H}_A)^2:\quad& \\
h(0)=gf,\ h(1)=\mathrm{id}_{C_A},\quad k(0)=fg,\ k(1)=\mathrm{id}_{C_A}&\}.
\end{split}\end{equation}
Two quadruples are equivalent when their first coordinates are homotopic.
This relation is analytic, since a homotopy is a witness in the Polish path
space. The quotient represents exactly the units of $[C_A,C_A]$: $g$
represents the inverse of $f$. It is a semidefinable group. Inversion swaps $f,g$ and $h,k$. For multiplication, the representatives are $ff^{\prime }$
and $g^{\prime }g$; concatenating the paths $g^{\prime }hf^{\prime }$ and $h^{\prime }$, and the paths $fk^{\prime }g$ and $k$, gives the required
inverse witnesses. These operations have Borel, in fact continuous, lifts.

\begin{proposition}\label{prop:st-definable-kk-units}
 The unit group $KK(A,A)^\times$, with the
witnessed Cuntz presentation $\mathcal{W}_A/{\sim}$, is definable. The
natural isomorphism 
\begin{equation*}
\mathsf G_A/J_A\longrightarrow KK(A,A)^\times 
\end{equation*}
is a definable isomorphism.
\end{proposition}

\begin{proof}
\uses{prop:st-kk-path-cover, thm:st-definable-transfer}
An automorphism $\alpha$ induces the automorphism $f_\alpha=\mathrm{id}_{\mathcal{K}}\otimes q\alpha$ of $C_A$. The map 
\begin{equation*}
\alpha\longmapsto  (f_\alpha,f_{\alpha^{-1}},\mathrm{const}_{\mathrm{id}},  
\mathrm{const}_{\mathrm{id}}) 
\end{equation*}
is Borel, indeed continuous. Naturality of \eqref{eq:st-cuntz-endomorphism}
identifies its class with $KK(\alpha)$. Proposition~\ref{prop:st-kk-path-cover}
shows that the induced map on quotients is bijective and that its source is
definable. Apply Theorem~\ref{thm:st-definable-transfer}.
\end{proof}

\begin{proof}[Proof of Theorem~\protect\ref{thm:st-kk}]
\proves{thm:st-kk}
\uses{lem:st-kl-realization, lem:st-stable-kl-uniqueness, prop:st-definable-kk-units, prop:st-kk-path-cover, prop:st-transfer, thm:st-kernel-ea}
Lemmas~\ref{lem:st-stable-kl-uniqueness} and~\ref{lem:st-kl-realization} make $\kappa_A$ a continuous surjective homomorphism of Polish groups with kernel $N_A$. It is open by the Polish open mapping theorem. Theorem~\ref{thm:st-kernel-ea} gives whirly amenability of $N_A$, and
Proposition~\ref{prop:st-transfer} gives the universal-minimal-flow and Bohr
formulas. Propositions~\ref{prop:st-kk-path-cover} and~\ref{prop:st-definable-kk-units} prove the assertions
about the definable unit group and its Hausdorff quotient.
The whirly-amenability equivalence is the additional assertion of
Proposition~\ref{prop:st-transfer}.
\end{proof}

Throughout the paper $Q_A=\mathsf G_A/N_A$ denotes this Hausdorff
quotient. Theorem~\ref{thm:st-kk} identifies it with
$R_A^\times=KL(A,A)^\times$, and Section~\ref{sec:st-total} identifies it,
under UCT, with $\operatorname{Aut}_\Lambda(\underline K(A))$; the single symbol $Q_A$
is used for all three descriptions. The kernel $N_A$ is the closure of
the multiplier-inner automorphisms; the companion unital paper
\cite{LupiniUnital2026} instead divides by the closure of
$\operatorname{Ad} U_0(A)$, and its quotient is written $Q_A^0$ there.

\begin{corollary}\label{cor:st-kirchberg-fd-representations}
 For every stable Kirchberg algebra $A$, without a UCT assumption, pullback along $\kappa_A$ gives an equivalence 
\begin{equation*}
\operatorname{Rep}_{\mathrm{fd}}(KL(A,A)^\times)  \simeq  \operatorname{Rep}_{\mathrm{fd}}(\operatorname{Aut}(A)) 
\end{equation*}
between continuous finite-dimensional unitary representation categories. In
particular, $\operatorname{Aut}(A)$ is minimally almost periodic if and only if $KL(A,A)^\times$ is. Under UCT, the latter group is $\operatorname{Aut}_\Lambda(\underline
K(A))$.
\end{corollary}

\begin{proof}
\uses{cor:st-finite-dimensional-representations}
Apply Corollary~\ref{cor:st-finite-dimensional-representations} to $\kappa_A$,
and then the UCT identification of Section~\ref{sec:st-total} under UCT.
\end{proof}

View $KK(A,A)^\times$ as the quotient topological group $\mathsf G_A/J_A$ in
Proposition~\ref{prop:st-kk-path-cover}. Although this topology need not be
Hausdorff, amenability still has its usual compact-dynamical meaning: every
continuous action on a compact Hausdorff space admits an invariant
probability measure.

\begin{corollary}[UCT-free amenability criterion]\label{cor:st-kirchberg-amenability-no-uct}
 For every stable Kirchberg algebra $A$, without a UCT assumption, the following are equivalent:

\begin{enumerate}
\item $\operatorname{Aut}(A)$ is amenable;

\item $KL(A,A)^\times$ is amenable;

\item the definable group $KK(A,A)^\times$, with its
group-with-a-Polish-cover topology, is amenable.
\end{enumerate}

Equivalently, every compact Hausdorff action of the definable $KK$-unit
group has an invariant probability measure.
\end{corollary}

\begin{proof}
\uses{prop:st-transfer, thm:st-kk}
The equivalence of (i) and (ii) follows from Theorem~\ref{thm:st-kk}, or
directly from Proposition~\ref{prop:st-transfer}. The closure of the identity
in $\mathsf G_A/J_A$ is $N_A/J_A$. It acts trivially in every continuous action on a
Hausdorff space: for each point, continuity of its orbit map identifies the
action of every element in that closure with the action of the identity.
Hence every compact Hausdorff action of $\mathsf G_A/J_A$ factors through its
maximal Hausdorff quotient $\mathsf G_A/N_A\cong KL(A,A)^\times$, and conversely
every action of the latter pulls back to the former. This proves the
equivalence of (ii) and (iii).
\end{proof}

\begin{corollary}[Locally compact \(KL\)-quotients]\label{cor:st-locally-compact-kl}
 Let $A$ be a stable Kirchberg algebra,
without a UCT assumption, and suppose that $KL(A,A)^\times$ is locally
compact. Then the following are equivalent:

\begin{enumerate}
\item $M(\operatorname{Aut}(A))$ is metrizable;

\item $KL(A,A)^\times$ is compact.
\end{enumerate}

When these conditions hold, 
\begin{equation*}
M(\operatorname{Aut}(A))\cong\kappa_A^*KL(A,A)^\times,  \qquad b\operatorname{Aut}(A)\cong
KL(A,A)^\times, 
\end{equation*}
where the quotient acts on itself by left translation. In particular, $\operatorname{Aut} (A)$ is amenable and its universal minimal flow is uniquely ergodic, with
Haar measure on the displayed quotient.
\end{corollary}

\begin{proof}
\uses{lem:st-locally-compact-umf, thm:st-kk}
Theorem~\ref{thm:st-kk} and Lemma~\ref{lem:st-locally-compact-umf} prove the
equivalence. When it holds, the displayed formulas follow from Theorem~\ref{thm:st-kk} and the compact-group facts in Section~\ref{sec:st-dynamics}; Haar measure is the unique invariant
probability measure on the left-translation flow.
\end{proof}


\subsection{Total \(K\)-theory and the automorphism quotient}

\label{sec:st-total} Throughout this section $A$ is a stable UCT Kirchberg
algebra. We identify the quotient of Section~\ref{sec:st-definable-kk} in terms
of total $K$-theory and then describe its phantom kernel. The UCT of
Rosenberg and Schochet \cite{RosenbergSchochet1987} gives the degree-zero exact sequence 
\begin{equation}  \label{eq:st-uct}
\begin{split}
0\longrightarrow \operatorname{Ext}^1_\mathbb{Z}(K_0(A),K_1(A))\oplus \operatorname{Ext}^1_\mathbb{Z}(K_1(A),K_0(A)) \longrightarrow KK(A,A) \\
\longrightarrow \operatorname{Hom}(K_0(A),K_0(A))\oplus \operatorname{Hom}(K_1(A),K_1(A))
\longrightarrow0.
\end{split}\end{equation}
We retain the composition convention for the Kasparov product.

Total $K$-theory retains the coefficient exact sequences 
\begin{equation}  \label{eq:st-bockstein}
K_i(A)\xrightarrow{n}K_i(A) \xrightarrow{\rho_n}K_i(A;\mathbb{Z}/n) \xrightarrow{\beta_n}K_{i-1}(A) \xrightarrow{n}K_{i-1}(A),
\end{equation}
as well as all coefficient-change operations. The notation $\Lambda$
includes these operations and their compositions, not only the maps in a
single displayed sequence.

The universal multicoefficient theorem of Dadarlat--Loring \cite{DadarlatLoring1996}
identifies 
\begin{equation}  \label{eq:st-umct}
KL(A,A)\cong\operatorname{Hom}_\Lambda(\underline K(A),\underline K(A)),
\end{equation}
where 
\begin{equation*}
KL(A,A)=KK(A,A)/P_A,\qquad  P_A=\operatorname{Pext}^1_\mathbb{Z}(K_0(A),K_1(A))  \oplus \operatorname{Pext}^1_\mathbb{Z}(K_1(A),K_0(A)). 
\end{equation*}
Here pure extensions are those whose restrictions to finitely generated
subgroups of the quotient split; $P_A$ is embedded by the UCT. We use
exactness, not a splitting of the multicoefficient sequence: contrary to the
statement in \cite{DadarlatLoring1996}, it need not split, as recorded in the discussion
following \cite[Theorem~1.3]{EilersRestorffRuiz2011}. Dadarlat identifies the topology in \eqref{eq:st-umct} with pointwise
convergence and, for $A$ in the UCT class, $P_A$ with the closure of zero
in the natural $KK$-topology \cite[Theorem~4.1 and p.~411]{Dadarlat2005}. On
units the topology of \eqref{eq:st-inverse-pair} is exactly the pointwise
automorphism topology: if bijections converge pointwise to a bijection,
their inverses also converge pointwise. Thus the $KL$-unit quotient in
Theorem~\ref{thm:st-kk} is precisely the group appearing in Theorem~\ref{thm:cu-stable-general}.

For the UCT presentation we use the topological identification from
Theorem~\ref{thm:cu-stable-general}:
\[
 q_A:\operatorname{Aut}(A)\longrightarrow Q_A\cong\operatorname{Aut}_\Lambda(\underline K(A)),
 \qquad \ker q_A=N_A.
\]
The map is continuous, surjective, and open. Its realisation and
uniqueness proof, including continuity on coefficient classes, belongs
to the cited theorem. The strengthened whirly-kernel conclusion follows
from the norm-unitary theorem and strict approximation, as proved in
Section~\ref{sec:st-kernel}.

\begin{proposition}\label{prop:st-pure-state-quotient}
 Let $A$ be a stable UCT Kirchberg algebra,
let $\varphi$ be a pure state, and put $H_\varphi=\operatorname{Aut}(A,\varphi)$.
Restriction gives a continuous open surjection 
\begin{equation*}
q_A|_{H_\varphi}:H_\varphi\longrightarrow  \operatorname{Aut}_\Lambda(\underline
K(A)) 
\end{equation*}
with whirly amenable kernel $H_\varphi\cap N_A$. Consequently,
\begin{equation*}
\begin{split}
M(H_\varphi)&\cong (q_A|_{H_\varphi})^*M(\operatorname{Aut}_\Lambda(\underline
K(A))), \\
bH_\varphi&\cong b\operatorname{Aut}_\Lambda(\underline K(A)).
\end{split}
\end{equation*}
The restriction of $M(\operatorname{Aut}(A))$ to $H_\varphi$ is universal minimal.
Amenability and extreme amenability of these two groups are equivalent to
the corresponding properties of the total-$K$ quotient.
Both groups are whirly amenable exactly when that quotient is trivial.
\end{proposition}

\begin{proof}
\uses{cor:st-stable-stabilizer-umf, lem:st-nonarch-whirly, prop:st-transfer}
Apply Corollary~\ref{cor:st-stable-stabilizer-umf}, proved in Section~\ref{sec:st-stable-stabilizers}, and the topological total-$K$ identification
above. Its kernel theorem is independent of UCT and does not use the present
proposition.
The whirly criterion follows from Proposition~\ref{prop:st-transfer} and
Lemma~\ref{lem:st-nonarch-whirly}.
\end{proof}

\begin{corollary}\label{cor:st-aut-bohr}
For a stable UCT Kirchberg
algebra $A$, extreme amenability, amenability, and metrizability of the
universal minimal flow hold for $\operatorname{Aut}(A)$ exactly when they hold for
$\operatorname{Aut}_\Lambda(\underline K(A))$, respectively, and
\begin{equation*}
b\operatorname{Aut}(A)\cong b\operatorname{Aut}_\Lambda(\underline K(A)).
\end{equation*}
Furthermore, $\operatorname{Aut}(A)$ is whirly amenable exactly when
$\operatorname{Aut}_\Lambda(\underline K(A))=1$, equivalently when $\operatorname{Aut}(A)=N_A$.
\end{corollary}

\begin{proof}
\uses{lem:st-nonarch-whirly, prop:st-transfer}
Apply Proposition~\ref{prop:st-transfer} and \eqref{eq:st-transfer-bohr} to
$q_A$; for the whirly assertion also use Lemma~\ref{lem:st-nonarch-whirly},
since the total-$K$ quotient is a closed permutation group on a countable
discrete invariant.
\end{proof}

\begin{corollary}\label{cor:st-closed-preimages}
 Let $L$ be a closed subgroup of $Q_A=\operatorname{Aut} _\Lambda(\underline K(A))$, and put 
\begin{equation*}
G_{A,L}=q_A^{-1}(L). 
\end{equation*}
Then the restriction $q_A:G_{A,L}\to L$ is a continuous open surjection with
whirly amenable kernel $N_A$. Consequently
\begin{equation*}
M(G_{A,L})\cong M(L),\qquad bG_{A,L}\cong bL, 
\end{equation*}
where the first action is pulled back along $q_A$.
Moreover, $G_{A,L}$ is whirly amenable if and only if $L=1$.
\end{corollary}

\begin{proof}
\uses{lem:st-nonarch-whirly, prop:st-transfer}
Only openness requires comment. If $U=V\cap G_{A,L}$, where $V$ is open in $\operatorname{Aut}(A)$, then 
\begin{equation*}
q_A(U)=q_A(V)\cap L, 
\end{equation*}
which is open in $L$. The kernel is $N_A$, so Proposition~\ref{prop:st-transfer}
applies.
For the additional assertion, apply its whirly-amenability part and
Lemma~\ref{lem:st-nonarch-whirly} to the Polish non-Archimedean group $L$.
\end{proof}


\subsubsection{The extension term invisible to ordinary \(K\)-theory}

Put $G_i=K_i(A)$ and 
\begin{equation*}
E_A=\operatorname{Ext}^1_\mathbb{Z}(G_0,G_1)\oplus\operatorname{Ext}^1_\mathbb{Z}(G_1,G_0). 
\end{equation*}
Restriction gives a continuous open surjection 
\begin{equation*}
r:Q_A\longrightarrow\operatorname{Aut}(G_0)\times\operatorname{Aut}(G_1). 
\end{equation*}
Surjectivity follows by lifting a graded isomorphism through the UCT and
using classification; openness follows from the Polish open mapping theorem.

We recall the algebraic extension from
\cite[Section~8.1]{LupiniCuntz2026}:
\begin{equation}\label{eq:st-ext-kernel}
 1\longrightarrow E_A/P_A\longrightarrow Q_A
 \xrightarrow{r}\operatorname{Aut}(G_0)\times\operatorname{Aut}(G_1)\longrightarrow1.
\end{equation}
The kernel is additive, with its inherited topology. No general
splitting or discreteness is assumed. The square-zero argument belongs
to the cited proof. We use this sequence to obtain the following
amenability consequence.

\begin{corollary}[Ordinary \(K\)-theory amenability
criterion]\label{cor:st-amenability-k-groups}
 For every stable UCT Kirchberg algebra $A$, 
\begin{equation*}
\operatorname{Aut}(A)\text{ is amenable}  \quad\Longleftrightarrow\quad  \operatorname{Aut}(K_0(A))\text{
and }\operatorname{Aut}(K_1(A))\text{ are amenable}. 
\end{equation*}
The automorphism groups of the $K$-groups carry pointwise convergence on the
underlying countable discrete groups.
\end{corollary}

\begin{proof}
\uses{cor:st-aut-bohr, lem:st-amenable-extensions}
The kernel $E_A/P_A$ in \eqref{eq:st-ext-kernel} is an abelian topological
group and hence is amenable, by the Markov--Kakutani fixed-point theorem.
Therefore Lemma~\ref{lem:st-amenable-extensions} shows that $Q_A$ is amenable
whenever $\operatorname{Aut}(G_0)\times\operatorname{Aut}(G_1)$ is amenable. Conversely, amenability
passes to continuous quotients, so amenability of $Q_A$ implies amenability
of both factors. Finally apply Corollary~\ref{cor:st-aut-bohr}.
\end{proof}

\begin{corollary}[Vanishing Hausdorff extension kernel]\label{cor:st-vanishing-ext}
If \(E_A=P_A\), restriction induces a topological isomorphism
\[
 Q_A\cong\operatorname{Aut}(G_0)\times\operatorname{Aut}(G_1).
\]
\end{corollary}
\begin{proof}
In \eqref{eq:st-ext-kernel}, restriction is then a continuous bijective
homomorphism of Polish groups, hence an isomorphism by the open mapping
theorem. This criterion permits both degrees to be nonzero.
\end{proof}

When one \(K\)-group vanishes, the reduction to automorphisms of the
other is Corollary~\ref{cor:cu-k1zero}; the four-point
\(K_0=K_1=\mathbb Z/2\) example is its Example~8.7. These existing computations
are not reproved here. Phantom information, which disappears from the
Hausdorff quotient, is considered next.


\subsection{Explicit \(K\)-theoretic computations}

\label{sec:st-computations}


\subsubsection{The existing finitely generated computation}
The discreteness, finite-rank criterion, and finite cardinality formula
for finitely generated \(K\)-groups are
Proposition~\ref{prop:cu-finite-k}. The stabilised Cuntz flows and
their table are Corollary~\ref{cor:cu-cuntz-stable}; we do not reproduce
them. The additional amenability criterion here is
Corollary~\ref{cor:st-amenability-k-groups}: in the finitely generated
family amenability is equivalent to both free ranks being at most one.
We now turn to infinite invariants and extensions of the existing
Pr\"ufer and vector-space examples.


\subsection{Pure-state stabilisers of stable Kirchberg algebras}\label{sec:st-stable-stabilizers}

We now return to a stable Kirchberg algebra $A$, without UCT. Write $H=\operatorname{Aut} (A,\varphi)$, $N=N_A$, and $L=H\cap N$. The norm-unitary theorem gives a direct proof once the relative inner density is established.

\begin{lemma}\label{lem:st-stable-u0-density}
 For every stable Kirchberg algebra $A$, 
\begin{equation*}
N_A=\overline{\operatorname{Ad} U_0(1+A)}, 
\end{equation*}
where $U_0(1+A)$ denotes the component of one in the scalar-part-one
unitaries of the minimal unitisation.
\end{lemma}

\begin{proof}
\uses{lem:st-adjoint, thm:st-multiplier-whirly}
As in the proof of Theorem~\ref{thm:st-kernel-ea}, write
$A\cong B\otimes\mathcal K$ with $B$ a unital Kirchberg algebra.
The strict-density assertion of Theorem~\ref{thm:st-multiplier-whirly}
makes $U_0(1+A)=U_1(A)_0$ strictly dense in $U(M(A))$.
Continuity of the adjoint map, Lemma~\ref{lem:st-adjoint}, then gives
the equality of the two point-norm closures. Only the strict-density
part of that theorem is used here.
\end{proof}

\begin{lemma}\label{lem:st-state-projections}
 For every pure state $\varphi$ of $A$, there
is an increasing projection approximate identity $(e_n)$ with $\varphi(e_n)=1
$.
\end{lemma}

\begin{proof}
Choose a nonzero projection $p\in A$. The irreducible GNS representation is
faithful. Kadison unitary transitivity in $A^+$ therefore gives a unitary
moving the cyclic vector into the range of $\pi_\varphi(p)$. Conjugating $p$
back gives a projection $e_0\in A$ with $\varphi(e_0)=1$. The hereditary
algebra $(1-e_0)A(1-e_0)$ has real rank zero and an increasing projection
approximate identity $(f_n)$. Set $e_n=e_0+f_n$. These projections
approximate the diagonal corners of every element. For an off-diagonal
corner $x=e_0a(1-e_0)$, one has 
\begin{equation*}
\left\|x(1-f_n)\right\|^2 
=\left\|(1-f_n)x^*x(1-f_n)\right\|\longrightarrow0. 
\end{equation*}
The other corner is similar. Thus $(e_n)$ is an approximate identity of $A$,
and every $e_n$ fixes the cyclic vector.
\end{proof}

\begin{lemma}\label{lem:st-stable-marked-inner}
 Put 
\begin{equation*}
D_\varphi=\{a\in A:\varphi(aa^*)=\varphi(a^*a)=0\},\qquad  W^{(0)}=\{u\in
U_1(D_\varphi):[u]=0\in K_1(D_\varphi)\}. 
\end{equation*}
Then $L=\overline{\operatorname{Ad} W^{(0)}}$.
\end{lemma}

\begin{proof}
\uses{lem:st-stable-u0-density, lem:st-state-projections}
Fix $\alpha\in L$, a finite test set $F\subset A$, and a tolerance. Apply
Property~1.2 of \cite{KishimotoOzawaSakai2003}, which holds for every $C^*$-algebra by
Theorems~3.1 and~4.3 of that paper, to $(F,\varphi)$, obtaining additional tests and a
tolerance for exact correction of equivalent pure states. The essentiality
hypothesis holds: the faithful irreducible image of a simple purely infinite
algebra contains no nonzero compact operators. By Lemma~\ref{lem:st-stable-u0-density}, choose $v\in U_0(1+A)$ whose conjugation
approximates $\alpha$ on both sets of tests. Property~1.2 supplies $w\in
U_0(1+A)$ such that 
\begin{equation*}
\varphi\operatorname{Ad} v=\varphi\operatorname{Ad} w,\qquad  \operatorname{Ad} w\approx\mathrm{id}\text{ on }F. 
\end{equation*}
A unitary path in $A^+$ can be normalised continuously by its scalar
quotient, so the scalar part can be kept equal to one. Then $z=vw^*$ lies in 
$U_0(1+A)$, its conjugation fixes $\varphi$, and $\operatorname{Ad} z$ approximates $\alpha
$ on $F$. Purity gives $\pi_\varphi(z)\xi_\varphi=e^{it}\xi_\varphi$ for
some real $t$.

Use the projections from Lemma~\ref{lem:st-state-projections} and put 
\begin{equation*}
q_n=\exp(-it e_n),\qquad u_n=zq_n. 
\end{equation*}
The unitaries $u_n$ have scalar part one, belong to $U_0(1+A)$, and fix $\xi_\varphi$. Moreover, $q_n\to e^{-it}1$ strictly, so $\operatorname{Ad} u_n\to\operatorname{Ad} z$
point-norm. Both $u_n-1$ and its adjoint annihilate $\xi_\varphi$; hence $u_n-1\in D_\varphi$. The nonzero hereditary algebra $D_\varphi$ is full in $A
$, and its inclusion induces a $K_1$-isomorphism. Thus $[u_n]=0$ also in $K_1(D_\varphi)$, proving the required density. Conversely, every element of $\operatorname{Ad} W^{(0)}$ fixes $\varphi$ and is multiplier inner, so belongs to $L$.
\end{proof}

\begin{theorem}\label{thm:st-stable-marked-kernel}
 For every stable Kirchberg algebra $A$ and
every pure state $\varphi$, the group $\operatorname{Aut}(A,\varphi)\cap N_A$ is whirly
amenable.
\end{theorem}

\begin{proof}
\uses{lem:st-stable-marked-inner, thm:st-v}
Kadison transitivity makes \(D_\varphi\) nonzero, and simplicity
makes this hereditary algebra full. It is \(\sigma\)-unital by
separability. If it had an identity \(e\), its bidual support would be
\(1-s_\varphi\), where \(s_\varphi\) is the minimal support projection
of the pure state. For \(a\in A\) with \(\varphi(a)\ne0\),
\[
 (1-e)a(1-e)=\varphi(a)(1-e)\in A.
\]
Thus \(1-e\in A\) and \(1\in A\), contradicting nonzero stability.
So \(D_\varphi\) is nonunital. Zhang's dichotomy \cite[Theorem~1.2(i)]{Zhang1992} makes it stable, and
Brown's theorem \cite[Theorem~2.8]{Brown1977} identifies it stably with
\(A\). Choosing a
nonzero projection in \(D_\varphi\) expresses it as \(C\otimes\mathcal K\)
for a unital nuclear Kirchberg corner \(C\). The stable component
identification makes \(W^{(0)}=U_1(D_\varphi)_0\). Theorem~\ref{thm:st-v}
therefore makes \(W^{(0)}\) whirly amenable in norm. Its continuous
adjoint image is point-norm dense in \(L\) by
Lemma~\ref{lem:st-stable-marked-inner}. Dense-image permanence proves the
assertion. No extension of a corner automorphism is required.
\end{proof}

\begin{corollary}\label{cor:st-stable-stabilizer-umf}
 Let $A$ be a stable Kirchberg algebra, $Q_A=\operatorname{Aut}(A)/N_A$, and $H=\operatorname{Aut}(A,\varphi)$ for a pure state $\varphi$.
Restriction of the quotient map is an open continuous surjection $H\to Q_A$
with whirly amenable kernel. In particular,
\begin{equation*}
M(H)\cong(q_A|_H)^*M(Q_A),\qquad bH\cong bQ_A. 
\end{equation*}
The restriction of $M(\operatorname{Aut}(A))$ to $H$ is itself universal minimal. The
groups $H$, $\operatorname{Aut}(A)$, and $Q_A$ are simultaneously extremely amenable, and
simultaneously amenable.
They are also simultaneously whirly amenable; if $Q_A$ is
non-Archimedean, in particular under UCT, this occurs exactly when $Q_A=1$.
\end{corollary}

\begin{proof}
\uses{lem:st-nonarch-whirly, prop:st-transfer, thm:st-stable-marked-kernel}
For $\alpha\in\operatorname{Aut}(A)$, pure-state homogeneity \cite[Theorem~2.1]{KishimotoOzawaSakai2003}
gives $\beta\in\operatorname{AInn}_0(A)\subseteq N_A$ with $\varphi\alpha\beta=\varphi$. Thus $\alpha\beta\in H$ represents
the same quotient element as $\alpha$, proving surjectivity. The Polish open
mapping theorem gives openness. Apply Theorem~\ref{thm:st-stable-marked-kernel}
and Proposition~\ref{prop:st-transfer}, including its whirly-amenability
assertion. Lemma~\ref{lem:st-nonarch-whirly} gives the criterion when $Q_A$
is non-Archimedean. The same transfer theorem for $\operatorname{Aut}(A)$
identifies both minimal flows with the indicated pullbacks. Surjectivity of $H\to Q_A$ gives the assertion about restriction.
\end{proof}

\begin{proof}[Proof of Theorem~\protect\ref{thm:st-stable-stabilizer}]
\proves{thm:st-stable-stabilizer}
\uses{cor:st-stable-stabilizer-umf, thm:st-stable-marked-kernel}
Combine Theorem~\ref{thm:st-stable-marked-kernel} and Corollary~\ref{cor:st-stable-stabilizer-umf}.
\end{proof}


\subsection{Direct sums and zero-dimensional fields}

\label{sec:st-geometric} Throughout this section, $B$ is a stable
Kirchberg algebra such that $E=\operatorname{Aut}(B)$ is extremely amenable. By Theorem~\ref{thm:st-kk}, this is equivalent to extreme amenability of $KL(B,B)^\times$; under UCT it is equivalent to extreme amenability of $\operatorname{Aut} _\Lambda(\underline K(B))$. The hypothesis holds, in particular, for $B=\mathcal{O}_2\otimes\mathcal{K}$ and $B=\mathcal{O}_3\otimes\mathcal{K}$.
These two examples have $E=N_B$ whirly amenable. More generally, whenever
$E$ is whirly amenable the coordinate and field kernels below have that
stronger property. We retain the extreme-amenability hypothesis for the
general UMF formulas.


\subsubsection{Permutations of summand ideals}

\begin{proposition}\label{prop:st-direct-sum}
 For a nonempty finite or countable set $I$, 
\begin{equation*}
\operatorname{Aut}(c_0(I,B))\cong E^I\rtimes\operatorname{Sym}(I) 
\end{equation*}
as topological groups, with the product topology on $E^I$ and the pointwise
topology on $\operatorname{Sym}(I)$. Consequently 
\begin{equation*}
M(\operatorname{Aut}(c_0(I,B)))\cong M(\operatorname{Sym}(I)).
\end{equation*}
If $E$ is whirly amenable, then the normal coordinate kernel $E^I$ is
whirly amenable.
\end{proposition}

\begin{proof}
\uses{lem:st-products, prop:st-transfer}
The coordinate ideals are exactly the minimal nonzero ideals, so an
automorphism permutes them. After specifying a permutation $\sigma$, the
restrictions are automorphisms $\gamma_j\in E$, and the formula is 
\begin{equation}  \label{eq:st-sum-action}
(\alpha a)_j=\gamma_j(a_{\sigma^{-1}(j)}).
\end{equation}
Conversely this formula defines an automorphism for any such data.

For a norm-one element supported in a single coordinate, images supported in
different coordinates have distance one. Thus point-norm convergence forces
eventual agreement on the image of each coordinate, proving continuity of
the permutation map. Once the relevant permutation coordinates agree,
testing supported elements gives point-norm convergence of each $\gamma_j$.
Conversely, coordinate convergence and eventual agreement of the
permutations give convergence on finitely supported elements. These are
dense in $c_0(I,B)$, and the maps are isometric, so convergence holds
everywhere. This proves the topological identification.

The kernel $E^I$ is extremely amenable by Lemma~\ref{lem:st-products}. Apply
Proposition~\ref{prop:st-transfer}.
The whirly assertion follows from the stronger part of Lemma~\ref{lem:st-products}.
\end{proof}

\begin{corollary}\label{cor:st-direct-sum-bohr-general}
 For every nonzero simple stable $C^*$-algebra $C
$ and every nonempty set $I$, put $E=\operatorname{Aut}(C)$. With the
coordinate-permutation action on $(bE)^I$, one has 
\begin{equation*}
b\operatorname{Aut}(c_0(I,C))\cong  
\begin{cases}
(bE)^I\rtimes\operatorname{Sym}(I), & I\text{ finite}, \\ 
1, & I\text{ infinite}.\end{cases}
\end{equation*}
Thus every infinite direct sum is minimally almost periodic, with no
separability, nuclearity, or amenability assumption. If $E$ is
amenable, then $\operatorname{Aut}(c_0(I,C))$ is amenable for every $I$. When $C$ is a
stable Kirchberg algebra, this last hypothesis is equivalent to amenability
of $KL(C,C)^\times$.
\end{corollary}

\begin{proof}
\uses{cor:st-kirchberg-amenability-no-uct, cor:st-wreath-minap, lem:st-amenable-extensions, lem:st-amenable-products, lem:st-bohr-finite-products, prop:st-direct-sum}
The algebraic and topological parts of the proof of Proposition~\ref{prop:st-direct-sum} use only that $C$ is nonzero and simple; using nets in the
possibly nonmetrizable case, they give 
\begin{equation*}
\operatorname{Aut}(c_0(I,C))\cong E^I\rtimes\operatorname{Sym}(I) 
\end{equation*}
for every nonempty $I$. If $I$ is infinite, apply Corollary~\ref{cor:st-wreath-minap}.

Suppose that $I$ is finite. The canonical homomorphism 
\begin{equation*}
E^I\rtimes\operatorname{Sym}(I)\longrightarrow(bE)^I\rtimes\operatorname{Sym}(I) 
\end{equation*}
has dense image and compact target. To verify its universal property, let $\rho:E^I\rtimes\operatorname{Sym}(I)\to K$ be a continuous homomorphism to a compact
group. By Lemma~\ref{lem:st-bohr-finite-products}, the restriction of $\rho$ to 
$E^I$ extends uniquely to a continuous homomorphism $\bar\rho:(bE)^I\to K$.
The conjugacy relations with the finite coordinate-permutation subgroup
continue to hold on $(bE)^I$, by density. Consequently 
\begin{equation*}
(x,\sigma)\longmapsto  \bar\rho(x)\rho(1,\sigma) 
\end{equation*}
is the unique extension of $\rho$ to $(bE)^I\rtimes\operatorname{Sym}(I)$, proving the
finite formula.

Finally, if $E$ is amenable, then $E^I$ is amenable by Lemma~\ref{lem:st-amenable-products}. The finitary permutations form a dense directed
union of finite groups in $\operatorname{Sym}(I)$, so that group is amenable by the same
lemma. Lemma~\ref{lem:st-amenable-extensions} now proves amenability of the
semidirect product. The last assertion is Corollary~\ref{cor:st-kirchberg-amenability-no-uct}.
\end{proof}

For finite $I$, the UMF is the regular $|I|!$-point flow of $\operatorname{Sym}(I)$. For $I=\mathbb{N}$, Glasner--Weiss \cite{GlasnerWeiss2002,Glasner2002Orders} identify $M(\operatorname{Sym}(\mathbb{N}))$ with the compact subspace $\operatorname{LO}(\mathbb{N})\subseteq2^{\mathbb{N}\times\mathbb{N}}$ of strict linear orders. The action is 
\begin{equation*}
i<_{\alpha\cdot L}j  \quad\Longleftrightarrow\quad 
\sigma_\alpha^{-1}(i)<_L\sigma_\alpha^{-1}(j). 
\end{equation*}
This also specifies the action of the full algebra automorphism group.

For any closed subgroup $P\leq\operatorname{Sym}(\mathbb{N})$, its preimage under the
permutation map is a closed Polish group $E^\mathbb{N}\rtimes P$ with UMF $M(P)$. Thus the construction transfers other known permutation-group
computations.


\subsubsection{Continuous fields over zero-dimensional spaces}

\begin{lemma}\label{lem:st-continuous-maps}
 If $X$ is compact and zero-dimensional and $E$
is extremely amenable, respectively whirly amenable, then $C(X,E)$, with
the compact-open group topology, has the corresponding property.
\end{lemma}

\begin{proof}
\uses{lem:st-products}
For each finite clopen partition $\mathcal{P}$ of $X$, the maps constant on
its parts form a subgroup isomorphic to $E^{|\mathcal{P}|}$. These subgroups
are extremely amenable and directed by refinement. They have dense union:
uniform continuity and a finite clopen refinement allow a continuous map to
be approximated by choosing its value at one point of each part. Apply Lemma~\ref{lem:st-products}.
If $E$ is whirly amenable, those finite products are whirly amenable, so
the stronger part of the same lemma applies.
\end{proof}

\begin{proposition}\label{prop:st-field}
 For nonempty compact metrizable zero-dimensional $X$, 
\begin{equation*}
\operatorname{Aut}(C(X)\otimes B)  \cong C(X,E)\rtimes\operatorname{Homeo}(X) 
\end{equation*}
topologically. Hence 
\begin{equation*}
M(\operatorname{Aut}(C(X)\otimes B))\cong M(\operatorname{Homeo}(X)).
\end{equation*}
If $E$ is whirly amenable, then the normal field kernel $C(X,E)$ is
whirly amenable.
\end{proposition}

\begin{proof}
\uses{lem:st-continuous-maps, prop:st-transfer}
Write $C(X)\otimes B=C(X,B)$. Since $B$ is simple, the primitive ideals are
the kernels of the evaluations at points of $X$. An automorphism induces a
homeomorphism $h$ of $X$. Its fiber maps give 
\begin{equation*}
(\alpha f)(x)=\gamma_x(f(h^{-1}(x))),\qquad  \gamma_x\in E. 
\end{equation*}
For a constant section $a$, the function $x\mapsto\gamma_x(a)=\alpha(a)(x)$
is continuous. Thus $x\mapsto\gamma_x$ is point-norm continuous. Conversely
this formula for continuous $\gamma$ and $h$ defines an automorphism.

We check the topology. If $U$ is clopen and $\left\|a\right\|=1$, then 
\begin{equation*}
\left\|\alpha(1_Ua)(x)\right\|=1_{h(U)}(x). 
\end{equation*}
Point-norm closeness less than one forces equality of these supports.
Agreement on finite clopen partitions describes the topology of $\operatorname{Homeo}(X)$,
so the quotient map is continuous. The pullback map $h\mapsto(f\mapsto
f\circ h^{-1})$ is a continuous section.

On the kernel, convergence is uniform on $X$ in each point-norm test $a\in B$, by constant sections. These tests describe the compact-open topology on $C(X,E)$. They also imply convergence on arbitrary sections: approximate the
compact range of a section by finitely many values in $B$, and use the
isometry of all fiber maps. The semidirect product identification is
therefore topological. Lemma~\ref{lem:st-continuous-maps} makes the kernel
extremely amenable, so Proposition~\ref{prop:st-transfer} applies.
The same lemma gives whirly amenability of the kernel when $E$ is whirly
amenable.
\end{proof}

If $X$ is the Cantor set, let $\mathcal{F}(X)$ be the Vietoris hyperspace of
nonempty closed subsets. Let $\Phi(X)\subseteq\mathcal{F}(\mathcal{F}(X))$
be the space of chains maximal under inclusion. Glasner--Weiss \cite[Theorem
6.1]{GlasnerWeiss2003} prove that $\Phi(X)$ is the universal minimal flow of $\operatorname{Homeo} (X)$. Thus the flow in Proposition~\ref{prop:st-field} is 
\begin{equation*}
\alpha\cdot\mathcal{C}=\{h_\alpha(F):F\in\mathcal{C}\}. 
\end{equation*}
This is a compact metrizable flow. Zero-dimensionality was used essentially
in the kernel argument; no extension to arbitrary compact bases is asserted.

\begin{corollary}\label{cor:st-geometric-bohr}
 Under the hypotheses of this section, one has 
\begin{equation*}
b\operatorname{Aut}(c_0(I,B))\cong\operatorname{Sym}(I)  \quad\text{for finite nonempty }I, 
\end{equation*}
while 
\begin{equation*}
b\operatorname{Aut}(c_0(\mathbb{N},B))=1,  \qquad  b\operatorname{Aut}(C(2^{\mathbb{N}})\otimes B)=1. 
\end{equation*}
Thus the last two automorphism groups have nontrivial universal minimal
flows but no nontrivial continuous homomorphisms to compact groups.
\end{corollary}

\begin{proof}
\uses{cor:st-direct-sum-bohr-general, prop:st-transfer}
Proposition~\ref{prop:st-transfer} reduces the three Bohr compactifications to
those of $\operatorname{Sym}(I)$, $\operatorname{Sym}(\mathbb{N})$, and $\operatorname{Homeo}(2^{\mathbb{N}})$,
respectively. The finite case is immediate.

The countable direct-sum assertion also follows, without the extreme-amenability hypothesis, from Corollary~\ref{cor:st-direct-sum-bohr-general}. For the Cantor-field group, use Tsankov's description of the irreducible
continuous unitary representations of the automorphism group of an
$\omega$-categorical structure with weak elimination of imaginaries: each is
induced from an irreducible representation of the automorphism group of a
finite algebraically closed substructure \cite[Corollary~5.2]{Tsankov2012}. The countable atomless Boolean algebra, whose
automorphism group is $\operatorname{Homeo}(2^{\mathbb{N}})$, is one of the structures
covered there, as the discussion following that corollary records. An
induced representation is finite dimensional only if the inducing subgroup
has finite index; every finite subalgebra larger than $\{0,1\}$ has
infinitely many conjugates, while the subalgebra $\{0,1\}$ gives only the
trivial representation. Hence the group has no nontrivial finite-dimensional
continuous unitary representation. By the Peter--Weyl theorem its Bohr
compactification is trivial.
\end{proof}

\begin{corollary}\label{cor:st-geometric-amenability}
\uses{cor:st-geometric-bohr, def:lo}
 For every nonempty finite or countable set 
$I$, the group $\operatorname{Aut}(c_0(I,B))$ is amenable and its universal minimal flow
is uniquely ergodic. For $I=\mathbb{N}$, the unique measure on $\operatorname{LO}(\mathbb{N})$ is the uniform random-order measure, characterised by 
\begin{equation*}
\mu\{L:L|_F=\ell\}=\frac1{|F|!} 
\end{equation*}
for every finite $F\subseteq\mathbb{N}$ and every linear order $\ell$ on $F$. In contrast, $\operatorname{Aut}(C(2^{\mathbb{N}})\otimes B)$ is nonamenable. Thus the
two infinite geometric examples in Corollary~\ref{cor:st-geometric-bohr} are
both minimally almost periodic, but lie on opposite sides of amenability.
\end{corollary}

\begin{proof}
\uses{prop:st-direct-sum, prop:st-field, prop:st-transfer}
The finitary permutations form a dense directed union of finite groups in $\operatorname{Sym}(\mathbb{N})$, so $\operatorname{Sym}(\mathbb{N})$ is amenable. Invariance forces all $|F|!$ orders on each finite set $F$ to have equal mass; these compatible
finite distributions define the random-order measure and determine it
uniquely. The finite symmetric-group case is Haar measure on the regular
finite flow. The group $\operatorname{Homeo}(2^{\mathbb{N}})$ is nonamenable \cite[remark
following Theorem~6.7]{Tsankov2012}. Apply Proposition~\ref{prop:st-transfer}, including its unique-measure assertion, to Propositions~\ref{prop:st-direct-sum} and~\ref{prop:st-field}.
\end{proof}


\subsubsection{Unitisation}

\begin{lemma}\label{lem:st-unitization}
 If $B$ is nonunital and has no characters,
restriction induces a topological isomorphism 
\begin{equation*}
\operatorname{Aut}(B^+)\cong\operatorname{Aut}(B). 
\end{equation*}
\end{lemma}

\begin{proof}
The scalar quotient $B^+\to\mathbb{C}$ is the unique character: every
character restricts to either zero or a character on $B$, and the latter is
excluded. Its kernel $B$ is therefore characteristic. Restriction is inverse
to $\alpha(b+\lambda1)=\alpha(b)+\lambda1$. Both operations are point-norm
continuous.
\end{proof}

Stable Kirchberg algebras, their sums and fields here, and nonzero
stable AF algebras have no characters. Their automorphism conclusions
therefore persist for minimal unitisations. In particular this verifies
the unitisation assertion in Theorem~\ref{thm:st-geometric}. Weighted
states on sums of unital fibers are treated in \cite{LupiniUnital2026}.

\begin{proof}[Proof of Theorem~\ref{thm:st-geometric}]
\proves{thm:st-geometric}
\uses{prop:st-direct-sum, prop:st-field, cor:st-geometric-bohr, cor:st-geometric-amenability, lem:st-unitization, cor:st-direct-sum-bohr-general}
The hypothesis holds for
$B=\mathcal{O}_2\otimes\mathcal{K}$ and $B=\mathcal{O}_3\otimes\mathcal{K}$, as recorded at the
start of this section. By Propositions~\ref{prop:st-direct-sum} and~\ref{prop:st-field}, the
universal minimal flows of $\operatorname{Aut}(c_0(I,B))$ and $\operatorname{Aut}(C(X)\otimes B)$ are those of
$\operatorname{Sym}(I)$ and $\operatorname{Homeo}(X)$, through the permutation of summand ideals and the
induced homeomorphism of $X$. For finite $I$ this is the regular flow of
$\operatorname{Sym}(I)$. For $I=\mathbb{N}$ and for the Cantor set it is $\operatorname{LO}(\mathbb{N})$ and
$\Phi(X)$, by the results of Glasner and Weiss recalled above
\cite{GlasnerWeiss2002,Glasner2002Orders} and \cite[Theorem~6.1]{GlasnerWeiss2003}.
The Bohr compactifications are computed in Corollary~\ref{cor:st-geometric-bohr}, and
amenability, unique ergodicity and nonamenability in
Corollary~\ref{cor:st-geometric-amenability}. Lemma~\ref{lem:st-unitization} and the paragraph
following it give the statements after unitisation. The final assertions, for
nonzero simple stable $C$, are Corollary~\ref{cor:st-direct-sum-bohr-general}.
\end{proof}


\subsection{Stable nuclear norm-unitary groups}\label{sec:st-unitaries}

For this section \(A\) denotes a unital coefficient algebra and
\(B=A\otimes\mathcal K\) its stabilisation. No separability of \(A\) is
assumed. In the notation of Section~1 of~\cite{LupiniStable2026}, the group under
study is \(V_B=U_1(B)_0\); the coefficient
algebra supplies the finite matrix corners, not a claim about norm
amenability of \(U(A)\).

We use throughout two standard identifications, valid for every stable
$C^*$-algebra $B$. Writing $\lambda:B^+\to\mathbb C$ for the scalar quotient,
evaluation and normalization give the topological direct product
\begin{equation}  \label{eq:st-scalar-splitting}
 U(B^+)\cong\mathbb T\times W_B,\qquad
 u\longmapsto(\lambda(u),\lambda(u)^{-1}u),
\end{equation}
with inverse $(z,w)\mapsto zw$; and stability identifies the component
group of $W_B$ with $K$-theory,
\begin{equation}  \label{eq:st-component-identification}
 W_B/V_B\cong K_1(B)_{\mathrm{disc}}
\end{equation}
\cite[Section~8.1]{Blackadar1998}, the quotient being discrete because $V_B$
is open: every sufficiently small norm neighbourhood of the identity
consists of unitaries admitting logarithms in $B$. When $B=A\otimes\mathcal K$ we
write $K_1(B)=K_1(A)$.

For the norm metric \(d(u,v)=\left\|u-v\right\|\), multiplication on either
side and inversion are isometries. A bounded family \(\mathcal F\) of
functions is uniformly equicontinuous and bounded (UEB) if it is
bounded in supremum norm and has a common modulus of uniform continuity.
For a general topological group, UEB means uniformly bounded and
equicontinuous for the right uniformity: for each \(\varepsilon>0\)
there is an identity neighbourhood \(V\) such that
\(\left|f(vg)-f(g)\right|<\varepsilon\) for all \(f\in\mathcal F\),
\(v\in V\), and \(g\in G\).
A net \(\mu_i\) of Borel probabilities is asymptotically left invariant
in the UEB topology if, for each \(g\) and each such \(\mathcal F\),
\[
 \sup_{f\in\mathcal F}
 \left|\int f(gx)\,d\mu_i(x)-\int f(x)\,d\mu_i(x)\right|\longrightarrow0.
\]
A net \((\mu_i)\) is a L\'evy net if, for every open identity
neighbourhood \(V\) and every net of Borel sets \(B_i\) with
\(\liminf_i\mu_i(B_i)>0\), one has \(\mu_i(VB_i)\to1\).
In a metric group this is the usual concentration property at every
fixed positive radius. A L\'evy group is a group admitting an increasing
directed family of compact subgroups with dense union whose normalised
Haar probabilities form a L\'evy net; for Polish groups and sequences
this is the notion already used in Section~\ref{sec:st-dynamics}.

\begin{theorem}[Concentration criterion]\label{thm:st-schneider}
\uses{def:cu-definition-1}
If a Hausdorff topological group admits a L\'evy net of Borel
probability measures converging to invariance in the UEB topology,
then it is whirly amenable.
\end{theorem}
\begin{proof}
The general assertion is \cite[Corollary 5.7]{Schneider2019}.
We give a direct proof for the compactly supported regular probabilities
used in this paper, including the assertion about all invariant measures.

First, compactness and joint continuity of a compact \(G\)-flow \(X\)
imply that, for every real \(f\in C(X)\), the orbit functions
\(g\mapsto f(gx)\), as \(x\) varies over \(X\), form a UEB family.
Fix \(x_0\in X\) and push \(\mu_i\) forward by \(g\mapsto gx_0\).
UEB convergence to left invariance makes every weak-* cluster point
of these probabilities \(G\)-invariant. Such a cluster point exists
by compactness of the space of regular probabilities on \(X\).
This proves amenability.

We next record the consequence of concentration that, for every
real UEB family \(\mathcal F\),
\begin{equation}\label{eq:st-levy-variance}
 \sup_{a\in\mathcal F}
 \int\left|a-\int a\,d\mu_i\right|^2d\mu_i\longrightarrow0.
\end{equation}
Indeed, choose a median \(m_i(a)\) of each \(a\). The two sets
\(\{a\leq m_i(a)\}\) and \(\{a\geq m_i(a)\}\) have measure at least
\(1/2\). Given \(\varepsilon>0\), choose an identity neighbourhood
\(V\) for which the oscillation of every \(a\in\mathcal F\) on each
pair \((vg,g)\) is below \(\varepsilon\). The L\'evy property makes
the measures of both \(V\)-neighbourhoods tend to one, uniformly
over all these sets. This convergence is uniform: otherwise, for some
\(c>0\), choose at a cofinal set of indices Borel sets \(B_i\) with
\(\mu_i(B_i)\geq1/2\) and \(\mu_i(VB_i)\leq1-c\), and put
\(B_i=G\) at the remaining indices, contradicting the L\'evy-net
condition. Hence
\(\mu_i\{|a-m_i(a)|>\varepsilon\}\to0\) uniformly in \(a\).
Uniform boundedness gives convergence of the second moment about
the median to zero. The variance is no larger than that second
moment, proving \eqref{eq:st-levy-variance}.

Now let \(\nu\) be any invariant regular probability measure on \(X\),
fix \(h\in G\) and real \(f\in C(X)\), and put
\[
 a_x(g)=f(gx),\qquad b_x(g)=f(hgx).
\]
These functions, for all \(x\in X\), form a single UEB family.
By left UEB invariance,
\[
 \sup_{x\in X}\left|\int b_x\,d\mu_i-\int a_x\,d\mu_i\right|\longrightarrow0.
\]
Equation~\eqref{eq:st-levy-variance} and
\[
 \int|b_x-a_x|^2d\mu_i
 \leq2\operatorname{Var}_{\mu_i}(a_x)
      +2\operatorname{Var}_{\mu_i}(b_x)
      +\left|\int b_x\,d\mu_i-\int a_x\,d\mu_i\right|^2
\]
therefore give \(\sup_x\int|f(hgx)-f(gx)|^2d\mu_i(g)\to0\).
Integrate over \(\nu\). Integration of the continuous function on
\(\operatorname{supp}\mu_i\times X\) can be interchanged, and
\(G\)-invariance of \(\nu\) yields
\[
 \int_X|f(hx)-f(x)|^2d\nu(x)
 =\int_G\int_X|f(hgx)-f(gx)|^2d\nu(x)d\mu_i(g)
 \longrightarrow0.
\]
The left side is independent of \(i\), so it is zero.
A continuous nonnegative function with zero integral vanishes on
\(\operatorname{supp}\nu\). Since real continuous functions separate
points and \(h\) was arbitrary, every point of this support is fixed
by \(G\). This proves whirly amenability.
Extreme amenability follows because every compact flow has an
invariant probability with nonempty support.
\end{proof}


\subsubsection{A finite-dimensional correction}
We use the following form of published stable uniqueness.

\begin{theorem}[Dadarlat--Eilers]\label{thm:st-de}
Let \(E\) be separable, unital and residually finite dimensional, and
let \(D\) be unital. Choose separating finite-dimensional unital
representations \(\pi_k\) of \(E\), each repeated infinitely often,
and put \(\gamma_n=\bigoplus_{k\leq n}\pi_k:E\to M_{r_n}(\mathbb C)\).
If unital nuclear homomorphisms \(\varphi,\psi:E\to D\) have equal
classes in \(KK_{\mathrm{nuc}}(E,D)\), there are
\(t_n\in U(M_{r_n+1}(D))\) such that
\[
 \left\|t_n(\varphi(a)\oplus(\gamma_n(a)\otimes1_D))t_n^*
       -(\psi(a)\oplus(\gamma_n(a)\otimes1_D))\right\|\longrightarrow0
\]
for every \(a\in E\).
\end{theorem}
\begin{proof}
This is \cite[Corollary 4.6]{DadarlatEilers2002}, in its published formulation.
Nuclearity of the source is not a hypothesis. Nuclear homotopies give
equal \(KK_{\mathrm{nuc}}\)-classes; when either algebra is nuclear,
\(KK_{\mathrm{nuc}}\) agrees with \(KK\), as explained in
\cite[Section 3.1]{DadarlatEilers2002}.
\end{proof}

\begin{lemma}[Simultaneous scalar-matrix correction]\label{lem:st-correction}
Let \(D\) be unital and nuclear. Given \(u_1,\ldots,u_m\in U_0(D)\),
\(\varepsilon>0\), and \(R\in\mathbb N\), there exist \(r\geq R\),
\(q_1,\ldots,q_m\in U(r)\), and \(t\in U(M_{r+1}(D))\) such that,
on writing
\[
 \lambda_r(P)=\operatorname{diag}(1_D,P\otimes1_D),\qquad
 U_j=\operatorname{diag}(u_j,1_r\otimes1_D),
\]
one has
\begin{equation}\label{eq:st-correction}
 \max_j\left\|t^*\lambda_r(q_j)tU_j-\lambda_r(q_j)\right\|<\varepsilon.
\end{equation}
\end{lemma}
\begin{proof}
\uses{thm:st-de}
Take the \emph{full} free-group algebra \(E=C^*(\mathbb F_m)\), with
canonical generators \(g_j\). It is residually finite dimensional:
\cite[Theorem~7]{Choi1980} proves this for \(m=2\), and the general case
reduces to it: \(\mathbb F_m\) embeds as a subgroup of \(\mathbb F_2\) for every finite
\(m\), full group $C^*$-algebras of subgroups embed, and residual finite
dimensionality passes to $C^*$-subalgebras. For \(m=1\) the algebra is
\(C(\mathbb T)\), whose evaluations separate points. Define
\(\varphi(g_j)=u_j^*\) and \(\psi(g_j)=1_D\).
The universal property gives unital homomorphisms, and nuclearity
of the target makes both maps nuclear. Paths \(u_j(s)\) from \(1\)
to \(u_j\) define a homotopy
\[
 \Phi:E\longrightarrow C([0,1],D),\qquad
 \Phi(g_j)(s)=u_j(s)^*.
\]
Its target is nuclear, so it is a nuclear homotopy.
Apply Theorem~\ref{thm:st-de} to the generators. With
\(q_j=\gamma_n(g_j)\) and \(t=t_n\), the conclusion is
\[
 \left\|t\operatorname{diag}(u_j^*,q_j\otimes1_D)t^*-\lambda_r(q_j)\right\|
 <\varepsilon.
\]
Conjugate back and multiply on the right by \(U_j\) to obtain
\eqref{eq:st-correction}. Finally, replacing \(q_j\) by \(q_j\oplus1_s\)
and \(t\) by \(t\oplus(1_s\otimes1_D)\) preserves this estimate.
Thus \(r\) can be made arbitrarily large.
\end{proof}


\subsubsection{Haar averaging by commutators}
For the data in Lemma~\ref{lem:st-correction}, define
\begin{equation}\label{eq:st-z}
 z_P=t^*\lambda_r(P)t\lambda_r(P)^*\qquad(P\in U(r)).
\end{equation}
The map is continuous and \(z_1=1\). Connectedness of \(U(r)\)
therefore puts its entire image in \(U_0(M_{r+1}(D))\).
No connectedness assumption on \(t\) is needed.

\begin{lemma}[Haar coupling]\label{lem:st-haar}
Let \(\mu\) be the image of Haar probability on \(U(r)\) under
\eqref{eq:st-z}. For every \(j\), \(\mu\) and its right \(U_j\)-translate
admit a coupling supported on pairs at norm distance less than
\(\varepsilon\). Also
\begin{equation}\label{eq:st-lipschitz}
 \left\|z_P-z_Q\right\|\leq 2\left\|P-Q\right\|_{\mathrm{op}}
                  \leq2\left\|P-Q\right\|_{\mathrm{HS}}.
\end{equation}
The Hilbert--Schmidt norm here is unnormalized.
\end{lemma}
\begin{proof}
\uses{lem:st-correction}
Each \(U_j\) commutes with every \(\lambda_r(P)\). Direct multiplication
therefore gives
\[
 \begin{split}
 z_P^*z_{Pq_j}U_j
 &=\lambda_r(P)t^*\lambda_r(q_j)t\lambda_r(Pq_j)^*U_j\\
 &=\lambda_r(P)\bigl(t^*\lambda_r(q_j)tU_j\bigr)\lambda_r(Pq_j)^*.
 \end{split}
\]
Equation~\eqref{eq:st-correction} makes this element
\(\varepsilon\)-close to \(1\). Consequently
\begin{equation}\label{eq:st-coupling}
 \sup_P\left\|z_{Pq_j}U_j-z_P\right\|<\varepsilon.
\end{equation}
Since \(P\) and \(Pq_j\) have the same Haar law,
\((z_P,z_{Pq_j}U_j)\) is the required coupling.
For \eqref{eq:st-lipschitz}, change the two occurrences of \(P\)
one at a time and use the norm-one unitary factors.
Scalar matrix inclusion in \(M_r(D)\) is isometric, independently
of \(D\).
\end{proof}

\begin{lemma}\label{lem:st-unitary-levy}
Haar probability on \(U(r)\), with unnormalized Hilbert--Schmidt
distance, is a L\'evy family as \(r\to\infty\).
\end{lemma}
\begin{proof}
The assertion for \(SU(r)\) is the classical concentration theorem;
see \cite[Example~2.5]{GiordanoPestov2007}, and \cite[Theorem~5.17]{Meckes2019} for a
quantitative form that covers \(U(r)\) directly. Let \(S\) be Haar in \(SU(r)\) and let
\(\theta\) be independent and uniform on \([-\pi,\pi]\).
Then \(P=e^{i\theta/r}S\) is Haar in \(U(r)\). Indeed its determinant
is uniform and, conditionally on it, its law is Haar on the
corresponding determinant coset. Moreover
\[
 \left\|P-S\right\|_{\mathrm{HS}}
 =\sqrt r\,|e^{i\theta/r}-1|\leq\pi/\sqrt r.
\]
This coupling transfers concentration to \(U(r)\).
\end{proof}


\subsubsection{The stable group}
\begin{theorem}\label{thm:st-v}
For every unital nuclear \(A\), the norm group
\(V_B=U_1(A\otimes\mathcal K)_0\) is whirly amenable.
\end{theorem}
\begin{proof}
\uses{lem:st-correction, lem:st-haar, lem:st-unitary-levy, thm:st-de, thm:st-schneider}
Put \(B=A\otimes\mathcal K\) and
\(e_k=1_A\otimes\sum_{\ell=1}^k e_{\ell\ell}\).
Fix \(v_1,\ldots,v_m\in V_B\), and choose paths \(v_j(s)\) from
\(1\) to \(v_j\). The set of all \(v_j(s)-1\) is norm compact in \(B\).
It follows that
\[
 x_{j,k}(s)=1+e_k(v_j(s)-1)e_k\longrightarrow v_j(s)
\]
uniformly in \(j,s\). For large \(k\) these elements are invertible,
and their polar parts have the form
\[
 b_{j,k}(s)=x_{j,k}(s)(x_{j,k}(s)^*x_{j,k}(s))^{-1/2}
          =1-e_k+a_{j,k}(s).
\]
They converge uniformly to \(v_j(s)\), start at \(1\), and
\(a_{j,k}(s)\in U(M_k(A))\). In particular, the endpoints
\(a_j=a_{j,k}(1)\) lie in \(U_0(M_k(A))\).

Given \(\eta>0\), choose \(k\) so that
\(\left\|1-e_k+a_j-v_j\right\|<\eta/2\).
Apply Lemmas~\ref{lem:st-correction} and \ref{lem:st-haar} with
\(D=M_k(A)\) and error below \(\eta/2\).
The matrix algebra \(M_{r+1}(D)\) is a larger corner of \(B\), with
its first \(D\)-block equal to the chosen \(e_k\)-corner.
Extend each \(z_P\) by the identity outside this larger corner.
Its compact image lies in \(V_B\).
Equation~\eqref{eq:st-coupling} now gives, for each \(j\), a coupling
of its Haar image \(\mu\) with the right \(v_j\)-translate of \(\mu\)
at distance below \(\eta\).

Direct the construction by triples \((F,\eta,R)\), with finite
\(F\subset V_B\) increasing, \(\eta\downarrow0\), and \(R\to\infty\),
and require \(r\geq R\). For a UEB family with common modulus
\(\omega\), the preceding couplings bound the translation error
by \(\omega(\eta)\). Thus the resulting net is asymptotically right
invariant in the UEB topology.

The maps defining its measures are uniformly \(2\)-Lipschitz for
unnormalized Hilbert--Schmidt distance by
\eqref{eq:st-lipschitz}. Lemma~\ref{lem:st-unitary-levy} and \(r\to\infty\)
show that their images form a L\'evy net for the norm uniformity
of \(V_B\). Push the measures forward by inversion.
Inversion is an isometry and converts right invariance into left
invariance (with \(g\) replaced by \(g^{-1}\)).
Theorem~\ref{thm:st-schneider} applies.

The auxiliary source \(C^*(\mathbb F_m)\) in Lemma~\ref{lem:st-correction}
is separable even when \(A\) is not; Theorem~\ref{thm:st-de} permits an
arbitrary unital target \(D=M_k(A)\). The corner approximation uses
norm compactness of finitely many paths, and the concentrating
probabilities are indexed by a net. Thus the proof requires no
separability assumption on \(A\).
\end{proof}

\begin{remark}\label{rem:st-stable-essential}
The argument takes place in arbitrarily large finite corners of
\(A\otimes\mathcal K\). It does not give norm amenability of \(U(A)\) or
of a fixed \(U(M_k(A))\). Also the matrices \(q_j\) represent
\(C^*(\mathbb F_m)\); they need not respect additional joint relations
of the original \(u_j\). For example, a null-homotopic symmetry
\(1-2p\) need not be supplemented by another symmetry.
There is consequently no assertion that projection \(K_0\)-classes
can be discarded in a fixed matrix corner.
\end{remark}

\begin{corollary}[Nuclearity of stabilisations of unital algebras]\label{cor:st-nuclearity}
\uses{def:ea}
Let \(B=A\otimes\mathcal K\), with \(A\) unital. The following are equivalent:
\begin{enumerate}
\item \(B\) is nuclear;
\item \(V_B\) is norm amenable;
\item \(V_B\) is extremely amenable in norm;
\item \(V_B\) is whirly amenable in norm;
\item \(U((A\otimes\mathcal K)^+)\) is norm amenable.
\end{enumerate}
\end{corollary}
\begin{proof}
\uses{thm:st-v}
Nuclearity passes between \(A\) and \(A\otimes\mathcal K\), so Theorem~\ref{thm:st-v} proves (i)\(\Rightarrow\)(iv), and
(iv)\(\Rightarrow\)(iii)\(\Rightarrow\)(ii) are immediate.
By \eqref{eq:st-scalar-splitting} and \eqref{eq:st-component-identification}
the two quotients of \(U(B^+)\) by \(V_B\) are abelian, so amenability of
\(V_B\) implies (v) by extension stability.
Put \(D=(A\otimes\mathcal K)^+\). Norm amenability of \(U(D)\) makes \(D\)
strongly amenable by \cite[Theorem~2.6]{AlekseevSchmidtThom2025}, hence nuclear by Ozawa's
theorem \cite[Theorem~2.2]{AlekseevSchmidtThom2025}. Its ideal \(A\otimes\mathcal K\) is nuclear, proving
(v)\(\Rightarrow\)(i). Alternatively, and directly: the norm
metric on \(U(D)\) is bi-invariant, so inversion turns a left-invariant
mean on bounded norm-uniformly continuous functions into a
right-invariant mean on the same space; Paterson's invariant function
space \(B_{22}(D)\) consists of functions uniformly continuous for the
relative weak topology \cite[Proposition~3(d)]{Paterson1991}, hence also for
the norm topology, and restricting the right-invariant mean to
\(B_{22}(D)\) and applying \cite[Theorem~1]{Paterson1991} makes \(D\)
nuclear. Paterson's criterion is stated for the weak topology on the
unitary group; the norm topology is finer, so the restriction is
legitimate.
\end{proof}

\begin{remark}\label{rem:st-ast}
\uses{cor:st-cuntz-norm, cor:st-nuclearity}
Corollary~\ref{cor:st-nuclearity} proves the implication from nuclearity to
norm amenability of the unitary group, which is the open direction of
\cite[Conjecture~1.1]{AlekseevSchmidtThom2025}, for every algebra of the form
\(D=(A\otimes\mathcal K)^+\) with \(A\) unital and nuclear. For such a \(D\) the
quotient tracial state property is automatic. Put \(B=A\otimes\mathcal K\),
and let \(I\) be a proper closed ideal of \(D\). If \(I\not\subseteq B\),
then \(D=B+I\), so \(D/I\cong B/(I\cap B)\) is stable. This is
impossible for the nonzero unital algebra \(D/I\). Hence \(I\subseteq B\),
and the scalar character \(D\to\mathbb C\) descends to a tracial state on
\(D/I\). Strong amenability of \((\mathcal O_n\otimes\mathcal K)^+\) was
already observed by Rosenberg \cite{Rosenberg1977}; the norm
amenability of its unitary group is asked in \cite[Section~4.1]{AlekseevSchmidtThom2025} and in
\cite[Question~9.3]{LupiniCuntz2026}, and
Corollary~\ref{cor:st-cuntz-norm} answers both.
\end{remark}


\subsubsection{The scalar-and-component flow}
\label{sec:st-unitary-flow}

Applied to \(B=A\otimes\mathcal K\), the identifications
\eqref{eq:st-scalar-splitting} and \eqref{eq:st-component-identification} read
\[
 U(B^+)\cong\mathbb T\times W_B,\qquad W_B/V_B\cong K_1(A)_{\mathrm{disc}};
\]
the first is also recalled in \cite[Section~5]{LupiniCuntz2026}, and the
second can be seen directly by compressing compact paths into finite
matrix corners as in the proof of Theorem~\ref{thm:st-v}.

\begin{proof}[Proof of Theorem~\ref{thm:st-nuclear-norm}]
\proves{thm:st-nuclear-norm}
\uses{cor:st-nuclearity, prop:st-transfer, thm:st-v}
These two identifications give the continuous open map
\[
 q_U:U((A\otimes\mathcal K)^+)\longrightarrow\mathbb T\times K_1(A),\qquad
 q_U(u)=(\lambda(u),[\lambda(u)^{-1}u]),
\]
with kernel \(V_B\). Theorem~\ref{thm:st-v}, the kernel reduction in
Proposition~\ref{prop:st-transfer}, and the compact-factor rule give
\[
 M(U((A\otimes\mathcal K)^+))=\mathbb T\times M(K_1(A)).
\]
The action is \(u\cdot(z,x)=(\lambda(u)z,[\lambda(u)^{-1}u]\cdot x)\).
Corollary~\ref{cor:st-nuclearity} gives the nuclearity equivalences.
\end{proof}

\begin{corollary}\label{cor:st-unitary-consequences}
Let \(A\) be unital and nuclear.
\begin{enumerate}
\item \(W_B\) is whirly amenable exactly when it is extremely amenable,
and exactly when \(K_1(A)=0\).
\item \(M(U((A\otimes\mathcal K)^+))\) is metrizable exactly when \(K_1(A)\)
is finite. In that case it is \(\mathbb T\times K_1(A)\), with translation
through \(q_U\).
\item \(bU((A\otimes\mathcal K)^+)\cong\mathbb T\times b(K_1(A))\).
\item The full group \(U((A\otimes\mathcal K)^+)\) is amenable and is never
extremely amenable or whirly amenable.
\item For every compact \(U((A\otimes\mathcal K)^+)\)-flow \(X\) and every
invariant regular probability measure \(\nu\),
\(\operatorname{supp}\nu\subseteq X^{V_B}\); the action on its support
descends through \(q_U\) to \(\mathbb T\times K_1(A)\).
\end{enumerate}
\end{corollary}
\begin{proof}
\uses{prop:st-transfer, thm:st-v}
For (i), if \(K_1(A)=0\), then \(W_B=V_B\), so use
Theorem~\ref{thm:st-v}. Conversely, use the quotient by \(V_B\);
a nontrivial discrete abelian
group has a nontrivial compact character and cannot be extremely
amenable. For (ii), a finite group is its own UMF, whereas an
infinite discrete group has nonmetrizable UMF \cite{Pestov2006}.
Statement (iii) follows from
Proposition~\ref{prop:st-transfer}. Amenability in (iv) follows by
extensions, and the nontrivial scalar quotient \(\mathbb T\) excludes
extreme amenability and hence whirly amenability.
Statement (v) is the whirly-kernel part of Proposition~\ref{prop:st-transfer}.
\end{proof}

\begin{example}\label{ex:st-example-1}
For a UCT Kirchberg algebra with \(K_1(A)=\mathbb Z/2\), the stabilised
norm-unitary UMF is
\(\mathbb T\times\mathbb Z/2\); with \(K_1(A)=\mathbb Z\) it is
\(\mathbb T\times M(\mathbb Z_{\mathrm{disc}})\).
The latter is not the Bohr compactification of \(\mathbb Z\).
\end{example}


\subsubsection{The real-rank-zero framework beyond nuclearity}
The conditional component-and-flow theorem for stabilisations of
Kirchberg algebras is \cite[Theorem~9.1]{LupiniCuntz2026}. The theorem below
extends that framework to every separable stable real-rank-zero algebra,
including nonnuclear algebras, and adds the whirly-amenability equivalence
together with separate metrizability and Bohr conclusions. Its whirliness
input is recalled from Lemma~\ref{lem:cu-stable-whirly}.

Here $A$ is again a stable algebra, and $W_A=U_1(A)$, $V_A=(W_A)_0$ in
the notation fixed in Section~1 of~\cite{LupiniStable2026}; the scalar splitting
\eqref{eq:st-scalar-splitting} and the component identification
\eqref{eq:st-component-identification} apply to it verbatim.

Polish L\'evy groups are whirly by Glasner--Tsirelson--Weiss
\cite[Theorem 1.1 and Remark 1.4]{GlasnerTsirelsonWeiss2005}. Whirliness passes to topological
groups generated by continuous images of whirly groups: for an invariant
probability $\mu$, the same closed set $\operatorname{supp}\mu$ is fixed by each image and
then by the generated closure. This support argument avoids intersections of
arbitrarily many measurable sets of full measure.

\begin{proposition}\label{prop:st-whirly}
 If $A$ is a separable stable $C^*$-algebra of real rank
zero, then the norm group $V_A=U_1(A)_0$ is whirly.
\end{proposition}

\begin{proof}
\uses{lem:cu-stable-whirly}
The proof of Lemma~\ref{lem:cu-stable-whirly} uses only stability to
place each projection circle in a copy of \(U_1(\mathcal K)\), and real rank
zero to generate the identity component by those circles. These
hypotheses hold for the present algebra; nuclearity and simplicity
are not used. The same fixed-support argument applies.
\end{proof}

\begin{remark}\label{rem:st-alternative-whirly}
\uses{lem:st-stable-u0-density, prop:st-whirly}
 Proposition~\ref{prop:st-whirly}, together
with the strict-density argument in Section~\ref{sec:st-kernel}, gives a
second proof of the whirliness half of Theorem~\ref{thm:st-kernel-ea}.
Let $A$ be a stable Kirchberg algebra. It has
real rank zero \cite[Theorem~1.2(ii)]{Zhang1992}, so
Proposition~\ref{prop:st-whirly} makes $V_A$ whirly in norm, and
Lemma~\ref{lem:st-stable-u0-density} gives $N_A=\overline{\operatorname{Ad}(V_A)}$ in the
point-norm topology. The adjoint map $V_A\to N_A$ is continuous, since
norm convergence of unitaries implies point-norm convergence of their
conjugations. If $X$ is a compact Hausdorff $N_A$-flow with invariant
regular Borel probability measure $\mu$, pulling the action back along
$\operatorname{Ad}$ gives a continuous $V_A$-action preserving $\mu$, so whirliness of
$V_A$ gives
\begin{equation*}
\operatorname{supp}\mu\subseteq\operatorname{Fix}_{\operatorname{Ad}(V_A)}(X)=\operatorname{Fix}_{N_A}(X),
\end{equation*}
the equality by density and closedness of point stabilisers. Amenability
of $N_A$ is not supplied by this argument, and is taken from
Theorem~\ref{thm:st-kernel-ea}.
\end{remark}

\begin{theorem}\label{thm:st-norm-dichotomy}
 Let $A$ be a separable stable $C^*$-algebra of
real rank zero. The following are equivalent:

\begin{enumerate}
\item $V_A$ is whirly amenable;

\item $V_A$ is extremely amenable;

\item $V_A$ is amenable;

\item $W_A$ is amenable;

\item $U(A^+)$ is amenable.
\end{enumerate}

When these conditions hold, 
\begin{align}
M(W_A)&\cong M(K_1(A)_{\mathrm{disc}}),  \notag \\
M(U(A^+))&\cong \mathbb{T}\times M(K_1(A)_{\mathrm{disc}}),
\label{eq:st-norm-umf-k1}
\end{align}
with action through the scalar and component quotients. Moreover, 
\begin{equation}  \label{eq:st-va-dichotomy}
\begin{split}
M(V_A)\text{ is metrizable} &\quad\Longleftrightarrow\quad V_A\text{ is
amenable} \\
&\quad\Longleftrightarrow\quad V_A\text{ is extremely amenable}
\\
&\quad\Longleftrightarrow\quad V_A\text{ is whirly amenable}
\quad\Longleftrightarrow\quad M(V_A)=\{*\},
\end{split}\end{equation}
and 
\begin{equation}  \label{eq:st-wa-metrizability}
M(W_A)\text{ is metrizable} \quad\Longleftrightarrow\quad \bigl[K_1(A)\text{
is finite and }W_A\text{ is amenable}\bigr].
\end{equation}
Finally, 
\begin{equation}  \label{eq:st-norm-metrizability}
M(U(A^+))\text{ is metrizable} \quad\Longleftrightarrow\quad \bigl[K_1(A)\text{ is finite and }U(A^+)\text{ is amenable}\bigr].
\end{equation}
In the metrizable case the flow in \eqref{eq:st-norm-umf-k1} is the compact
left-translation flow $\mathbb{T}\times K_1(A)$.
Furthermore, $W_A$ is whirly amenable if and only if it is extremely
amenable, and these conditions are equivalent to $K_1(A)=0$ and
amenability of $V_A$.
\end{theorem}

\begin{proof}
\uses{lem:st-locally-compact-umf, prop:st-whirly}
The scalar splitting \eqref{eq:st-scalar-splitting} shows that $U(A^+)$ is amenable
exactly when $W_A$ is amenable. Amenability passes to open subgroups and is
preserved under extensions by amenable groups \cite[Theorem~2.8]{AlekseevSchmidtThom2025}. Since 
$V_A$ is open in $W_A$ and the quotient in \eqref{eq:st-component-identification} is
abelian, $W_A$ is amenable exactly when $V_A$ is amenable. Proposition~\ref{prop:st-whirly} makes the latter equivalent to whirly amenability of $V_A$,
which implies extreme amenability. Conversely, extreme amenability implies
amenability, proving all five equivalences.

Under these conditions, the whirly amenable normal-kernel reduction and \eqref{eq:st-component-identification} give 
\begin{equation*}
M(W_A)\cong M(K_1(A)_{\mathrm{disc}}),  \qquad  M(U(A^+))\cong M(\mathbb{T}\times K_1(A)_{\mathrm{disc}}). 
\end{equation*}
The compact-factor rule recalled in Section~\ref{sec:st-dynamics} gives \eqref{eq:st-norm-umf-k1}.

The norm metric makes both $V_A$ and $W_A$ Polish SIN groups. If $M(V_A)$ is
metrizable, then $V_A$ is amenable by \cite[Corollary 3.4]{BYMT2017}; the
equivalences already proved and the definition of extreme amenability give \eqref{eq:st-va-dichotomy}. If $M(W_A)$ is metrizable, the same result makes $W_A$ amenable, while the universal minimal flow of its discrete quotient $K_1(A)$ is a metrizable factor and hence $K_1(A)$ is finite. Conversely,
amenability of $W_A$ makes $V_A$ whirly amenable, and finiteness of the
quotient makes the displayed formula for $M(W_A)$ metrizable. This proves \eqref{eq:st-wa-metrizability}.

For \eqref{eq:st-norm-metrizability}, first suppose the universal minimal flow
is metrizable. The norm metric makes $U(A^+)$ a Polish SIN group, so it is
amenable by \cite[Corollary 3.4]{BYMT2017}. The universal minimal flow of the
discrete quotient $K_1(A)$ is a factor. Since $K_1(A)$ is discrete, Lemma~\ref{lem:st-locally-compact-umf} shows that metrizability forces it to be
finite. Conversely, amenability and finiteness of $K_1(A)$ make \eqref{eq:st-norm-umf-k1} the compact metrizable flow $\mathbb{T}\times K_1(A)$.
Finally, if $W_A$ is extremely amenable, it is amenable and its discrete
quotient $K_1(A)$ is extremely amenable. A discrete group with this property
is trivial, by Lemma~\ref{lem:st-locally-compact-umf} and the compact-group
case. Thus $K_1(A)=0$ and $V_A$ is amenable. Conversely, those conditions
give $W_A=V_A$, which is whirly amenable by the equivalences above.
Whirly amenability implies extreme amenability, completing the last claim.
\end{proof}

\begin{corollary}\label{cor:st-norm-bohr}
\uses{thm:st-norm-dichotomy}
 Under the hypotheses of Theorem~\ref{thm:st-norm-dichotomy}, the Bohr compactifications are 
\begin{equation*}
bV_A=1,\qquad  bW_A\cong b(K_1(A)_{\mathrm{disc}}),\qquad  bU(A^+)\cong b(K_1(A)_{\mathrm{disc}})\times\mathbb{T}. 
\end{equation*}
These conclusions do not require amenability.
\end{corollary}

\begin{proof}
Let $f:V_A\to C$ be a continuous homomorphism into a compact group, and
replace $C$ by the closure of its image. Haar probability on $C$, for the
left-translation action pulled back along $f$, has support all of $C$.
Whirliness therefore forces every point of $C$ to be fixed, so $f$ is
trivial. Thus every compact group image of $W_A$ factors through \eqref{eq:st-component-identification}. The first two formulas follow from the
universal property of the Bohr compactification, and the last follows from
the scalar splitting \eqref{eq:st-scalar-splitting}.
\end{proof}

\begin{corollary}[The remaining stabilised Cuntz cases]\label{cor:st-cuntz-norm}
For \(3\leq n<\infty\), and for \(n=\infty\), the norm group
\(U_1(\mathcal O_n\otimes\mathcal K)\) is whirly amenable and
\[
 M(U((\mathcal O_n\otimes\mathcal K)^+))=\mathbb T.
\]
\end{corollary}
\begin{proof}
\uses{thm:cu-main, thm:st-nuclear-norm}
Each coefficient algebra is unital nuclear with \(K_1=0\), so apply
Theorem~\ref{thm:st-nuclear-norm}. The \(n=2\) calculation remains the
existing result Theorem~\ref{thm:cu-main}.
\end{proof}

Every stable Kirchberg algebra has a full unital nuclear corner.
Brown's theorem therefore lets Theorem~\ref{thm:st-nuclear-norm} apply
to all of them, without UCT. Norm amenability is no longer an unresolved
input for this class. The conditional real-rank-zero formulation above
retains its purpose for algebras outside the nuclear class.

\subsection{Proofs of statements in Section~\ref{sec:operator-algebras}}\label{sec:st-encyclopaedia-proofs}

\begin{proof}[Proof of Theorem~\ref{thm:lupini-stable}]
\proves{thm:lupini-stable}
\uses{thm:st-kk, cor:st-closed-preimages, thm:cu-stable-general, prop:cu-finite-k, cor:st-amenability-k-groups, prop:ea-normal}
Part~(1) is Theorem~\ref{thm:st-kk}; whirly amenability of $N_A$ implies its
extreme amenability. Part~(2) is Corollary~\ref{cor:st-closed-preimages} with
$L=Q_A$. For part~(3), the quotient in~(2) is the group $Q_B$ of
Theorem~\ref{thm:cu-stable-general}, and Proposition~\ref{prop:cu-finite-k}
shows that it is countable and discrete, finite exactly when both free ranks
are at most one, and gives the order of the finite group and the
nonmetrizability in the other case. When $Q_B$ is finite,
$M(\mathrm{Aut}(A))$ is the translation flow of $Q_B$, by
Theorem~\ref{thm:cu-stable-general} and Proposition~\ref{prop:ea-normal}.
Finally, $\mathrm{Aut}(A)$ is amenable if and only if
$\mathrm{Aut}(K_0(A))$ and $\mathrm{Aut}(K_1(A))$ are, by
Corollary~\ref{cor:st-amenability-k-groups}; for finitely generated groups
this means that both free ranks are at most one, as recorded in
Section~\ref{sec:st-computations}.
\end{proof}

\begin{proof}[Proof of Theorem~\ref{thm:lupini-stable-geometric}]
\proves{thm:lupini-stable-geometric}
\uses{prop:st-direct-sum, prop:st-field, prop:ea-normal, thm:glasner-weiss, thm:gw-cantor}
By Proposition~\ref{prop:st-direct-sum},
$M(\mathrm{Aut}(c_0(I,B)))\cong M(\mathrm{Sym}(I))$. For finite $I$ this is
$\mathrm{Sym}(I)$, as for every compact group
(Proposition~\ref{prop:ea-normal}), and for $I=\mathbb{N}$ it is
$\mathrm{LO}(\mathbb{N})$ by Theorem~\ref{thm:glasner-weiss}. By
Proposition~\ref{prop:st-field},
$M(\mathrm{Aut}(C(X)\otimes B))\cong M(\mathrm{Homeo}(X))$, which is $\Phi(X)$
for the Cantor set $X$ by Theorem~\ref{thm:gw-cantor}.
\end{proof}

\begin{proof}[Proof of Theorem~\ref{thm:lupini-stable-norm}]
\proves{thm:lupini-stable-norm}
\uses{thm:st-nuclear-norm}
The first two assertions are Theorem~\ref{thm:st-nuclear-norm}. For
$C=\mathcal{O}_n$ one has $K_1(\mathcal{O}_n)=0$ \cite{Cuntz1981}, so the
second factor is a point.
\end{proof}

\section{Proofs: unital Kirchberg and AF algebras}\label{sec:proofs-unital-kirchberg-af}

This section reproduces, from the preprint of the editor \cite{LupiniUnital2026} in its version of 15 September 2026, the
proofs of the results stated in Section~\ref{sec:operator-algebras}. The text is that of the
preprint, with British spelling: its sections appear as subsections, its macros
are expanded and its labels renamed, and of its introduction only the
statements of the main results are kept. Citations of the companion preprints
point to their own proof pages where the cited statement is reproduced there.
The preprint has not been refereed. The last subsection derives the statements
of Section~\ref{sec:operator-algebras} from those of the preprint.

\subsection{Main results of the preprint}\label{sec:un-main-results}

\begin{theorem}[Unital Kirchberg automorphisms]\label{thm:un-main-kirchberg}
For every unital Kirchberg algebra \(A\),
\begin{equation}\label{eq:un-main-kirchberg}
 M(\operatorname{Aut}(A))\cong S(A)\times M(Q_A^0),
\end{equation}
where \(\alpha\cdot(\rho,x)=(\rho\circ\alpha^{-1},q_A^0(\alpha)\cdot x)\).
Moreover \(M(N_A^0)=S(A)\), and \(S(A)\) is the Furstenberg boundary
of \(N_A^0\). If \(\varphi\) is pure, the restriction
\[
 \operatorname{Aut}(A,\varphi)\longrightarrow Q_A^0
\]
is an open surjection with whirly amenable kernel. If \(A\) satisfies
the UCT and \(K_1(A)=0\), then
\[
 Q_A^0\cong\operatorname{Aut}(K_0(A),[1_A])
\]
with pointwise convergence on the countable discrete group \(K_0(A)\).
\end{theorem}

\begin{theorem}[Unital AF automorphisms]\label{thm:un-main-af}
Suppose that \(A\) is separable, unital and AF, and has no nonzero
finite-dimensional representations. Then \(U(A)\) is whirly amenable
in norm, \(\overline{\operatorname{Inn}(A)}\) is whirly amenable in the point-norm
topology, and
\begin{equation}\label{eq:un-main-af}
 M(\operatorname{Aut}(A))\cong q_A^*M(\Gamma_A),
 \qquad q_A(\alpha)=K_0(\alpha).
\end{equation}
In particular \eqref{eq:un-main-af} applies to every infinite-dimensional
simple separable unital AF algebra.
Moreover, \(\operatorname{Aut}(A)\) is whirly amenable if and only if \(\Gamma_A=1\).
\end{theorem}

\subsection{Dynamical inputs}\label{sec:un-prelim}

We use continuous actions on nonempty compact Hausdorff spaces. Extreme
amenability means that every such action has a fixed point. Whirly
amenability adds amenability to the requirement that every invariant
regular probability measure have support in the fixed-point set.
The following permanence and quotient principles are used in the forms established
in \cite[Section~2]{LupiniCuntz2026} and
\cite[Section~2]{LupiniStable2026}; they are recalled as inputs.

\begin{proposition}[Permanence and quotient principles]\label{prop:un-dynamical-inputs}
\begin{enumerate}
\item\label{item:un-dense-permanence} Continuous homomorphisms with dense image preserve extreme and
whirly amenability. Products with the product topology preserve these
properties, as do dense directed unions of subgroups with the property.
\item\label{item:un-kernel-principle} An open continuous quotient \(q:G\to Q\) with extremely amenable
kernel satisfies \(M(G)=q^*M(Q)\) and \(bG=bQ\); amenability transfers
in both directions. If the kernel is whirly amenable, whirly amenability
also transfers in both directions. When the kernel is whirly amenable,
any invariant measure on a compact \(G\)-flow has support fixed by the
kernel, and the action on that support descends to \(Q\).
\item\label{item:un-compact-principle} For a compact group \(K\),
\(M(K\times L)=K\times M(L)\). A locally compact Polish group has
metrizable universal minimal flow exactly when it is compact.
\item\label{item:un-nonarch-principle} A non-Archimedean Polish group is whirly amenable exactly when it
is trivial.
\item\label{item:un-wreath-principle} An infinite permutation wreath product
\(E^I\rtimes\operatorname{Sym}(I)\), with product and pointwise topologies, has
trivial Bohr compactification.
\end{enumerate}
\end{proposition}

The right uniformity of \(G\) compares \(vg\) with \(g\) for small
\(v\). We use the corresponding quotient uniformity on \(G/H\).
The Samuel compactification is defined by bounded uniformly continuous
functions; for a group it is the greatest right ambit \(S_R(G)\).
Here a UEB family consists of bounded right uniformly continuous
functions that are uniformly bounded and share a modulus of uniform
continuity. UEB convergence to left invariance means that, for each
\(g\in G\), the differences between integration against \(\mu_i\) and
against \(g\mu_i\) tend to zero uniformly over each such family.
Schneider's criterion \cite[Corollary~5.7]{Schneider2019} implies that every
Polish L\'evy group is whirly amenable: the Haar measures of its increasing
compact subgroups concentrate and UEB-converge to invariance.

\begin{lemma}[Samuel-space criterion]\label{lem:un-samuel}
Let \(H\leq G\) be extremely amenable. Every minimal subflow of
\(\operatorname{Samuel}(G/H)\) is a model of \(M(G)\).
In particular, if the right-uniform completion of \(G/H\) is compact
and minimal, that completion is \(M(G)\).
\end{lemma}
\begin{proof}
Every compact \(G\)-flow \(X\) has an \(H\)-fixed point \(x\).
The map \(gH\mapsto gx\) is right uniformly continuous: for each
entourage of \(X\), compactness and joint continuity supply an
identity neighbourhood moving every point by that entourage.
The orbit map extends to the Samuel compactification.
If \(X\) is minimal, the image of any nonempty minimal subflow is
all of \(X\). This proves universality. In the precompact case the
Samuel compactification is the compact uniform completion.
\end{proof}

Polish automorphism groups and their closed quotients obey the open
mapping theorem: a continuous surjective homomorphism is open
\cite{BeckerKechris1996}. This will determine the topology of each
classification quotient.


\subsection{The alignment input from the Cuntz paper}\label{sec:un-alignment}

The uniform alignment theorem
Theorem~\ref{thm:cu-alignment} applies to every simple, separable,
unital, infinite-dimensional nuclear algebra. For each finite
\(F\subset A\) and \(\varepsilon>0\), it supplies finite
\(E\subset A\) and \(\delta>0\), independent of the pure states,
such that pure states \(\rho_0,\rho_1\) that are \(\delta\)-close on \(E\)
can be aligned by an automorphism moving \(F\) by less than
\(\varepsilon\).

We use the identity-component refinement of this input:
\[
 \rho_1=\rho_0\circ v^{-1},\qquad
 \max_{a\in F}\left\lVert v(a)-a\right\rVert<\varepsilon,
 \qquad v\in N_A^0.
\]
It follows from the cited construction: the corrections produced by
Proposition~\ref{prop:cu-approx-alignment} are endpoints of unitary paths
starting at one, hence lie in \(U_0(A)\), and the automorphisms of
Theorem~\ref{thm:cu-alignment} are point-norm limits of adjoint
actions of finite products of such unitaries. Its limiting adjoint
actions therefore lie in \(N_A^0\). This observation records the subgroup needed below;
the UCP-row, Gram-matrix, and alternating-construction proofs remain
in the Cuntz paper. For unitarily equivalent states we also use the
exact unitary correction of Proposition~\ref{prop:cu-approx-alignment},
or equivalently Property~1.2 of \cite{KishimotoOzawaSakai2003}, which holds for every
$C^*$-algebra by Theorems~3.1 and~4.3 of that paper.


\subsection{Pure-state stabilisers of unital Kirchberg algebras}
\label{sec:un-kirchberg-stabilizer}

Fix a unital Kirchberg algebra \(A\) and a pure state \(\varphi\), with
GNS triple \((\pi,\mathcal H,\xi)\). Use the notation
\[
 \begin{gathered}
 G=\operatorname{Aut}(A),\quad N^0=\overline{\operatorname{Ad} U_0(A)},\quad Q=G/N^0,\\
 q=q_A^0:G\longrightarrow Q,\\
 H=\operatorname{Aut}(A,\varphi),\qquad H^0=H\cap N^0,\\
 D_\varphi=\{a\in A:\varphi(a^*a)=\varphi(aa^*)=0\}.
 \end{gathered}
\]
The notation \(H^0\) denotes this intersection; it is not a definition
of the identity component of \(H\).

For \(A=\mathcal O_2\), the following hereditary-algebra identification
is Lemma~\ref{lem:cu-hereditary}; the point here is its form for every
unital Kirchberg algebra.

\begin{lemma}\label{lem:un-hereditary}
The hereditary algebra \(D_\varphi\) is a full, nonunital, stable
Kirchberg algebra, and
\[
 D_\varphi\cong A\otimes\mathcal K.
\]
Its inclusion in \(A\) induces isomorphisms on \(K\)-theory.
Furthermore, viewing \(D_\varphi^+=D_\varphi+\mathbb C1_A\) inside \(A\),
\[
 U_1(D_\varphi)
   =\{u\in U(A):\pi(u)\xi=\xi\},\qquad
 U_1(D_\varphi)_0=U_1(D_\varphi)\cap U_0(A).
\]
\end{lemma}
\begin{proof}
The algebra \(D_\varphi\) is hereditary. Kadison transitivity applied
on a two-dimensional subspace gives a nonzero selfadjoint element
annihilating \(\xi\), so \(D_\varphi\neq0\); simplicity makes it full.
The support of a pure state is a minimal projection in \(A^{**}\).
If \(D_\varphi\) had an identity \(e\), its bidual support would give
\(e=1-s_\varphi\). Thus \(s_\varphi=1-e\) would be a nonzero minimal
projection in \(A\), contrary to pure infiniteness.
Zhang's dichotomy \cite[Theorem 1.2]{Zhang1992} makes the nonunital,
\(\sigma\)-unital, purely infinite simple algebra \(D_\varphi\) stable.
Brown's full-hereditary theorem \cite[Theorem 2.8]{Brown1977} then gives
the displayed isomorphism and the \(K\)-theory equivalence.

For a unitary, \(\pi(u)\xi=\xi\) is equivalent to
\(u-1\in D_\varphi\); the condition also gives \(\pi(u^*)\xi=\xi\).
The stable component identification gives
\[
 U_1(D_\varphi)/U_1(D_\varphi)_0\cong K_1(D_\varphi).
\]
On the other hand, purely infinite \(K_1\)-injectivity gives
\(U(A)/U_0(A)\cong K_1(A)\) \cite{Cuntz1981}; see also \cite[Proposition~8.1.4]{Blackadar1998}.
Since the inclusion induces a \(K_1\)-isomorphism, the final equality
follows.
\end{proof}

\begin{lemma}[Relative inner approximation]\label{lem:un-marked-inner}
One has
\[
 H^0=\overline{\operatorname{Ad} U_1(D_\varphi)_0}.
\]
\end{lemma}
\begin{proof}
\uses{lem:un-hereditary, prop:cu-approx-alignment}
Let \(\alpha\in H^0\), and fix a finite test set \(F\subset A\) and a
tolerance. Choose the state tests \(E,\delta\) in the exact unitary
correction recalled in Section~\ref{sec:un-alignment}, with small movement
on \(F\). Take \(v\in U_0(A)\) such that \(\operatorname{Ad} v\) approximates
\(\alpha\) on \(F\cup E\). Since \(\varphi\alpha=\varphi\),
the equivalent pure states \(\varphi\circ\operatorname{Ad} v\) and \(\varphi\) are
\(\delta\)-close on \(E\).
That correction Proposition~\ref{prop:cu-approx-alignment} supplies
\(w\in U_0(A)\), almost central on \(F\),
such that \(\varphi\circ\operatorname{Ad} v=\varphi\circ\operatorname{Ad} w\).
Then \(z=vw^*\in U_0(A)\) preserves \(\varphi\), and \(\operatorname{Ad} z\)
approximates \(\alpha\) on \(F\).

Equality of the vector states of \(\xi\) and \(\pi(z^*)\xi\) in an
irreducible representation makes these vectors scalar multiples:
the intertwiner they define lies in the scalar commutant.
After multiplying \(z\) by a scalar, which does not change
\(\operatorname{Ad} z\) or its component, we have \(\pi(z)\xi=\xi\).
Lemma~\ref{lem:un-hereditary} puts \(z\) in \(U_1(D_\varphi)_0\).
The reverse inclusion in the asserted equality is immediate.
\end{proof}

\begin{theorem}\label{thm:un-marked-whirly}
The group \(H^0=\operatorname{Aut}(A,\varphi)\cap\overline{\operatorname{Ad} U_0(A)}\)
is whirly amenable in the point-norm topology, and hence extremely
amenable.
\end{theorem}
\begin{proof}
\uses{lem:un-hereditary, lem:un-marked-inner, prop:un-dynamical-inputs, thm:st-v}
By Lemma~\ref{lem:un-hereditary} and the stable nuclear norm-unitary theorem
Theorem~\ref{thm:st-v},
\(U_1(D_\varphi)_0\) is whirly amenable in norm.
Its adjoint map to the point-norm automorphism group is continuous,
since
\(\left\lVert uau^*-vav^*\right\rVert\leq2\left\lVert a\right\rVert\left\lVert u-v\right\rVert\).
Its image has dense closure \(H^0\) by Lemma~\ref{lem:un-marked-inner}.
Proposition~\ref{prop:un-dynamical-inputs}\ref{item:un-dense-permanence} gives whirly amenability of
\(H^0\).
\end{proof}

The hereditary algebra \(D_\varphi\) identifies the connected
unitaries fixing the chosen GNS vector with a stabilised norm-unitary
group. Its dense adjoint image gives the pure-state stabiliser in
\(N^0\).


\subsection{The state space and the Kirchberg product formula}
\label{sec:un-kirchberg-product}

\begin{lemma}\label{lem:un-state-minimal}
The action \(N^0\curvearrowright S(A)\),
\(\alpha\cdot\rho=\rho\circ\alpha^{-1}\), is minimal and strongly
proximal. Its pure states form a comeager orbit.
\end{lemma}
\begin{proof}
\uses{lem:cu-glimm-avoid}
Fix \(\rho\in S(A)\), a pure state \(\sigma\), a finite subset \(F\)
of the unit ball, and \(\varepsilon>0\). Pure-state excision
\cite{AkemannAndersonPedersen1986} and real rank zero \cite[Theorem 1.2]{Zhang1992} give a nonzero proper projection \(p\)
such that
\[
 \left\lVert pap-\sigma(a)p\right\rVert<\varepsilon\qquad(a\in F).
\]
For detail, excise by a norm-one positive contraction \(h\), and take
a nonzero projection in the hereditary algebra generated by
\((h-(1-\gamma))_+\). Then \(\left\lVert hp-p\right\rVert\leq\gamma\).
Choosing the excision and \(\gamma\) errors small gives the estimate;
a proper nonzero subprojection preserves it.

Choose \(m\geq2\) mutually orthogonal nonzero projections in \(A\).
One of them, \(z\), has \(\rho(z)\leq1/m\).
Pure infiniteness gives \(r\leq z\) with \(r\sim1-p\).
Set \(q=1-r\). Both \(q,r\) are nonzero and
\([q]=[p]\). Classification of nonzero projections in a purely
infinite simple algebra gives \(q\sim p\)
\cite[Corollary 6.11.9]{Blackadar1998}.
The two equivalences combine to a unitary \(v\) with \(vpv^*=q\).

We can choose \(v\in U_0(A)\). Indeed the full corner \(pAp\) has
the same \(K_1\)-group as \(A\), and its unitary group surjects onto
\(K_1\). Choose \(w\in U(pAp)\) with
\([w+(1-p)]=-[v]\) and replace \(v\) by \(v(w+(1-p))\).
This keeps \(vpv^*=q\), and \(K_1\)-injectivity puts the new
unitary in \(U_0(A)\).

For \(\rho'=\rho\circ\operatorname{Ad} v\), we have
\(t=\rho'(1-p)=\rho(r)\leq1/m\). The two off-diagonal corners are
bounded by Cauchy--Schwarz, so
\[
 |\rho'(a)-\sigma(a)|
 \leq \varepsilon+2\sqrt t+2t
 \leq \varepsilon+2m^{-1/2}+2m^{-1}\qquad(a\in F).
\]
Every \(N^0\)-orbit therefore approaches every pure state.
Pure states are weak-* dense by Glimm approximation, as used in Lemma~\ref{lem:cu-glimm-avoid} in a
faithful essential irreducible representation. This proves minimality.

For strong proximality, move the barycenter of a probability measure
on \(S(A)\) towards a pure state. Any cluster point of the translated
measures has that pure barycenter, so it must be the corresponding
Dirac measure. Finally, pure-state homogeneity by asymptotically
inner automorphisms implemented by paths starting at \(1\)
\cite[Theorem 2.1]{KishimotoOzawaSakai2003} makes the pure states one \(N^0\)-orbit.
Extreme points of a compact metrizable convex set form a \(G_\delta\).
Their density then makes this orbit comeager.
\end{proof}

\begin{proposition}\label{prop:un-state-completion}
The orbit map
\[
 N^0/H^0\longrightarrow P(A),\qquad
 nH^0\longmapsto\varphi\circ n^{-1},
\]
is a uniform isomorphism for quotient right uniformity and weak-*
uniformity. Its compact completion is \(S(A)\).
Consequently \(M(N^0)=S(A)\).
\end{proposition}
\begin{proof}
\uses{lem:un-samuel, lem:un-state-minimal, thm:un-marked-whirly}
Pure-state homogeneity gives bijectivity. For every state \(\rho\)
and every automorphism \(v\),
\[
 |(\rho\circ v^{-1})(a)-\rho(a)|
 \leq\left\lVert v^{-1}(a)-a\right\rVert=\left\lVert v(a)-a\right\rVert.
\]
This gives forward uniform continuity.
The alignment input in Section~\ref{sec:un-alignment}, including \(v\in N^0\), gives
inverse uniform continuity with tests independent of the base state.
Weak-* density of pure states identifies the completion with \(S(A)\).
Now use Theorem~\ref{thm:un-marked-whirly},
Lemma~\ref{lem:un-state-minimal}, and Lemma~\ref{lem:un-samuel}.
\end{proof}

\begin{proof}[Proof of Theorem~\ref{thm:un-main-kirchberg}]
\proves{thm:un-main-kirchberg}
\uses{lem:un-samuel, prop:un-dynamical-inputs, thm:un-marked-whirly, thm:un-uct}
The preceding proposition gives \(M(N^0)=S(A)\).
Since this flow is strongly proximal, it is also the Furstenberg
boundary of \(N^0\).
Pure-state transitivity under \(N^0\) gives \(G=N^0H\). Hence
\(q|_H:H\to Q\) is a continuous surjection, with kernel \(H^0\).
Both groups are Polish, so the Polish open mapping theorem makes
it open. Proposition~\ref{prop:un-dynamical-inputs}\ref{item:un-kernel-principle} gives \(M(H)=q^*M(Q)\).

We verify the additional uniformity needed for the full-group product.
The map
\begin{equation}\label{eq:un-quotient-product}
 G/H^0\longrightarrow P(A)\times Q,\qquad
 gH^0\longmapsto(g\cdot\varphi,q(g))
\end{equation}
is an equivariant bijection. For surjectivity, first choose an element
with the prescribed quotient and then correct its state using \(N^0\).
For injectivity, equality of both coordinates gives a difference
in \(H\cap N^0\).
The state estimate above and uniform continuity of the group quotient
give forward uniform continuity.

For the converse fix an identity neighbourhood \(V\) in \(G\), and
choose \(V_1V_2\subseteq V\). Uniform alignment gives a weak-*
entourage \(\mathcal U\) such that pairs of pure states in
\(\mathcal U\) can be aligned by \(V_1\cap N^0\).
Choose a symmetric \(\mathcal U_0\) with
\(\mathcal U_0\circ\mathcal U_0\subseteq\mathcal U\).
Compactness of \(S(A)\) gives \(V_3\subseteq V_2\) moving every
state by \(\mathcal U_0\). Since \(q\) is open, \(q(V_3)\) is an
identity neighbourhood in \(Q\).
If \((\rho',\rho)\in\mathcal U_0\) and \(t't^{-1}\in q(V_3)\),
lift the latter difference to \(v\in V_3\), and then choose
\(n\in V_1\cap N^0\) with \(nv\cdot\rho=\rho'\).
The element \(nv\in V\) carries one corresponding coset to the
other. This proves inverse uniform continuity.

It follows that
\begin{equation}\label{eq:un-samuel-product}
 \operatorname{Samuel}(G/H^0)\cong S(A)\times S_R(Q).
\end{equation}
Here is a justification of the compactification step. A bounded
uniformly continuous function on \(P(A)\times Q\) extends to
\(S(A)\times Q\), uniformly in the second coordinate.
A finite partition of unity on \(S(A)\) approximates it, uniformly,
by sums \(h_i(\rho)f_i(q)\) with \(f_i\) right uniformly continuous.
These sums extend to \(S(A)\times S_R(Q)\).
Conversely, Stone--Weierstrass gives uniform approximation of every
continuous function on the compact product by such sums.

Theorem~\ref{thm:un-marked-whirly} and Lemma~\ref{lem:un-samuel} make every
minimal subflow of \eqref{eq:un-samuel-product} universal.
The group \(N^0\) acts minimally on \(S(A)\) and trivially on
\(S_R(Q)\). Thus every nonempty closed \(G\)-invariant subset of
the product is \(S(A)\times Y\), with \(Y\) closed and
\(Q\)-invariant: each nonempty fiber contains a dense \(N^0\)-orbit
in the first coordinate. It is minimal exactly when \(Y\) is
\(Q\)-minimal. This proves the product formula and its action.
The UCT assertion follows from Theorem~\ref{thm:un-uct} below.
\end{proof}

\begin{corollary}[Invariant measures for pure-state stabilisers]\label{cor:un-marked-measures}
Let \(A\) be a unital Kirchberg algebra and \(\varphi\) a pure state.
Put \(H=\operatorname{Aut}(A,\varphi)\) and \(H^0=H\cap N_A^0\).
In every compact \(H\)-flow \(X\), any invariant regular probability
measure has support in \(X^{H^0}\), and the action on its support
descends to \(Q_A^0\) through \(q_A^0|_H\). Moreover,
\[
 H\text{ is whirly amenable}
 \quad\Longleftrightarrow\quad Q_A^0\text{ is whirly amenable}.
\]
In particular, if \(Q_A^0=1\), the entire pure-state stabiliser \(H\)
is whirly amenable.
If \(Q_A^0\) is non-Archimedean, then \(H\) is whirly amenable
exactly when \(Q_A^0=1\). In particular, when \(A\) satisfies the UCT
and \(K_1(A)=0\), this is equivalent to
\(\operatorname{Aut}(K_0(A),[1_A])=1\).
\end{corollary}
\begin{proof}
\uses{prop:un-dynamical-inputs, thm:un-marked-whirly}
The preceding proof shows that \(q_A^0|_H\) is an open surjection
with kernel \(H^0\). Apply Theorem~\ref{thm:un-marked-whirly} and
Proposition~\ref{prop:un-dynamical-inputs}\ref{item:un-kernel-principle}.
For the final assertions use Proposition~\ref{prop:un-dynamical-inputs}\ref{item:un-nonarch-principle} and the
quotient identification in Theorem~\ref{thm:un-main-kirchberg}.
\end{proof}

\begin{corollary}\label{cor:un-kirchberg-dynamics}
For a unital Kirchberg algebra \(A\):
\begin{enumerate}
\item \(\operatorname{Aut}(A)\) and \(N_A^0\) are not amenable.
\item \(bN_A^0=1\), and \(b\operatorname{Aut}(A)\cong bQ_A^0\).
\item \(M(\operatorname{Aut}(A))\) is metrizable if and only if \(M(Q_A^0)\) is.
If \(Q_A^0\) is locally compact, this is equivalent to its
compactness, and then \(M(\operatorname{Aut}(A))=S(A)\times Q_A^0\).
\end{enumerate}
\end{corollary}
\begin{proof}
A nontrivial minimal strongly proximal flow has no invariant
probability measure: translating one would leave it unchanged,
while strong proximality would make it a Dirac mass at a fixed point.
The state space here is nontrivial, so (i) follows for either group.
Every compact group image of \(N_A^0\) is a factor of its UMF,
hence strongly proximal; a compact group translation flow is
strongly proximal only when trivial. This proves (ii), using the
open quotient to factor continuous homomorphisms.
For (iii), use the product formula and compact metrizability of
\(S(A)\). A noncompact locally compact group has nonmetrizable
UMF, whereas a compact group is its own UMF \cite{Pestov2006}.
\end{proof}


\subsection{The residual unital quotient under the UCT}\label{sec:un-uct}

Assume in this section that \(A\) is a unital UCT Kirchberg algebra.
Write \(G_i=K_i(A)\), \(u=[1_A]\), and
\[
 N_A=\overline{\operatorname{Ad} U(A)},\qquad N_A^0=\overline{\operatorname{Ad} U_0(A)}.
\]
For a subgroup \(F\subseteq G_0\) containing \(u\), write
\[
 H_u(F,G_1)=\{f(u):f\in\operatorname{Hom}(F,G_1)\}.
\]
Choose increasing finitely generated subgroups \(F_n\subseteq G_0\)
containing \(u\), with union \(G_0\), and define
\begin{equation}\label{eq:un-l}
 H_n=H_u(F_n,G_1),\qquad
 L_A=\varprojlim_n G_1/H_n.
\end{equation}
Restriction gives \(H_{n+1}\subseteq H_n\), so the bonding maps
are the quotient maps. Each coordinate is discrete. Thus \(L_A\)
is an abelian Polish group. The map
\[
 \delta:G_1\longrightarrow L_A,\qquad
 x\longmapsto(x+H_n)_n
\]
has dense image: a compatible finite list of coordinates is realised
by any representative of its last coordinate.

Let \(\underline K(A)\) denote the ordinary graded \(K\)-groups
and the finite-coefficient \(K\)-groups, with their coefficient
and Bockstein operations. Put
\[
 T_A=\operatorname{Aut}_\Lambda(\underline K(A),u).
\]
Its topology is pointwise convergence on this countable collection
of countable discrete groups. Preservation of operations and of the
unit is closed, so \(T_A\) is a non-Archimedean Polish group.

\begin{theorem}[The unit obstruction]\label{thm:un-uct}
There is an exact sequence of Polish groups
\begin{equation}\label{eq:un-uct-sequence}
 1\longrightarrow L_A\longrightarrow Q_A^0
 \longrightarrow T_A\longrightarrow1.
\end{equation}
The last map is open. The first identifies \(L_A\) homeomorphically
with the closed subgroup \(N_A/N_A^0\), by
\[
 \delta([v])\longmapsto \operatorname{Ad} v\,N_A^0\qquad(v\in U(A)).
\]
If \(\alpha\in\operatorname{Aut}(A)\), conjugation by its class on \(L_A\) is
\begin{equation}\label{eq:un-l-action}
 (x_n+H_n)_n\longmapsto
 (K_1(\alpha)(x_n)+H_n)_n.
\end{equation}
In particular, if \(G_1=0\), then
\(Q_A^0\cong\operatorname{Aut}(G_0,u)\) topologically.
\end{theorem}
\begin{proof}
Ng--Ruiz prove Property~(C) for unital simple nuclear separable UCT
purely infinite algebras \cite[Theorem~3]{NgRuiz2008}, and identify
\(N_A/N_A^0\) with \eqref{eq:un-l} through precisely the displayed dense map
\cite[Theorem~2]{NgRuiz2008}. The latter assumes in addition that the natural map
\(U(A)/U_0(A)\to K_1(A)\) is an isomorphism; this holds here by pure
infiniteness, as they note, by \cite{Cuntz1981} and
\cite[Proposition~8.1.4]{Blackadar1998}.
The identification of \(\operatorname{Aut}(A)/N_A\) with \(T_A\) is the unital
Kirchberg classification theorem \cite[Theorem~4.2.1 and
Corollary~4.2.2]{Phillips2000} together with the universal multicoefficient
theorem \cite{DadarlatLoring1996} and the uniqueness theorem
\cite[Theorem~5.1(a)]{Dadarlat2005}; the positivity condition appearing in
the stably finite formulations is vacuous in the purely infinite simple
case.
These statements give algebraic exactness in \eqref{eq:un-uct-sequence}.
The action on total \(K\)-theory is continuous, since sufficiently
close projection and unitary representatives have the same class.
The induced surjection of Polish groups is open.
The topology on \(N_A/N_A^0\) agrees with its subspace topology in \(Q_A^0\),
since \(N_A\) is closed and contains \(N_A^0\).

Every automorphism \(a_1\) of \(G_1\) preserves every \(H_n\), by
composing homomorphisms \(F_n\to G_1\) with \(a_1\) and its inverse.
The identity
\(\alpha\circ\operatorname{Ad} v\circ\alpha^{-1}=\operatorname{Ad}(\alpha(v))\) proves
\eqref{eq:un-l-action} on \(\delta(G_1)\), and continuity extends it.
The intrinsic identification with \(N_A/N_A^0\) also proves independence
of the exhaustion \(F_n\).
When \(G_1=0\), \(L_A=0\), and both even UCT extension terms
vanish. The multicoefficient invariant therefore has precisely the
unit-preserving \(G_0\)-automorphisms. Equivalently this follows
directly from unital Kirchberg classification \cite{RosenbergSchochet1987,Phillips2000}.
\end{proof}

\begin{remark}\label{rem:un-remark-1}
\uses{thm:un-uct}
We cite the theorem numbering of the published Ng--Ruiz paper:
Theorems~2 and~3 correspond to Theorems~2.7 and~3.1, respectively,
in the linked 2007 preprint. In that paper the notation for
the approximately inner group denotes a closure; it does not mean only
the algebraically inner automorphisms.
Sequence~\eqref{eq:un-uct-sequence} is an extension with specified
topology and conjugation action. These data alone do not supply a
splitting or determine its extension class.
\end{remark}

\begin{proposition}[Computing \(L_A\)]\label{prop:un-l-computation}
If \(G_0\) is finitely generated, then
\[
 L_A=G_1/H_u(G_0,G_1)
\]
discretely. More explicitly, if
\[
 G_0=\mathbb Z^r\oplus\bigoplus_{j=1}^s\mathbb Z/n_j,\qquad
 u=(a_1,\ldots,a_r;b_1,\ldots,b_s),
\]
then
\begin{equation}\label{eq:un-h-formula}
 H_u(G_0,G_1)=\sum_{i=1}^r a_iG_1+
             \sum_{j=1}^s b_jG_1[n_j],
\end{equation}
where \(G_1[n]=\{x:nx=0\}\).
For arbitrary countable \(G_0\), \(L_A=0\) if \(u\) generates
a free direct summand, and \(L_A=G_1\) discretely if \(u=0\).
\end{proposition}
\begin{proof}
A finite set of generators eventually lies in \(F_n\), after which
the inverse system is constant. Evaluation of homomorphisms on
the indicated generators gives \eqref{eq:un-h-formula}.
If a homomorphism \(r:G_0\to\mathbb Z\) satisfies \(r(u)=1\), the maps
\(g\mapsto r(g)x\) realise every \(x\in G_1\), already after
restriction to each \(F_n\). If \(u=0\), all \(H_n=0\).
\end{proof}

\begin{proposition}[Finite \(K\)-groups]\label{prop:un-finite-k}
If \(G_0,G_1\) are finite, \(Q_A^0\) is finite, with
\[
 \begin{split}
 |Q_A^0|={}&|G_1/H_u(G_0,G_1)|\,
 |\operatorname{Aut}(G_0,u)|\,|\operatorname{Aut}(G_1)|\\
 &{}\cdot|\operatorname{Ext}^1(G_0,G_1)|\,|\operatorname{Ext}^1(G_1,G_0)|.
 \end{split}
\]
In this case
\[
 M(\operatorname{Aut}(A))=S(A)\times Q_A^0.
\]
The order formula is not a description of the multiplication in \(Q_A^0\).
\end{proposition}
\begin{proof}
\uses{thm:un-uct}
The even UCT kernel is the square-zero ideal
\[
 E=\operatorname{Ext}^1(G_0,G_1)\oplus\operatorname{Ext}^1(G_1,G_0).
\]
Pure extensions from finitely generated groups split, so the
pure-extension subgroup is zero here. The multicoefficient theorem
\cite{DadarlatLoring1996,Dadarlat2005} identifies the kernel of
\(T_A\to\operatorname{Aut}(G_0,u)\times\operatorname{Aut}(G_1)\) with \(1+E\), whose group
law is addition in \(E\).
Every lift of a graded automorphism is invertible: lift its inverse,
and correct the errors \(1+e\) by \(1-e\), using \(E^2=0\).
Counting this exact sequence and \eqref{eq:un-uct-sequence} gives
the formula. The UMF assertion follows from Theorem~\ref{thm:un-main-kirchberg}.
\end{proof}

\begin{example}[Two different component effects]\label{ex:un-two-components}
The range theorem for UCT Kirchberg algebras \cite[Theorem~4.2.5(1)]{Phillips2000} realises
\((G_0,u,G_1)=(0,0,\mathbb Z/2)\). Here \(T_A=1\), while \(L_A=\mathbb Z/2\).
All automorphisms are approximately inner, but \(N_A^0\) has index two:
\[
 M(\operatorname{Aut}(A))=S(A)\times\mathbb Z/2.
\]
The example shows why \(\overline{\operatorname{Inn}(A)}\) cannot simply replace
\(\overline{\operatorname{Ad} U_0(A)}\) in the unital Kirchberg product formula.
\end{example}

\begin{example}[A profinite unit obstruction]\label{ex:un-profinite}
Let \(p\) be prime, and take
\[
 (G_0,u,G_1)=(\mathbb Z[1/p],1,\mathbb Z).
\]
For \(F_n=p^{-n}\mathbb Z\), a homomorphism \(F_n\to\mathbb Z\) takes \(1\)
to a multiple of \(p^n\). Hence
\[
 L_A=\varprojlim_n\mathbb Z/p^n\mathbb Z=\mathbb Z_p.
\]
The group \(\mathbb Z[1/p]\) is torsion free, so every extension from it
by \(\mathbb Z\) is pure: the quotient has no torsion.
Also \(\operatorname{Ext}^1(\mathbb Z,\mathbb Z[1/p])=0\).
Thus the Hausdorff UCT kernel vanishes and \(T_A=\{\pm1\}\),
acting by signs on \(G_1\).
Conjugation on \(L_A\) is inversion for the nontrivial element.
If \(s\) is any lift of that element, \(s^2\in L_A\).
It commutes with \(s\), so \(s^2=-s^2\) in the additive notation
on \(L_A\). The group \(\mathbb Z_p\) has no two-torsion, whence \(s^2=1\).
Therefore
\[
 Q_A^0\cong\mathbb Z_p\rtimes\{\pm1\}.
\]
It is compact. The automorphism UMF is the metrizable space
\(S(A)\times(\mathbb Z_p\rtimes\{\pm1\})\). The unit obstruction is compact even though its defining dense group \(\mathbb Z\) is discrete and infinite.
\end{example}


\subsection{Unital AF algebras and dimension groups}\label{sec:un-af}


\subsubsection{Concentration with arbitrarily many matrix blocks}
The AF argument uses the following estimate, whose independence of the
number of factors is essential.

\begin{lemma}[Product concentration]\label{lem:un-product-concentration}
For positive integers \(d_1,\ldots,d_k\), put
\[
 K=\prod_{j=1}^k U(d_j),\quad d_*=\min_jd_j,\quad
 d_2(P,Q)=\left(\sum_j\left\lVert P_j-Q_j\right\rVert_{\mathrm{HS}}^2\right)^{1/2}.
\]
If \(f:K\to\mathbb R\) is \(1\)-Lipschitz for \(d_2\), Haar probability \(h\)
satisfies, for \(d_*>2\) and \(t>0\),
\begin{equation}\label{eq:un-product-concentration}
 h\{|f-\textstyle\int f\,dh|\geq t\}
 \leq2\exp\bigl(-(d_*-2)t^2/24\bigr).
\end{equation}
Consequently, as \(d_*\to\infty\), these groups with Haar measure
are a L\'evy family for the maximum operator-norm metric, with
no restriction on \(k\).
\end{lemma}
\begin{proof}
The estimate follows by applying \cite[Theorem 5.17]{Meckes2019}
to \(f\) and \(-f\). Its proof uses the logarithmic
Sobolev inequality for each matrix group and its tensorization
with the maximum of the factor constants.
To obtain the assertion about operator norm, observe
\[
 \max_j\left\lVert P_j-Q_j\right\rVert_{\mathrm{op}}\leq d_2(P,Q).
\]
Thus distance from any set in the maximum norm is \(1\)-Lipschitz
for \(d_2\). If the set has measure at least \(a>0\), the lower
tail estimate forces the mean of that distance to tend to zero;
the upper tail estimate makes the measure outside any fixed
positive neighbourhood tend to zero.
\end{proof}

\begin{lemma}[Growing block sizes]\label{lem:un-block-growth}
Let \(A=\overline{\bigcup_n F_n}\) be a unital AF algebra, with
\(F_n\subseteq F_{n+1}\) finite dimensional and containing \(1_A\).
Write \(F_n=\bigoplus_{j=1}^{k_n}M_{d_{n,j}}\).
Then \(A\) has no nonzero finite-dimensional representations if
and only if
\[
 \min_j d_{n,j}\longrightarrow\infty.
\]
\end{lemma}
\begin{proof}
A unital representation of \(A\) into \(M_d\) restricts to a unital
representation of each \(F_n\), so at least one block of \(F_n\)
has size at most \(d\).
Conversely, if the displayed limit fails, pass to a sequence of
levels with a block of the same size \(d\), for some fixed \(d\).
Projection onto that block gives unital homomorphisms
\(\pi_\ell:F_{n_\ell}\to M_d\).
For each fixed \(m\), the unital homomorphisms \(F_m\to M_d\)
form a compact space: there are finitely many multiplicity types,
and each type is a compact unitary orbit.
A diagonal subsequence makes \(\pi_\ell|_{F_m}\) converge for
every \(m\). The limits are compatible unital homomorphisms.
Their union extends by continuity to a unital homomorphism \(A\to M_d\),
a contradiction. A nonzero finite-dimensional representation can
always be restricted to its support and regarded as unital.
\end{proof}

\begin{theorem}\label{thm:un-af-unitary}
For a separable unital AF algebra \(A\), the following are equivalent:
\begin{enumerate}
\item \(A\) has no nonzero finite-dimensional representations;
\item the norm group \(U(A)\) is a L\'evy group;
\item \(U(A)\) is whirly amenable in norm;
\item \(U(A)\) is extremely amenable in norm.
\end{enumerate}
\end{theorem}
\begin{proof}
\uses{lem:un-block-growth, lem:un-product-concentration}
Choose \(F_n\) as in Lemma~\ref{lem:un-block-growth}. The compact groups
\[
 K_n=U(F_n)\cong\prod_jU(d_{n,j})
\]
are increasing. Their union is norm dense in \(U(A)\), by approximating
a unitary in \(F_n\) and taking its polar part.
Under (i), the minimum block size tends to infinity.
Lemma~\ref{lem:un-product-concentration} shows that the Haar measures
of \(K_n\), viewed in \(U(A)\), form a L\'evy sequence. This is
(ii). For completeness, any increasing sequence of compact groups
with dense union has asymptotically invariant Haar measures in the
UEB topology: given \(u\in U(A)\), approximate it by \(v\in K_m\).
For \(n\geq m\), Haar invariance under \(v\), together with
\(\left\lVert ux-vx\right\rVert=\left\lVert u-v\right\rVert\), bounds the error by the common
continuity modulus. The concentration criterion \cite[Corollary~5.7]{Schneider2019} proves
(ii)\(\Rightarrow\)(iii), and (iii) implies (iv).

If \(\pi:A\to M_d\) is a nonzero irreducible finite-dimensional
representation, its image is \(M_d\). Its unitary map is onto:
every matrix unitary is \(e^{ih}\), and a selfadjoint lift of \(h\)
can be exponentiated in \(A\).
Thus \(U(A)\) has a continuous surjection onto the nontrivial
compact group \(U(d)\). Its translation action has no fixed point,
contradicting (iv). Any nonzero finite-dimensional representation
has an irreducible summand. This proves (iv)\(\Rightarrow\)(i).
\end{proof}

The concentration arguments for the UHF and fixed-width AF examples
treated in \cite[Proposition~5.1]{EFHKKL2016} are special cases.
Here the number of summands at the \(n\)-th stage may grow arbitrarily
fast. An estimate obtained by a union bound over those summands
would not suffice; the tensorized concentration bound does.


\subsubsection{Automorphisms and the scaled dimension group}
For every separable unital AF algebra \(A\), Elliott classification and
finite-dimensional uniqueness give a surjection
\begin{equation}\label{eq:un-af-quotient}
 q_A:\operatorname{Aut}(A)\longrightarrow\Gamma_A=\operatorname{Aut}(\mathsf D_A),\qquad
 q_A(\alpha)=K_0(\alpha),
\end{equation}
with kernel
\[
 N_A=\overline{\operatorname{Inn}(A)}
 =\overline{\{\operatorname{Ad}(u):u\in U(A)\}}
\]
\cite{Elliott1976}; see also \cite[Exercise~7.7.5]{Blackadar1998}. In particular the inner groups used in
this formula are not equipped with their discrete topology.

\begin{proposition}\label{prop:un-af-quotient}
The map \eqref{eq:un-af-quotient} is continuous and open.
Its target is a non-Archimedean Polish group. Its kernel is always
amenable. If \(A\) has no nonzero finite-dimensional representations,
that kernel is whirly amenable in the point-norm topology.
\end{proposition}
\begin{proof}
\uses{prop:un-dynamical-inputs, thm:un-af-unitary}
Represent \(K_0\)-classes by projections in matrix algebras over \(A\).
Sufficiently close projections are equivalent, so their classes are
locally constant under point-norm changes in the automorphism.
This proves continuity.
The countable dimension group, with its addition, order and distinguished unit,
is a countable structure; its automorphism group is a closed
permutation group. Surjectivity and the kernel description are
the existence and uniqueness parts of AF classification. The
Polish open mapping theorem gives openness.

The unitary group \(U(A)\) is the norm closure of an
increasing union of compact finite-dimensional unitary groups,
and hence is amenable. Passing through the adjoint map and taking its
closure proves amenability of \(N_A\). When \(A\) has no nonzero
finite-dimensional representations, use
Theorem~\ref{thm:un-af-unitary} and the continuous adjoint map
\(U(A)\to N_A\), whose image is dense. Proposition~\ref{prop:un-dynamical-inputs}\ref{item:un-dense-permanence}
gives whirly amenability of \(N_A\).
\end{proof}

\begin{proof}[Proof of Theorem~\ref{thm:un-main-af}]
\proves{thm:un-main-af}
\uses{prop:un-af-quotient, prop:un-dynamical-inputs}
Combine Proposition~\ref{prop:un-af-quotient} with
Proposition~\ref{prop:un-dynamical-inputs}\ref{item:un-kernel-principle};
for the whirly-amenability criterion also use
Proposition~\ref{prop:un-dynamical-inputs}\ref{item:un-nonarch-principle}.
If \(A\) is simple, unital and infinite dimensional, any nonzero
finite-dimensional representation would be injective, which is
impossible; thus the representation hypothesis holds.
\end{proof}

\begin{corollary}\label{cor:un-af-consequences}
\uses{thm:un-main-af}
For the unital AF algebras in Theorem~\ref{thm:un-main-af},
\[
 b\operatorname{Aut}(A)\cong b\Gamma_A,\qquad
\operatorname{Aut}(A)\text{ amenable}\ \Longleftrightarrow\
\Gamma_A\text{ amenable}.
\]
In every compact \(\operatorname{Aut}(A)\)-flow \(X\), an invariant regular
probability measure has support in \(X^{N_A}\), and the action on its
support descends to \(\Gamma_A\) through \(q_A\).
If \(\Gamma_A\) is compact, \(M(\operatorname{Aut}(A))\) is its translation
flow through \(q_A\). If \(\Gamma_A\) is infinite discrete, that
UMF is nonmetrizable.
\end{corollary}
\begin{proof}
\uses{prop:un-dynamical-inputs}
Apply Proposition~\ref{prop:un-dynamical-inputs}\ref{item:un-kernel-principle}--\ref{item:un-compact-principle}.
\end{proof}


\subsubsection{Finite-dimensional algebras and finite sums of simple algebras}
\begin{lemma}\label{lem:un-finite-af}
Let \(A\) be finite dimensional and write
\(A=\bigoplus_{d\in J}M_d(\mathbb C)^{m_d}\), where \(J\) is a finite set of
distinct positive integers. Then
\[
 \operatorname{Aut}(A)\cong
 \prod_{d\in J}\left(\operatorname{PU}(d)^{m_d}
                    \rtimes\operatorname{Sym}(m_d)\right)
\]
topologically, and in particular \(\operatorname{Aut}(A)\) is compact.
\end{lemma}
\begin{proof}
An automorphism permutes isomorphic simple summands. On each
matrix summand it is inner, with unitary implementer unique up
to a scalar. This gives the indicated continuous isomorphism;
both sides are compact Hausdorff. 
\end{proof}

The finite-dimensional case of the next proposition then returns
\(M(\operatorname{Aut}(A))=\operatorname{Aut}(A)\) itself, a compact group acting on itself by left
translation.

\begin{proposition}[Finite sums of simple unital AF algebras]\label{prop:un-finite-simple-sum}
Let
\[
 A=\bigoplus_{t=1}^s B_t^{m_t}
\]
where the \(B_t\) are pairwise nonisomorphic simple unital AF
algebras. Define
\[
 R_t=
 \begin{cases}
 \operatorname{PU}(d),& B_t=M_d(\mathbb C),\\
 \operatorname{Aut}(\mathsf D_{B_t}),& B_t\text{ infinite dimensional}.
 \end{cases}
\]
There is an open continuous quotient with whirly amenable kernel
from \(\operatorname{Aut}(A)\) onto
\[
 R_A=\prod_{t=1}^s (R_t^{m_t}\rtimes\operatorname{Sym}(m_t)).
\]
Thus \(M(\operatorname{Aut}(A))=q^*M(R_A)\).
If all the \(R_t\) are compact, the answer is \(R_A\) itself.
\end{proposition}
\begin{proof}
\uses{lem:un-finite-af, prop:un-dynamical-inputs}
Permuting isomorphic simple summands identifies
\(\operatorname{Aut}(A)=\prod_t(\operatorname{Aut}(B_t)^{m_t}\rtimes\operatorname{Sym}(m_t))\).
For an infinite-dimensional \(B_t\), take the dimension-group
quotient from Theorem~\ref{thm:un-main-af}; for a matrix algebra,
take its identity quotient, \(\operatorname{Aut}(M_d(\mathbb C))=\operatorname{PU}(d)\) being
Lemma~\ref{lem:un-finite-af}.
Their product respects the permutations. Its kernel is a finite
product of whirly amenable groups, hence whirly amenable by
Proposition~\ref{prop:un-dynamical-inputs}\ref{item:un-dense-permanence}.
It is open, and Proposition~\ref{prop:un-dynamical-inputs}\ref{item:un-kernel-principle} applies.
\end{proof}


\subsection{What remains for arbitrary separable unital AF automorphisms}
\label{sec:un-af-gap}

The representation hypothesis in Theorem~\ref{thm:un-af-unitary} has an
exact meaning for unitary groups. For automorphism groups it is a
sufficient hypothesis for discarding the approximately inner kernel,
but it is not necessary. Commutative AF algebras have trivial inner
groups, despite having many characters. At the other extreme,
matrix algebras retain nontrivial compact inner groups.
In general these behaviours occur over a nontrivial space of ideals.

\begin{proposition}[A finite-dimensional obstruction in the kernel]\label{prop:un-fd-obstruction}
If a unital AF algebra \(A\) has an irreducible representation
\(\pi:A\to M_d\) with \(d\geq2\), then \(N_A=\overline{\operatorname{Inn}(A)}\)
has a continuous surjection onto \(\operatorname{PU}(d)\).
In particular \(N_A\) is not extremely amenable.
\end{proposition}
\begin{proof}
Every inner automorphism preserves \(I=\ker\pi\).
A point-norm limit of inner automorphisms, and its inverse, preserve
\(I\), so every element of \(N_A\) induces an automorphism of \(A/I\).
This gives a continuous map \(N_A\to\operatorname{Aut}(M_d)=\operatorname{PU}(d)\):
to test convergence on a matrix, choose any lift in \(A\).
It is onto because every matrix unitary lifts to a unitary of \(A\),
by lifting a selfadjoint logarithm and exponentiating.
A nontrivial compact translation quotient prevents extreme amenability.
\end{proof}

Even this obstruction does not describe all residual dynamics of
\(N_A\). A basic test class is a continuous matrix bundle over a
zero-dimensional compact space.

\begin{proposition}[The homogeneous AF test class]\label{prop:un-gauge}
Let \(X\) be a nonempty compact metrizable zero-dimensional space
and \(d\geq2\). For \(A=C(X,M_d)\), one has topological isomorphisms
\[
 \operatorname{Aut}(A)\cong C(X,\operatorname{PU}(d))\rtimes\operatorname{Homeo}(X),
 \qquad
 \Gamma_A\cong\operatorname{Homeo}(X),
\]
and \(N_A=C(X,\operatorname{PU}(d))\).
The function group has the uniform topology, and the homeomorphism
group has the compact-open topology.
\end{proposition}
\begin{proof}
The action of an automorphism on the centre \(C(X)\) specifies a
homeomorphism. After composing with its inverse action on the base,
the automorphism acts fiberwise by a continuous map
\(X\to\operatorname{Aut}(M_d)=\operatorname{PU}(d)\).
Continuity can be tested on a finite matrix-unit basis, which also
identifies the uniform topology on this kernel with its point-norm
topology on \(A\).
Conversely these data give an automorphism, continuously in both
coordinates, with multiplication given by the indicated semidirect
product.
The dimension group is \(C(X,\mathbb Z)\) with pointwise order and order
unit the constant function \(d\). Its order-unit automorphisms are
the homeomorphisms of \(X\): they preserve the unique element \(1\)
with \(d1=u\), and hence the Boolean algebra of characteristic
functions in the interval \([0,1]\). Stone duality gives the claim.
Finally, finite locally constant \(\operatorname{PU}(d)\)-valued
maps lift to unitary-valued maps. They are uniformly dense in
\(C(X,\operatorname{PU}(d))\), by zero dimensionality and compactness.
They act by inner automorphisms, proving the kernel assertion.
\end{proof}

When \(X\) is infinite, the uniform function group in
Proposition~\ref{prop:un-gauge} is not compact. Choose infinitely
many pairwise disjoint nonempty clopen subsets of \(X\), and
functions taking a fixed nonidentity value on one of them and
the identity elsewhere. They are uniformly separated.
Thus the residual issue cannot in general be handled merely by
adjoining a finite-dimensional compact inner factor.

\begin{remark}[Question]\label{rem:un-general-af}
Compute the universal minimal flow of \(\operatorname{Aut}(A)\) for every separable
unital AF algebra that has finite-dimensional representations.
Already for \(A=C(X,M_d)\), find an effective description of
\[
 M\bigl(C(X,\operatorname{PU}(d))
             \rtimes\operatorname{Homeo}(X)\bigr)
\]
for arbitrary compact metrizable zero-dimensional \(X\).
Describe how this information combines with the dimension-group
quotient in a general Bratteli system.
\end{remark}

Thus Problem~\ref{rem:un-general-af} is precisely the part of the
arbitrary-unital-AF case not covered by Theorem~\ref{thm:un-main-af};
the general realisation of a UMF as a minimal Samuel subflow does
not resolve this AF-specific problem.

\subsection{Proofs of statements in Section~\ref{sec:operator-algebras}}\label{sec:un-encyclopaedia-proofs}

\begin{proof}[Proof of Theorem~\ref{thm:lupini-unital}]
\proves{thm:lupini-unital}
\uses{thm:un-main-kirchberg}
This is part of Theorem~\ref{thm:un-main-kirchberg}.
\end{proof}

\begin{proof}[Proof of Theorem~\ref{thm:lupini-af}]
\proves{thm:lupini-af}
\uses{thm:un-main-af, prop:un-finite-simple-sum}
The first assertions, including the case of infinite-dimensional simple
separable unital AF algebras, are Theorem~\ref{thm:un-main-af}. Let
$A=M_d(\mathbb{C})$ with $d\geq2$. Its scaled dimension group is
$(\mathbb{Z},\mathbb{Z}_{\geq0},d)$, whose only automorphism is the identity,
so $\Gamma_A$ is trivial. Every automorphism of $M_d(\mathbb{C})$ is inner, so
$\mathrm{Aut}(A)=PU(d)$, and Proposition~\ref{prop:un-finite-simple-sum}, with
a single summand, gives $M(\mathrm{Aut}(A))\cong PU(d)$, which is not a
point.
\end{proof}

\section{Proofs: inductive limits of compact Lie groups}\label{sec:proofs-inductive-limits}

This section reproduces, from the preprint of Codenotti and Lupini
\cite{LupiniInductive2026} in its version of 13 September 2026, the proofs of
the results stated in Section~\ref{sec:inductive-limits}. The text is that of
the preprint, with British spelling: its sections appear as subsections, its
macros are expanded and its labels renamed, and its introduction is reduced to
the statements of the main results and the notation they use. The preprint has
not been refereed. The last subsection derives the statements of
Section~\ref{sec:inductive-limits} from those of the preprint.

\subsection{Main results of the preprint}\label{sec:ind-main-results}

\begin{theorem}\label{thm:ind-main}
\uses{def:final-topology}
Let $K_1\subseteq K_2\subseteq\cdots$ be compact Lie groups, and give
$G=\bigcup_nK_n$ the final topology. Suppose that, for some $j$,
$K_j^0$ is nonabelian or $K_j/K_j^0$ is nonsolvable. Then $G$ has a
jointly continuous action on a nonempty compact metrizable space $X$
such that $X^{K_j}=\varnothing$. In particular, $G$ is not extremely
amenable.
\end{theorem}

Here $K^0$ denotes the identity component. For a compact Lie group,
$K/K^0$ is finite, and an abelian identity component is a torus, with
the trivial group allowed as a zero-dimensional torus. The hypothesis on
$K_j$ therefore says exactly that $K_j$ is not solvable.

\begin{corollary}\label{cor:ind-classical-intro}
\uses{def:ea}
With the standard block embeddings and final topologies, none of
\[
\varinjlim\operatorname{SU}(n),\qquad \varinjlim\operatorname{SO}(n),\qquad
\varinjlim\operatorname{Sp}(n),\qquad \varinjlim\operatorname{U}(n),\qquad
\varinjlim\operatorname{O}(n)
\]
is extremely amenable. The notation $\operatorname{Sp}(n)$ denotes the compact group
of quaternionic unitary matrices.
\end{corollary}

Throughout, $\operatorname{SU}(\infty)$ denotes $\varinjlim\operatorname{SU}(n)$, and analogously
for $\operatorname{SO}(\infty)$, $\operatorname{Sp}(\infty)$, $\operatorname{U}(\infty)$, and $\operatorname{O}(\infty)$.
Unless another topology is explicitly specified, these symbols carry
the final topology for the standard block embeddings.

The proof of Theorem~\ref{thm:ind-main} uses the full fixed-set prescription
in Oliver's Theorem~4 and Proposition~8, combined in his Theorem~5
\cite{Oliver1976}. Its role is
to provide equivariant maps between successive disks. An inverse limit
then carries the action of the whole union. Ordinary existence of a
fixed-point-free disk for each stage would not by itself provide those
maps. Section~\ref{sec:ind-observables} gives a finite-observable version
of the obstruction.

Write $M(G)$ for the universal minimal flow and $bG$ for the Bohr
compactification of a topological group $G$; the definitions are recalled
in Sections~\ref{sec:ind-bohr} and~\ref{sec:ind-umf}.

\begin{theorem}\label{thm:ind-dynamics-intro}
\uses{cor:ind-classical-intro, thm:ellis}
For each of the five classical limits in
Corollary~\ref{cor:ind-classical-intro}, $M(G)$ contains a homeomorphic copy
of the Stone--\v{C}ech compactification $\beta\mathbb N$ of the discrete natural
numbers. In particular, $M(G)$ is nonmetrizable. Moreover,
\begin{gather*}
b\operatorname{SU}(\infty)=b\operatorname{SO}(\infty)=b\operatorname{Sp}(\infty)=\{1\},\\
b\operatorname{U}(\infty)\cong\mathbb T,\qquad b\operatorname{O}(\infty)\cong\{1,-1\},
\end{gather*}
where the last two Bohr maps are the determinants.
\end{theorem}

The Bohr calculation uses only finite-dimensional representations.
The assertion about $M(G)$ uses the compact metrizable actions above,
a separation lemma specific to the final topology, and a criterion of
Basso--Zucker valid for arbitrary topological groups
\cite{BassoZucker2021}. A compact metrizable action without fixed points
alone would not imply nonmetrizability of $M(G)$.

The universal minimal actions also exhibit a distinction between
freeness everywhere and freeness on a large subset. In
Section~\ref{sec:ind-free-locus} we prove the following result, whose first
assertion answers Pestov's Question~4.

\begin{theorem}\label{thm:ind-free-intro}
\uses{thm:ellis}
For each of the five classical limits, every jointly continuous action
on a nonempty compact Hausdorff space has a point with nontrivial
stabiliser. Consequently, none of these $k_\omega$-groups admits a free
action on a compact space. Nevertheless, the free locus in $M(G)$ is a
proper dense $G_\delta$ and has full measure for every invariant Radon
probability measure.
\end{theorem}

Question~4 asks whether Veech's theorem, that every locally compact
group admits a free action on a compact space, extends to
$k_\omega$-groups~\cite[p.~440]{Pestov2007};
Theorem~\ref{thm:ind-free-intro} shows that it does not.
Nonfreeness follows from McGibbon's theorem on maps from infinite
loop spaces and the Zabrodsky lemma. The assertion about the free
locus uses the stage stabilisers supplied by the Oliver construction
and a generic-vector argument.

Sections~\ref{sec:ind-relative}--11 of~\cite{LupiniInductive2026} and
Appendix~\ref{sec:ind-separation} compare these final-topology results with
a coarser trace-norm topology. Let
$G_{\mathrm{alg}}=\bigcup_n\operatorname{SU}(n)$ denote the canonical algebraic union. We write
$G_{\mathrm{fin}}$ for $G_{\mathrm{alg}}$ with the final topology and $G_{\mathrm{tr}}$ for
$G_{\mathrm{alg}}$ with the topology of the unnormalized trace-norm metric
$d_1(g,h)=\operatorname{Tr}\lvert g-h\rvert$. The completion of $G_{\mathrm{tr}}$ is
\begin{equation}\label{eq:ind-suone}
\operatorname{SU}_1=\{u\in\operatorname{U}(\ell^2):u-I\in\mathcal S_1,\ \det_F(u)=1\},
\end{equation}
where $\det_F$ is the Fredholm determinant.

Throughout we write $e$ for the identity element of a group and $I$ for
the identity operator on a Hilbert space. Identity components of compact
Lie groups are written $K^0$. The symbol $\mathcal F(K)$ is reserved for
Oliver's family of closed solvable subgroups. Averaging over a compact
subgroup $K$ is written $\operatorname{Avg}_K$; the orthogonal projection onto
$K$-fixed vectors in a unitary representation is written $\operatorname{P}^{\mathrm{fix}}_K$.

For a compact group $K$ with normalised Haar measure, we use
\begin{equation}\label{eq:ind-variance-convention}
\operatorname{Var}_{k\in K}\phi(k)
:=\int_K\lvert\phi(k)-\int_K\phi(k')\,\mathrm d k'\rvert^2\,\mathrm d k.
\end{equation}

\begin{theorem}\label{thm:ind-trace-intro}
\uses{def:ea}
The groups $G_{\mathrm{tr}}$ and $\operatorname{SU}_1$, each with its trace-norm topology, are
extremely amenable. More precisely, if
$N>m\ge2$ and $f:\operatorname{SU}(N)\to\mathbb R$ is $L$-Lipschitz for $d_1$, then
\begin{equation}\label{eq:ind-variance-intro}
\int_{\operatorname{SU}(N)}\operatorname{Var}_{k\in\operatorname{SU}(m)}f(kh)\,\mathrm d h
\le \frac{2mL^2}{N-m}.
\end{equation}
All Haar measures are normalised to have mass one.
\end{theorem}

The estimate concerns variance along a fixed subgroup. It does not
assert concentration of $f$ around its mean over the entire large
group. The final topology in Theorem~\ref{thm:ind-main} is strictly finer;
Appendix~\ref{sec:ind-separation} gives an explicit function separating
the two right uniformities.

\subsection{Inverse limits of compact actions}\label{sec:ind-inverse}

We first isolate the topological construction. No smoothness is needed
in this section.

\begin{lemma}\label{lem:ind-inverse-action}
Let $G=\bigcup_{n\ge1}K_n$ be a topological group with the final
topology for an increasing sequence of topological groups. Suppose
that $X_n$ is a nonempty compact Hausdorff $K_n$-space and that
\[
p_n:X_{n+1}\longrightarrow X_n
\]
is continuous and $K_n$-equivariant, where the action on $X_{n+1}$ is
restricted from $K_{n+1}$. Then $X=\varprojlim(X_n,p_n)$ is a nonempty
compact Hausdorff $G$-space. The projection $X\to X_n$ is
$K_n$-equivariant. If every $X_n$ is metrizable, so is $X$.
\end{lemma}

\begin{proof}
The compatibility equations $p_n(x_{n+1})=x_n$ define closed subsets
of the compact product $\prod_nX_n$. Every finite collection of these
equations is satisfiable: choose the largest relevant coordinate and
apply the bonding maps downwards, then fill later coordinates
arbitrarily. The finite intersection property gives $X\ne\varnothing$.
This argument does not require any bonding map to be surjective.

For $n<m$, write $p_{n,m}=p_n\circ\cdots\circ p_{m-1}$, and set
$p_{m,m}=\operatorname{id}$. If $g\in K_m$, define
\begin{equation}\label{eq:ind-tail-action}
(g\cdot x)_n=
\begin{cases}
g\cdot x_n,&n\ge m,\\
p_{n,m}(g\cdot x_m),&n<m.
\end{cases}
\end{equation}
For $n\ge m$, equivariance of $p_n$ preserves the compatibility
equations. At earlier coordinates compatibility follows by composition.
If $r\ge m$, then for $m\le n\le r$,
\[
p_{n,r}(g\cdot x_r)=g\cdot p_{n,r}(x_r)=g\cdot x_n.
\]
Thus the formula is independent of the chosen stage containing $g$.
To check the action law, choose a common stage containing two given
group elements and use the action law on all later coordinates;
earlier coordinates follow by applying the bonding maps.

For a fixed $m$, every coordinate of the map $K_m\times X\to X$ in
\eqref{eq:ind-tail-action} is continuous. Hence each stage action is
jointly continuous and induces a continuous homomorphism
\[
\rho_m:K_m\longrightarrow\operatorname{Homeo}(X),
\]
where $\operatorname{Homeo}(X)$ has the compact-open topology. For completeness,
if $C\subseteq X$ is compact, $U\subseteq X$ is open, and
$g_0C\subseteq U$, joint continuity and compactness of $C$ give a
neighbourhood $V$ of $g_0$ with $VC\subseteq U$. This proves
continuity for the subbasic compact-open sets.

The maps $\rho_m$ agree on overlaps, and finality makes the assembled
homomorphism $\rho:G\to\operatorname{Homeo}(X)$ continuous. Evaluation
$\operatorname{Homeo}(X)\times X\to X$ is continuous because $X$ is compact
Hausdorff. Therefore the action of $G$ is jointly continuous.
Projection to $X_n$ is $K_n$-equivariant by the formula. Finally, a
closed subspace of a countable product of compact metrizable spaces
is compact metrizable.
\end{proof}

\begin{remark}\label{rem:ind-remark-1}
Joint continuity above follows by passing through $\operatorname{Homeo}(X)$.
It is not inferred merely from continuity of individual orbit maps.
\end{remark}


\subsection{Oliver disks and the general obstruction}\label{sec:ind-oliver}

For a compact Lie group $K$, define
\begin{equation}\label{eq:ind-family}
\mathcal F(K)=\{H\le K:H\text{ is closed},\ H^0\text{ is a torus},
\ H/H^0\text{ is solvable}\}.
\end{equation}
Equivalently, $\mathcal F(K)$ is the family of closed solvable subgroups of
$K$: a compact Lie group $H$ is solvable if and only if $H^0$ is a torus
and $H/H^0$ is solvable. The condition is intrinsic to $H$. In
particular, if $K\le L$ is a closed subgroup inclusion, then for every
closed $H\le K$,
\begin{equation}\label{eq:ind-intrinsic}
H\in\mathcal F(K)\quad\Longleftrightarrow\quad H\in\mathcal F(L).
\end{equation}

\begin{theorem}[Oliver]\label{thm:ind-oliver}
For every compact Lie group $K$, there is a smooth action on a compact
disk $D_K$ such that, for every closed subgroup $H\le K$,
\begin{equation}\label{eq:ind-fixed-sets}
D_K^H\text{ is a disk if }H\in\mathcal F(K),
\qquad D_K^H=\varnothing\text{ otherwise}.
\end{equation}
\end{theorem}

\begin{proof}[Reference]
We recall Oliver's terminology. A compact Lie group is \emph{$p$-toral}
if its identity component is a torus and its component group is a
$p$-group, and \emph{prime toral} if it is $p$-toral for some prime $p$
\cite[p.~80]{Oliver1976}. A family $\mathcal G$ of closed subgroups is
\emph{separating} if, for every pair $H\trianglelefteq L$ of closed
subgroups with $L/H$ prime toral, either both $H$ and $L$ lie in
$\mathcal G$ or neither does; \emph{closed} refers to the Hausdorff
topology on the space of closed subgroups, not to closure under passage
to subgroups~\cite[pp.~92--93]{Oliver1976}. Proposition~8 of
\cite[p.~95]{Oliver1976} states that \eqref{eq:ind-family} is a closed
separating family. It is nonempty because it contains the trivial
subgroup, and it is conjugation invariant because its defining condition
is intrinsic. Theorem~4 of~\cite[p.~94]{Oliver1976} gives precisely
\eqref{eq:ind-fixed-sets}. For $K$ as in Theorem~\ref{thm:ind-main}, the two
statements are combined in Theorem~5 of~\cite[p.~96]{Oliver1976}.
\end{proof}

\begin{lemma}\label{lem:ind-bonding}
\uses{thm:ind-oliver}
If $K\le L$ are compact Lie groups and $D_K,D_L$ are chosen as in
Theorem~\ref{thm:ind-oliver}, there is a continuous $K$-equivariant map
$D_L\to D_K$.
\end{lemma}

\begin{proof}
\uses{lem:ind-inverse-action}
The restricted smooth action of $K$ on $D_L$ admits a $K$-CW
structure, by equivariant triangulation for smooth compact Lie group
actions, including the boundary case
\cite[Section~7, especially Theorem~7.1 and Corollary~7.2]{Illman1983}.
Illman uses his own notion of an equivariant CW complex, which differs in
general from Matumoto's~\cite[Remark~6.2]{Illman1983}; the two agree for a
smooth compact manifold, and all that is used below is the decomposition
into equivariant cells $K/H\times B^q$ with closed stabilisers $H$.
For a cell stabiliser $H$, the fixed set $D_L^H$ is nonempty. Hence
$H\in\mathcal F(L)$, and \eqref{eq:ind-intrinsic} gives $H\in\mathcal F(K)$.

Construct the map cell by cell. On a zero-dimensional equivariant
cell $K/H$, choose a point of the nonempty space $D_K^H$.
For a cell $K/H\times B^q$ with $q>0$, an equivariant map already
defined on its boundary is determined by a map
$\partial B^q\to D_K^H$. Since $D_K^H$ is contractible, this map
extends over $B^q$. The weak topology of the equivariant CW structure
gives a continuous map on the whole space. The resulting map need not
be surjective; Lemma~\ref{lem:ind-inverse-action} does not require
surjective bonding maps.
\end{proof}

\begin{remark}\label{rem:ind-remark-2}
\uses{lem:ind-bonding}
The two $K$-disks in Lemma~\ref{lem:ind-bonding} have exactly the same
fixed-set prescription. Oliver observes, without proof, that two
disk actions corresponding to the same family are equivariantly homotopy
equivalent~\cite[p.~96]{Oliver1976}. The cellular proof above spells out the
existence of the map needed here.
\end{remark}

\begin{proof}[Proof of Theorem~\ref{thm:ind-main}]
\proves{thm:ind-main}
\uses{lem:ind-bonding, lem:ind-inverse-action}
Discarding the finitely many stages before $K_j$ does not change the
union or its final topology. Choose $D_n=D_{K_n}$ for $n\ge j$.
Lemma~\ref{lem:ind-bonding} provides $K_n$-equivariant maps
$p_n:D_{n+1}\to D_n$. Lemma~\ref{lem:ind-inverse-action} yields a
nonempty compact metrizable $G$-space
\[
X=\varprojlim_{n\ge j}(D_n,p_n).
\]
By hypothesis $K_j\notin\mathcal F(K_j)$, so
$D_j^{K_j}=\varnothing$. The $K_j$-equivariant projection $X\to D_j$
implies $X^{K_j}=\varnothing$, and therefore $X^G=\varnothing$.
\end{proof}

\begin{figure}[htbp]
\centering
\includegraphics[width=\linewidth]{inverse-limit-disks}
\caption{The proof of Theorem~\ref{thm:ind-main}. Lemma~\ref{lem:ind-bonding} gives $K_n$-equivariant maps $p_n\colon D_{n+1}\to D_n$ between the disks $D_n=D_{K_n}$ of Theorem~\ref{thm:ind-oliver}, and Lemma~\ref{lem:ind-inverse-action} makes $X=\varprojlim_{n\ge j}(D_n,p_n)$ a nonempty compact metrizable $G$-space whose projection to $D_j$ is $K_j$-equivariant. Since $D_j^{K_j}=\varnothing$, also $X^{K_j}=\varnothing$, and therefore $X^G=\varnothing$. (Illustration added for the encyclopaedia.)}\label{fig:inverse-limit-disks}
\illustrates{thm:ind-main}
\end{figure}

\begin{remark}\label{rem:ind-sharpness}
\uses{thm:ind-main, thm:ind-oliver}
The hypothesis of Theorem~\ref{thm:ind-main} says exactly that some stage is
not solvable, and it cannot be dropped from the present method. Let $K$ be
a solvable compact Lie group. The trivial subgroup lies in $\mathcal F(K)$, and
$K^0$ is a torus and therefore prime toral, so the separating property
places $K^0$ in $\mathcal F(K)$. The finite solvable group $K/K^0$ has a chief
series with elementary abelian factors; pulling it back gives a chain
$K^0=N_0\trianglelefteq\cdots\trianglelefteq N_m=K$ with each $N_{i+1}/N_i$
a finite elementary abelian $p$-group, hence prime toral. Applying the
separating property along the chain gives $K\in\mathcal F(K)$, so the disk of
Theorem~\ref{thm:ind-oliver} satisfies $D_K^K\ne\varnothing$. This is the
observation recorded after Theorem~4 in~\cite[p.~94]{Oliver1976}. It does
not say that a solvable compact Lie group admits no fixed-point-free smooth
action on a disk, only that no such action has all of its subgroup fixed
sets disks or empty.
\end{remark}

When every stage is solvable, extreme amenability can still fail for
unrelated reasons. For $G=\bigcup_n\mathbb T^n$ with the standard inclusions the
first coordinate $z\mapsto z_1$ is compatible with the inclusions and
defines a continuous surjective homomorphism $G\to\mathbb T$, so $G$ acts on $\mathbb T$
without a fixed point. This obstruction is absent as soon as $G$ is
minimally almost periodic.

\begin{remark}[Question]\label{rem:ind-solvable-stages}
Let $G=\bigcup_nK_n$ be an increasing union of solvable compact Lie groups
with the final topology, and suppose that $G$ is minimally almost periodic.
Is $G$ extremely amenable?
\end{remark}

The coordinate projection $\bigcup_n\mathbb T^n\to\mathbb T$ is excluded by the
minimal almost periodicity hypothesis, so the question is not settled
by that elementary quotient.


\subsection{The classical groups}\label{sec:ind-classical}

The embeddings are the standard block inclusions over
$\mathbb R$, $\mathbb C$, or the quaternions, as appropriate.

\begin{corollary}\label{cor:ind-semisimple-classical}
The groups $\varinjlim\operatorname{SU}(n)$, $\varinjlim\operatorname{SO}(n)$, and
$\varinjlim\operatorname{Sp}(n)$, with their final topologies, are not extremely
amenable. They admit compact metrizable actions with no
$\operatorname{SU}(2)$-fixed point, no $\operatorname{SO}(3)$-fixed point, and no
$\operatorname{Sp}(1)$-fixed point, respectively.
\end{corollary}

\begin{proof}
\uses{thm:ind-main}
The compact Lie groups $\operatorname{SU}(2)$, $\operatorname{SO}(3)$, and
$\operatorname{Sp}(1)\cong\operatorname{SU}(2)$ are connected and nonabelian. Apply
Theorem~\ref{thm:ind-main} at these stages. The real and quaternionic block
embeddings are smooth closed subgroup embeddings, so they satisfy the
same hypotheses.
\end{proof}

The special orthogonal and symplectic cases are instances of the
same obstruction. For the full unitary and orthogonal groups there
are simpler obstructions.

\begin{proposition}\label{prop:ind-determinants}
\uses{def:ea}
The final-topology groups $\varinjlim\operatorname{U}(n)$ and
$\varinjlim\operatorname{O}(n)$ are not extremely amenable.
\end{proposition}

\begin{proof}
Compatibility with the embeddings gives surjective homomorphisms
\[
\det:\varinjlim\operatorname{U}(n)\longrightarrow\mathbb T,
\qquad
\det:\varinjlim\operatorname{O}(n)\longrightarrow\{1,-1\}.
\]
Both are continuous by finality, because their restrictions to every
finite stage are continuous. Acting by multiplication through the
first homomorphism on $\mathbb T$, or through the second on $\{1,-1\}$,
gives a jointly continuous compact action without a global fixed
point. In either case an element of determinant $-1$ fixes no point
of the indicated space.
\end{proof}

Theorem~\ref{thm:ind-main} also applies directly to these two groups by
using $\operatorname{U}(2)$ and $\operatorname{O}(3)$. The determinant proof records the
elementary reason specific to the standard embeddings. The unitary
case was already noted in~\cite{CacciatoriUrsino2023}; these quotient
arguments are included for context rather than as novelty claims.


\subsection{Bohr compactifications}\label{sec:ind-bohr}

The \emph{Bohr compactification} of a topological group $G$ is a compact
Hausdorff group $bG$ with a continuous homomorphism $\eta_G:G\to bG$
having dense image, such that every continuous homomorphism from $G$
to a compact Hausdorff group factors uniquely through $\eta_G$ by a
continuous homomorphism. A group with $bG=\{1\}$ is called
\emph{minimally almost periodic}.

\begin{lemma}\label{lem:ind-no-finite-representations}
Every continuous finite-dimensional unitary representation of
$\operatorname{SU}(\infty)$, $\operatorname{SO}(\infty)$, or $\operatorname{Sp}(\infty)$ is trivial.
The same conclusion holds for any topology on the algebraic union
for which all stage inclusions are continuous.
\end{lemma}

\begin{proof}
Let $\rho:G\to\operatorname{U}(d)$ be such a representation. Its restriction to a
finite stage $K_n$ is a continuous homomorphism of Lie groups, hence
is smooth. For all sufficiently large $n$, the real Lie algebra
$\mathfrak k_n$ is simple and
\[
\dim\mathfrak k_n>d^2=\dim\mathfrak u(d).
\]
Here the dimensions of $\mathfrak{su}(n)$, $\mathfrak{so}(n)$, and
$\mathfrak{sp}(n)$ are respectively $n^2-1$, $n(n-1)/2$, and
$n(2n+1)$; in the orthogonal case we take $n\ge5$ to avoid the
low-dimensional exceptions. The kernel of
$d\rho_n:\mathfrak k_n\to\mathfrak u(d)$ is an ideal, so simplicity
makes $d\rho_n$ either zero or injective. The dimension inequality
excludes injectivity. Since $K_n$ is connected, $\rho|_{K_n}=1$.
Every element of $G$ belongs to an arbitrarily large stage, proving
the assertion.
\end{proof}

\begin{proposition}\label{prop:ind-bohr-classical}
For the standard classical limits with their final topologies,
\begin{gather*}
b\operatorname{SU}(\infty)=b\operatorname{SO}(\infty)=b\operatorname{Sp}(\infty)=\{1\},\\
b\operatorname{U}(\infty)\cong\mathbb T,\qquad b\operatorname{O}(\infty)\cong\{1,-1\}.
\end{gather*}
The last two Bohr maps are the determinants.
\end{proposition}

\begin{proof}
\uses{lem:ind-no-finite-representations}
If $G$ is one of the first three groups and $\phi:G\to C$ is a
continuous homomorphism to a compact Hausdorff group, compose $\phi$
with every finite-dimensional continuous unitary representation of
$C$. Lemma~\ref{lem:ind-no-finite-representations} makes all these
compositions trivial. The Peter--Weyl theorem says that these
representations separate points of $C$, so $\phi$ is trivial.

For $G=\operatorname{U}(\infty)$, the same argument applied to the special unitary
subgroup shows that every such $\phi$ kills $\operatorname{SU}(\infty)$. The
inclusion of this subgroup, equipped with its own final topology, is
continuous by finality. Since
$\ker(\det:G\to\mathbb T)=\operatorname{SU}(\infty)$, the homomorphism factors uniquely
through $\det$. The factor is continuous because
$s(z)=\operatorname{diag}(z,1,1,\ldots)$ is a continuous section: the factor is
$\phi\circ s$. The determinant is continuous by finality and is
surjective, so it has the required universal property.

The orthogonal case is identical, using
$\ker(\det:\operatorname{O}(\infty)\to\{1,-1\})=\operatorname{SO}(\infty)$ and the section
$s(\varepsilon)=\operatorname{diag}(\varepsilon,1,1,\ldots)$.
\end{proof}


\subsection{Nonmetrizable universal minimal flows}\label{sec:ind-umf}

A \emph{minimal flow} is a nonempty compact Hausdorff $G$-space in
which every orbit is dense. The \emph{universal minimal flow} $M(G)$
is the minimal flow admitting a continuous equivariant surjection onto
every minimal $G$-flow; it exists and is unique up to isomorphism.
We prove that the classical groups considered here have universal
minimal flows containing $\beta\mathbb N$. The embedding is topological;
no prescribed action on $\beta\mathbb N$ is asserted.


\subsubsection{Finite-stage orbits and commuting tails}

\begin{lemma}\label{lem:ind-compact-centralizer}
Suppose a compact Hausdorff group $K$ acts transitively on a compact
Hausdorff space $Y$, and write $Y\cong K/L$ for a closed subgroup
$L\le K$. The centraliser of $K$ in $\operatorname{Homeo}(Y)$, with the compact-open
topology, is a compact group isomorphic to $N_K(L)/L$.
\end{lemma}

\begin{proof}
A $K$-equivariant homeomorphism is determined by the image $aL$ of
$L$. Since a bijective equivariant map preserves stabilisers, we have
$aLa^{-1}=L$, so $a\in N_K(L)$. Conversely, each $a\in N_K(L)$
defines a $K$-equivariant homeomorphism by $kL\mapsto ka^{-1}L$.
This convention gives a homomorphism from $N_K(L)$ with kernel $L$.
The resulting map $N_K(L)/L\to\operatorname{Homeo}(Y)$ is continuous: the displayed
action is jointly continuous, and compactness of $Y$ gives continuity
into the compact-open topology. It is a bijection onto the centraliser.
Its domain is compact and its codomain Hausdorff, so it is a
homeomorphism onto that centraliser.
\end{proof}

For the rest of this subsection let $(G,G_{\mathrm{ss}},L_0)$ be one of the triples
in the following table. The middle column is a closed subgroup of $G$
isomorphic to one of $\operatorname{SU}(\infty)$, $\operatorname{SO}(\infty)$, or $\operatorname{Sp}(\infty)$.
All subgroups use the standard coordinate
blocks; $K_n$ denotes the corresponding full finite stage of $G$.
\[
\begin{array}{c|c|c}
G&G_{\mathrm{ss}}&L_0\\\hline
\operatorname{SU}(\infty)&\operatorname{SU}(\infty)&\operatorname{SU}(2)\\
\operatorname{SO}(\infty)&\operatorname{SO}(\infty)&\operatorname{SO}(3)\\
\operatorname{Sp}(\infty)&\operatorname{Sp}(\infty)&\operatorname{Sp}(1)\\
\operatorname{U}(\infty)&\operatorname{SU}(\infty)&\operatorname{SU}(2)\\
\operatorname{O}(\infty)&\operatorname{SO}(\infty)&\operatorname{SO}(3)
\end{array}
\]

\begin{lemma}\label{lem:ind-stage-orbits-nowhere-dense}
\uses{def:minimal}
Let $Y$ be a minimal $G$-flow on which $G_{\mathrm{ss}}$ acts nontrivially.
Then $K_ny$ is nowhere dense in $Y$ for every $n$ and $y\in Y$.
\end{lemma}

\begin{proof}
\uses{lem:ind-compact-centralizer, prop:ind-bohr-classical}
Suppose $A=\operatorname{Int}(K_ny)$ is nonempty. Minimality implies
that the translates $gA$, $g\in G$, cover $Y$. Choose a finite
subcover $g_1A,\ldots,g_rA$, and a stage $K_N$ containing $K_n$ and
all $g_i$. Then $Y=K_Ny$.

Let $T_N$ be the copy of $G_{\mathrm{ss}}$ acting on coordinates strictly after $N$
and fixing the first $N$ coordinates. Give it its own final topology.
Its inclusion into $G$ is continuous, since this holds on every stage.
The groups $T_N$ and $K_N$ commute, so the action of $T_N$ takes
values in the compact centraliser supplied by
Lemma~\ref{lem:ind-compact-centralizer}. Proposition~\ref{prop:ind-bohr-classical}
applies to this shifted copy of $G_{\mathrm{ss}}$ and makes its action trivial.

Every element of $G_{\mathrm{ss}}$ is conjugate in $G$ to an element of $T_N$:
move its finite coordinate block into the tail by a finitary
orthogonal, unitary, or quaternionic unitary matrix. In the special
orthogonal and special unitary cases, correct the determinant on an
additional unused coordinate if necessary. The kernel of the action
is normal in $G$, and hence contains all of $G_{\mathrm{ss}}$. This contradicts
the hypothesis. Finally, $K_ny$ is compact and therefore closed, so
empty interior implies that it is nowhere dense.
\end{proof}

\begin{lemma}\label{lem:ind-semisimple-witness}
\uses{def:flow}
For every triple above there is a nonempty compact metrizable minimal
$G$-flow $Y$ with $Y^{L_0}=\varnothing$.
\end{lemma}

\begin{proof}
\uses{cor:ind-semisimple-classical, thm:ind-main}
For $\operatorname{SU}(\infty)$, $\operatorname{SO}(\infty)$, and $\operatorname{Sp}(\infty)$, use the
actions in Corollary~\ref{cor:ind-semisimple-classical} and take a minimal
subflow. For $\operatorname{U}(\infty)$, apply Theorem~\ref{thm:ind-main} to the chain
\[
\operatorname{SU}(2)\subset\operatorname{U}(2)\subset\operatorname{U}(3)\subset\cdots.
\]
For $\operatorname{O}(\infty)$ use
$\operatorname{SO}(3)\subset\operatorname{O}(3)\subset\operatorname{O}(4)\subset\cdots$ instead.
Adding the first compact subgroup does not change the union or its
final topology: continuity of its inclusion into the next stage
makes its open-set condition redundant. The initial subgroup is
connected and nonabelian, so the resulting action has no $L_0$-fixed
point. Every nonempty compact flow has a minimal subflow; this
subflow remains metrizable and has no $L_0$-fixed point.
\end{proof}


\subsubsection{Separation in the final topology}

The next lemma isolates the topological step. It does not require
minimality or any Lie structure on the stages.

\begin{lemma}\label{lem:ind-final-open-separation}
Let $G=\bigcup_{n\ge1}K_n$ be a topological group with the final
topology for an increasing sequence of compact subgroups. Suppose
$X$ is a compact metrizable $G$-space and $x\in X$ satisfies
$\operatorname{Int}(K_nx)=\varnothing$ for every $n$. Then there
exist an open identity neighbourhood $W\subseteq G$ and nonempty open
sets $A_n\subseteq X$ such that the sets $WA_n$ are pairwise disjoint.
\end{lemma}

\begin{proof}
Fix a compatible metric on $X$. Being compact metrizable, $X$ is locally
compact, and the action is jointly continuous, so $K_nB$ is compact for
every compact $B\subseteq X$. Recursively choose nested open
neighbourhoods $V_n$ of $x$ with $\operatorname{diam}(V_n)\le2^{-n}$,
and compact sets $Q_n$ with nonempty interior, so that
\begin{equation}\label{eq:ind-separation-recursion}
Q_n\subset V_n\setminus K_nx,\qquad
V_{n+1}\subset V_n\setminus K_nQ_n.
\end{equation}
Indeed, $K_nx$ is closed with empty interior, so
$V_n\setminus K_nx$ contains a compact neighbourhood of some point.
Choose this as $Q_n$. Since $K_n$ is a group,
$Q_n\cap K_nx=\varnothing$ implies $x\notin K_nQ_n$; the latter
set is compact. Thus an appropriate $V_{n+1}$ exists. We obtain
\begin{equation}\label{eq:ind-separation-compact-sets}
\sup_{z\in Q_n}d(z,x)\longrightarrow0,\qquad
K_iQ_i\cap Q_j=\varnothing\quad(i<j).
\end{equation}
In particular, the $Q_n$ are pairwise disjoint and none contains $x$.

Define the set of forbidden group elements
\begin{align*}
E={}&\{g\in G:\text{there exist }i\ne j\text{ with }gQ_i\cap Q_j\ne\varnothing\}\\
&\quad\cup\{g\in G:\text{there exists }i\text{ with }gx\in Q_i
                    \text{ or }x\in gQ_i\}.
\end{align*}
Then $e\notin E$. We claim that $E\cap K_m$ is closed for every $m$.
For this purpose put
\[
C_i=\{x\}\cup\bigcup_{j\ne i}Q_j.
\]
Each $C_i$ is compact: any neighbourhood of $x$ contains all but
finitely many $Q_j$, and the remaining sets are compact. For compact
$B,C\subseteq X$, the set
\[
\mathcal T_m(B,C)=\{g\in K_m:gB\cap C\ne\varnothing\}
\]
is closed, as the projection of a closed subset of $K_m\times B$.
We have the finite-union identity
\begin{equation}\label{eq:ind-forbidden-stage}
E\cap K_m=\bigcup_{i<m}
\bigl(\mathcal T_m(Q_i,C_i)\cup\mathcal T_m(C_i,Q_i)\bigr).
\end{equation}
To check it, a transporter $gQ_i\cap Q_j\ne\varnothing$ with
$g\in K_m$ cannot have both $i,j\ge m$: use
\eqref{eq:ind-separation-compact-sets} with the smaller index, replacing
$g$ by $g^{-1}$ when necessary. Also, $gx\in Q_i$ or $x\in gQ_i$
is impossible for $i\ge m$, because $Q_i\cap K_ix=\varnothing$.
All remaining cases are exactly the terms in
\eqref{eq:ind-forbidden-stage}. Thus $E\cap K_m$ is closed.

The final topology now makes $E$ closed in $G$. Choose an open
identity neighbourhood $W$ such that $W^{-1}W\subseteq G\setminus E$,
and let $A_n=\operatorname{Int}(Q_n)$. If $w_i a_i=w_j a_j$ with
$w_i,w_j\in W$, $a_i\in A_i$, $a_j\in A_j$, and $i\ne j$, then
$w_j^{-1}w_iQ_i\cap Q_j\ne\varnothing$, contrary to the choice of
$W$. Hence the sets $WA_n$ are pairwise disjoint.
\end{proof}

\begin{proposition}[Consequence of Basso--Zucker]\label{prop:ind-basso-zucker}
\uses{thm:ellis}
Let $G$ be any topological group. Suppose some minimal $G$-flow $Y$
has nonempty open sets $A_n$ and an open identity neighbourhood $W$
such that the sets $WA_n$ are pairwise disjoint. Then $M(G)$ contains
a homeomorphic copy of $\beta\mathbb N$.
\end{proposition}

\begin{proof}
Let $\pi:M(G)\to Y$ be an equivariant surjection. The sets
$B_n=\pi^{-1}(A_n)$ are nonempty and open, and the sets $WB_n$ are
pairwise disjoint, because $\pi(WB_n)\subseteq WA_n$.
Shrinking $W$, we may assume that it belongs to the base of symmetric
open identity neighbourhoods used in~\cite{BassoZucker2021}.
We follow the numbering of the accepted preprint
\texttt{arXiv:2008.08471v2}: Definition~5.1, Proposition~5.6 and
Corollary~5.7. Those statements give the conclusion. In their
terminology, the separated open sets imply that $G$ is not UEB, and
the corollary then embeds $\beta\mathbb N$ in $M(G)$. The statements
apply to arbitrary topological groups.
\end{proof}

\begin{theorem}\label{thm:ind-classical-umf}
With their standard block embeddings and final topologies, each of
\[
\operatorname{SU}(\infty),\quad\operatorname{SO}(\infty),\quad\operatorname{Sp}(\infty),\quad
\operatorname{U}(\infty),\quad\operatorname{O}(\infty)
\]
has a universal minimal flow containing a homeomorphic copy of
$\beta\mathbb N$. In particular, all five universal minimal flows are
nonmetrizable.
\end{theorem}

\begin{proof}
\uses{lem:ind-final-open-separation, lem:ind-semisimple-witness, lem:ind-stage-orbits-nowhere-dense, prop:ind-basso-zucker}
Choose the compact metrizable minimal flow $Y$ from
Lemma~\ref{lem:ind-semisimple-witness}. Its subgroup $G_{\mathrm{ss}}$ acts nontrivially,
since $L_0\le G_{\mathrm{ss}}$ and $Y^{L_0}=\varnothing$. Thus every finite-stage
orbit is nowhere dense by Lemma~\ref{lem:ind-stage-orbits-nowhere-dense}.
Apply Lemma~\ref{lem:ind-final-open-separation} at any point of $Y$, then
Proposition~\ref{prop:ind-basso-zucker}. The compact space $\beta\mathbb N$ is
nonmetrizable, so no compact metrizable space can contain it.
\end{proof}

Together, Proposition~\ref{prop:ind-bohr-classical} and
Theorem~\ref{thm:ind-classical-umf} prove
Theorem~\ref{thm:ind-dynamics-intro}. The Bohr calculation is independent
of the Oliver construction. In the unitary and orthogonal cases,
determinant actions alone would not suffice for the nonmetrizability
proof: the witnesses in Lemma~\ref{lem:ind-semisimple-witness} detect
$\operatorname{SU}(2)$ and $\operatorname{SO}(3)$ inside the determinant kernels.


\subsection{Nonfreeness and the free locus}\label{sec:ind-free-locus}

For a $G$-flow $Z$, write $G_z=\{g\in G:g\cdot z=z\}$ and
\[
\operatorname{Free}(Z)=\{z\in Z:G_z=\{e\}\}.
\]
We call a Borel subset of $Z$ \emph{universally invariant-conull} if
it has measure one for every $G$-invariant Radon probability measure
on $Z$. The invariance qualification is essential: a proper subset
cannot have full measure for every Radon probability measure.

\begin{theorem}\label{thm:ind-classical-free-locus}
\uses{cor:ind-classical-intro, thm:ellis}
Let $G$ be any of the five classical limits in
Corollary~\ref{cor:ind-classical-intro}, with its standard final topology.
Then $\operatorname{Free}(M(G))$ is a proper dense $G_\delta$ and is
universally invariant-conull. Invariant Radon probability measures
on $M(G)$ exist. The nonfree locus is an invariant dense meager
$F_\sigma$ and is null for every such measure.
\end{theorem}


\subsubsection{An involution with fixed points in every compact flow}

We first prove nonfreeness. For a nilpotent finite CW complex $P$,
write $\widehat P$ for its Sullivan profinite completion.
McGibbon's theorem states that
\begin{equation}\label{eq:ind-mcgibbon-null}
\operatorname{Map}_*(W,\widehat P)\simeq_w *
\end{equation}
whenever $W$ is a connected infinite loop space with torsion
fundamental group~\cite[Theorem~3]{McGibbon1996}; the precise
formulation is also recorded in
\cite[Example~1.6(iii), p.~68]{KiharaOda2012}.
Here $\operatorname{Map}_*$ denotes the pointed mapping space,
and weak contractibility includes the assertion that it has only
one path component. We also use the Zabrodsky lemma in the following
form: if $F\to E\to B$ is a fibration with $B$ of connected CW
homotopy type and $\operatorname{Map}_*(F,Y)$ is weakly contractible,
then
\begin{equation}\label{eq:ind-zabrodsky-mapping}
\operatorname{Map}_*(B,Y)\longrightarrow\operatorname{Map}_*(E,Y)
\end{equation}
is a weak equivalence~\cite[Lemma~2.1]{CenceljEtAl2007}.
This formulation requires neither $B$ nor $Y$ to be simply connected.

\begin{lemma}\label{lem:ind-no-odd-stable-su}
Put $a=\operatorname{diag}(-1,-1,1,\ldots)\in\operatorname{SU}(\infty)$. There is no
continuous map $f:\operatorname{SU}(\infty)\to S^d$, for any finite $d\ge0$,
such that $f(ag)=-f(g)$ for every $g\in\operatorname{SU}(\infty)$.
\end{lemma}

\begin{proof}
Let
\[
T=\operatorname{Tel}\bigl(\operatorname{SU}(2)\hookrightarrow\operatorname{SU}(3)
\hookrightarrow\cdots\bigr)
\]
be the mapping telescope. The compatible left actions of
$C_2=\langle a\rangle$ give a free action on $T$ and a canonical
continuous equivariant map $T\to\operatorname{SU}(\infty)$. The telescope has
the usual stable special-unitary homotopy type. In particular, it
is simply connected and has the homotopy type of an infinite loop
space: the stable unitary space is an infinite loop space by Bott
periodicity, and its simply connected cover is again an infinite
loop space. The telescope allows us to work with CW homotopy types
without making any assertion about the literal quotient of the
final-topology group.

Suppose $f$ exists. Compose with the map from $T$ and an equivariant
equatorial inclusion to obtain an odd map $T\to S^m$, where
$m\ge3$ is odd. Set
\[
E=EC_2\times_{C_2}T,\qquad P=\mathbb{RP}^m.
\]
The finite CW complex $P$ is nilpotent. Its universal cover is
$S^m$, and the antipodal deck transformation has degree $+1$, hence
is homotopic to the identity. Thus $\pi_1(P)=C_2$ acts trivially
on every higher homotopy group.

Both $T$ and $BC_2=K(C_2,1)$ satisfy the hypothesis of
\eqref{eq:ind-mcgibbon-null}. Apply \eqref{eq:ind-zabrodsky-mapping} to
the Borel fibration
\[
T\longrightarrow E\longrightarrow BC_2.
\]
It follows that
\begin{equation}\label{eq:ind-borel-completion-null}
\operatorname{Map}_*(E,\widehat P)\simeq_w *.
\end{equation}
On the other hand, the odd map induces a map of Borel fibrations
over $BC_2$ and hence a map
\[
\bar f:E\longrightarrow EC_2\times_{C_2}S^m\simeq P.
\]
Both projections to $BC_2$ induce isomorphisms on fundamental
groups, because their fibres $T$ and $S^m$ are simply connected.
Therefore $\bar f_*$ is an isomorphism $C_2\to C_2$.

Completion preserves the finite fundamental group of the nilpotent
complex $P$. Consequently the composite $E\to P\to\widehat P$
still induces a nontrivial map on fundamental groups, contradicting
\eqref{eq:ind-borel-completion-null}. Equivalently, completion
naturality factors the covering classifier $P\to BC_2$ through
$\widehat P\to\widehat{BC_2}\simeq BC_2$; a null composite
$E\to\widehat P$ would make the projection $E\to BC_2$ null.
\end{proof}

\begin{remark}\label{rem:ind-index-three}
\uses{lem:ind-no-odd-stable-su}
Lemma~\ref{lem:ind-no-odd-stable-su} cannot be obtained from the usual
cohomological index of the involution. Let $w$ generate
$H^1(BC_2;\mathcal F_2)$ and let $i:C_2\to\operatorname{SU}(\infty)$ send the generator to
$a$. The Borel fibration $T\to E\to BC_2$ maps to the universal
fibration over $B\operatorname{SU}$ along $Bi$, and $c_2$ transgresses to the
generator $x_3$ of $H^3(T;\mathcal F_2)$ there, so the transgression of $x_3$
is $(Bi)^*c_2$ by naturality. The representation underlying $i$ is
$\sigma\oplus\sigma$, with $\sigma$ the sign character, so its total
Chern class is $(1+t)^2$ for $t=c_1(\sigma)$, giving
$(Bi)^*c_2=t^2$, whose mod~$2$ reduction is $w^4$. Hence $w^4=0$ in
$H^4(E;\mathcal F_2)$, while $w^3\ne0$ because $\operatorname{SU}(2)=S^3$ sits inside
$\operatorname{SU}(\infty)$ equivariantly for the antipodal action. The largest
$k$ with $w^k\ne0$ is therefore $3$, and no argument based on the
height of $w$ can exclude an odd map to a finite-dimensional sphere.
The homotopy-theoretic input above is not a convenience.
\end{remark}

\begin{proposition}\label{prop:ind-involution-fixed}
\uses{def:flow}
Let $G$ be any of the five classical limits. Every nonempty compact
Hausdorff $G$-flow has a point fixed by a nonidentity involution. For
$\operatorname{SU}(\infty)$ one may take $a=\operatorname{diag}(-1,-1,1,\ldots)$, and for the
other groups one may take its image under the embeddings used below.
Consequently, none of the five groups admits a free compact action, and
their universal minimal actions are nonfree.
\end{proposition}

\begin{proof}
\uses{lem:ind-no-odd-stable-su}
Let $Z$ be a nonempty compact $\operatorname{SU}(\infty)$-flow and suppose $Z^a$ is
empty. Continuous real functions separate points. Compactness
therefore supplies $u_1,\ldots,u_r\in C(Z,\mathbb R)$ such that
\[
D(z)=\bigl(u_1(az)-u_1(z),\ldots,u_r(az)-u_r(z)\bigr)\ne0
\qquad(z\in Z).
\]
The map $F(z)=D(z)/\lVert D(z)\rVert$ takes values in $S^{r-1}$ and
satisfies $F(az)=-F(z)$. For any $z_0\in Z$, the continuous map
$g\mapsto F(gz_0)$ contradicts Lemma~\ref{lem:ind-no-odd-stable-su}.

For the other groups, use the compatible stagewise embeddings
\[
\operatorname{SU}(n)\hookrightarrow\operatorname{U}(n),\qquad
\operatorname{SU}(n)\hookrightarrow\operatorname{Sp}(n),\qquad
\operatorname{SU}(n)\hookrightarrow\operatorname{SO}(2n)\hookrightarrow\operatorname{O}(2n).
\]
The second views complex matrices as quaternionic matrices, and
the third is realification. They induce continuous injective
homomorphisms on the final-topology unions. Restricting any compact
flow along one of these homomorphisms gives a fixed point for the
nonidentity image of $a$. In the real groups this image is
$\operatorname{diag}(-I_4,I,\ldots)$. Apply the assertion to the universal
minimal flow of each group.
\end{proof}

\begin{remark}\label{rem:ind-pestov-question-four}
\uses{prop:ind-involution-fixed}
The final-topology group $\operatorname{SU}(\infty)=\bigcup_n\operatorname{SU}(n)$ is a
$k_\omega$-space, since the stages are compact and determine its closed
sets. Proposition~\ref{prop:ind-involution-fixed} therefore gives a negative
answer to Pestov's Question~4~\cite[p.~440]{Pestov2007}: this
$k_\omega$-group admits no free action on a nonempty compact Hausdorff
space. The proof gives the same conclusion for all five classical
limits. Thus the paper answers Question~4 in addition to Question~19.
\end{remark}


\subsubsection{Generic vectors and stage stabilisers}

\begin{lemma}\label{lem:ind-generic-free-vectors}
Let $L\le\operatorname{U}(V)$ be a compact subgroup whose identity component is
a torus, where $V$ is a finite-dimensional complex inner-product
space. Almost every vector of $V$, and almost every vector of its
unit sphere, has trivial $L$-stabiliser. If $L\le\operatorname{O}(W)$ has the
same property and $\dim_{\mathbb R}W\ge2$, then almost every ordered
orthonormal two-frame in $W$ has trivial pointwise $L$-stabiliser.
The measures on spheres and frame spaces are the invariant
probability measures.
\end{lemma}

\begin{proof}
Put $T=L^0$ and decompose $V$ into its $T$-weight spaces,
$V=\bigoplus_\chi V_\chi$. Choose representatives $b$ of the
finite quotient $L/T$. Outside a set of Lebesgue measure zero,
the components $v_\chi$ of a vector $v$ are all nonzero and have
pairwise distinct norms. If $b$ permutes the weight spaces
nontrivially, these norm conditions prevent $btv=v$ for every
$t\in T$. If $b$ preserves every weight space but is nonscalar
on some $V_\chi$, the equation $btv=v$ forces $v_\chi$ into an
eigenspace of $b|_{V_\chi}$. Exclude these finitely many proper
eigenspaces for each such representative. Finally, if $b$ acts
scalarly on every weight space, $btv=v$ forces $bt$ to act as the
identity on all of $V$, so $bt=e$. The same argument covers the
identity coset. All excluded sets are cones of measure zero, so
their intersections with the unit sphere also have measure zero.

For the real assertion, apply the complex assertion to the faithful
complexification of the representation on $W$. If $u,v$ are
independent standard real Gaussian vectors, then $u+iv$ has trivial
stabiliser almost surely. Its stabiliser is precisely the common
stabiliser of $u$ and $v$. The pair is linearly independent almost
surely, and Gram--Schmidt preserves this common stabiliser while
giving the invariant distribution on ordered orthonormal
two-frames.
\end{proof}

\begin{lemma}\label{lem:ind-compact-free-open}
For a continuous action of a compact Lie group $K$ on a compact
Hausdorff space $Z$, the $K$-free locus is open.
\end{lemma}

\begin{proof}
Let $z$ be $K$-free and choose a faithful unitary representation
$\rho:K\to\operatorname{U}(d)$. The matrix function $kz\mapsto\rho(k)$ is
continuous on the closed orbit $Kz$. Extend its real and imaginary
entries by the Tietze extension theorem to a continuous
$A:Z\to M_d(\mathbb C)$, and put
\[
B(y)=\int_K\rho(k)^{-1}A(ky)\,\mathrm d k.
\]
Then $B$ is continuous, $B(z)=I$, and
$B(hy)=\rho(h)B(y)$. If $B(y)$ is invertible and $hy=y$, this
identity implies $\rho(h)=I$, hence $h=e$. Invertibility therefore
provides the required open neighbourhood of $z$.
\end{proof}

\begin{proposition}\label{prop:ind-generic-free-locus}
\uses{def:flow}
Let $G$ be one of the five classical final-topology unions, with
standard stages $K_n$. Suppose that a nonempty compact Hausdorff
$G$-flow $Z$ satisfies
\[
(G_z\cap K_m)^0\text{ is a torus}\qquad(z\in Z,\ m\in\mathbb N).
\]
Then $\operatorname{Free}(Z)$ is a dense $G_\delta$ and is
universally invariant-conull.
\end{proposition}

\begin{proof}
\uses{lem:ind-compact-free-open, lem:ind-generic-free-vectors}
Put $F_n=\{z\in Z:G_z\cap K_n=\{e\}\}$. These sets are open
by Lemma~\ref{lem:ind-compact-free-open}. Fix $z\in Z$ and $n$.
For the unitary and symplectic chains take $m=n+1$ and let $b=e_m$
be the last coordinate vector. For the orthogonal chains take
$m=n+2$ and let $b=(e_{m-1},e_m)$ be the last ordered two-frame.
In each case $K_n$ fixes $b$ pointwise and $K_m$ is transitive on
the relevant sphere or frame space. Haar probability under
$k\mapsto k^{-1}b$ is its invariant probability measure.

Apply Lemma~\ref{lem:ind-generic-free-vectors} to $L=G_z\cap K_m$.
For the symplectic chain use its standard complex unitary
representation on $\mathbb C^{2m}$, whose unit sphere is the unit sphere
of $\mathbb H^m$. For Haar-almost every $k\in K_m$, the vector
or frame $k^{-1}b$ has trivial $L$-stabiliser. If
$h\in G_{kz}\cap K_n$, then $k^{-1}hk\in L$ fixes $k^{-1}b$,
so $h=e$. We have proved
\begin{equation}\label{eq:ind-haar-stage-free}
\int_{K_m}\mathbf1_{F_n}(kz)\,\mathrm d k=1\qquad(z\in Z).
\end{equation}
Every identity neighbourhood in $K_m$ has positive Haar measure.
Thus, by continuity of the orbit map at the identity, every
neighbourhood of $z$ meets $F_n$. Hence $F_n$ is dense. Since $Z$
is compact Hausdorff, it is Baire, and
\[
\operatorname{Free}(Z)=\bigcap_n F_n
\]
is a dense $G_\delta$.

If $\mu$ is any $G$-invariant Radon probability measure, Fubini and
\eqref{eq:ind-haar-stage-free} give
\[
1=\int_Z\int_{K_m}\mathbf1_{F_n}(kz)\,\mathrm d k\,\mathrm d\mu(z)
=\int_{K_m}\mu(k^{-1}F_n)\,\mathrm d k=\mu(F_n).
\]
The integrand is product measurable even if $Z$ is nonmetrizable:
the inverse image of $F_n$ under the action map is open in
$K_m\times Z$, and every such open set is a countable union of
open rectangles because $K_m$ has a countable base. Taking the
countable intersection gives $\mu(\operatorname{Free}(Z))=1$.
\end{proof}

\begin{proof}[Proof of Theorem~\ref{thm:ind-classical-free-locus}]
\proves{thm:ind-classical-free-locus}
\uses{prop:ind-generic-free-locus, prop:ind-involution-fixed, thm:ind-oliver, thm:ind-main}
Let $X=\varprojlim D_m$ be the Oliver flow constructed in the proof
of Theorem~\ref{thm:ind-main}. A minimal subflow of $X$ is a factor of
$M(G)$, so there is a $G$-map $q:M(G)\to X$. For $z\in M(G)$,
the compact subgroup $L=G_z\cap K_m$ fixes $\pi_m(q(z))$.
The fixed-set prescription \eqref{eq:ind-fixed-sets} implies that
$L$ is solvable, and hence $L^0$ is a torus. If the construction
discarded some initial stages, handle them inside a later stage.
Proposition~\ref{prop:ind-generic-free-locus} gives density, the
$G_\delta$ property, and universal invariant-conullness.

For completeness, invariant Radon probability measures exist on
every nonempty compact $G$-flow. For each $n$, Haar averaging gives
a $K_n$-invariant probability measure. The sets of such measures
are nonempty weak-* compact subsets of the space of Radon
probabilities, and are nested as $n$ increases. Their intersection
is nonempty and consists of $G$-invariant measures. On $M(G)$ every
such measure has full support by minimality.

Proposition~\ref{prop:ind-involution-fixed} shows that the free locus
is proper. Its complement is invariant, meager, $F_\sigma$, and
null for every invariant Radon probability. It is dense as well:
the orbit of a nonfree point consists of nonfree points and is
dense by minimality.
\end{proof}

Together, Proposition~\ref{prop:ind-involution-fixed} and
Theorem~\ref{thm:ind-classical-free-locus} prove
Theorem~\ref{thm:ind-free-intro}.

\begin{remark}\label{rem:ind-remark-3}
\uses{thm:ind-classical-free-locus, thm:ind-main}
Theorem~\ref{thm:ind-classical-free-locus} concerns the five classical
chains and their specified topologies. It does not hold for every
chain covered by Theorem~\ref{thm:ind-main}: for a constant chain at
a compact Lie group $K$, the universal minimal flow is $K$ with
its free translation action. Also, the existence of a fixed point
for each individual involution does not give a common fixed point
for different involutions.
\end{remark}


\subsection{A finite family of obstructing observables}\label{sec:ind-observables}

We use the convention
\begin{equation}\label{eq:ind-ruc-convention}
f\in\operatorname{RUCB}(G)
\quad\Longleftrightarrow\quad
f\text{ is bounded and }
\sup_{h\in G}\lvert f(uh)-f(h)\rvert\longrightarrow0\quad(u\to e).
\end{equation}
All functions in this section are real-valued.

\begin{proposition}\label{prop:ind-obstructing-observables}
\uses{thm:ind-main}
Under the hypotheses of Theorem~\ref{thm:ind-main}, there are finitely
many $f_1,\ldots,f_d\in\operatorname{RUCB}(G)$ and a constant $\delta>0$ such that
\begin{equation}\label{eq:ind-variance-obstruction}
\sum_{i=1}^d\operatorname{Var}_{k\in K_j}f_i(kh)\ge\delta
\qquad(h\in G).
\end{equation}
\end{proposition}

\begin{proof}
Identify the ordinary compact disk $D_j$ with a closed ball in
$\mathbb R^d$, and let $a_1,\ldots,a_d$ be its coordinate functions. Define
\[
v(y)=\sum_{i=1}^d\operatorname{Var}_{k\in K_j}a_i(k\cdot y),\qquad y\in D_j.
\]
This is a continuous nonnegative function. If $v(y)=0$, each
continuous orbit function $k\mapsto a_i(k\cdot y)$ is constant,
because Haar measure has full support. The coordinates separate
points, so the orbit of $y$ is a singleton. This contradicts
$D_j^{K_j}=\varnothing$. Thus compactness gives
$\delta=\min_{D_j}v>0$.

For the inverse-limit action constructed above, fix $x_*\in X$, put
$A_i=a_i\circ\pi_j\in C(X)$, and define $f_i(g)=A_i(g\cdot x_*)$.
Joint continuity and compactness of $X$ imply
\[
\sup_{h\in G}\lvert f_i(uh)-f_i(h)\rvert
\le\sup_{x\in X}\lvert A_i(u\cdot x)-A_i(x)\rvert\longrightarrow0.
\]
Hence $f_i\in\operatorname{RUCB}(G)$. Equivariance of $\pi_j$ gives
\[
\sum_i\operatorname{Var}_{k\in K_j}f_i(kh)
=v\bigl(\pi_j(h\cdot x_*)\bigr)\ge\delta.
\]
\end{proof}

The obstruction is uniform over every translate $K_jh$. In particular,
failure of extreme amenability is witnessed by a fixed finite family
of observables, rather than by a family whose size grows with the
ambient stage.


\subsection{A relative variance estimate in trace norm}\label{sec:ind-relative}

We now work in the compact group $\operatorname{SU}(N)$ with
$d_1(g,h)=\operatorname{Tr}\lvert g-h\rvert$, without normalising by $N$. Fix $N>m\ge2$
and write $K=\operatorname{SU}(m)$ for the upper-left block subgroup.


\subsubsection{A representation inequality}

Let $e_{ij}$ be matrix units. In a finite-dimensional unitary
representation $\pi$ of $\operatorname{SU}(N)$, extend the differential complex
linearly to $\mathfrak{sl}_N(\mathbb C)$ and write
$E_{ij}=d\pi(e_{ij})$ for $i\ne j$. Then
$E_{ij}^*=E_{ji}$. Define
\begin{equation}\label{eq:ind-cross-casimir}
Q=\sum_{i=1}^m\sum_{a=m+1}^N
\bigl(E_{ia}E_{ai}+E_{ai}E_{ia}\bigr).
\end{equation}
The operator is positive and commutes with $K$: its summands form the
sum of squares over an orthonormal basis of the cross-block subspace,
which is invariant under conjugation by $K$.

\begin{lemma}\label{lem:ind-relative-gap}
Let $\operatorname{P}^{\mathrm{fix}}_K$ be the orthogonal projection onto the $K$-fixed vectors.
Then
\begin{equation}\label{eq:ind-relative-gap}
Q\ge(N-m)(I-\operatorname{P}^{\mathrm{fix}}_K).
\end{equation}
More precisely, on a $K$-isotypic component of highest weight
$\lambda$, $Q\ge(N-m)R$, where
$R=\lambda_1-\lambda_m$ is the highest-weight spread.
\end{lemma}

\begin{proof}
Set $Z=d\pi(e_{11}-e_{mm})$. Matrix-unit commutators give the
following identity of self-adjoint operators:
\begin{align}\label{eq:ind-positive-identity}
Q-(N-m)Z
={}&2\sum_{a>m}\bigl(E_{a1}E_{1a}+E_{ma}E_{am}\bigr)\notag\\
&+\sum_{i=2}^{m-1}\sum_{a>m}
\bigl(E_{ia}E_{ai}+E_{ai}E_{ia}\bigr)\ge0.
\end{align}
The second sum is empty if $m=2$. Each term on the right is of the
form $T^*T$ or $TT^*$. For example, the identity follows by subtracting
\[
\sum_{a>m}\bigl([E_{1a},E_{a1}]-[E_{ma},E_{am}]\bigr)=(N-m)Z
\]
from \eqref{eq:ind-cross-casimir}.

On a highest-weight vector of an irreducible $K$-type, $Z$ acts by
$R=\lambda_1-\lambda_m$. For a nontrivial type, $R$ is a positive
integer. Since $Q$ is self-adjoint and commutes with $K$, each
eigenspace of $Q$ is $K$-invariant. Apply
\eqref{eq:ind-positive-identity} to a highest-weight vector in each
irreducible summand of each eigenspace. It follows that the
corresponding eigenvalue of $Q$ is at least $(N-m)R$. This also
handles multiplicities. Positivity on the trivial type completes
\eqref{eq:ind-relative-gap}.
\end{proof}


\subsubsection{Cross-block energy}

For $i\le m<a$, put
\[
X_{ia}=\frac{e_{ia}-e_{ai}}{\sqrt2},\qquad
Y_{ia}=\frac{\mathrm i(e_{ia}+e_{ai})}{\sqrt2}.
\]
These are Hilbert--Schmidt orthonormal in $\mathfrak{su}(N)$ for
$\langle X,Y\rangle=-\operatorname{Tr}(XY)$. Write
$D_Xf(g)=\left.\frac{\mathrm d}{\mathrm dt}\right|_{t=0}f(e^{tX}g)$.

\begin{lemma}\label{lem:ind-energy}
If $f$ is smooth, real, and $L$-Lipschitz for $d_1$, then at every
$g\in\operatorname{SU}(N)$,
\begin{equation}\label{eq:ind-energy-bound}
\sum_{i\le m<a}\bigl(\lvert D_{X_{ia}}f(g)\rvert^2+
\lvert D_{Y_{ia}}f(g)\rvert^2\bigr)\le2mL^2.
\end{equation}
\end{lemma}

\begin{proof}
For traceless Hermitian $A$,
$\lvert df_g(\mathrm iAg)\rvert\le L\lVert A\rVert_1$.
Extend this real linear functional from the traceless Hermitian
subspace by Hahn--Banach. Trace-norm duality gives a Hermitian
matrix $B$ with
\[
df_g(\mathrm iAg)=\operatorname{Tr}(BA),\qquad \lVert B\rVert_{\mathrm{op}}\le L.
\]
The extension is determined only up to adding a real multiple of the
identity, since the trace annihilates the traceless Hermitian subspace;
the entries $B_{ia}$ with $i\le m<a$, which are the only ones used
below, are the same for every admissible extension. The choice of $B$ is
pointwise; no measurable choice is required.
Let $P_m$ be projection onto the first $m$ coordinates. Direct
calculation gives
\begin{align*}
\sum_{i\le m<a}\bigl(\lvert D_{X_{ia}}f(g)\rvert^2+
\lvert D_{Y_{ia}}f(g)\rvert^2\bigr)
&=2\sum_{i\le m<a}\lvert B_{ia}\rvert^2\\
&=2\lVert P_mB(I-P_m)\rVert_{\mathrm{HS}}^2\le2mL^2.
\end{align*}
\end{proof}

\begin{theorem}\label{thm:ind-variance}
Every real $L$-Lipschitz $f$ on $(\operatorname{SU}(N),d_1)$ satisfies
\begin{equation}\label{eq:ind-variance}
\int_{\operatorname{SU}(N)}\operatorname{Var}_{k\in K}f(kh)\,\mathrm d h\le\frac{2mL^2}{N-m}.
\end{equation}
For a finite collection $f_1,\ldots,f_r$ with respective Lipschitz
constants $L_1,\ldots,L_r$, some $h\in\operatorname{SU}(N)$ satisfies
\begin{equation}\label{eq:ind-finite-variance}
\sum_{i=1}^r\operatorname{Var}_{k\in K}f_i(kh)
\le\frac{2m}{N-m}\sum_{i=1}^rL_i^2.
\end{equation}
\end{theorem}

\begin{proof}
\uses{lem:ind-energy, lem:ind-relative-gap}
Use the left regular representation $U(k)f(g)=f(k^{-1}g)$.
Its differential is $-D_X$, so the operator in
\eqref{eq:ind-cross-casimir} is
\[
Q=-\sum_{i\le m<a}(D_{X_{ia}}^2+D_{Y_{ia}}^2).
\]
Peter--Weyl decomposition applies Lemma~\ref{lem:ind-relative-gap} to
each finite-dimensional summand. Summing the resulting quadratic-form
inequality and integrating by parts gives, for smooth $f$,
\begin{equation}\label{eq:ind-form-bound}
\lVert f-\operatorname{Avg}_Kf\rVert_2^2\le\frac{1}{N-m}
\int_{\operatorname{SU}(N)}\sum_{i\le m<a}
\bigl(\lvert D_{X_{ia}}f\rvert^2+\lvert D_{Y_{ia}}f\rvert^2\bigr)\,\mathrm d g.
\end{equation}
Here $(\operatorname{Avg}_Kf)(g)=\int_Kf(kg)\,\mathrm d k$, which is the projection
$\operatorname{P}^{\mathrm{fix}}_K$ of Lemma~\ref{lem:ind-relative-gap} for the left regular
representation. Haar invariance shows that
\[
\lVert f-\operatorname{Avg}_Kf\rVert_2^2=\int_{\operatorname{SU}(N)}\operatorname{Var}_{k\in K}f(kg)\,\mathrm d g.
\]
Lemma~\ref{lem:ind-energy} proves \eqref{eq:ind-variance} for smooth functions.
The quadratic-form passage can equivalently be obtained first for
finite Peter--Weyl sums and then by their convergence in the relevant
first-order Sobolev norm.

For general Lipschitz $f$, pass to the right convolutions
$f_\varepsilon(g)=\int_{\operatorname{SU}(N)}f(gu)\varphi_\varepsilon(u)\,\mathrm d u$ with a
smooth nonnegative Haar approximate identity $\varphi_\varepsilon$.
Right convolution is smooth in the left translation directions
differentiated above. Right invariance of $d_1$ preserves the Lipschitz
constant, and the convolutions converge uniformly. Their variances
therefore converge uniformly, proving \eqref{eq:ind-variance}. Sum the
inequality over $i$ and choose a point where the sum is no greater
than its average to obtain \eqref{eq:ind-finite-variance}.
\end{proof}

In particular, if
\[
\Delta_{m,r}(N)=
\sup_{\substack{f_i:\operatorname{SU}(N)\to[0,1]\\\operatorname{Lip}_{d_1}(f_i)\le1}}
\min_{h\in\operatorname{SU}(N)}\sum_{i=1}^r\operatorname{Var}_{k\in\operatorname{SU}(m)}f_i(kh),
\]
then $\Delta_{m,r}(N)\le2mr/(N-m)$. The functions in the supremum
may depend on $N$.


\subsection{Extreme amenability in trace norm}\label{sec:ind-trace-ea}

The following elementary dynamical lemma makes the passage from
finite observables to fixed points explicit.

\begin{lemma}\label{lem:ind-variance-to-fixed}
\uses{def:ea}
Let a metrizable topological group $\Gamma$ have a compatible right-invariant
metric and contain an increasing sequence of compact subgroups $K_m$
with dense union. Suppose that for each $m$ and every finite family
of bounded real Lipschitz functions $b_1,\ldots,b_r$ on $\Gamma$,
\begin{equation}\label{eq:ind-flatness-hyp}
\inf_{h\in\Gamma}\sum_i\operatorname{Var}_{k\in K_m}b_i(kh)=0.
\end{equation}
Then $\Gamma$ is extremely amenable.
\end{lemma}

\begin{proof}
Consider a continuous action of $\Gamma$ on a nonempty compact Hausdorff
space $Z$, and fix $z_*\in Z$. If $A\in C(Z,\mathbb R)$, the orbit
function $a(g)=A(gz_*)$ is bounded and right-uniformly continuous,
since
\[
\sup_g\lvert a(ug)-a(g)\rvert\le
\sup_{z\in Z}\lvert A(uz)-A(z)\rvert\longrightarrow0.
\]
For a compatible right-invariant metric this is ordinary uniform
continuity, since $d(ug,g)=d(u,e)$ for all $g$. Every bounded uniformly continuous real function on a
metric space is uniformly approximable by bounded Lipschitz
functions. For example, the functions
$a_L(x)=\inf_y\{a(y)+L d(x,y)\}$ are bounded and $L$-Lipschitz
and converge uniformly to $a$ as $L\to\infty$. Uniform continuity
controls points with $d(x,y)$ small, and boundedness excludes all
other points once $L$ is large.

Fix a finite subset $S\subseteq K_m$, finitely many
$A_1,\ldots,A_r\in C(Z,\mathbb R)$, and $\varepsilon>0$. Approximate the
associated orbit functions $a_i$ uniformly by bounded Lipschitz
functions $b_i$. Given $h$, let $c_i=\int_{K_m}b_i(kh)\,\mathrm d k$.
For $s\in S$, Haar invariance and $\lvert v-w\rvert^2\le
2\lvert v-c_i\rvert^2+2\lvert w-c_i\rvert^2$ give
\[
\int_{K_m}\sum_{s\in S}\sum_i
\lvert b_i(skh)-b_i(kh)\rvert^2\,\mathrm d k
\le4\lvert S\rvert\sum_i\operatorname{Var}_{k\in K_m}b_i(kh).
\]
By \eqref{eq:ind-flatness-hyp}, choose $h$ making the right side
arbitrarily small and then a $k$ where the integrand is small.
Choosing the uniform approximations sufficiently accurate produces
$z=khz_*$ such that
\[
\lvert A_i(sz)-A_i(z)\rvert<\varepsilon\qquad(s\in S,\ 1\le i\le r).
\]
Every finite subset of $\bigcup_mK_m$ lies in a common stage.
Compactness of $Z$, applied to the closed finite sets of approximate
fixed-point conditions, now gives a point at which
$A(sz)=A(z)$ for every $A\in C(Z,\mathbb R)$ and every
$s\in\bigcup_mK_m$. Continuous real functions separate points of
$Z$, so this point is fixed by the dense union. Continuity makes it
fixed by all of $\Gamma$.
\end{proof}

\begin{lemma}\label{lem:ind-density}
The trace-norm completion of $G_{\mathrm{tr}}$ is $\operatorname{SU}_1$.
\end{lemma}

\begin{proof}
The set $\operatorname{SU}_1$ is closed in $I+\mathcal S_1$ for the trace norm:
unitarity is closed under operator-norm convergence, and the
Fredholm determinant is trace-norm continuous. Thus it is complete.

Let $u\in\operatorname{SU}_1$. Functional calculus gives $u=e^{\mathrm iA}$
with $A=A^*\in\mathcal S_1$: choose principal arguments for the eigenvalues
of $u$, which accumulate only at $1$ because $u-I$ is compact
(so only finitely many eigenvalues can lie near $-1$), using
$\lvert t\rvert\le(\pi/2)\lvert e^{\mathrm it}-1\rvert$ for $\lvert t\rvert\le\pi$.
The identity $\det_F(e^{\mathrm iA})=e^{\mathrm i\operatorname{Tr} A}$ gives
$\operatorname{Tr} A\in2\pi\mathbb Z$.

Let $P_n$ be the standard coordinate projection and put
$A_n=P_nAP_n$. Define $h_n$ on the first $n+1$ coordinates by
\[
h_n=e^{\mathrm iA_n|_{P_n\ell^2}}
\oplus e^{-\mathrm i\operatorname{Tr} A_n},
\]
and extend it by the identity. Then $h_n\in\operatorname{SU}(n+1)$.
Since $A_n\to A$ in trace norm, the integral formula for the
difference of two exponentials gives
$\lVert e^{\mathrm iA_n}-e^{\mathrm iA}\rVert_1\le\lVert A_n-A\rVert_1$.
Also $e^{-\mathrm i\operatorname{Tr} A_n}\to1$. The triangle inequality therefore
gives $h_n\to u$ in trace norm.
\end{proof}

\begin{proof}[Proof of Theorem~\ref{thm:ind-trace-intro}]
\proves{thm:ind-trace-intro}
\uses{lem:ind-density, lem:ind-variance-to-fixed, thm:ind-variance}
The estimate is Theorem~\ref{thm:ind-variance}. For either
$\Gamma=G_{\mathrm{tr}}$ or $\Gamma=\operatorname{SU}_1$, restrict any finite
family of bounded Lipschitz functions on $\Gamma$ to $\operatorname{SU}(N)$.
For each fixed $m$, inequality \eqref{eq:ind-finite-variance} tends to
zero as $N\to\infty$. Thus \eqref{eq:ind-flatness-hyp} holds.
The trace-norm metric is bi-invariant, in particular right-invariant,
and the compact subgroups
$\operatorname{SU}(m)$ have dense union in $\Gamma$, by Lemma~\ref{lem:ind-density} in
the second case. Lemma~\ref{lem:ind-variance-to-fixed} proves extreme
amenability.
\end{proof}


\subsection{An explicit separation of the right uniformities}\label{sec:ind-separation}

The difference between the final and trace-norm topologies on $G_{\mathrm{alg}}$
can be seen directly at
the level of bounded right-uniformly continuous functions. On
\[
E=\left\{x:\sum_{j\ge1}2^{-j}\lvert x_j\rvert^2<\infty\right\},
\]
let $T_g$ be the finitary matrix action of $g\in G_{\mathrm{alg}}$.
Each $T_g$ is a bounded invertible operator on $E$. On every finite
stage, $g\mapsto T_g$ is continuous in operator norm, so it is
continuous for the final topology. Put
\[
\ell(g)=\log\max\{\lVert T_g\rVert_{\mathrm{op}},\lVert T_g^{-1}\rVert_{\mathrm{op}}\},
\qquad f(g)=\min\{1,\ell(g)\}.
\]
The function $\ell$ is nonnegative, symmetric, continuous, and
subadditive. Hence
$\lvert f(ug)-f(g)\rvert\le\ell(u)$ uniformly in $g$, and
$f\in\operatorname{RUCB}(G_{\mathrm{fin}})$.

For $n\ge2$, let $u_n$ be a real determinant-one rotation on the
plane of coordinates $1,n$, through the angle
$\theta_n=4\,2^{-(n-1)/2}$, and the identity elsewhere. Then
\[
\lVert u_n-I\rVert_1=4\sin(\theta_n/2)\longrightarrow0.
\]
In an orthonormal basis of this plane for the weighted norm of $E$,
one off-diagonal entry of $T_{u_n}$ has absolute value
\[
2^{(n-1)/2}\sin\theta_n\longrightarrow4.
\]
Thus $\ell(u_n)>1$ eventually, and $f(u_n)=1$ while $f(e)=0$.
The function $f$ is not trace-norm continuous. In particular, the final
topology is strictly finer than the trace-norm topology, and final-RUC
functions need not be uniformly approximable by trace-Lipschitz
functions.

This example separates uniformities; the actual fixed-point
obstruction is Proposition~\ref{prop:ind-obstructing-observables}.

\subsection{Proofs of the statements in Section~\ref{sec:inductive-limits}}\label{sec:ind-encyclopaedia-proofs}

\begin{proof}[Proof of Corollary~\ref{cor:ind-classical-intro}]
\proves{cor:ind-classical-intro}
\uses{cor:ind-semisimple-classical, prop:ind-determinants}
Corollary~\ref{cor:ind-semisimple-classical} covers $\varinjlim\operatorname{SU}(n)$,
$\varinjlim\operatorname{SO}(n)$ and $\varinjlim\operatorname{Sp}(n)$, and
Proposition~\ref{prop:ind-determinants} covers $\varinjlim\operatorname{U}(n)$ and
$\varinjlim\operatorname{O}(n)$.
\end{proof}

\begin{proof}[Proof of Theorem~\ref{thm:ind-dynamics-intro}]
\proves{thm:ind-dynamics-intro}
\uses{prop:ind-bohr-classical, thm:ind-classical-umf}
Together, Proposition~\ref{prop:ind-bohr-classical} and
Theorem~\ref{thm:ind-classical-umf} prove the theorem.
\end{proof}

\begin{proof}[Proof of Theorem~\ref{thm:ind-free-intro}]
\proves{thm:ind-free-intro}
\uses{prop:ind-involution-fixed, thm:ind-classical-free-locus}
Together, Proposition~\ref{prop:ind-involution-fixed} and
Theorem~\ref{thm:ind-classical-free-locus} prove the theorem.
\end{proof}

\begin{proof}[Proof of Theorem~\ref{thm:lupini-inductive}]
\proves{thm:lupini-inductive}
\uses{thm:ind-main}
This is Theorem~\ref{thm:ind-main}. The two forms of the hypothesis agree: as
noted after~\eqref{eq:ind-family}, a compact Lie group $H$ is solvable if and
only if $H^0$ is a torus and $H/H^0$ is solvable, and a compact connected Lie
group is abelian exactly when it is a torus.
\end{proof}

\begin{proof}[Proof of Corollary~\ref{cor:lupini-inductive}]
\proves{cor:lupini-inductive}
\uses{cor:ind-classical-intro}
This is Corollary~\ref{cor:ind-classical-intro}, with
$\operatorname{SU}(\infty)=\varinjlim\operatorname{SU}(n)$ and so on, as in
Section~\ref{sec:ind-main-results}.
\end{proof}

\begin{proof}[Proof of Theorem~\ref{thm:lupini-inductive-umf}]
\proves{thm:lupini-inductive-umf}
\uses{thm:ind-dynamics-intro}
This is Theorem~\ref{thm:ind-dynamics-intro}.
\end{proof}

\begin{proof}[Proof of Theorem~\ref{thm:lupini-inductive-free}]
\proves{thm:lupini-inductive-free}
\uses{prop:ind-involution-fixed, thm:ind-classical-free-locus}
By Proposition~\ref{prop:ind-involution-fixed}, every action of $G$ on a nonempty
compact Hausdorff space has a point fixed by an element of order two; so no
such action is free, and in particular the action on $M(G)$ is not free. The
assertions about the points with trivial stabiliser, and the existence of
invariant Radon probability measures on $M(G)$, are
Theorem~\ref{thm:ind-classical-free-locus}.
\end{proof}

\begin{proof}[Proof of Theorem~\ref{thm:lupini-trace-class}]
\proves{thm:lupini-trace-class}
\uses{thm:ind-trace-intro}
This is the first assertion of Theorem~\ref{thm:ind-trace-intro}: the group
$G_{\mathrm{tr}}$ is $\bigcup_n\operatorname{SU}(n)$ with the trace-norm metric,
and $\operatorname{SU}_1$ is its completion~\eqref{eq:ind-suone}.
\end{proof}


\bibliographystyle{plain}
\bibliography{references}
\end{document}
